\documentclass[10pt,a4paper]{article}

% ----------------- Paquetes -------------------%
\usepackage[T1]{fontenc}
\usepackage[utf8]{inputenc}
\usepackage[a4paper,margin=2.5cm]{geometry}
\usepackage{mathptmx}             
\usepackage{changepage} 
\usepackage{setspace}
\usepackage[spanish]{babel}
\usepackage[labelfont=bf,labelsep=space]{caption}
\usepackage{amssymb}
\usepackage{amsmath}
\usepackage{ebgaramond}
\usepackage{placeins}
\usepackage{etoolbox}
\usepackage{setspace}
\usepackage{parskip} 
\usepackage{tcolorbox}
\usepackage{needspace}
\usepackage{xparse}
\usepackage{ocgx2}
\usepackage{hyperref}
\usepackage{hypcap}


%%%%%%%%%%%%%%%%%%%%%%%%%%%%%%%%%%%%%%%%%%%%%%%%%%%%%%%%%%%%%%%%
% --- Fija primero tus valores globales "buenos" ---
\setstretch{1.2}                 % Interlineado global
\setlength{\parskip}{0.7\baselineskip}
\setlength{\parindent}{0pt}
\newlength{\GlobalParskip}
\setlength{\GlobalParskip}{\parskip}

\newcounter{eqnum}[subsection]
\renewcommand{\theeqnum}{\arabic{eqnum}}
\newcounter{alphaitem}
\newcounter{numitem}

%% ------------------- Recuadros----------------------%%
\newcommand{\puntosclave}[1]{%
  \begin{tcolorbox}[
    colframe=black,
    title={Puntos clave},
    before upper={\setlength{\parindent}{0.5cm}\setlength{\parskip}{0.5\baselineskip}}
  ]
  #1
  \end{tcolorbox}
}

\newcommand{\modificaciones}[1]{
  \begin{tcolorbox}[
    colback=white,
    colframe=red,
    title={Modificaciones a realizar},
    before upper={\setlength{\parindent}{0.5cm}\setlength{\parskip}{0.5\baselineskip}}
  ]
  #1
  \end{tcolorbox}
}

%% ------------------- Preguntas TEST----------------------%%
% Opciones: radio (single-choice)
\NewDocumentCommand{\OptRadio}{m m m}{%
  \noindent\CheckBox[radio,name=#1,value=#2,width=1.2ex,height=1.2ex]{}%
  \;\textbf{#2.}~#3\par
}

% Opciones: checkbox (multi-choice)
\NewDocumentCommand{\OptCheck}{m m}{%
  \noindent\CheckBox[name=#1,width=1.2ex,height=1.2ex,bordercolor={0 0 0}]{}%
  \;#2\par
}

% Pregunta single-choice (radio)
% Args: 1=id, 2=enunciado, 3=opciones, 4=solución+justificación, 5=imagen (o {}), 6=origen
\NewDocumentCommand{\PreguntaTestSingle}{m m m m m m}{%
  \par\medskip
  \noindent\textbf{#1.}~#2\par\smallskip

    % Imagen (si no es {})
    \IfBlankTF{#5}{}{%
      \begin{center}
        \includegraphics[width=0.75\linewidth]{#5}
      \end{center}
    }

    \begin{Form}
      #3
    \end{Form}

    \par\smallskip
    \switchocg{sol-#1}{\fbox{Mostrar / ocultar solución}}%
    \hfill{\footnotesize #6}\par\smallskip

    \begin{ocg}{Solución}{sol-#1}{0}
      \noindent #4
    \end{ocg}
  \par\medskip
}

% Pregunta multi-choice (checkboxes)
\NewDocumentCommand{\PreguntaTestMulti}{m m m m m m}{%
  \par\medskip
  \noindent\textbf{#1.}~#2\par\smallskip

    % Imagen (si no es {})
    \IfBlankTF{#5}{}{%
      \begin{center}
        \includegraphics[width=0.75\linewidth]{#5}
      \end{center}
    }
    \begin{Form}
      #3
    \end{Form}

    \par\smallskip
    \switchocg{sol-#1}{\fbox{Mostrar / ocultar solución}}%
    \hfill{\footnotesize #6}\par\smallskip

    \begin{ocg}{Solución}{sol-#1}{0}
      \noindent #4
    \end{ocg}
  \par\medskip
}

%% ------------------- Listas----------------------%%
    % --- Lista alfabética ---
    \newenvironment{la}
      {\begin{list}{\alph{alphaitem})}{%
          \usecounter{alphaitem}%
          \setlength{\leftmargin}{1cm}%
          \setlength{\parskip}{3pt}% 
          \setlength{\itemsep}{0pt}% ← espacio entre ítems
          \setlength{\parsep}{0pt}%  ← párrafos dentro de un ítem
      }}
      {\end{list}}
      
    % --- Lista numérica ---
    \newenvironment{lnum}
      {\begin{list}{\arabic{numitem}.}{%
          \usecounter{numitem}%
          \setlength{\leftmargin}{1cm}%
          \setlength{\parskip}{3pt}% 
          \setlength{\itemsep}{0pt}% ← espacio entre ítems
          \setlength{\parsep}{0pt}%  ← párrafos dentro de un ítem
      }}
      {\end{list}}
      
% ------------------- Imágenes----------------------%
    \graphicspath{{Img/}}% Rutas por defecto 
    % --- Imagen adaptativa ---
    % Registros de dimensión definidos una sola vez
    \newdimen\imgcap
    \newdimen\imgmin
    \newdimen\imgmax
    \newdimen\imgremain
    \newdimen\imgh
    \newdimen\imgneed
    
    \newcommand{\imagenfit}[3]{%
      \addvspace{10pt}% separación antes
      \par\begingroup
        % Parámetros de composición (asignaciones, no \newdimen)
        \imgcap=1\baselineskip            % alto estimado de título+fuente
        \imgmin=0.15\textheight
        \imgmax=0.15\textheight
        \imgremain=\pagegoal
        \ifdim\pagegoal=\maxdimen \imgremain=0pt\fi
        \advance\imgremain by -\pagetotal
        \advance\imgremain by -\imgcap
        % Decisión de altura destino
        \ifdim\imgremain<\imgmin
          \newpage
          \imgh=0.20\textheight
        \else
          \imgh=\imgremain
          \ifdim\imgh>\imgmax \imgh=\imgmax \fi
        \fi
        % Reservar espacio para evitar cortes del bloque completo
        \imgneed=\imgh \advance\imgneed by \imgcap
        \Needspace{\imgneed}
        % Cuerpo
        \noindent\begin{minipage}{\textwidth}
          \centering
          \captionof{figure}{\textemdash\ #2}\par\vspace{0.3em}% título arriba
          \includegraphics[width=\linewidth,height=\imgh,keepaspectratio]{\detokenize{#1}}\par
          \vspace{0.3em}\small\textit{Fuente: }#3% fuente abajo
        \end{minipage}%
      \endgroup
      \par\addvspace{10pt}%
    }

%------------ Bloque de RECUADRO ESPECIAL -------------%
    \newenvironment{specialblock}[1]{%
      \begin{quote} % crea sangría a izquierda y derecha
      \small % texto más pequeño
      \noindent\textbf{#1}\\[-0.3em] % título en negrita
      \rule{\linewidth}{0.4pt}\\[0.6em]}{
      \end{quote}}
      
%-------------------------- Portada -----------------------%
\newcommand{\oldtitle}[1]{\def\OldTitle{#1}}
\newcommand{\changerationale}[1]{\def\ChangeWhy{#1}}

\newcommand{\changebox}{%
\begin{tcolorbox}[title={Historial de cambios del tema}]
\textbf{Título anterior:} \OldTitle\par
\textbf{Motivo del cambio:} \ChangeWhy
\end{tcolorbox}
}

%-------------------------- Author Font -----------------------%
    \newcommand{\authorfont}[1]{{\fontfamily{EBGaramond-TLF}\selectfont #1}}

%-------------------------- Anexos -----------------------%
    \makeatletter
    \newcounter{anexo}
    \renewcommand{\theanexo}{\Roman{anexo}}
    \newcommand{\anexo}[1]{%
      \clearpage
      \phantomsection
      \refstepcounter{anexo}%
      \noindent\textbf{Anexo \theanexo.- }#1\par\medskip
      \addcontentsline{stc}{section}{Anexo \theanexo.- #1}%
    }
    \makeatother

    %-------------------------- Lista de figuras (personal) -----------------------%
\makeatletter
\newcommand{\MyListOfFigures}{%
  \clearpage
  \phantomsection
  {\centering\bfseries\Large Índice de figuras\par}%
  \medskip
  % Entrada en nuestro índice de contenido (stc)
  \addcontentsline{stc}{section}{Índice de figuras}%
  % Recuperación del listado (sin cabecera propia)
  \begingroup
    \parindent=0pt\parskip=2pt
    \@starttoc{lof}%
  \endgroup
}
\makeatother

%-------------------------- Bibliografía -----------------------%
\makeatletter
% Contador para las referencias
\newcounter{myref}
% Cabecera de la bibliografía, entra en el índice 'stc'
\newcommand{\MyBibliography}{%
  \clearpage
  \phantomsection
  {\centering\bfseries\Large Bibliografía y referencias\par}%
  \medskip
  \addcontentsline{stc}{section}{Bibliografía y referencias}%
  \setcounter{myref}{0}%
}
\newcommand{\MyBibItem}[4]{%
  \refstepcounter{myref}%
  \noindent[\themyref]\ #1.\ \emph{#2}. #3. #4\par
}
\makeatother

%--------- Establecer cuatro niveles de índice --------%
    \setcounter{tocdepth}{4}   
    \setcounter{secnumdepth}{4}% numera hasta \paragraph

%%%%%%%%%%%%%%%%%%%%%%%%%%%%%%%%

\makeatletter

    %--------- Interlineado Global --------%
    \edef\GlobalBaseLineStretch{\baselinestretch}
    
    %--------- Espaciado de seccinones --------%
    \renewcommand\section{\@startsection{section}
      {1}{\z@}
      {2.5ex \@plus 1ex \@minus .2ex} % espacio antes
      {1.5ex \@plus .2ex}             % espacio después
      {\normalfont\Large\bfseries}    % formato del título
    }
    \renewcommand\subsection{\@startsection{subsection}{2}{\z@}
      {3ex \@plus 0.8ex \@minus .2ex}
      {1.5ex \@plus .2ex}
      {\normalfont\large\bfseries}
    }
    \renewcommand\subsubsection{\@startsection{subsubsection}{3}{\z@}
      {3ex \@plus 0.5ex \@minus .2ex}
      {1ex \@plus .2ex}
      {\normalfont\normalsize\bfseries}
    }
    \renewcommand\paragraph{\@startsection{paragraph}{4}{\z@}%
    {-2ex \@plus -0.5ex \@minus -.2ex}% espacio antes
    {0.8ex \@plus .2ex}% espacio después
    {\normalfont\normalsize\itshape}}% ← cursiva en lugar de negrita

    %--------- Ecuaciones --------%   
    % Variables globales para almacenar títulos y notas
    \gdef\leq@current@section{}%
    \gdef\leq@current@subsection{}%
    \gdef\leq@last@printed@section{}%
    \gdef\leq@last@printed@subsection{}%
    
    % Comando para añadir nota (invisible en el cuerpo)
    \newcommand{\eqnote}[1]{%
      \addcontentsline{leq}{leqnote}{#1}%
    }
    
    % Comando para ecuaciones
    \newcommand{\eqblock}[2]{%
      % Verificar si necesitamos imprimir el título de sección
      \ifx\leq@current@section\leq@last@printed@section\else
        \addcontentsline{leq}{leqsection}{\leq@current@section}%
        \global\let\leq@last@printed@section\leq@current@section
        \gdef\leq@last@printed@subsection{}% Reset subsección al cambiar sección
      \fi
      % Verificar si necesitamos imprimir el título de subsección
      \ifx\leq@current@subsection\leq@last@printed@subsection\else
        \ifx\leq@current@subsection\@empty\else
          \addcontentsline{leq}{leqsubsection}{\leq@current@subsection}%
          \global\let\leq@last@printed@subsection\leq@current@subsection
        \fi
      \fi
      % Ecuación
      \refstepcounter{eqnum}%
      \begin{equation*}
        #1\tag{E.\theeqnum}
      \end{equation*}%
      \addcontentsline{leq}{leqt}{[E.\theeqnum] -- #2}%
      \addcontentsline{leq}{leqm}{#1}%
    }
    
    %%% ----- Índice de ecuaciones ----------%%%
    % Tamaño para []
    \newcommand{\leqIndexFont}{\normalsize}
    \newcommand{\leqTitleFont}{\tiny}
    \newcommand{\leqNoteFont}{\tiny}
    % Cabecera de sección 
    \newcommand*\l@leqsection[2]{%
      \addvspace{25pt}% Mayor separación antes de nueva sección
      \noindent\leqIndexFont
      \makebox[\linewidth]{\rule{.38\linewidth}{0.4pt}\ \ \textbf{  \MakeUppercase{#1}}\ \ \rule{.38\linewidth}{0.4pt}}%
      \par\vspace{-10pt}
    }
    % Cabecera de subsección 
    \newcommand*\l@leqsubsection[2]{%
      \par\addvspace{15pt}%
      \noindent\leqIndexFont #1
      \par\vspace{-7pt}
    }
    % Título de ecuación 
    \newcommand*\l@leqt[2]{%
    \par\addvspace{10pt}%
      \par\noindent\leqTitleFont #1\par
    }
    % Miniatura de ecuación
    \newcommand*\l@leqm[2]{%
      \vspace{0.3\baselineskip}%
      \par\scriptsize\noindent\hspace*{1.5em}\(#1\)\par%
    }
    % Nota de ecuación 
    \newcommand*\l@leqnote[2]{%
      \vspace{0.2\baselineskip}%
      \par\noindent\leqNoteFont\hspace*{1.5em}\textbf{Nota.--} #1\par\vspace{0.5\baselineskip}%
    }
    % Generación del índice
    \newcommand{\mylistofequations}{%
      \clearpage
      \phantomsection
      % Título visual de la lista (ajústalo a tu gusto):
      {\centering\bfseries\Large Índice de ecuaciones\par}
      % Entrada en nuestro índice 'stc':
      \addcontentsline{stc}{section}{Índice de ecuaciones}%
      % Llamada a tu lista real (usa el nombre exacto de tu macro):
      \@starttoc{leq}%
    }
    
    %--------- Captura de títulos de sección --------%
    \let\leq@original@section\section
    \let\leq@original@subsection\subsection
    % Flag para saber si estamos en una sección asterisco
    \newif\ifLeqStarSection
    \LeqStarSectionfalse
    
    % Redefinir \section CON CONTROL DE ASTERISCO
    \renewcommand{\section}{%
      \@ifstar{\leq@section@star}{\@dblarg\leq@section@nostar}%
    }
    \def\leq@section@star#1{%
      \global\LeqStarSectiontrue
      \gdef\leq@current@section{#1}%
      \gdef\leq@current@subsection{}%
      \leq@original@section*{#1}%
      \global\LeqStarSectionfalse
    }
    \def\leq@section@nostar[#1]#2{%
      \global\LeqStarSectionfalse
      \gdef\leq@current@section{#2}%
      \gdef\leq@current@subsection{}%
      \leq@original@section[#1]{#2}%
    }
    % Redefinir \subsection
    \renewcommand{\subsection}{%
      \@ifstar{\leq@subsection@star}{\@dblarg\leq@subsection@nostar}%
    }
    \def\leq@subsection@star#1{%
      \global\LeqStarSectiontrue
      \gdef\leq@current@subsection{#1}%
      \leq@original@subsection*{#1}%
      \global\LeqStarSectionfalse
    }
    \def\leq@subsection@nostar[#1]#2{%
      \global\LeqStarSectionfalse
      \gdef\leq@current@subsection{#2}%
      \leq@original@subsection[#1]{#2}%
    }

    % ----- Restauración de valores Globales de estilo ----------%
    % Flags
    \newif\ifInFootnote
    \InFootnotefalse
    \pretocmd{\@footnotetext}{\InFootnotetrue}{}{}
    \apptocmd{\@footnotetext}{\InFootnotefalse}{}{}
    % Profundidad de listas (soporta anidadas)
    \newcount\MyListDepth \MyListDepth=0
    \AtBeginEnvironment{la}{\advance\MyListDepth by 1}
    \AfterEndEnvironment{la}{\advance\MyListDepth by -1}
    \AtBeginEnvironment{lnum}{\advance\MyListDepth by 1}
    \AfterEndEnvironment{lnum}{\advance\MyListDepth by -1}
    \AtBeginEnvironment{itemize}{\advance\MyListDepth by 1}
    \AfterEndEnvironment{itemize}{\advance\MyListDepth by -1}
    \AtBeginEnvironment{enumerate}{\advance\MyListDepth by 1}
    \AfterEndEnvironment{enumerate}{\advance\MyListDepth by -1}
    % Restauración diferida: hazlo al INICIO del próximo párrafo real
    \newif\if@pendingRestore
    \@pendingRestorefalse
    \newcommand{\RequestRestore}{\global\@pendingRestoretrue}
    \everypar{%
      \if@pendingRestore
        \global\@pendingRestorefalse
        \ifInFootnote\else
          \ifnum\MyListDepth=0
            \setlength{\parskip}{\GlobalParskip}%
            \setstretch{\GlobalBaseLineStretch}%
          \fi
        \fi
      \fi
    }
    % Al final de bloques:
    \AfterEndEnvironment{equation}{\RequestRestore}
    \AfterEndEnvironment{equation*}{\RequestRestore}
    \AfterEndEnvironment{la}{\RequestRestore}
    \AfterEndEnvironment{lnum}{\RequestRestore}
    % Al final de macros:
    \apptocmd{\eqblock}{\RequestRestore}{}{}
    \apptocmd{\imagenfit}{\RequestRestore}{}{}
    
\makeatother

% ===================== ToC propio con 4 niveles (único bloque) =====================
\makeatletter

% --- Escritores: añadimos una línea al .stc en cada cabecera numerada ---
% Guardamos las versiones actuales para llamar al comportamiento normal
\let\MyTOC@orig@section\section
\let\MyTOC@orig@subsection\subsection
\let\MyTOC@orig@subsubsection\subsubsection
\let\MyTOC@orig@paragraph\paragraph

% SECTION
\renewcommand{\section}{%
  \@ifstar{\MyTOC@section@star}{\@dblarg\MyTOC@section@nostar}%
}
\def\MyTOC@section@star#1{%
  \MyTOC@orig@section*{#1}% sin entrada al .stc
}
\def\MyTOC@section@nostar[#1]#2{%
  \MyTOC@orig@section[{#1}]{#2}% comportamiento normal (escribe al .toc)
  % Escribimos también nuestra línea al .stc (texto con \numberline)
  \addcontentsline{stc}{section}{\protect\numberline{\thesection}#2}%
}

% SUBSECTION
\renewcommand{\subsection}{%
  \@ifstar{\MyTOC@subsection@star}{\@dblarg\MyTOC@subsection@nostar}%
}
\def\MyTOC@subsection@star#1{%
  \MyTOC@orig@subsection*{#1}%
}
\def\MyTOC@subsection@nostar[#1]#2{%
  \MyTOC@orig@subsection[{#1}]{#2}%
  \addcontentsline{stc}{subsection}{\protect\numberline{\thesubsection}#2}%
}

% SUBSUBSECTION
\renewcommand{\subsubsection}{%
  \@ifstar{\MyTOC@subsubsection@star}{\@dblarg\MyTOC@subsubsection@nostar}%
}
\def\MyTOC@subsubsection@star#1{%
  \MyTOC@orig@subsubsection*{#1}%
}
\def\MyTOC@subsubsection@nostar[#1]#2{%
  \MyTOC@orig@subsubsection[{#1}]{#2}%
  \addcontentsline{stc}{subsubsection}{\protect\numberline{\thesubsubsection}#2}%
}

% PARAGRAPH
\renewcommand{\paragraph}{%
  \@ifstar{\MyTOC@paragraph@star}{\@dblarg\MyTOC@paragraph@nostar}%
}
\def\MyTOC@paragraph@star#1{%
  \MyTOC@orig@paragraph*{#1}%
}
\def\MyTOC@paragraph@nostar[#1]#2{%
  \MyTOC@orig@paragraph[{#1}]{#2}%
  % Si no numeras los \paragraph, cambia \theparagraph por texto plano sin \numberline
  \addcontentsline{stc}{paragraph}{\protect\numberline{\theparagraph}#2}%
}

% --- Impresor del índice propio (lee el .stc) ---
\newcommand{\MyTOC}{%
  \par\bigskip
  {\centering\bfseries\Large Índice de contenido\par}%
  \medskip
  \begingroup
    \parindent=0pt\parskip=2pt
    \@starttoc{stc}%
  \endgroup
}

\makeatother
% ===================== fin ToC propio =====================


%%%%%%%%%%%%%%


% --- Datos del tema ---
\title{Unión Europea (I)\\
\large Tratados, orden jurídico e instituciones de la Unión Europea}
\author{Marcelo Ortega}
\date{Marzo 2026}
\newcommand{\TemaCode}{3B36}

\oldtitle{
    B.37. La Unión Europea: Instituciones y Orden jurídico. Los tratados de la Unión Europea
}
\changerationale{
    Ajuste de redacción.
}

\usepackage{soul}\sethlcolor{yellow}% resaltado amarillo: \hl{texto} (Panel TCEE, ⌃H)
\begin{document}


\maketitle
\changebox

\newpage

% --- Página de resumen y modificaciones ---
\puntosclave{ %% Escribir puntos clave Apuntes CECO
     El último apartado (Gestión de crisis con el Tratado de Lisboa) permitiría introducir las cuestiones de actualidad que el opositor estime relevantes; las novedades han de seguirse por prensa y consultando las páginas web institucionales.
    }
\vspace{1cm}

\modificaciones{
    
    }

\newpage
\MyTOC
\newpage

\mylistofequations
\clearpage

\MyListOfFigures
\clearpage


\pagenumbering{arabic}
\setcounter{page}{1}


\section{Introducción}



\subsection{Contextualización}


\subsubsection{Relevancia}




\subsubsection{Problemática}



\subsection{Estructura de la exposición}







\clearpage
\section{Unión Europea: Evolución histórica}
\puntosclave{

}

En este apartado, se distinguen cinco fases que ayudan a entender la construcción europea.


\subsection{Origen y expansión: 1951-1973}
El punto de partida de la Unión Europea se sitúa en el 9 de mayo de 1950, cuando Robert Schuman (ministro francés de Asuntos Exteriores) pronunció un discurso en el que proponía crear una organización entre Francia y Alemania, y los países que se unieran, para la producción y el consumo del carbón y el acero. En él, destacó dos elementos que aún a día de hoy guían el proyecto europeo:
La idea de integración europea nació de los efectos desastrosos de la Segunda Guerra Mundial y de la amenaza de enfrentamiento en el continente. La ¡juesta en común de los recursos franceses y alemanes suponía el final del conflicto histórico a propósito del control del sector siderúrgico y afianzaba así la paz en Europa.
Desde el comienzo, se han utilizado metas económicas pragmáticas y graduales para tratar de conseguir objetivos políticos. De fonna elocuente, Schuman declaraba: "Europa no se hará de golpe ni en una construcción de conjunto. Se construirá mediante logros concretos, creando primero una solidaridad de hecho". Según los "padres fundadores" de la Unión, la fusión de los intereses económicos de los Estados miembros crearía las condiciones para lograr más adelante una integración política de mayor envergadura.

La iniciativa de Schuman prosperó y, en 1951, Francia, la República Federal de Alemania, Italia, Bélgica, Países Bajos y Luxemburgo finnaron el Tratado de París, constitutivo de la Comunidad Europea del Carbón y del Acero (CECA). Con una duración prevista de 50 aflos, este Tratado estableció un mercado común para estos productos (libre circulación. respeto a la competencia,apoyo a la reconversión, etc.) y sentó la base de la arquitectura comunitaria, al contar con un órgano ejecutivo (la Alta Autoridad, presidida por primera vez por Jean Monnet), un Consejo de Ministros, una Asamblea Parlamentaria y un Tribunal de Justicia.

El Tratado de París representó solo el primer paso de un proceso más ambicioso. Tras su puesta en marcha, se siguieron estudia!1do las posibilidades de una unión económica general, así como una unión en el terreno nuclear en Europa, lo que condujo a la negociación de los Tratados de Roma. El Tratado constitutivo de la Comunidad Económica Europea (CEE) y el Tratado constitutivo de la Comunidad Europea de la Energía Atómica ("Euratom") fueron suscritos por los Estados miembros de la CECA en 1957 por un período de tiempo ilimitado. 

El objetivo de la CEE era establecer un mercado común que garantizara cuatro libertades (libre circulación de mercancías, servicios, personas y capitales). Para ello, el Tratado preveía la supresión de los derechos de aduana, la adopción de políticas comunes 1 (comercial, agrícola, pesca, transportes y competencia) y la constitución de diversos instrumentos para atenuar los costes de ajuste (Fondo Social Europeo, Banco Europeo de Inversiones, Comité Económico y Social).\footnote{%
    Además de otorgar competencias a la Comunidad en estos ámbitos. d Tratado introdujo una ·•cláusula de cierre"' o "comodín". que permitía a la Comunidad actuar e;:n competencias no mencionadas expresamente en el Tr.itado sin tener que refonnarlo. El objetivo de esta cláusula era no limitar demasiado la capacidad de actuación de la Comunidad a medida que avanzase la integración. Esta cláusula pennitió. por ejemplo. la incorporación ·'extraoficiar· al sistema comunitario de la política regional. la política industrial y tecnológica, la política de medio ambiente y la política social en 1972.
} Por su parte. el objetivo de Euratom era desarrollar de forma común la energía nuclear para fines pacíficos, coordinando el suministro de materiales y los programas de investigación. Desde el punto de vista institucional, el Tratado CEE y el Tratado Euratom desarrollaron esquemas paralelos al de la CECA. Los "ejecutivos" y los Consejos de Ministros de las tres Comunidades Europeas se mantuvieron separados hasta 1967, cuando entró en vigor el ·Tratado de Fusión" que unificó el Consejo y la Comisión, y estableció el principio de unidad presupuestaria. Aunque siguieron existiendo las Comunidades en plural. se impuso la lógica de la CEE, demostrando la preferencia que existía por la integración del conjunto de la actividad económica sobre la sectorial.

Hasta 1973, la evolución de la CEE coincidió con un periodo de notable expansión económica en el mundo occidental. Ello facilitó que la unión aduanera se completara en I 968. al tiempo que ésta reforzó el crecimiento económico europeo en una especie de círculo virtuoso. A su vez, la Política Agrícola Común (PAC) tuvo un desarrollo e impacto político muy significativo a lo largo de los años sesenta.

Los progresos en el funcionamiento de la CEE atrajeron muy pronto a otros países, en particular a los que formaban parte la Asociación Europea de Libre Comercio (EFTA) constituida en 1960. Reino Unido, Irlanda y Dinamarca solicitaron su adhesión en 1961, y Noruega en 1962. La primera ampliación culminó cuando entraron en vigor los Tratados de Adhesión del Reino Unido, Irlanda y Dinamarca, el I de enero de 1973. Pese a haber concluido las negociaciones, Noruega no entró porque las mayoría de su población votó en contra en referéndum.

La CEE no solo recibió solicitudes de adhesión, sino también de asociación comercial. Grecia y Turquía solicitaron su asociación con la CEE en 1959, y España y Portugal en 1962. Aunque en el caso de España se solicitaba "una asociación susceptible de llegar en su día a la plena , integración", la falta de un sistema democrático impedía su ingreso a la Comunidad. España habría de esperar a 1964 para que se iniciasen conversaciones con la CEE con vistas a negociar un Acuerdo Preferencial, el cual se terminó firmando en octubre de 1970.

\subsection{Crisis: 1973-1985}
En el contexto de la crisis económica internacional de los años setenta. hubo un estancamiento en el proceso de integración europea, al que se denominó "euro-esclerosis". Por una parte, surgieron o se mantuvieron toda una serie de barreras no arancelarias al comercio que perpetuaron la fragmentación del mercado común, mermando la posición de la Comunidad frente a Estados Unidos y Japón en cuanto a productividad, innovación y empleo. Por otra. se recrudecieron los conflictos entre los Estados miembros relacionados con la financiación comunitaria y se hicieron patentes varias disfuncionalidades institucionales, que se fueron solucionando, aunque siempre
con cierto aire de provisionalidad.
El Consejo funcionaba con decisiones por unanimidad en exclusiva, incluso en asuntos en los que no era necesaria. Hasta finales de los ochenta, imperó de facto el "compromiso de Luxemburgo" de 1966, según el cual, cuando estuvieran en juego intereses vitales de uno o varios países, los miembros del Consejo se esforzarían por llegar a soluciones aceptables por todos. Este acuerdo había resuelto la llamada crisis de la "silla vacía", durante la que Francia había dejado de participar en las reuniones para evitar que se tomaran decisiones (sobre financiación de la PAC o control presupuestario) que podían adoptarse por mayoría. 

Aunque externas al marco institucional de la Comunidad, las cumbres de jefes de Estado y de Gobierno de los Estados miembros comenzaron a desempeñar una función de impulso político y a resolver situaciones de impasse el Consejo (por ejemplo, relativas a la financiación o la ampliación). En 1974, se decidió que los jefes de Estado y de Gobierno se reuniesen periódicamente como "Consejo Europeo". 

Durante los setenta, también aumentó el peso institucional del Parlamento Europeo. Al tiempo que se creó un sistema de recursos propios en sustitución de las contribuciones financieras nacionales, en 1970, se concedió al Parlamento la potestad de dar su opinión sobre los gastos obligatorios y decidir sobre los gastos no obligatorios, y, en 1975, el derecho a rechazar el presupuesto y aprobar la gestión de la Comisión en la ejecución del mismo. Este aumento de competencias exigió conferir al Parlamento una legitimidad nueva, por lo que se introdujo la elección de los diputados por sufragio universal directo en 1976.

Por otro lado, en el contexto de inestabilidad cambiaria provocada por la quiebra de Bretton Woods, se llevaron a cabo los primeros esfuerzos de integración monetaria. En 1972, los Estados miembros establecieron la "Serpiente Monetaria Europea" para mantener los tipos de cambio bilaterales dentro de un margen restringido, la cual subsistió al inicio de la flotación libre y general de monedas frente al dólar, pero por poco tiempo. Tras el fracaso de la Serpiente, s�� logró establecer el Sistema Monetario Europeo (SME) en 1979, si bien éste mantuvo el carácter de mero acuerdo voluntario entre Bancos Centrales.

A pesar del contexto general de "euro-pesimismo", la Comunidad siguió ofreciendo un gran atractivo para otros países. En particular, los menos desarrollados del sur de Europa (Grecia, España y Portugal) solicitaron su ingreso en cuanto abandonaron sus regímenes políticos autoritarios, deseosos de estabilizar sus nuevos sistemas democráticos, impulsar su desarrollo y acabar con el aislamiento respecto a Europa. Grecia presentó su solicitud en 1975, y España y Portugal en 1977. El Tratado de Adhesión de Grecia entró en vigor el I de enero de 1981. España y Portugal firmaron sus Tratados de Adhesión en Madrid y Lisboa, respectivamente, el 12 de junio de 1985, entrando en vigor ambos el I de enero de 1986.

\subsection{Relanzamiento: 1985-1995}
El relanzamiento de la integración europea coincidió con una mejora del ciclo económico y el nombramiento de Jacques Delors como Presidente de la Comisión Europea en 1985. En enero de dicho año, la Comisión presentó un Libro Blanco que proponía medidas concretas para culminar, antes del 1 de enero de 1993. el \textit{gran mercado interior} previsto en el Tratado de Roma. En junio, se convocó una Conferencia lntergubernamental para estudiar la reforma institucional necesaria para lograr dicho objetivo. El resultado fue la firma del Acta Única Europea en 1986 por parte de los doce Estados miembros.

Este Tratado reconoció nuevas competencias comunitarias en ámbitos que se habían demostrado insuficientes los años anteriores (política monetaria, social, cohesión. l+D. medio ambiente, cooperación en política exterior). Desde el punto de vista institucional. "reimplantó" la regla de votación por mayoría cualificada en el Consejo y ro·rtaleció el papel del Parlamento Europeo introduciendo un "procedimiento de cooperación" legislativa entre el Consejo y el Parlamento. El Parlamento pasaba a contar con auténticas competencias legislativas en algunas áreas, aunque limitadas (por ejemplo. si después de examinar una determinada propuesta, el Parlamento la rechazaba, el Consejo podía seguir adoptándola por unanimidad). 

El clima favorable para la integración europea se extendió con el seísmo político provocado por la caída del muro de Berlín en noviembre de 1989. Por un lado, este contexto pudo servir para vencer las reticencias alemanas a seguir adelante con el Plan Delors para el establecimiento de una unión monetaria, el cual era un objetivo prioritario para Francia. Por otro, se hizo patente la necesidad de que el proyecto europeo contara con suficiente legitimidad democrática y eficacia institucional para responder a los desafios internacionales que se estaban viviendo. Los trabajos realizados por dos Conferencias lntergubernamentales (una relativa a la Unión Económica y Monetaria, y otra dedicada a la unión política) desembocaron en el Tratado de la Unión Europea (TUE) que los doce finnaron en 1992 en Maastricht. El Tratado de Maastricht supuso un avance fundamental. Contenía el acto fundacional de la UE (si bien no le dotó de personalidad jurídica) y se proponía la instauración progresiva de la unión monetaria (lo que darla lugar al nacimiento del euro en 1999). Este Tratado estableció tres pilares en la Unión.\footnote{%
    El TUE también instauró la .. ciudadanla europea", que se añadía a la ciudadanía nacional y suponía: derecho de libre circulación y residencia en la Unión: derecho a votar y a ser candidato en las elecciones europea��: derecho a protección diplomútica y consular de un Estado miembro distinto al de origen en un país tercero en el que el de origen no esté representado: derecho de petición ante el Parlamento Europeo y a presentar una denuncia allle el recién creado Defensor del Pueblo Europeo.
}

Primer pilar: La Comunidad Europea (CE). Este pilar estaba compuesto por la CEE, la Euratom y la CECA hasta su extinción. El primer pilar se distinguía por seguir el denominado método comunitario, caracterizado por: (i) el derecho exclusivo de iniciativa legislativa de la Comisión; (ii) el poder de decisión del Consejo y el Parlamento mediante un nuevo procedimiento legislativo, la "codecisión"; y (iii) la votación por mayoría cualificada en el Consejo. En este pilar se incluyeron: • Nuevas competencias en redes transeuropeas, política industrial. salud y protección del consumidor. educación, juventud y cultura. • Otras ya existentes pero ampliadas, tales como el medio ambiente, la investigación y la política social. En el ámbito de la política social, también se incorporó, anejo al Tratado, un Protocolo, del que Reino Unido estaba excluido, relativo a la promoción del empleo, la mejora de las condiciones de vida y trabajo, el diálogo social, etc.

Segundo pilar: La Política Exterior y de Seguridad Común (PESC). Este pilar se creó para permitir a los países emprender acciones comunes en materia de política exterior. Se previó que siguiera un método de funcionamiento intergubernamental, en el que:  (i) la Comisión comparte el derecho de iniciativa con los países; (ii) el Consejo Europeo suele desempeñar un papel fundamental; (iii) el Consejo decide, en general, por unanimidad; (iv) el Parlamento Europeo tiene un papel meramente consultivo; y (v) la jurisdicción del Tribunal de Justicia no aplica.

Tercer pilar: Cooperación en · 1os ámbitos de Justicia y Asuntos de Interior (fronteras exteriores, m igración y asilo, cooperación judicial, policial, etc.). Este pilar se creó con el objetivo de establecer un espacio de libertad, seguridad y justicia. También seguía un método intergubernamental.

En lo que concierne a la ampliación, rápidamente tras la caída del muro de Berlín, Alemania se reunificó y la antigua República Democrática Alemana se incorporó a la UE en noviembre de 1990, siendo ésta una ampliación de hecho pero no formal de la Unión. En este tiempo, los cambios acelerados que se estaban produciendo en el centro y el este de Europa provocaron que algunos países de la EFTA solicitaran su adhesión. Tras las negociaciones y los referéndums pertinentes, la cuarta ampliación de la UE, a Austria, Finlandia y Suecia, se produjo oficialmente el I de enero de 1 995, elevándose a qu ince el número de Estados miembros (en Noruega, que había solicitado su adhesión en 1992, sus ciudadanos volvieron a pronunciarse en contra).

\subsection{Proceso de reforma institucional: 1995-2014}
La cuarta ampliación hizo ya patente la necesidad de una importante reforma en un sistema que se había diseñado en ·principio para seis países. Ésta se convirtió en inevitable ante la perspectiva  de una ampliación de la Unión hacia el Este, posibilidad que se abrió oficialmente en el Consejo Europeo de Copenhague en junio de 1 993.\footnote{%
    En este Consejo Europeo, se establecieron unos ��riterios (los•· Criterios de Copenhague .. ) que han de cumplir los países candidatos para poder adherirse a la Unión Europea: (i) existencia de instituciones estables que garanticen la de111ocracia, el Estado de derecho, el respeto de los derechos humanos. y el respeto y la protección de las minorías: (ii) existencia de una economía de mercado en funcionamiento. y la capacidad de hacer frente a la presión competitiva y las fueras del mercado dentro de la UE: (iii) capacidad para asumir las obligaciones que se derivan de la adhesión. incluida la capacidad para poner en práctica de manera eficaz las nom1as. estándares y políticas que forman el acervo comunitario, y aceptar los objetivos de la unión política. económica y 111onetaria. Las decisiones sobre ampliación se to111an por unanimidad.
} Cada uno d�� los países de Europa Central y Oriental fueron solicitando su adhesión a la UE los años posteriores, uniéndose a Chipre y a Malta que lo habían hecho en 1990. Dado que esta ampliación no iba a tener precedentes en cuanto a dimensión y complejidad, mientras avanzaban las negociaciones con los países candidatos se puso en marcha el proceso de reforma institucional de la Unión. 

El Tratado de Ámsterdam, finnado en 1997 y que entró en vigor el I de mayo de 1999, fue el primer resultado de estos esfuerzos. En cuanto a las políticas, este Tratado: (i) supuso un avance en el ámbito social al establecer un mecanismo de coordinación de las políticas de empleo de los Estados miembros e incorporar al Tratado el Protocolo sobre política social; (ii) trasladó al primer pilar ("comunitarizó") importantes ámbitos del tercero, incluyendo la incorporación al Tratado del Convenio Schengen a fin de garantizar la libre circulación de los ciudadanos;\footnote{%
    Bélgica, Francia, Alemania, Luxemburgo y los Países Bajos suscribieron el Acuerdo de Schengcn en 1985 y elConvenio de Schcngen en 1990 con el fin de eliminar progresivamente los controles en sus fronteras interiores yestablecer un régimen de libre circulación para todos los nacionales de los países signatarios. En la actualidad. 26 países conforman el ·•espacio Scheng.:n··. Irlanda disfruta de e:xencion.:s (también lo hacía Reino Unido). Los países EFTA (Islandia. Liechtenstein. Noniega y Suiza) forman parte. Los países candidatos deben aceptar íntegramente el acervo de Schengen antes de su adhesión. No obstante. los controles solo se eliminan (por decisión unánime del Consejo) tras una evaluación que confirme que se han adoptado todas las medidas para dio.
} y (iii) trató de consolidar la política exterior de la Unión, mediante un nuevo cargo. el Alto Representante de la PESC (ocupado por el Secretario General del Consejo). No obstante, el Tratado de Ámsterdam apenas logró avances en el marco institucional.

Como la presión para preparar a la Unión para la gran ampliación iba en aumento (en 1999 se acordó que la UE debía estar lista a partir de finales de 2002), se volvió a convocar una Conferencia lntergubemamental, que desembocó en la finna del Tratado de Niza en 2001 , el cual entró en vigor el 1 de febrero de 2003. Los principales ajustes que introdujo este Tratado fueron: (i) limitar el número de Comisarios a uno por Estado miembro (eliminando el privilegio que hasta dicha fecha tenían los grandes, entre ellos España, de designar dos); (ii) establecer un nuevo sistema de votación por mayoría cualificada en el Consejo (doble mayoría, de votos y de población); (iii) algunas áreas pasaron total o parcialmente de la unanimidad a la mayoría cualificada, entre ellas la libre circulación de ciudadanos, la cooperación judicial civil, la política industrial, la aprobación de los fondos estructurales. etc. Sin embargo, en materias sensibles como fiscalidad. política social, asilo e inmigración se mantuvo la unanimidad.

Los logros alcanzados en el Tratado de Niza fueron en cierto modo decepcionantes ya que no procedió a la profunda reforma a la que había aspirado. El Consejo Europeo celebrado en diciembre de 2001 en Laeken convocó una Convención presidida por Valéry Giscard d'Estaing para llevar a cabo un debate más amplio y profundo sobre el futuro de la UE, incluyendo l a tarea de elaborar un proyecto de Constitución para Europa. La Convención terminó esta tarea en 2004 y los Jefes de Estado y de Gobierno firmaron el proyecto de Constitución en Roma. No obstante, éste fue descartado tras el resultado negativo de los referéndums celebrados en Francia y Países Bajos en 2 005.

A pesar del fiasco del proyecto de Constitución. los ajustes del Tratado de Niza fueron suficientes para iniciar la andadura de una UE ampliada. Después de unas intensas negociaciones, la quinta ampliación tuvo lugar en dos etapas. El I de mayo de 2004 se incorporaron a la UE diez nuevospaíses: dos pequeños mediterráneos ( M alta y Chipre) y ocho de Europa Central y Oriental (Estonia, Letonia. Lituania, Polonia, la República Checa, Eslovaquia, Hungría y Eslovenia). Bulgaria y Rumanía lo hicieron después, el I de enero de 2 007, debido a retrasos en las reformas de sus sistemas judiciales y lucha contra la corrupción. Por otra parte, Croacia solicitó el ingreso a la UE en 2003, produciéndose su adhesión el I de julio de 2 0 1 3. Tras esta sexta ampliación, la UE llegó a contar con veintiocho Estados miembros.

\subsection{Base actual de la Unión Europea} 
Tras un periodo de reflexión, gran parte del contenido del proyecto de Constitución se incorporó al Tratado de L isboa, pero eliminando todos los rasgos, símbolos y términos de ésta que podían hacer recelar que la Unión se asimilara a una estructura cuasi-estatal. El Tratado se firmó el 13 de diciembre de 2007 en Lisboa y entró en vigor el I de diciembre de 2009, tras ser ratificado por todos los Estados miembros (en un proceso no exento de dificultades; rechazo en referéndum en Irlanda, problemas legales en Alemania y República Checa). 

El Tratado de Lisboa supone la última modificación de los Tratados y ha introducido importantes novedades: 
Ha otorgado personalidad jurídica a la Unión, lo que refuerza su papel en el plano internacional al permitirle ser parte contratante en los tratados internacionales. Esto también permite que se suprima la estructura en pilares y el funcionamiento de la UE pase a basarse en el principio de atribución de competencias o cesión de soberanía.\footnote{%
    El funcionamiento por cooperación intergubemamental solo se mantiene en materia de defensa dando lugar a la Política Común de Seguridad y Defensa ( PCSD), que es una política específica dentro de la Política Exterior y de Seguridad Común (PESC).
} Es decir, la Unión actúa dentro de los límites de las competencias que le  atribuyen los Estados miembros en los Tratados, y toda competencia no atribuida a laUnión en los Tratados corresponde a los Estados miembros.\footnote{%
    El Tratado de Lisboa mantiene la \textit{cláusula de cierre} denominada ahora '•cláusula de tlexibilidad", pero impone dos limitaciones: (i) para permitir la acción de la Unión se requerirá la previa autorización del Parlamento Europeo (antes bastaba la simple consulta); y (ii) nunca se podrá utilizar para adoptar medidas de annonización cuando así se prohíba en otros artículos de los Tratados.
}

Se distingue entre tres ámbitos de competencia: (i) exclusiva, cuando solo la Unión es capaz de legislar y adoptar actos vinculantes;\footnote{%
    Ámbitos de competencia exclusiva (art. 3 TUE): unión aduanera: nonnas sobre competencia necesarias para el funcionamiento del mercado interior: política monetaria de los países de la zona euro: la conservación de los recursos biológicos marinos: política comercial común; celebración de acuerdos internacionales bajo ciertas condiciones.
} (ii) compartida, cuando los Estados miembros pueden adoptar actos jurídicamente vinculantes en la medida en que la Unión no haya ejercido su competencia;\footnote{%
    Ámbitos de competencia compartida (art. 4 TUE): mercado interior: los aspectos de política social incluidos en el Tratado: cohesión económica, social y territorial: agricultura y pesca: medio ambiente: protección de los consumidores: transporte: redes trans-europeas: energía: espacio de libertad. seguridad y justicia: problemas de sobre la seguridad compartidos en materia de salud pública: investigación. desarrollo tecnológico y espacio: cooperación al desarrollo y ayuda humanitaria.
} y (iii) de apoyo, cuando la Unión solo puede intervenir para apoyar, coordinar o complementar la acción de los Estados.\footnote{%
    Ámbitos de competencia de apoyo (art. 6 TUE): protección y mejora de la salud humana: industria: cultura: turismo: educación, fomiación profesional. juventud y el deporte: protección civil: cooperación administrativa.
} 

También para aclarar el reparto y el ejercicio de las competencias, el Tratado de Lisboa establece otros principios fundamentales: 
• Subsidiariedad: La Unión solo intervendrá en los ámbitos que le son propios "en la medida en que los objetivos de la acción pretendida no puedan ser alcanzados de manera suficiente por los Estados miembros, ni a nivel central ni a nivel regional y local.". Este principio había sido introducido en el Tratado de Maastricht, pero ahora se refuerza estableciéndose la obligación de informar a los Parlamentos Nacionales de cada propuesta legislativa europea.
• Proporcionalidad: "El contenido y la forma de acción de la Unión no excederán lo necesario para alcanzar los objetivos de los Tratados." Para facilitar los avances en la integración, el Tratado de Lisboa ha generalizado el procedimiento de codecisión (que pasa a ser el •·procedimiento legislativo ordinario'') y la votación por mayoría cualificada.\footnote{%
    El Tratado de Lisboa establece un nuevo sistema de doble mayoría condicionada. según el cual se alcanza la mayoría cualificada si votan a favor: (i) al menos un 55\% de los Estados miembros; y, (ii) dichos Estados miembros representan como mínimo el 65\% de la población total de la Unión. Para impedir que un grupo pequeño de países con un gran número de habitantes obstruya la adopción de decisiones, está previsto que una minoría de bloqueo debe estar formada. como mínimo. por cuatro Estados miembros que representen más de un 35\% de la población. Se sigue exigiendo launanimidad para ciertos asuntos políticamente sensibles.
} No obstante, teniendo en cuenta que existen ámbitos especialmente sensibles en los que unos pocos países pueden frenar la integración de la mayoría que sí lo desea, el Tratado de Lisboa ha previsto algunos mecanismos de flexibil idad (que, de hecho, proporcionan la base para llevar a cabo una ''integración a varias velocidades"). 

En este sentido:
• Ha ordenado las disposiciones relativas a la cooperación reforzada, la cual ofrece la posibilidad de que un grupo de Estados miembros avance en la integración en un determinado ámbito sobre el que la UE tenga competencia, sin que vean obstaculizado su avance por las dudas o negativas de otros Estados miembros.\footnote{%
    Las condiciones y los procedimientos para la utilización de este instrumento que inicialmente habían sido muy estrictas (Tratado de Ámsterdam). se flexibilizaron en cierta medida con miras a la an1pliación de la UE (Tratado de Niza). Este procedimiento se ha utilizado en los ámbitos de la legislación en materia de divorcio. las patentes y el impuesto sobre las transacciones financieras. así como para proteger los intereses financieros de la UE mediante la creación de la Fiscalía Europea.
} La cooperación reforzada debe estar abierta a todos los Estados miembros e iniciarse solo como último recurso. El umbral mínimo para emprenderla es de nueve Estados miembros. Los actos adoptados en el marco de una cooperación reforzada no forman parte del acervo de la UE.
• Ha introducido unas "cláusulas pasarela", que evitan tener que modificar los Tratados para: (i) sustituir la votación por unanimidad por reglas de votación por mayoría, o (ii) pasar de un procedimiento legislativo especial al procedimiento legislativo ordinario para adoptar un acto dt:lt:rminado. La activación de una cláusula pasarela sietl)pre depende de  una decisión adoptada por unanimidad del Consejo o del Consejo Europeo.

Por otro lado, el Tratado de Lisboa facilita la revisión de los Tratados al prever un procedimiento de revisión simplificado, que evita la necesidad de convocar una Convención Europea y de una Conferencia l ntergubemamental. Este procedimiento se puede aplicar para la revisión de las políticas y acciones internas de la UE, aunque no puede aumentar las competencias de la Unión. Debe ser aprobado por unanimidad por el Consejo Europeo y las modificaciones de los Tratados solo entran en vigor si han sido ratificadas por el conjunto de los países de la UE.

El Tratado de Lisboa revisa asimismo los tipos de actos jurídicos '.de la UE y la lista de instituciones europeas, cuestiones que repasaremos a continuación en su configuración actual.


\clearpage
\section{Unión Europea: Ordenamiento Jurídico}
\puntosclave{

}


Mediante el Tratado de la Unión Europea, adoptado en Maastricht en 1992, los Estados miembros de las Comunidades Europeas constituyeron entre sí una Unión Europea (art. 1 del Tratado de la UE).

El proyecto de Unión Europea se remonta casi a los orígenes mismos del proceso de la integración; por ejemplo, la propuesta del ministro francés, Aristide Briand, ante la Sociedad de Naciones en 1929. Pero, sin perjuicio de esos orígenes, algo lejanos y conceptualmente distintos, el objetivo de transformar cualitativamente las relaciones emprendidas en 1951 y 1957 se puede situar en los años setenta del pasado siglo, siendo el «Informe Tindemans» sobre la Unión Europea de 1975 el documento que preveía una evolución progresiva consolidando y desarrollando las Comunidades Europeas y ampliándolas a nuevas políticas.

De la Declaración Solemne sobre la Unión Europea, adoptada por el Consejo Europeo de 1983 en Stuttgart, se deducía que elconcepto es una idea de síntesis: mediante la creación de la UE los Estados miembros estaban decididos a alcanzar «una concepción común, global y coherente» del proceso.

El núcleo de la Unión Europea en esos años se situaba exclusivamente en las Comunidades Europeas; se estimaba que se progresaría en la vía de la Unión en la medida en que se desarrollasen y se profundizase en las políticas existentes así como se emprendiesen nuevas políticas de cooperación en el marco de los Tratados de París y de Roma. La Unión Europea en 1983 era
todavía exclusivamente comunitaria.

El Acta Única Europea de 1986 (en vigor desde el 1 de julio de 1987) fue el primer texto de derecho originario que recogió el objetivo de la Unión Europea y lo hizo en el primer párrafo de su Preámbulo: «animados por la voluntad de proseguir la obra emprendida a partir de los Tratados constitutivos de las Comunidades Europeas y de transformar el conjunto de las relaciones entre sus Estados en una Unión Europea». Sin embargo, no se adoptó la decisión de constituir la Unión Europea si bien se acordó introducir, junto al pilar comunitario, aunque completamente separado de éste, el pilar intergubernamental de la Cooperación Política Europea (CPE). La Unión se perfilaba como un híbrido político.

Por fin, la decisión de constitución cobró realidad en el primer artículo del Tratado de la Unión Europea de 7 de febrero de 1992 firmado en Maastricht (en vigor desde el 1 de noviembre de 1993). La Unión Europea, tal como se ideó en el inicial Tratado de Maastricht de 1992, era un ente político, era el todo basada en tres partes o pilares sobre los que se sustentaba el proceso en su conjunto: las Comunidades Europeas y las dos formas de cooperación intergubernamental —la Política Exterior y de Seguridad Común PESC y la cooperación judicial penal y policial—. Podríamos decir que la Unión Europea era un ente de la razón con fundamento in re (en las Comunidades). Era un ideal político, que tiene una entidad política pero cuya realidad —política, jurídica, económica, social— se sustentaba en las realizaciones materiales conseguidas a través de las Comunidades Europeas.

Los tres pilares conducían a los objetivos últimos de la Unión y participaban parcialmente del sistema institucional; pero se distinguían por utilizar instrumentos de integración diferentes y por el método comunitario (basado, en breves palabras, en la iniciativa exclusiva de la Comisión y la aprobación por mayoría cualificada en el Consejo y por mayoría en el PE).

Pero esta idea de una Unión Europea basada en las tres Comunidades (CECA-CEE-Euratom) como pilar central comunitario, junto a los pilares o formas de cooperación intergubernamentales, resultó muy compleja y difícil de entender por la ciudadanía y por los propios políticos que daban por desaparecidas aquéllas. Tras la reforma introducida por el Tratado de Lisboa de 2007 al Tratado de la Unión Europea desaparece formalmente la estructura de pilares; se extingue la antigua Comunidad Europea (pero no su tratado que pasa a denominarse Tratado de Funcionamiento de la UE) y hace que la renovada Unión Europea le suceda de forma unificada.

La Unión Europea no partía ex novo en 1992 ni es un proceso nuevo de integración tras la reforma hecha en Lisboa en 2007; es una organización internacional que asume todo el acervo de integración y sucede a la Comunidad Europea (como se explicita tras la reforma del Tratado de Lisboa). Continúa el proceso iniciado en 1951 y 1957. Por ello, el artículo 1 del Tratado de la Unión Europea (TUE), tal como está en vigor, declara que la Unión «se fundamenta» en el Tratado de la UE y en el Tratado de Funcionamiento de la Unión Europea (TFUE), continuadores del proceso. En efecto, la realidad —política, jurídica, económica, social — de la UE se sustenta en las realizaciones conseguidas mediante las «antiguas» Comunidades Europeas en más de sesenta años de integración económica, social y política por medio del Derecho y en los avances a través de los dos antiguos pilares intergubernamentales de la PESC y los asuntos de interior y justicia. 

Los Estados (y, por medio de ellos, la ciudadanía) crearon la Unión Europea sobre el edificio jurídico-político de las Comunidades Europeas. La Unión no es la meta misma de la integración sino «una nueva etapa en el proceso creador de una Unión» (primer párrafo del Preámbulo y art. 1 TUE).

La Unión mantiene la idea y objetivo tradicional de una unión «cada vez más estrecha entre los pueblos de Europa», que avanzalenta y continuamente hacia una mayor integración sin definir la meta final.

El compromiso de una «unión cada vez más estrecha» apareció originalmente en el preámbulo el Tratado de Roma, 1957. En el Tratado CECA de 1951 no se utilizó, pero sí la idea de «más amplia y profunda». Desde la entrada en vigor de la reforma de Maastricht en 1993, este enunciado se recoge en el artículo 1 del Tratado de la UE. Fue uno de los puntos de enfrentamiento entre Reino Unido y la Unión en 2016; aunque se aceptó exonerar al Reino Unido de cooperar en la ambición de «más» Europa no fue suficiente para su permanencia en la Unión. 

Esa expresión responde a una concepción filosófica y jurídica sobre la idea de proceso, de construcción intergeneracional, deevolución, de adaptación a los tiempos y necesidades. Y no supone ningún automatismo ni una ampliación de competencias que siempre exigiría una reforma de los tratados bajo control parlamentario nacional.

Responde a la idea de progresividad; a la idea de que el futuro se construye poco a poco y entre todos, que no hay nada acabado ni cerrado y que las generaciones futuras están llamadas a ser protagonistas del proceso, condensando toda una filosofía política.

La Unión es una organización internacional sui generis que se constituye por Estados democráticos y su ciudadanía, de los que recibe las competencias para alcanzar los objetivos comunes que aquéllos quieren lograr. Se enuncia así el principio básico y clásico de todos los Tratados constitutivos de organizaciones internacionales: el principio de atribución de competencias (art. 5 TUE). Son los Estados, y por medio de ellos los ciudadanos, quienes han dotado de competencias a la Unión. La Unión Europea no ha sido ni es un poder originario creado a partir de un acto constituyente popular; no es un Estado ni tiene como horizonte serun Estado federal unificado.

La UE no es un Estado federal, aunque sólo sea por el hecho de que ni tan siquiera es un Estado. Es una organización internacional intergubernamental, ciertamente muy distinta a las Organizaciones clásicas; es una organización internacional original, única. La UE es un ente jurídico-internacional y político atípico; esta asociación de Estados soberanos no tiene en el horizonte sustituir a los Estados soberanos.

La creación de la Unión Europea no cambia ni transforma la naturaleza política y jurídica del proceso iniciado en 1950; sigue siendo una asociación voluntaria de Estados soberanos a la que se dota de competencias concretas y limitadas que puede ejercer en las condiciones establecidas en los tratados internacionales que la regulan. No podía ser de otra manera entre Estados que siguen siendo soberanos e independientes en el orden internacional. La voluntad de los Estados miembros ha nutrido siempre un proceso de integración que no comenzó con el fracasado Tratado constitucional. Esa voluntad de los Estados es una condición sine qua non del sistema. Sin un pacto entre los Estados no se da vida a normas originarias o constitutivas. Nuestros Estados, nuestros gobiernos democráticos, son el poder constituyente intergeneracional constante.

El Tratado de la Unión deja claro que los Estados miembros, soberanos e independientes, permanecen como tales y que esta entidad política internacional hunde sus raíces y sus límites en el Derecho Internacional. Que seguimos en el marco de una Organización internacional y de un modelo de federalismo internacional y, por tanto, en el laxo marco del Derecho Internacional. La verdadera naturaleza de la Unión Europea y su Tratado sigue anclada en el Derecho de las Organizaciones internacionales. Como dijo el Consejo Constitucional francés, en su Decisión de 19 de noviembre de 2004 sobre la compatibilidad del fracasado Tratado constitucional con la Constitución francesa, ni tan siquiera el uso del término «Constitución» cambiaba la naturaleza de «organización internacional» de la UE.

En efecto, el TUE mantiene su carácter de organización internacional sin definir ni orientarse hacia contornos políticos federales, máxime teniendo en cuenta las continuas ampliaciones, el fracaso del Tratado constitucional de 2004, las tensiones de la globalización y la tendencia a un fuerte control intergubernamental.



\subsection{Ordenamiento Jurídico: Derecho de la Unión Europea}

Este capítulo y el siguiente tienen por objeto el estudio de las normas y los actos que constituyen los instrumentos de regulación jurídica del sistema de la Unión Europea. La singularidad, permanente evolución y complejidad del mismo aconsejan algunas consideraciones previas a su tratamiento.

El singular sistema jurídico de la Unión Europea, fundamentado en los Tratados constitutivos, TUE y TFUE, es el complejo resultado de una larga evolución de los tratados originarios de los años cincuenta. El concepto de sistema jurídico resultante como tal, es decir, del Derecho de la Unión Europea es una construcción de la jurisprudencia del TJUE y de una inestimable aportación doctrinal. La precariedad de estructura y organización del modelo normativo europeo, donde ni siquiera existe una jerarquía normativa, no ha facilitado las cosas. Ha sido en el seno de las CCEE donde se forjaron los elementos centrales de este sistema jurídico al que, para mayor complejidad se le adicionaron después fórmulas de organización distinta —Política Exterior y de Seguridad Común (PESC) y la Cooperación en asuntos de Justicia y de Interior— donde, con un fundamento y estructuras diferentes, se estableció un dispositivo de producción de actos y de formación de obligaciones cuyo alcance siempre resultó difícil de identificar, pero que tienen un profundo significado para el modelo de regulación de la Unión Europea. El Tratado de Lisboa intentó reordenar los elementos esenciales del sistema aspirando a someterlo a una sola lógica jurídica. El resultado es algo deficitario en la medida en que, aunque, aparentemente al menos, solo pervive diferenciado como procedimiento la PESC (art. 40 TUE), lo cierto es que la pervivencia temporal de algunas competencias del ELSJ y la acusada diferenciación de obligaciones para algunos Estados miembros, dificultan una consideración unitaria del sistema jurídico de la Unión. Pese a todo, y asumiendo los riesgos de una aproximación que tenga en cuenta todas estas circunstancias, estos temas pretenden situarnos desde el punto de vista de la configuración normativa del sistema en esa perspectiva unitaria y en una visión de conjunto.




2.1. NORMAS Y ORDENAMIENTO JURÍDICO
Hace tiempo que la teoría general abandonó la pretensión de entender el Derecho como un conjunto de normas individualizadas y exigió su comprensión desde la noción de «ordenamiento jurídico». Ello permite, como dijo Bobbio, abarcar la consideración del modo con el cual una determinada norma es eficaz por una compleja organización que determina la naturaleza y entidad de las sanciones, las personas que deben aplicarlas y su ejecución [1995: 155]. 

Hace tiempo que la quiebra de la simbolización del Derecho en el Estado permite plantearse la existencia de ordenamientos jurídicos no estatales (externos al Estado como el Derecho Internacional o internos como los sistemas federales) y las complejas relaciones entre ellos en una comprensión distinta del fenómeno jurídico. Y hace tiempo que el TJ, en una de las más importantes sentencias de su historia, la sentencia Van Gend en Loos de 5 de febrero de 1963, afirmó esta cualidad del sistema comunitario cuando advirtió que en los Tratados constitutivos de las CCEE había algo más que un mero agregado de normas:
— que, este Tratado (CEE) constituye algo más que un acuerdo que se limitara a crear obligaciones mutuas entre los Estados contratantes;
— que, esta concepción se encuentra confirmada por el preámbulo del Tratado que más allá de los Gobiernos, contempla a los pueblos y de manera más concreta, por la creación de órganos que institucionalizan derechos soberanos cuyo ejercicio afecta tanto a los Estados miembros como a sus nacionales;
— que, por otro lado hay que señalar que los nacionales de los Estados reunidos en la Comunidad están llamados a colaborar por medio del Parlamento Europeo y del Comité Económico y Social al funcionamiento de esta Comunidad;
— que, además, el papel del Tribunal de Justicia [...], cuya finalidad es asegurar la unidad de interpretación del Tratado por los órganos jurisdiccionales nacionales confirma que los Estados han reconocido al Derecho comunitario una autoridad susceptible de ser invocada por sus nacionales ante sus jurisdicciones;
— que de esta situación hay que concluir que la Comunidad constituye un nuevo ordenamiento jurídico de Derecho Internacional [cursiva añadida]. 

Poco después, en la no menos célebre sentencia Costa c. ENEL de 15 de julio de 1964, el Tribunal lo calificó como un ordenamiento jurídico propio, excluyendo la mención al Derecho Internacional. En ambos casos el TJ entendió el Derecho de las entonces CCEE, no como un mero agregado de normas, sino como un ordenamientojurídico. Todo ello reconduce la comprensión de las normas y actos jurídicos de la UE a un tratamiento de conjunto referido a un sistema normativo, específico y particular, el Derecho de la Unión, que constituye un verdadero ordenamiento y cuya esencia jurídica emerge de su profunda y compleja interrelación con los ordenamientos jurídicos internos.

2.2. ESPECIFICIDAD DEL SISTEMA Y TEORÍA DE LAS FUENTES 
Como es sabido, el concepto de «fuentes del derecho» es en sí mismo difícil de precisar y ha originado en la ciencia jurídica serios problemas derivados de la polisemia del término, de que alberga concepciones jurídicas distintas y, en ocasiones, contradictorias y de que su extensión a ordenamientos no estatales desarraiga un concepto surgido, en cualquier caso, en el seno del Estado. Acoger la teoría de las fuentes sin, al menos, advertir que se está manejando un concepto muy complejo y para algunos, incluso, uno de los más confusos y perturbadores de la ciencia jurídica (Santamaría, 1988: 286), puede resultar poco conveniente. Y si, desde luego, puede serlo en el Estado constitucional, con mucho más motivo resulta inconveniente en el ámbito de otros ordenamientos jurídicos como el internacional o el de la Unión. Sus insuficiencias han quedado claras en el Derecho Internacional público, donde el monopolio del consentimiento del Estado en la  formación de normas elimina toda utilidad del concepto de fuentes en un sentido material y donde, además, la carencia de una normaestructuradora del sistema de producción normativa y de jerarquización entre los diversos tipos reduce a mínimos su utilidad formal.

Desde luego, siempre es posible mantener el término «fuentes» para referirse a la existencia de un conjunto de «categorías normativas» formalmente diferenciadas. Pero si la carencia de una norma sobre producción normativa en un ordenamiento concreto nos priva de canalizar a través de la forma la información acerca de los poderes que las producen, traduciendo así una jerarquización o determinación de zonas de competencia y su sustanciación dentro del complejo modelo de relación de un ordenamiento jurídico, es posible que no compense la confusión que puede seguirse de una utilización del concepto de «sistema de fuentes». El Derecho de la Unión Europea es, ciertamente, un caso de excepcional complejidad en este sentido. Como ya se ha estudiado, se trata de un orden jurídico que partió de un modelo relativamente sencillo diseñado en normas convencionales internacionales que proporcionaban «datos relativamente poco consistentes» (Kovar, 1990: 4) desde el punto de vista de la organización del sistema de producción de normas. En realidad, aparte de las propias normas constitutivas (Tratados internacionales) y una tipología de actos sin una relación de jerarquía y diferenciados esencialmente en función de sus efectos, su alcance respecto de la mediatización estatal y la identificación de sus destinatarios, no había nada más. El carácter marcadamente evolutivo del sistema y, de manera muy especial, la relevante labor del TJ, apoyada en un importante cuadro contencioso previsto en los Tratados, han permitido, en cierto modo, superar sus deficiencias, elaborando un marco muy complejo y peculiar que, si por las razones antedichas no puede asimilarse en sentido estricto a un «sistema de fuentes», bien merece calificarse  como un verdadero sistema de normas y actos jurídicos.

2.3. NORMAS ORIGINARIAS Y NORMAS DERIVADAS
El sistema de normas de la Unión Europea se apoya en una summa divisio entre las normas originarias y derivadas que estaba presente desde el origen de la construcción europea. La norma originaria se identifica sustancialmente con los Tratados constitutivos y las normas convencionales que los han modificado a lo largo del tiempo, y cuyo último exponente general es el Tratado de Lisboa que modificó el TUE y el TCE (denominado ahora TFUE). Sus características básicas pueden agruparse en torno a tres ejes fundamentales: a) Tanto material como formalmente se trata de normas jurídicointernacionales y, en consecuencia, regidas por las normas de Derecho Internacional general aplicables a los tratados internacionales. b) Como elemento de normatividad esencial, revisten una dimensión «constitucional» que se manifiesta en sus contenidos (establecen los principios, determinan los poderes atribuidos y sus límites, estructuran el sistema institucional y la distribución de sus poderes y funciones así como los procedimientos para su ejercicio y control) y en la garantía de su preeminencia sobre cualquier otra norma.
c) Como elemento de regulación, las normas originarias de la Unión Europea contienen regulaciones materiales específicas (así, las relativas a las libertades comunitarias, a las reglas de competencia o a las políticas comunes). Se trata de un núcleo de normas que podrían ser objeto de regulación por normas no originarias, pero los Estados decidieron su inclusión en esta categoría por la importancia que les otorgan a sus contenidos que, de este modo, se ven protegidos por la supremacía de la norma primaria y por asegurar la necesidad de su consentimiento para una eventual modificación. Aunque este fenómeno no es ajeno a las Constituciones modernas, que suelen regular ámbitos sociales importantes de forma directa, lo cierto es que en la Unión Europea estas normas son abrumadoramente mayoritarias y en muchos casos se encuentra poco justificada su inclusión en la categoría de las normas originarias. A pesar de todo, el nuevo TFUE reproduce la mayor parte de las disposiciones materiales contenidas en los anteriores Tratados.

Las normas derivadas en el sistema europeo se cualifican esencialmente por tratarse de un conjunto de modos de instrumentación jurídica con fundamento en la norma constitutiva. El grueso de estas normas proviene del sistema de atribución de competencias que, constituyendo la esencia misma del modelo de la Unión, otorga al sistema institucional los poderes jurídicos necesarios para la consecución de los fines y objetivos establecidos en la norma originaria. Sería demasiado prolijo enumerar aquí todas las competencias previstas de una manera u otra por los Tratados y, en cierto modo, desviarían nuestra atención del objeto de este capítulo. Baste recordar que ese amplio abanico de competencias se encuentran definidas en el TUE, la de PESC, y el resto, las catalogadas en los artículos 3 a 6 TFUE, a lo largo de los Tratados. Son competencias estas últimas que se instrumentan básicamente a través de la tipología prevista en el artículo 288 TFUE (reglamentos, directivas, decisiones, recomendaciones y dictámenes). Prevista su forma y caracterizados los actos jurídicamente en dicho artículo, la doctrina ha dado en denominarlos «actos típicos» para diferenciarlos del conjunto de instrumentos jurídicos a los que las Instituciones han recurrido en la práctica por razones diversas (imprevisión de la forma o debilidad de la base competencial, entre otras) y que no entran en estas categorías.

Junto al Derecho derivado institucional se encuentran, con un fundamento más o menos preciso en la norma originaria, otros modos de instrumentación jurídica, como las normas sustanciadas a través del Derecho Internacional (como consecuencia de la sumisión al Derecho Internacional general o a través de la actividad convencional de la UE en virtud del poder de concluir acuerdos que, para determinados ámbitos, les otorgan los Tratados) o los acuerdos entre los Estados miembros como es el caso de las decisiones de los representantes de los Estados miembros reunidos en el seno del Consejo.

La PESC es una competencia especial y sigue manteniendo unas peculiaridades en la producción de actos. El Tratado de Lisboa ha suprimido en términos generales las particularidades de las normas y actos de la cooperación policial y judicial penal, aunque durante algún tiempo hasta que sean sustituidos mantendrán esa peculiaridad.




A diferencia del malogrado proyecto de Constitución, el Tratado de Lisboa no sustituye los Tratados constitutivos de la Unión por un único texto. sino que supone la última modificación de los dos textos fundamentales:

\subsubsection{Normativa originaria: Tratados}
3.1. LA NORMA CONVENCIONAL CONSTITUTIVA
Como hemos visto con anterioridad, tanto material como formalmente, la norma que da origen al sistema jurídico de la Unión es de carácter netamente jurídico-internacional. En este sentido, hay que entender que la norma originaria del sistema se fundamenta, desde el punto de vista material, en el consentimiento estatal, formado a través de los procedimientos constitucionales de cada Estado miembro y encauzado a través de la manifestación del consentimiento prevista en el sistema jurídico internacional para los Tratados celebrados en forma solemne, esto es, a través de la ratificación. Desde el punto de vista formal se trata de normasconvencionales internacionales (Tratados con sus Protocolos y Anexos) sujetos a las reglas de Derecho Internacional en materia de Tratados, es decir, al Convenio de Viena sobre el Derecho de los Tratados de 23 de mayo de 1969. En el Dictamen 1/91 sobre el EEE1, el TJCE ha aludido de manera expresa al Convenio de Viena y de manera particular a sus reglas de interpretación. Uno de los aspectos básicos de su sujeción a estas reglas viene expresado por el artículo 39 del Convenio de Viena que, salvo que el Tratado disponga otra cosa, exige de la celebración de un nuevo Tratado para la enmienda del anterior, sujeto a un nuevo procedimiento decelebración y entrada en vigor. Tal es la última ratio del procedimiento ordinario de revisión del artículo 48 TUE. La consecuencia fundamental, desde el punto de vista que aquí nos interesa, es que cualquier modificación de las normas originarias, sea cual fuere su naturaleza o intensidad, exige, salvo en algunos casos muy menores y con cautelas, la celebración de un Tratado internacional.

Estas consideraciones tienen un doble efecto en relación con la identificación del Derecho originario: que toda modificación de ese Derecho ha de recogerse en una norma convencional y que toda norma convencional, concluida por los Estados miembros, con efectos modificativos de los Tratados constitutivos, forma parte del Derecho originario. Ello explica que el denominado Derecho originario albergue no sólo los Tratados constitutivos sino, además, todos los Tratados que de una manera u otra han modificado sus disposiciones. Al tratarse de modificaciones de los tratados, permanecen vigentes las normas originarias en tanto no hayan sido derogadas, modificadas y se deduce que, en tanto no sean incompatibles con el TUE o TFUE (o la propia CEEA). Se han configurado, así, en la evolución del proceso de construcción europea, dos vectores de creación del Derecho originario:
a) Los Tratados constitutivos y los Tratados concluidos para su modificación:
— Tratado constitutivo de la Comunidad Económica Europea, adoptado en Roma el 25 de marzo de 1957 (entrado en vigor el 1 de enero de 1958), denominado de la Comunidad Europea en Maastricht, y Tratado de Funcionamiento de la Unión Europea desde la entrada en vigor del Tratado de Lisboa el 1 de diciembre de 2009.
— Tratado constitutivo de la Comunidad Europea de la Energía Atómica, adoptado en Roma el 25 de marzo de 1957 (entrado en vigor el 1 de enero de 1958).
Estos Tratados fueron modificados en sus aspectos institucionales en el marco de un proceso de unificación institucional de gran relevancia por tratados de los años 1957 (Asamblea y TJCE) y 1965 (Comisión y Consejo) que fueron derogados por el Tratado de Ámsterdam (art. 9), aunque no los efectos jurídicos de los actos vigentes adoptados en virtud de aquellos tratados (art. 11 del Tratado de Ámsterdam), permaneciendo en vigor sólo el Acta relativa a la elección de los representantes en el Parlamento Europeo por sufragio universal directo firmada en Bruselas el 20 de septiembre de 1976.
Asimismo, algunas modificaciones parciales son llevadas a cabo por:
— Los dos Tratados, por los que se modifican determinadas disposiciones presupuestarias y financieras de los Tratados constitutivos de las Comunidades Europeas y del Tratado por el que se constituye un Consejo único y una Comisión única de las Comunidades Europeas, el primero de 22 de abril de 1970 hecho en Luxemburgo y el segundo de 22 de julio de 1975 en Bruselas.
— Las Decisiones relativas a los recursos propios. Actualmente en vigor se encuentra la Decisión 2014/335/UE, Euratom del Consejo, de 26 de mayo de 2014, relativa al sistema de recursos propios2.
— El Tratado por el que se modifican determinadas disposiciones del Protocolo sobre los Estatutos del Banco Europeo de Inversiones de 10 de julio de 1975 hecho en Bruselas. No obstante, las modificaciones generales por excelencia son las llevadas a cabo a través del procedimiento de revisión contemplado en los Tratados constitutivos y que, tras los trabajos de las correspondientes Conferencias Intergubernamentales, desembocaron en:
— El Acta Única Europea, que, tras su adopción en Luxemburgo el 27 de enero de 1986, fue firmada en Luxemburgo el 17 de febrero de 1986 por nueve Estados miembros y el 28 de febrero de 1986 en La Haya por los tres Estados restantes (entrada en vigor el 1 de julio de 1987).
— El Tratado de la Unión Europea, firmado en Maastricht el 7 de febrero de 1992 (entrado en vigor el 1 de noviembre de 1993).
— El Tratado de Ámsterdam, por el que se modifica el Tratado de la Unión Europea, los Tratados constitutivos de las Comunidades Europeas y determinados actos conexos de 2 de octubre de 1997 (entrado en vigor el 1 de mayo de 1999).
— El Tratado de Niza, por el que se modifican el Tratado de la Unión Europea, los Tratados constitutivos de las Comunidades Europeas y determinados actos conexos de 26 de febrero de 2001 (entrado en vigor el 1 de febrero de 2003). 
— El Tratado de Lisboa, por el que se modifican el Tratado de la Unión Europea y el Tratado constitutivo de la Comunidad Europea, firmado en Lisboa el 13 de diciembre de 2007 (entrado en vigor el 1de diciembre de 2009).
— Protocolo por el que se modifica el Protocolo sobre las disposiciones transitorias, anejo al Tratado de la Unión Europea, al Tratado de Funcionamiento de la Unión Europea y al Tratado constitutivo de la Comunidad Europea de la Energía Atómica, adoptado en Bruselas el 23 de junio de 2010. Entró en vigor el 1 de diciembre de 2011. Es fruto de la Decisión del Consejo Europeo de 17 de junio de 2010 relativa al examen, por una Conferencia de representantes de los Gobiernos de los Estados miembros, de las modificaciones de los Tratados propuestas por el Gobierno español en lo que se refiere a la composición del Parlamento Europeo, sin convocar una Convención (2010/350/UE) que adopta una modificación en virtud del procedimiento del artículo 48, apartado 3, TUE.
— Decisión del Consejo Europeo de 25 de marzo de 2011(2011/199/UE) que modifica el artículo 136 del Tratado de Funcionamiento de la Unión Europea en relación con un mecanismo de estabilidad para los Estados miembros cuya moneda es el euro  (modificación adoptada en virtud del procedimiento del art. 48, apartado 6, TUE), en vigor desde 1 de enero de 20133.

b) Un segundo núcleo de modificaciones del Derecho originario viene dado por los actos de naturaleza convencional que se producen con motivo de las adhesiones de nuevos Estados miembros o para poner término a su aplicación. El primero ha cobrado una relevancia excepcional como consecuencia de las seis ampliaciones de que han sido objeto las Comunidades Europeas (las tres primeras) y la Unión Europea (las tres últimas). El segundo caso, la terminación del Tratado, sólo ha tenido lugar de manera excepcional en relación con Groenlandia, pero en cualquier caso hubo de ser canalizado a través del procedimiento de revisión antes aludido.

c) El caso especial de la Carta de Derechos Fundamentales de la Unión Europea. En sí no tiene carácter convencional, pero a tenor del artículo 6.1 TUE tiene «el mismo valor jurídico de los Tratados»,  razón por la que los principios, derechos y libertades enunciados en la Carta y en las condiciones previstas en el TUE, deben considerarse derecho originario. La versión en vigor (2010/C 83/02) es la publicada en el DOUE C 83, de 30 de marzo de 2010.

3.2. CARACTERES DE LA NORMA CONSTITUTIVA
a) La primera característica de las normas constitutivas que debe resaltarse es la pluralidad y diversificación de obligaciones. El objetivo de unificación de la norma primaria por tantos defendido como un elemento de racionalización del sistema ha conseguido importantes avances en el Tratado de Lisboa. En cualquier caso, es preciso inmediatamente señalar que tales avances no han roto esacaracterística de diversificación de las normas originarias en el Derecho actualmente en vigor. En primer lugar, TUE y TFUE siguen siendo dos Tratados formalmente independientes que gozan del mismo valor jurídico (arts. 1 TUE y TFUE). La Unión «sustituye y sucede a la Comunidad Europea», según dispone el artículo 1 TUE. A estos Tratados hay que añadir el Tratado constitutivo de la Comunidad Europea de la Energía Atómica5 por más que razones políticas lo hagan vagar semi-oculto por los vericuetos del Derecho de la Unión. En efecto, ni la formulación misma es unitaria en su totalidad, ni los elementos más expresivos de esa unidad han permanecido al abrigo de diversificaciones introducidas, con mayor o menor ortodoxia jurídica, por situaciones particulares de los Estados manifiestas en protocolos, declaraciones y otros actos de indescifrable naturaleza jurídica. Por si fuera poco, el mecanismo de la «Cooperación reforzada» y el de «Cooperación estructurada permanente» abundan en una diversificación de núcleos obligacionales difícilmente compatibles con la pretensión de unidad del sistema y desde luego con la de su simplificación. Los muchos intentos y propuestas de «simplificación» han resultado un complicado viaje a ninguna parte. Esa pluralidad de normas originarias tiene importantes efectos en relación con el sistema normativo del orden jurídico de la Unión.

Junto a esta diversificación de las normas originarias de la Unión, hay que ahora los actos que se adoptan en la PESC. Aquí, la permanencia de un modelo menos integrado y la muy peculiar «atribución» de competencias, debilitan aún más aquellos factores de unificación. El artículo 40 TUE establece con claridad la diversidad y autonomía de los procedimientos previstos para la PESC (Título V del TUE) y las otras competencias de los artículos 3 a 6 TFUE. Adicionalmente, debe señalarse que el TJ puede incidir sobre la unificación del sistema habida cuenta de la más abierta redacción del artículo 24, que, pese a que excluye la competencia del Tribunal sobre el pilar PESC (Título V del TUE), la reconoce sobre el respeto del artículo 40 TUE y del control de legalidad de las medidas restrictivas que afecten a personas físicas6.

Pues bien, el efecto de notable complejidad y disfuncionalidad que se deriva de esta pluralidad y diversificación de las normas constitutivas de la Unión Europea se ve muy seriamente agravado, como advertíamos antes, por la introducción de ciertos grados de diversificación en las obligaciones. Destacan entre éstos los introducidos por el «Protocolo 15 sobre determinadas disposiciones relativas al Reino Unido de Gran Bretaña e Irlanda del Norte» —que, simplemente, excluye o modifica la aplicación a este Estado de varias e importantes disposiciones del Tratado, relativas a la Unión Económica y Monetaria— o por la aún menos ortodoxa fórmula de la Decisión de los Jefes de Estado o de Gobierno reunidos en el seno del Consejo Europeo de Edimburgo de 11 y 12 de diciembre de 1992, que introduce asimismo un núcleo sustancial de excepciones en varias materias (ciudadanía, UEM, ELSJ y PESC, entre otras) para Dinamarca, recogida en el Protocolo 22. Igualmente deben mencionarse el Protocolo 19 sobre el acervo de Schengen integrado en el marco de la Unión Europea, o el Protocolo de Polonia y Reino Unido excluyéndose de la aplicación de la Carta de Derechos Fundamentales de la UE, las incorporadas «interpretaciones» de Irlanda como condición para aceptar su ratificación a Lisboa o, las concesiones de última hora a la República Checa, respecto de las que ni siquiera se respetaron los procedimientos constitucionales de los restantes veintiséis Estados miembros ya que en sus ratificaciones éstos no han tenido conocimiento de ello. Y así hasta un sinfín de Protocolos, declaraciones y excepciones materiales o temporales, más o menos encubiertas, que, más allá de las insufribles complicaciones que generan, hacen simplemente ilusoria la pretensión de unidad del sistema jurídico de la UE.

Mención aparte merece la apertura de una especie de línea de desarrollo jurídico paralelo (incluso en este espacio normativo originario o constitutivo) que está suponiendo ciertos desarrollos de la Zona Euro. De especial importancia es retener aquí la formulación por los Jefes de Estado y de Gobierno de los EM de la Zona Euro de una especie de acuerdo que fue luego refrendado en una difícil reunión del Consejo Europeo de 24 y 25 de marzo de 2011 y finalmente recogidas en un Tratado internacional: El «Tratado constitutivo del Mecanismo Europeo de Estabilidad» (MEDE, firmado en Bruselas el 2 de febrero de 2012 y entrado en vigor el 27 de septiembre de 2012).

Mayores problemas se derivarán del Tratado de Estabilidad, Coordinación y Gobernanza (TECG) en la Unión Económica y Monetaria que los Jefes de Estado y Gobierno de todos los países miembros de la UE, a excepción del Reino Unido y la República Checa, firmaron el 2 de marzo de 2012. El Tratado de Estabilidad, Coordinación y Gobernanza, que entró en vigor el 1 de enero de 2013, queda abierto a la adhesión de otros países de la UE, además de las Partes contratantes. El objetivo es incorporar el contenido de este Tratado en el marco jurídico de la UE al cabo de cinco años de su entrada en vigor. Buena parte de sus normas deberánincorporarse a los ordenamientos jurídicos nacionales en disposiciones de rango constitucional como fórmula general y quedando sujeta a la competencia del Tribunal de Justicia de la Unión Europea la capacidad de imponer sanciones en caso de falta de trasposición. 

Tanto el MEDE como el TECG son normas convencionales internacionales. Su fundamento y efectos jurídicos derivan del Derecho internacional y son, en consecuencia, ajenos al sistema normativo de la UE, pero en ambos casos ciertos procedimientos, la implicación de instituciones de la UE y la competencia del TJUE advierten de que la lógica jurídica que los envuelva no permite considerarlos absolutamente extraños a la Unión Europea.

b) La supremacía de las normas constitutivas y la consecuente subordinación a éstas de las normas derivadas constituye el segundo elemento de caracterización general de las mismas. Esta supremacía se traduce en que es la norma constitutiva la que confiere al Derecho derivado su fundamento, alcance y límites. En sentido estricto, no existe en el ordenamiento jurídico de la Unión Europea una norma de estructuración general del sistema de producción de normas y, consecuentemente con ello, no existe una determinación expresa de la jerarquía normativa ni siquiera en larelación entre la norma constitutiva y la derivada. Pero, en estecaso, no por implícita resulta menos notoria su presencia en el ordenamiento jurídico de la Unión Europea. En efecto, esta supremacía puede inferirse de una norma de alcance general y, de manera particular, de los procedimientos previstos para controlar jurisdiccionalmente el respeto de la norma constitutiva, de laactividad de la Unión en la conclusión de acuerdos internacionales y del mecanismo mismo de revisión de las normas constitutivas. Para mayor complejidad las disposiciones del artículo 6.1 TUE reconocen a los principios, derechos y libertades fundamentales de la Carta de Derechos Fundamentales de la UE «el mismo valor jurídico que los Tratados», lo que los hace gozar de esta posición de supremacía. 

La norma de alcance general es la contenida en el artículo 5 TUE, más arriba trascrito, que al sujetar «el ejercicio de las competencias» de las Instituciones a las condiciones y fines previstos en las normas constitutivas, somete, naturalmente y de forma muy especial, a esas mismas condiciones y fines la competencia normativa. Junto a ella deben considerarse todas las disposiciones del artículo 2 TFUE.

En relación con el control jurisdiccional del TJ, prácticamente la totalidad de los recursos son directa o indirectamente aptos para asegurar esa jerarquía. Ello es lógico en tanto en cuanto corresponde al Tribunal garantizar el respeto del Derecho y, en consecuencia, el cumplimiento de las normas relativas a la sujeciónde la competencia normativa a las disposiciones constitutivas. Dichoesto, de manera más específica debe considerarse el control de legalidad a través de los recursos de nulidad, de omisión o el control de legalidad de determinados actos de la PESC, incluida la posibilidad de que se establezcan medidas cautelares a través de una solicitud de medidas provisionales7.

Particularmente expresivas, en relación con el lugar que ocupa en la jerarquía la norma constitutiva, son las disposiciones relativasa acuerdos internacionales celebrados por la Unión y, en particular, el sistema de control previo de la compatibilidad de dichos acuerdos con las normas constitutivas comunitarias previsto en el artículo 218.11 TFUE. El Tribunal de Justicia ha reiterado ya en varias ocasiones esa función constitucional suya. Así lo ha hecho en sus sucesivos dictámenes y, en concreto, en el Dictamen 3/94, de 13 de diciembre de 1995, en el que el Tribunal reitera su función de control «previo» a la prestación del consentimiento por la Unión, precisando que no por ello, una vez prestado el consentimiento, quedan las Instituciones o los Estados miembros desprotegidos frente a unaposible conclusión (material o formal) contraria a las normas constitutivas ya que siguen teniendo abiertas otras vías de recurso «control posterior» y, eventualmente, una solicitud de medidas provisionales8. Tal vez el pronunciamiento de mayor impacto en esta línea sea el Dictamen 2/13 de 18 de diciembre de 2014 en el que el TJUE ha declarado la incompatibilidad del Proyecto de Acuerdo de adhesión al CEDH con los Tratados, imposibilitando incluso el mandato del artículo 6.2 TUE que dispone dicha adhesión. 

Deben adicionarse a lo anterior las normas relativas a la preeminencia sobre Tratados anteriores celebrados entre Estados miembros, ya que se prevé la vigencia de estos últimos sólo en la medida en que sean compatibles con los Tratados (así, el art. 350 TFUE sobre el Tratado del BENELUX), o las normas sobre los tratados celebrados con anterioridad por los Estados miembros con terceros. En este caso, dentro del respeto a las obligaciones contraídas según las normas internacionales, los Estados están obligados a utilizar los procedimientos correspondientes para su adecuación a las exigencias del sistema de la Unión (art. 351 TFUE).

Por último, como hemos señalado, la preeminencia de las normas constitutivas se infiere de manera especial de su propio procedimiento ordinario y simplificado de revisión (salvo en algún caso muy excepcional), del artículo 48 TUE, que sigue manteniendo la necesidad para su materialización de las ratificaciones de «todos los Estados miembros» o hasta en los casos excepcionales la nooposición de un Parlamento nacional. Ninguna norma, pues, que no sea del mismo rango puede modificarlos. Nos ocuparemos de ello más adelante, pero conviene subrayar aquí, frente a quienes detectan en el propio TUE algunas posibilidades de modificación menos rígidas, que, si tales vías son posibles, es exclusivamente porque están previstas como tales en los propios Tratados por lo que en nada se atenúa la rigidez del sistema ni se debilita la posición de superioridad jerárquica de la norma constitutiva.

c) La eficacia directa es objeto de detenida atención en el capítulo referido a las relaciones entre Derecho de la Unión y Derecho interno. No obstante, quedaría sustancialmente incompleta la enumeración de los caracteres de la norma constitutiva si no se subrayara la cualidad de sus normas, en ciertas condiciones, de tener eficacia directa. La cuestión, desde esta perspectiva, hunde su raíz en el problema de los destinatarios de la norma. Una interpretación de su carácter convencional (por cierto, con poco fundamento en el actual Derecho Internacional) pudo, en principio, identificar como los destinatarios naturales de estas normas a los Estados y excluir a los particulares a los que, en todo caso, los efectos de las normas llegarían de manera indirecta a través de una especie de interpositio legislatoris de sus ordenamientos nacionales. Esta interpretación fue tempranamente corregida por el TJ en una importante línea jurisprudencial iniciada con la tantas veces aludida sentencia Van Gend en Loos de 5 de febrero de 1963.

A partir de esta sentencia, el TJ ha ido elaborando una construcción coherente y sin fisuras reconociendo a un gran número de normas de los Tratados la capacidad de producir efectos directos sobre los particulares, cuando naturalmente la norma tiene ese carácter incondicional y formula con claridad y precisión la obligación o el derecho. Abrió la brecha con aquella obligación de no hacer y reconoció los efectos entre el particular y el Estado, después incluyó las obligaciones positivas y el efecto entre particulares — como en el caso de la sentencia Defrenne (del hoy art. 157 TFUE)— 9. Así, el Tribunal ha ido desgranando un elevado número de disposiciones de los Tratados que gozan de estos efectos. Dos consecuencias parecen evidentes: la primera, que no siempre resulta tarea fácil determinar esa efectividad directa de una disposición, y la segunda, que esta dificultad sólo puede subsanarse acudiendo a la jurisprudencia del Tribunal ya establecida o, en caso de dificultad especial, acudiendo al propio Tribunal a través de los procedimientos adecuados, en particular a través de las cuestiones  prejudiciales10.


\paragraph{Tratado de la Unión Europea}
El Tratado de la Unión Europea (TUE): Creado por el Tratado de Maastricht. este texto incluye disposiciones sobre los principios democráticos, las instituciones, la cooperación reforzada y la acción exterior de la Unión.

\paragraph{Tratado de Funcionamiento de la Unión Europea}
El Tratado de Funcionamiento de la Unión Europea (TFUE): Heredero del Tratado de Roma constitutivo de la CEE, este texto incluye disposiciones detalladas sobre todos los ámbitos de política que lleva a cabo la Unión.

\paragraph{Compromisos internacionales: }
En la jerarquía del Derecho de la Unión. al Derecho primario le siguen los acuerdos internacionales. La UE celebra este tipo de acuerdos sola o con los Estados miembros. en función del reparto de competencias. En los ámbitos de competencia exclusiva. solo la Unión puede negociar y celebrar el acuerdo (por ejemplo, si es de carácter puramente comercial). En los de competencia compartida, el acuerdo es mixto, y debe ser suscrito y ratificado tanto por los Estados miembros como por la UE (por ejemplo, si los acu��rdos incorporan disposiciones judiciales, sociales o de energía).



\subsubsection{Normativa derivada: }
De conformidad con los criterios expuestos al principio del capítulo anterior, entendemos que, aparte del Derecho no escrito, todo el Derecho no comprendido en las normas constitutivas es un «Derecho derivado» en el sentido de que encuentra su «fundamento, alcance y límites» en la norma constitutiva, con independencia de su diversidad de naturaleza y forma. Tal es el objeto de este segundo capítulo relativo a los actos y normas de la Unión Europea. Se trata de un sistema de enorme complejidad sobre el que se quiso introducir un sistema de jerarquía para racionalizar el modelo. Finalmente, el Tratado de Lisboa tampoco estableció expresamente una jerarquía normativa entre los actos comunitarios, pero su reordenación, y en especial su distinción entre «actos legislativos adoptados por el procedimiento legislativoordinario o especial» (art. 289 TFUE) y actos no legislativos, puso las bases para la elaboración de una teoría más perfeccionada sobre la aplicación del principio de jerarquía en el sistema jurídico de la Unión. Particular atención merece desde este punto de vista la mejor formulación de los actos que provienen de una «delegación» del poder normativo. Un acto legislativo puede delegar en la Comisión ese poder para el desarrollo de «actos no legislativos» de alcance general que «completen o modifiquen elementos no esenciales del acto legislativo» (art. 290 TFUE), provengan del Consejo o del PE. El TFUE exige que en estos casos deben los actos adjetivarse con «delegados» o «delegadas» (así Directiva delegada o Reglamento delegado). Es posible que como dice Garzón Clariana haya dejado de ser evidente que la presentación del sistema normativo haya de seguir el tradicional modelo de los actos y exigir un tratamiento desde las categorías de los «actos legislativos y no legislativos» (2010: 757). Siempre conviene atender a las consideraciones de este autor, pero tal vez sea aun algo prematuro.

Nos ocuparemos, pues, en una perspectiva más tradicional en primer lugar, del conjunto de normas y actos de origen estrictamente «institucional», aludiendo a su característica esencial de provenir de la actividad «normativa» atribuida por los Tratados constitutivos a las Instituciones de forma exclusiva. Este conjunto de normas, el Derecho Institucional de la Unión, presenta una cierta complejidad. Persiste tras el Tratado de Lisboa, como hemos dicho en muchos lugares, una cierta dificultad para su propia identificación formal. Hay un cierto número de actos, que provienen de las antiguas estructuras de cooperación (PESC y CPJP), que ni por su naturaleza o alcance son asimilables a las categorías generales que vamos a estudiar aquí. Perviven, en primer lugar, porque son actos en vigor hasta que sean sustituidos por nuevos actos con las nuevas denominaciones1. Además en el caso de la CPJP, formalmente asumida en el sistema general con todo el ELSJ, sus actos deberán asimilarse al modelo general que vamos a estudiar aquí. Sin embargo, como consecuencia de los efectos aplazados por el «Protocolo sobre disposiciones transitorias» y por las excepciones particulares de algunos Estados (véase el Cap. 4) persistirán durante algún tiempo los actos y denominaciones anteriores. En el caso de la PESC no hay asimilación alguna con el sistema general y creemos más adecuado, al menos por el momento, estudiarlos en el ámbito que explica sus particularidades. Por estos motivos nos ha parecido mejor tratar estas cuestiones en los temas correspondientes a las estructuras ELSJ y PESC (Caps. 4 y 25, respectivamente), lo que debería permitirnos mantener aquí la relativa coherencia del sistema general.

El Derecho Institucional de la Unión es, así, nuestro primer objeto de estudio. Por su mayor importancia relativa y por constituir, sin duda de ningún género, el conjunto de regulación más amplio, será tratado con más detenimiento. 

Nos ocuparemos a continuación de la instrumentación normativa internacional, un conjunto de normas y actos cuya característica común es la presencia en todos ellos de un elemento de internacionalidad. En realidad, el criterio esencial que utilizamos no está tanto referido a la mayor o menor participación de las Instituciones, muy variable en intensidad como apreciaremos en los distintos tipos, sino al modelo de instrumentación normativa. Así, por oposición a aquel primer bloque que tiene unos caracteres más cercanos a la actividad «legislativa», este segundo se produce a través del cauce internacional general o convencional. Este sistema de instrumentación normativa internacional, cuya importancia es equiparable a la que ha adquirido en los Derechos internos, ofrece en el sistema de la Unión mayores peculiaridades que nos obligan a una diferenciación de tipos. Incluimos en él, desde luego, la actividad convencional de la Unión con terceros Estados u organizaciones internacionales en virtud del reconocimiento del derecho de concluir tratados, previsto en los Tratados constitutivos. Incluimos también aquellos «actos» que los Gobiernos de los Estados miembros acuerdan «reunidos en el seno» de una Institución (Consejo Europeo o Consejo) pero que no son actos de la Institución sino de los Estados miembros con independencia de su poco correcta denominación. Todos ellos se diferencian en muchos rasgos, y desde luego en la participación institucional que oscila entre su carácter exclusivo hasta prácticamente asumir un papel irrelevante en el acto, pero todos mantienen un rasgo común: su carácter convencional. De ahí que sean objeto de un tratamiento conjunto.



2.1.1. Naturaleza y tipología
El Derecho derivado de naturaleza institucional es el que alberga la actividad «normativa» o, si se prefiere, los «actos jurídicos de la Unión», como los califica la sección primera del Capítulo 2, Título I, Sexta Parte del TFUE. Son, en cualquier caso la expresión formal de la producción por las Instituciones de actos jurídicos vinculantes (y no vinculantes). Algo impropiamente suele identificarse con la «legislación europea», expresión que se viene aceptando en el inadecuado pero expresivo sentido de que en él se contienen loselementos formales más parecidos a la actividad legislativa de unEstado. El TJ para marcar estos caracteres no ha tenido reparo en utilizar tempranamente esta terminología en muchas ocasiones. Así, ha empleado las expresiones «actos cuasilegislativos»2, «ley comunitaria», «legislador comunitario»3 o «poder legislativo», entre otras. 

Estos actos se encuentran previstos con carácter general en el artículo 288 TFUE: Para ejercer las competencias de la Unión, las instituciones adoptarán reglamentos, directivas, decisiones, recomendaciones y dictámenes. En este artículo se determinan los caracteres jurídicos esenciales de los actos. En la tipología se identifican como vinculantes los reglamentos, directivas y decisiones, y como no vinculantes las recomendaciones y dictámenes. Sin embargo, como hemos dicho, esta clasificación debe completarse con otra categorización que el Tratado de Lisboa hace explícita para los actos vinculantes, distinguiendo entre actos legislativos, delegados y de ejecución. Todo ello responde a una evolución compleja (a la que no es ajena el fracasado proyecto constitucional pero en la que no nos podemos detener), cuyo resultado es un modelo de una complicación notable, sin plena coherencia y lleno de ambigüedades. Cierta conciencia de estos problemas tienen las Instituciones de la UE. De ahí que recientemente hayan adoptado el Acuerdo interinstitucional entre el Parlamento Europeo, el Consejo de la Unión Europea y la Comisión europea sobre la mejora de la legislación de 13 de abril de 20164.

Los actos legislativos se identifican por un criterio formal: han de haber sido adoptados por el procedimiento legislativo ordinario (PE y Consejo a propuesta de la Comisión) o por el procedimiento legislativo especial (adoptado por el PE con la participación del Consejo o por el Consejo con la participación del PE). Así, todo reglamento, directiva o decisión, cuyo procedimiento de adopción es calificado por los Tratados como un procedimiento «legislativo» constituye un acto legislativo. Es obviamente un criterio de identificación formal: las distintas bases jurídicas de los Tratados no reservan los procedimientos legislativos exclusivamente para normas materialmente legislativas y resulta claro que, desde la perspectiva de la legitimidad democrática de la norma legal, no puede equipararse el procedimiento legislativo ordinario con un procedimiento donde el PE es meramente consultado por el Consejo. Sin embargo, esta vía de identificación formal no puede confundirse con la carencia de efectos jurídicos: el carácter legislativo de un acto es relevante, por ejemplo, a efectos de la puesta en práctica del control de la subsidiariedad o de quiénes pueden instar el control de su validez. Aún imperfecta, esta categoría recoge actos cuya base jurídica se encuentra necesariamente en el Tratado que determina el procedimiento para su adopción por parte del peculiar «legislador europeo» y que, según establece la jurisprudencia, debe contener necesariamente las opciones políticas básicas que no son susceptibles de ninguna delegación normativa5.

Frente a éstos, todo acto que emane de las instituciones por un procedimiento no calificado como legislativo, será un «acto reglamentario», con independencia de que su base jurídica esté o no en los Tratados. Así, por ejemplo, el artículo 103 TFUE reserva al Consejo, a propuesta de la Comisión y previa consulta al Parlamento, la adopción de reglamentos o directivas para la aplicación de las reglas de competencia en el mercado interior, sin que dicho procedimiento sea «legislativo». De la misma manera, al Consejo Europeo se le otorga una capacidad decisoria que deberá materializarse en actos no legislativos que se distinguirán por su procedencia (Decisiones del Consejo Europeo) o al Banco Central Europeo al que se le confiere asimismo la capacidad de adoptar actos no legislativos, algunos de ellos coincidentes con la denominación de la tipología general (reglamentos) y otros no (orientaciones o instrucciones). Entre los actos reglamentarios, el TFUE distingue dos tipos que debemos reseñar, cuyo fundamento jurídico no se encuentra en el Derecho originario sino en un acto de Derecho derivado: los «actos delegados» y los «actos de ejecución». Como veremos, su régimen jurídico es distinto y esto ha generado ya un importante contencioso interinstitucional. Aunque el TJUE controlará si se han respetado las reglas respectivas de los Tratados, ha reconocido al legislador europeo un amplio margen de discrecionalidad para determinar qué tipo de actos, delegados o de ejecución, puede adoptar la Comisión6. Veámoslos.

Los actos delegados se derivan de una mejor regulación que el Tratado de Lisboa hace de la delegación normativa a la Comisión. Su identificación formal no es difícil pues el artículo 290.3 TFUE señala que se adjetivarán como «delegados», es decir: «Reglamento delegado», «Directiva delegada» o «Decisión delegada». Estos actos que, insistimos, tendrán aquellos efectos jurídicos propios de reglamentos, directivas y decisiones tal y como estudiaremos seguidamente, poseen una jerarquía inferior al acto legislativo, pues están fundamentados y limitados doblemente: por el acto legislativo delegante que fijará los objetivos, contenido, alcance y extensión de la delegación y por el Tratado, ya que tales actos solo podrán «completar o modificar aspectos no esenciales del acto legislativo» 7. Su control, como se verá en el Capítulo dedicado a la Comisión correrá a cargo del Consejo y del Parlamento en las condiciones que establezca el mismo acto legislativo delegante. 

Por su parte, los «actos de ejecución» se encuentra regulados en el artículo 291 TFUE. De nuevo su identificación formal no es compleja pues llevarán en su título la expresión «de ejecución» (así: «Reglamento de ejecución», «Directiva de ejecución» o «Decisión de ejecución»). Dichos actos emanarán generalmente de la Comisión y, en circunstancias excepcionales del Consejo (habitualmente en PESC, pero también existen, por ejemplo, en materia de gobernanza económica o de restablecimiento temporal de controles en las fronteras interiores). El dato central que identifica estos actos de ejecución es la «necesidad de condiciones uniformes para la aplicación de un acto de la Unión jurídicamente vinculante» al que se encuentran jerárquicamente subordinados. Será, en consecuencia, tal acto el que fijará los ámbitos y contenidos de los actos de ejecución, esto es, la medida en que el acto de ejecución precisará (sin modificarlo8) el contenido del acto jurídicamente vinculante, y señalará cuál es el procedimiento por el que los Estados miembros (que son los competentes en materia de ejecución del Derecho de la Unión) controlarán las competencias «de ejecución» otorgadas a la Comisión.

Para terminar este confuso panorama, es importante en la caracterización general distinguir además de entre vinculantes y no vinculantes y, dentro de los vinculantes, los que tienen un alcance general o individual. La distinción es fundamental, naturalmente, porque en función de ella se determinan su obligatoriedad y efectos desde el punto de vista de los sujetos o destinatarios de los mismos. Pero también se vincula a esta doble distinción, aparte su diferente régimen desde la perspectiva de su control jurisdiccional que será objeto de atención en el capítulo correspondiente, el régimen de publicación, notificación y vigencia.


A continuación en la jerarquía, encontramos una fuente jurídica esencial, el Derecho derivado, constituida por los actos jurídicos emanados de las instituciones comunitarias, con los que se llenan de contenido las competencias atribuidas en los Tratados. Según el Tratado de Lisboa, los actos jurídicos de la Unión son: el reglamento, la directiva y la decisión (vinculantes), la recomendación y el dictamen (no vinculantes). En muchos casos, los Tratados detenninan el tipo de acto jurídico que ha de adoptarse. Cuando no es así, las instituciones deben decidirlo según los principios de proporcionalidad y subsidiariedad.


\paragraph{Actos típicos} 
2.2.1.1. El reglamento 
El artículo 288 TFUE establece que: El reglamento tendrá un alcance general. Será obligatorio en todos sus elementos y directamente aplicable en cada Estado miembro.  Se trata indiscutiblemente del instrumento de regulación jurídica más acabado dentro del sistema de la Unión. Buena parte de la doctrina, en particular la más empeñada en hacer encajar este sistema dentro de los sistemas de fuentes estatales, resaltan aquellos caracteres de este tipo normativo más similares a los de las leyes (generalidad, abstracción, efectos erga omnes y otros caracteres con los que la teoría más tradicional cualifica la ley). Preferimos por razones que ya han sido expuestas caracterizarlo, como ha dicho Louis, como «el acto más completo y eficaz de la gama de instrumentos de que disponen las Instituciones» (1995: 109) y atenernos en tanto que tal a su cualidad jurídico-normativa en la que destaca su mayor intensidad normativa por relación a los otros instrumentos normativos. Esta intensidad se manifiesta en sus caracteres:

a) Alcance general: Esta característica del reglamento sirve básicamente para diferenciarlo de los actos de alcance individual como la decisión. Se trata de una característica esencial y que, en consecuencia, es requisito necesario para que una norma tenga el carácter de reglamento con independencia de su denominación o la forma de su adopción. Se trata de una característica objetiva del acto que tiene importantes efectos en su control de legalidad. Como tal carácter objetivo, el «alcance general» no proviene tanto de la determinación del número o identificación de sus destinatarios como de su previsión de aplicabilidad a «situaciones objetivamente determinadas» y con efectos jurídicos respecto de una categoría de sujetos contemplada «de manera general y abstracta», como advirtió el TJUE13. Así, en el asunto Calpak, el TJ estimó que la naturaleza reglamentaria de una disposición no puede cuestionarse simplemente por la posibilidad de «determinar el número o la identidad de (en este caso) los productores beneficiarios de la ayuda»14, cuando esa determinación o identificación proviene, como en este caso, de una situación objetivamente determinada y los destinatarios lo son como consecuencia de estar éstos incursos en la situación que la norma prevé con carácter general y abstracto. Por el contrario, naturalmente, cuando la determinación de los destinatarios y su identificación provienen de una «situación» concreta, el TJ ha rechazado que se trate de un Reglamento. Es el caso que se produjo en el asunto de la International Fruit Company, donde el Tribunal estimó que el «Reglamento» por el que la Comisión otorgaba unas licencias de importación a las empresas que habían depositado su solicitud en una fecha determinada no era en realidad tal reglamento en la medida en que tales empresas en la fecha en cuestión habían sido ya identificadas, por lo que pese a su denominación y forma de adopción, el Tribunal determinó que la naturaleza del acto era la de una decisión (colectiva) y no un reglamento15.

b) Obligatoriedad en todos sus elementos: Se trata de una característica que tiene una doble función de identificación del reglamento. En primer lugar, afirma su carácter obligatorio frente a los actos no vinculantes y en segundo lugar al incidir en que lo es «en todos sus elementos» lo diferencia de aquellas normas, como la directiva, que obligan en el resultado, pero no en los medios. Esta noción de plenitud del efecto obligatorio del reglamento fue claramente definida por el TJ en la sentencia más significativa a estos efectos, la sentencia de 7 de febrero de 1973 (Comisión c. Italia), en la que el Tribunal dijo que la consecuencia de esta característica es que: no se podría admitir que un Estado miembro aplique de manera incompletao selectiva las disposiciones de un reglamento de la Comunidad, de manera que frustrara la aplicación de ciertas disposiciones de la legislación comunitaria respecto de las cuales hubiera manifestado su oposición o las hubiera estimado contrarias a ciertos intereses nacionales16. Además, según el Tribunal, desde que un reglamento es aprobado válidamente, reténgase en este sentido que los reglamentos gozan de una «presunción de validez» y, que por tanto, su ilegalidad sólo puede ser establecida por el Tribunal a través de los procedimientos previstos a tal efecto17, el alcance objetivo de las reglas establecidas [...] no podría ser modificado por reservas u objeciones que los Estados miembros hubiesen formulado en el momento de su elaboración [...] [ni] dificultades de aplicación aparecidas en la etapa de la ejecución [...] podrían permitir a un Estado miembro liberarse unilateralmente del cumplimiento de sus obligaciones18.

c) Directamente aplicable en todos los Estados miembros: Este último rasgo del reglamento tiene un marcado carácter instrumental en la medida en que lo que hace es eliminar toda intermediación por parte de los Estados, asegurando que el reglamento despliegue de manera efectiva y uniforme sus efectos de alcance general y de plenitud de su obligatoriedad. La «aplicabilidad directa», así entendida, se traduce en dos caracteres sustanciales: uno de orden positivo, por el que se le atribuye a este tipo normativo la cualidad de generar derechos y obligaciones por sí mismo para los órganos y sujetos dependientes de un ordenamiento estatal (lo que se asocia al «efecto directo» de estas normas) y otro, que encierra un mandato negativo, que prohíbe cualquier actividad por parte del Estado miembro que pueda poner en cuestión la inmediatez de la efectividad del reglamento. Si la famosa sentencia Van Gend en Loos lo apreció en normas constitutivas, en la sentencia Simmenthal de 9 de marzo de 1978 lo aprecia tanto en disposiciones de los tratados como en las reglamentarias19.

En relación con el segundo de los aspectos señalados, es decir, la prohibición consecuente de toda intermediación del Estado, la jurisprudencia ha precisado también su alcance. Y es que, en efecto, el reglamento es un tipo normativo tendencialmente«autosuficiente», esto es, que no requiere de ninguna medidainterna (legislativa, administrativa o judicial) para producir sus efectos. Pero esto, pese a lo que en ocasiones se dice, no elimina que, en algunos casos se exija de la intervención del Estado para dar cumplimiento al mismo. Por esta razón, el TJ parte de la noción central de la aplicabilidad directa del reglamento, según la cual exige que su aplicación se realice sin medida alguna que pueda suponer una recepción en el Derecho nacional20. 

A partir de esta premisa, el TJ ha ido precisando en su jurisprudencia el sentido concreto de la misma. En primer lugar, que esta prohibición no opera, naturalmente, cuando el reglamento en cuestión deja la responsabilidad a los Estados miembros de adoptar ellos mismos las medidas legislativas, reglamentarias, administrativas y financieras necesarias para que las disposiciones del reglamento puedan ser efectivamente aplicadas21.

En segundo lugar, precisa que su transformación en una norma interna no es una medida necesaria y en consecuencia es incompatible con el Derecho de la Unión (sentencia de 7 de febrero de 1973, donde rechazó la actitud del Gobierno italiano de recoger en un decreto las disposiciones del reglamento en cuestión). En tercer lugar, que es, asimismo, «incompatible» con el Derecho de la Unión toda disposición de un ordenamiento jurídico nacional —legislativa, administrativa o judicial— que tuviese por efecto disminuir la eficacia del Derecho comunitario (Simmenthal). 

En cuarto lugar, el Tribunal también ha establecido que es igualmente incompatible no adoptar las medidas necesarias para permitir «la aplicación efectiva y en los plazos apropiados» de un reglamento22. Por último, y como protección y garantía de la inmediatez del reglamento, que sus efectos sobre los particulares (y, en consecuencia, su propia invocabilidad ante los tribunales) nunca depende de las medidas que pueda exigir su efectiva aplicación23.

2.2.1.2. La directiva

Frente a estos caracteres del reglamento, la directiva es definida apoyándose en los caracteres que la diferencian del reglamento. En efecto el artículo 288 TFUE dispone que: La directiva obligará al Estado miembro destinatario en cuanto al resultado que deba conseguirse, dejando, sin embargo, a las autoridades nacionales la elección de la forma y de los medios.

Se trata, en esta perspectiva, de un instrumento regulador sin alcance general (obliga sólo al Estado destinatario). Frente a la plenitud de la obligatoriedad del reglamento, la directiva tiene una obligatoriedad parcial que recae sobre el elemento «resultado» y, desde luego, en la medida en que deja en manos de las autoridades nacionales la elección de la «forma y los medios» de darle efectividad en el orden interno, es decir, que requiere la intermediación estatal, carece por definición de aplicabilidad directa (y, en consecuencia, de «efecto directo»).

Esta caracterización por referencia negativa al reglamento es particularmente útil para poner de manifiesto la razón última de ser de este instrumento normativo, de su compleja y poco pacífica evolución, a duras penas encauzada por una impresionante labor jurisprudencial que no ha conseguido evitar que siga siendo el instrumento normativo más problemático del modelo europeo. Así, frente a la exigencia de aplicación uniforme del Derecho de la Unión (pilar donde reside la garantía y seguridad del edificio jurídico europeo) no cabe más alternativa que la sustitución de las legislaciones nacionales por una verdadera «legislación comunitaria» (lógica a la que responde el reglamento) o la configuración de un instrumento que, sin sustituir el poder legislativo nacional, permita la armonización de las legislaciones nacionales. Es justamente la opción por esta segunda alternativa la que origina la creación del instituto de las directivas. Un instrumento que, dirigido al poder en el que reside la capacidad legislativa (los Estados), obliga a éstos a alcanzar un «resultado» común. La convergencia obligada de las legislaciones nacionales en ese «resultado común» se identifica con la función «armonizadora» de la directiva. La no sustitución de la competencia normativa de losEstados explica todos los caracteres de la directiva: la exigencia deintermediación normativa de los mismos para alcanzar el«resultado»; que sean las propias autoridades nacionales las quedeterminen «las formas y los medios» para su consecución; y la vinculación de los efectos jurídicos a las normas estatales de intermediación y no a la directiva, excluyendo en consecuencia su posibilidad de producir efecto directo.

En definitiva, pues, la directiva es un instrumento jurídico que responde, esencialmente, a los siguientes caracteres:

a) La directiva impone a los Estados una obligación de resultado que debe ser cumplida en el plazo determinado por la propia directiva. La naturaleza de esta obligación ha sido objeto de diversas lecturas. Se encuentran, por una parte, las que resaltan su carácter internacional atendiendo a su destinatario (los Estados) y al margen de discrecionalidad de los mismos en cuanto a su aplicación. Por otra parte, la lectura «autónoma» concibe la obligación dimanada de la directiva como una obligación de «ejecución» más que de «recepción», entendiendo la directiva como un instrumento normativo integrado y equiparable desde esaperspectiva al resto de las normas de la Unión de las que sediferencia, eso sí, por el sistema mediato de la producción de sus efectos al exigir la intervención del Estado. Esta obligación de ejecución es absoluta de manera que su inobservancia supone un incumplimiento del Derecho de la Unión que puede ser sancionado mediante las vías de recurso previstas en el Tratado y con las consecuencias jurídicas que de tal incumplimiento se derivan.

b) Concebida como tal obligación de ejecución, la directiva no sólo exige la consecución en abstracto de un resultado, sino que impone a los Estados la obligación de adoptar las medidas nacionales necesarias para alcanzar dicho resultado con la libertad de «elección de la forma y de los medios». La exigencia de intermediación normativa se sustancia, pues, en la obligación de los Estados de adoptar las normas internas necesarias para el cumplimiento del objetivo previsto en la directiva con el margen de discrecionalidad que supone la elección de la forma y los medios, lo que se conoce como transposición de la directiva al Derecho interno. 
Al tratarse de una obligación de ejecución, en realidad, esa elección de la forma y de los medios, no implica una libertad absoluta por parte del Estado, sino que, reconociendo su autonomía institucional, se traduce en una obligación más condicionada de lo que, en una primera lectura, sugiere el artículo 288. Además de estar condicionada, naturalmente, por el resultado previsto y por los plazos fijados en la propia directiva, cabe destacar otros dos condicionantes generales: primero, la exigencia de «adecuación» de la forma y los medios para la consecución del objetivo; y, segundo, las exigencias de seguridad jurídica respecto de las normas internas de transposición, al ser éstas el origen de los efectos jurídicos. 

c) La vinculación de los efectos jurídicos a la norma nacional de transposición constituye el tercero de los elementos esenciales en la configuración formal de la directiva en el Tratado. Ésta es la característica fundamental que traslada la ausencia de aplicabilidad y de eficacia directa de la directiva. Y es, también, la que introduce un alto grado de complejidad en este instrumento jurídico. El problema de fondo aparece básicamente en los supuestos en los que no se procede a su transposición en los plazos previstos o cuando esta operación se realiza de forma incompleta o incorrecta por parte del Estado. Esta disfunción constituye, naturalmente, un incumplimiento del Estado que, como tal, puede ser sancionado. Esta sanción, sin embargo, deja abierto el problema de la inaplicabilidad efectiva de la directiva que puede comportar una lesión de los derechos de los particulares. La solución global de ese problema no es, desde luego, fácil. La respuesta dada por el TJ se ha encaminado a dotar de efectividad a la directiva y a proteger los derechos de los particulares y ha consistido en asegurar el «efecto útil» de la directiva24. A tal efecto, ha obrado sobre dos elementos fundamentales: asegurando la obligación de resultado, y delimitando, en función de ella, el alcance de libertad de elección de las formas y medios. El Tribunal ha advertido que en todos los casos en que las disposiciones de una directiva parecen ser, desde el punto de vista de su contenido, incondicionales y suficientemente precisas, dichas disposiciones si no se han adoptado dentro del plazo prescrito medidas de aplicación, pueden ser invocadas contra cualquier disposición nacional no conforme a la directiva o, en la medida en que definen derechos que los particulares pueden alegar, frente al Estado25.

En definitiva, que la directiva, en estas condiciones, se constituye en un acto con la intensidad normativa suficiente para imponerse frente a la legislación o reglamentación interna contraria y para proteger, en su caso, el derecho o pretensión del particular frente a su Estado. Para conseguir esa oponibilidad frente a la norma interna contraria o al Estado, el Tribunal ha reconocido su invocabilidad ante los órganos jurisdiccionales internos, es decir, su efecto directo. Es lógico, naturalmente, que esta «invocabilidad» sólo pueda producirse cuando el contenido normativo de la directiva tenga la precisión suficiente para sustentar una pretensión jurídica o un derecho y esté formulado de manera incondicional. Se trata de un efecto directo vertical ya que se contempla en la relación del particular frente al Estado, pero, eso sí, en este solo sentido, puesto que el Estado en cambio no puede invocar o prevalerse de la directiva frente a los particulares. El Tribunal lo ha repetido en una extensa jurisprudencia26.

Esta jurisprudencia que ha ido paulatinamente forjando a la directiva con unos caracteres lejanos de los que inicialmente se le supusieron, ha hecho de ella un verdadero instrumento normativo del Derecho de la Unión. Llevada demasiado lejos esta tesis estuvo cerca de considerar un efecto directo horizontal felizmente reconducido en las sentencias Faccini Dori de 14 de julio de 1994 y El Corte Inglés de 7 de marzo 1996.

2.2.1.3. La decisión
La decisión es caracterizada por el artículo 288 TFUE como «obligatoria en todos sus elementos para todos sus destinatarios». Esta lacónica, y en cierto modo imprecisa caracterización, nos enfrenta, en principio y sólo en principio, a un acto jurídico individual dirigido a uno o varios destinatarios a los que obliga en la totalidad de sus elementos, lo que significa tanto en «la forma y medios» como en su «resultado».

La identificación concreta (y no abstracta como en el reglamento) del destinatario o destinatarios es la que le otorga ese carácter de acto jurídico individual. Unido esto a su carencia de alcance general ha hecho que en muchas ocasiones se presente a la decisión como algo similar al acto administrativo propio de los Derechos internos. Una vez más la particularidad del sistema de la Unión inutiliza este tipo de comparaciones.

La concepción funcional que pesa sobre estos actos jurídicos, incluida la decisión, hace que estas categorías se conciban en función de unas finalidades propias del sistema y que gocen de una gran adaptabilidad a las exigencias de esos objetivos y del contexto donde han de producir sus efectos. Así, cuando la decisión tiene como destinatario un particular, sea éste una persona física o jurídica, sus caracteres básicos se dan con nitidez y guarda aquella similitud con el acto administrativo. Sin embargo, cuando las decisiones se dirigen a un Estado o Estados miembros se desdibujan esos perfiles en función de la naturaleza del destinatario o destinatarios. Además, en muchas ocasiones, la decisión no alcanza a establecer el «medio», recordándonos estas decisiones, en este aspecto, a las directivas. Pero es más, cuando la decisión es dirigida al Estado o Estados miembros, la idea de que sus efectos jurídicos sólo son para el destinatario o destinatarios de la decisión pueden también alterarse. Parece evidente que, aun con todos los matices y precisiones necesarios, hay ciertas ocasiones en que una decisión puede gozar indirectamente, puesto que provendrá del Tratado o de una norma de alcance general en última instancia, del efecto directo. No es extraño así que el TFUE haya previsto la posibilidad de adoptar decisiones con arreglo al procedimiento legislativo ordinario. Una vez más, nos encontramos con la dificultad de encerrar un tipo normativo en una categoría formal y estática. Estos tipos, construidos y caracterizados en el Derecho de la Unión en función de su utilidad para determinados fines, evolucionan y se adaptan a la propia expansión y exigencias del sistema jurídico de la UE. El TJ ha entendido así su función y, por ello, prescinde de su «denominación» y forma para ejercer su control en función del contenido y tipo de efectos jurídicos del acto. La decisión es, sin duda, un caso paradigmático.

2.2.2. Actos atípicos

Se alude generalmente con ello a un número abierto de actos de naturaleza un tanto oscura y efectos jurídicos indeterminados cuya presencia en el orden jurídico de la Unión es lo único indiscutible. «Decisiones del Consejo» (distintas por sus caracteres a las que acabamos de analizar, «Decisiones del Consejo y de los representantes de los Gobiernos de los Estados miembros», «Resoluciones», «Programas» suelen encuadrarse en este albergue de los «actos atípicos» cuando no encuentran mejor acomodo con la única condición de que tengan un origen institucional y se enmarquen en los ámbitos de competencia de la Unión. Tan curiosa situación, y en ocasiones sospechosa creatividad institucional, debe explicarse, en mi opinión, por varias razones. A veces, los Tratados han previsto la competencia pero sin precisar el tipo de acto correspondiente a su ejercicio; en otras ocasiones se trata de actuaciones de las Instituciones en ámbitos sobre los que pesa una cierta debilidad en la determinación de la competencia pero que no excluye ciertas posibilidades de actuación institucional y en las que, por diversas razones, se prefiere eludir los actos típicos; o, por último, hay ocasiones en que en actos internos (reglamentos internos de las Instituciones) se prevé la adopción de ciertas regulaciones no contempladas en las categorías generales (un caso excepcionalmente claro fue la «Decisión» del Parlamento Europeo de 9 de marzo de 1994 sobre el Defensor del Pueblo Europeo). Es bastante probable que el Tratado de Lisboa haya agravado la cuestión del uso de estos actos atípicos en la medida en que en demasiadas ocasiones no prefigura el «tipo» normativo que se debe utilizar. Tan seria es esta falta de concreción que el artículo 296 TFUE establece que: Cuando los Tratados no establezcan el tipo de acto que deba adoptarse, las instituciones decidirán en cada caso conforme a los procedimientos aplicables y al principio de proporcionalidad. 

Esta «flexibilidad» unida a ciertas indeterminaciones como que las Instituciones adoptarán «principios», «reglas», «normas», «orientaciones», «instrucciones», entre otras, parecen abonar el campo donde fructifican los actos atípicos. 

En todas estas circunstancias, seguramente, resultará de interés la concepción no formal de su función de control, asumida por el TJ que ha extendido su control a todo acto destinado a producir efectos jurídicos con independencia de la «forma» o «denominación» que adopte, excluyendo solamente de su control de legalidad los actos no vinculantes (recomendaciones y dictámenes)28. En los tiempos que se avecinan será especialmente útil esta comprensión del Tribunal sobre la flexible construcción del sistema.

Mención aparte merece la técnica de los «Acuerdos Interinstitucionales», hasta ahora fundamentados en una mera práctica, (pero de importantes efectos en la vida de la Unión) que tras Lisboa ha encontrado su fundamento en el artículo 295 TFUE a cuyo tenor: El Parlamento Europeo, el Consejo y la Comisión llevarán a cabo consultas recíprocas y organizarán de común acuerdo la forma de su cooperación. A tal efecto y dentro del respeto de los Tratados, podrán celebrar acuerdos interinstitucionales que podrán tener carácter vinculante.

Es evidente que ese «podrán» tener carácter vinculante es una opción que no permite su definición previa ni su tipificación respecto a sus efectos. Entiendo, que de otorgarles ese valor las propias instituciones implicadas en el acuerdo su efecto vinculante se circunscribe a la relación interinstitucional exclusivamente. Se publican en el DOUE en la serie L entre los «actos no legislativos»29.



Reglamento: Es la nonna que más se parece al concepto habitual de ley por sus características, muy poco comunes en el Derecho internacional clásico.

• Alcance general: El reglamento se aplica de forma general a situaciones objetivamente determinadas, no a destinatarios identificables.

• Obligatorio en todos sus elementos: El reglamento prescribe tanto el resultado como los medios. Garantiza la aplicación uniforme del Derecho de la Unión en todos ellos. Excluye la posibilidad de aplicar normas nacionales incompatibles con sus cláusulas. por lo que puede ser necesario modificar la no,mativa nacional para evitar conflictos con un reglamento.

• Aplicabilidad directa en cada Estado miembro: Este acto se aplica en todos los Estados desde su entrada en vigor (en la fecha que el propio reglamento fije o, en su defecto, el vigésimo día siguiente al de su publicación en el Diario Oficial de la Unión Europea), sin que sea necesaria ninguna incorporación expresa ni transposición del mismo en los ordenamientos nacionales.

• Efecto directo: El reglamento genera derechos y obligaciones para los Estados, sus órganos y para sus particulares. Los reglamentos pueden ser invocados ante los jueces nacionales y aplicados a un caso concreto, tanto desde una perspectiva vertical (particulares frente a Estados u organismos públicos) como horizontal (entre particulares).

\paragraph{Directiva}
Es la norma que intenta annonizar (no unificar) los ordenamientos jurídicos de los Estados miembros. Obliga a los Estados miembros destinatarios en cuanto al resultado que debe conseguirse, dejándoles la elección de la forma y de los medios para alcanzarlo. Los Estados ·miembros han de trasponerla, es decir desarrollarla y concretarla mediante una norma interna, en el plazo establecido para ello. Dado que la directiva tradicional no es lo bastante concreta, no puede tener efecto directo. No obstante, la acción jurisprudencia! ha ido reconociendo el efecto directo vertical (invocar la aplicación de la directiva frente al Estado o ente público) a algunas directivas que sí llegan a desarrollar medios e instrumentos.\footnote{%
    El Tribunal de Justicia ha decidido que. excepcionalmente. las disposiciones de una Directiva también puedan ser directamente aplicables en un Estado Miembro. sin necesidad de un acto de transposición, siempre y cuando el plazo de transposición haya vencido y la Directiva no haya sido transpuesta o lo haya sido deficientemente
}

\paragraph{Decisiones}
Decisión: Es un tipo de acto que obliga en todos sus elementos a sus destinatarios (Estados miembros, personas fisicas o personas jurídicas). Viene a ser el acto de ejecución administrativa en la UE (por ejemplo, las decisiones se utilizan para pronunciarse sobre las propuestas de fusión de empresas). Las decisiones puede ser muy detalladas y prescribir los medios para alcanzar un resultado. Una decisión dirigida a un Estado miembro puede, además, producir un efecto directo para los ciudadanos de la Unión en las mismas condiciones que una directiva.



Recomendación: Es un tipo de acto que sugiere a los destinatarios un comportamiento determinado. Aunque no obliga jurídicamente a adoptarlo, puede ser muy relevante políticamente (por ejemplo, el Consejo adopta cada año recomendaciones específicas por país en el marco del Semestre Europeo).

Dictamen (avis): Es un tipo de acto en el que una institución u organismo comunitario (Comisión, Parlamento, Comité Económico y Social, Comité de las Regiones, etc.) expresa un punto de vista sobre un asunto o propuesta especifica. Aunque no es vinculante, pued�� generar indirectamente efectos jurídicos (por ejemplo, un reglamento puede ser anulado si en un proceso de elaboración no se ha pedido opinión al Parlamento cuando era preceptivo).


Para tenninar con el repaso de las fuentes de Derecho de la Unión, cabe señalar que también existen principios generales del Derecho europeo que carecen de fonnulación normativa concreta (no están escritos), pero expresan las concepciones elementales de Derecho y justicia a las que obedece el ordenamiento jurídico. Éstos han desempeñado un papel esencial en la construcción europea, al hacer posible unas interpretaciones extensivas de las nonnas comunitarias, consolidando su fuerza vinculante y la supranacionalidad normativa. El Tribunal de Justicia los ha desarrollado a través de su jurisprudencia, tomando como referencia básica los principios comunes a los ordenamientos jurídicos de los Estados miembros. Por ejemplo, en el ámbito de la integración económica, el Tribunal de Justicia ha reconocido los principios de equilibrio institucional, preferencia comunitaria, no discriminación, etc.




\paragraph{Actos no vinculantes}
2.3. LOS ACTOS NO VINCULANTES
Se engloban en esta categoría las recomendaciones y los dictámenes aludidos en el artículo 288 TFUE, que son caracterizados exclusivamente por su naturaleza no vinculante. Es habitual en la doctrina identificar el contenido de la recomendación con la indicación de una conducta a seguir o la modificación de una situación o comportamiento, mientras que el dictamen contendría una opinión o valoración de situaciones o conductas. Su diferenciación es, en la práctica, complicada. Su carencia de obligatoriedad lleva a parte de la doctrina a considerar que, simplemente, no son actos jurídicos. En nuestra opinión, la carencia de obligatoriedad no los priva del carácter jurídico, sino, exclusivamente, de sanción directa. Una vez más su incardinación en el sistema puede dar alguna respuesta. Tanto un instrumento como otro adquieren significado por relación a los contenidos normativos obligatorios y en función de ellos pueden tener vinculadas, de manera indirecta, consecuencias jurídicas, excluida naturalmente la sanción.

Así, en la sentencia Grimaldi de 13 de diciembre de 198930 en relación con una recomendación de la Comisión sobre la conveniencia de la adopción de una lista europea de enfermedades profesionales, el TJ afirmó que esta lista establecida en la recomendación debería ser considerada a efectos de interpretación de las normas nacionales y de la Unión por los jueces nacionales. Que a un dictamen se le pueden vincular «consecuencias jurídicas» parece indiscutible como en el caso de los «dictámenes motivados», cuya inobservancia por el Estado en el plazo fijado por el dictamen faculta a la Comisión para interponer el recurso por incumplimiento según prevé el artículo 258 TFUE.

\paragraph{Instrumentación normativa internacional}
3.1. EL DERECHO INTERNACIONAL
El alcance y valor del Derecho Internacional General en el Derecho de la Unión presenta unos perfiles problemáticos. Buena parte de las dificultades proviene de la compleja estructura del sistema jurídico de la Unión y de la ausencia de una norma que lo clarifique. A veces se agrava la situación por percepciones algo radicales entre los que no acaban de ver la UE más que como una organización internacional, en la que el Derecho Internacional sería la sustancia misma del modelo, y los que, más allá de lo que exige una razonable consideración de la especificidad y autonomía del Derecho de la Unión, pretenden hacer entrar esta cuestión a todo trance por una concepción constitucional de la relación Derecho Internacional/Derecho interno. La complejidad del tema no puede ni debe ser resuelta con fórmulas simplificadoras que, sencillamente, no responden a la realidad. Se trata, pues, de un problema que proviene fundamentalmente de una cierta indefinición misma del sistema de la Unión y no de errores conceptuales de la doctrina. Sea como fuere lo cierto es que el «valor» del Derecho Internacional en el Derecho de la Unión no es un tema resuelto. No va más allá de una expresión de intenciones la llamada de atención del artículo 3 TUE, en su apartado 5, en el sentido de que: En sus relaciones con el resto del mundo, la Unión afirmará y promoverá sus valores e intereses y contribuirá a la protección de sus ciudadanos. Contribuirá a la paz, la seguridad, el desarrollo sostenible del planeta, la solidaridad y el respeto mutuo entre los pueblos, el comercio libre y justo, la erradicación de la pobreza y la protección de los derechos humanos, especialmente los derechos del niño, así como al estricto respeto y al desarrollo del Derecho Internacional, en particular el respeto de los principios de la Carta de las Naciones Unidas.

Entendida y asumida esta especificidad y autonomía, no creemos que el sistema jurídico comunitario pueda calificarse especialmente de poco receptivo del Derecho Internacional, como acertadamente ha dicho Díez-Hochleitner (1998: 11-15). Una vasta jurisprudencia así lo confirma31. Sin embargo, en los últimos tiempos, en particular desde la sentencia Kadi, de 3 de septiembre de 200832, que sembró la duda y el desconcierto en la consideración del valor del Derecho Internacional en el Derecho de la Unión, el asunto no está tan claro. Según el Tribunal, las obligaciones impuestas por un acuerdo internacional [la Carta] no pueden tener por efecto menoscabar los principios constitucionales del TCE (ahora TFUE)33.

Pese a todo la situación no es ni más ni menos espinosa que en los sistemas constitucionales de sus Estados miembros. La ausencia de una norma de sumisión al Derecho Internacional General carece de significado particular ya que en este ámbito el Derecho Internacional se impone a los Estados y a las Organizaciones internacionales, sin requerir una formulación expresa ya que su validez se sitúa en el propio Derecho Internacional General y no en los ordenamientos constitucionales o de Organizaciones internacionales. El Estado español carece de una norma así y no por ello se considera poco receptivo. El sistema de la Unión está sujeto al Derecho Internacional General en la medida en que lo están todos los sujetos de este ordenamiento, naturalmente, en el ámbito de sus competencias.

Pero, además, el TJUE ha dado relevancia al Derecho Internacional a través de la vía de los principios generales (ya aludidos), incorporando criterios de interpretación de tratados de codificación o en materia de derechos humanos, entre otras. Es ya una línea tradicional de jurisprudencia del TJCE utilizar el Convenio de Viena de 1969 sobre Derecho de los Tratados como elemento de interpretación en tanto que norma de Derecho Internacional general34. Incluso, recientemente, ha mostrado su posición sobre los difíciles contornos de los «principios de derecho internacional consuetudinarios», que tienen un grado de precisión menor que las obligaciones convencionales, pero advirtiendo que su control jurisdiccional alcanza a analizar si las Instituciones de la Unión al adoptar un acto, han cometido «errores manifiestos de apreciación» en cuanto a la aplicación de estos principios35. 

Distinta es, sin embargo, la cuestión de valorar el alcance en el Derecho de la Unión de las obligaciones jurídico-internacionales de sus Estados miembros, en tanto que sujetos del Derecho Internacional en sus ámbitos de competencias no «atribuidos» a la Unión Europea. En este sentido, el Derecho de la Unión parte de algunas disposiciones que establecen los parámetros básicos de esa relación.

El artículo 351 TFUE establece que los Tratados (TUE/TFUE) no afectarán a los derechos y obligaciones que resulten de convenios celebrados con anterioridad al 1 de enero de 1958 o, para los Estados que se hayan adherido, con anterioridad a la fecha de su adhesión, entre uno o varios Estados miembros, por una parte, y uno o varios terceros Estados, por otra. Suma a ello esta misma disposición la obligación de los Estados miembros de recurrir «a todos los medios apropiados» (evidentemente son los que establece el Derecho Internacional) para eliminar la parcial o total incompatibilidad de aquéllos con los Tratados, instaurando una obligación de todos los Estados miembros de cooperar con el Estado que se encontrase en esta situación para subsanarla y la obligación de tener en cuenta las particulares implicaciones de la Unión derivadas de estos casos.

La cuestión, aparentemente clara en el sentido de respetar las obligaciones internacionalmente asumidas por un Estado o Estados miembros y los derechos de terceros Estados, ha presentado algunos perfiles problemáticos en relación con su alcance en el Derecho de la Unión. La cuestión básica es determinar si esos tratados vinculan a la Unión y con qué alcance. Dos situaciones pueden distinguirse con diferente tratamiento jurídico: a) la de aquellos convenios que no afectan a la competencia de la Unión (p. ej., los tratados de delimitación fronteras entre Estados), en los que la Unión queda obligada a su respeto, y b) aquellos tratados que afectan a ámbitos de competencias atribuidas a la Unión, donde el corazón del asunto se sitúa en si la Unión sustituye al Estado o Estados en sus obligaciones con respecto a terceros y en las consecuencias de la opción por una u otra alternativa.

En relación con esta segunda cuestión pesa una cierta confusión. En una importante sentencia del TJ, relativa al GATT de 1948, se consideró que la Unión «sustituía» a los Estados miembros36. Ahora bien, el TJ fundamentó esta consideración en varias circunstancias particulares del caso (eran parte casi todos los Estados miembros, se trataba de un ámbito de competencias transferido de manera muy clara, existía un acuerdo unánime en aceptarlo...) que permiten dudar de su traslación a todos los supuestos sin más. Así, por ejemplo, no las aprecia en el asunto Levy37. En el caso primero, es evidente que se producen en cadena una serie de efectos derivados de su inclusión en el orden jurídico de la Unión, en particular el de su jerarquización sobre el Derecho derivado. En el segundo, es decir, en la no sustitución, se mantiene el respeto de las obligaciones contraídas con terceros y para el Estado o Estados concernidos se crea lo que con razón ha calificado Louis de «oponibilidad» (1995: 123) frente a las Instituciones europeas que, como dice el TJ, «no deben obstaculizar la ejecución de estos compromisos por el Estado miembro de que se trate»38. La sentencia Levy mantiene este sentido y cuestiona, como ha dicho Roldán Barbero, el reproche al TJUE «de ser más sensible a los principios internos comunitarios que a las reglas de Derecho Internacional» (1994: 494).

3.2. LA ACTIVIDAD CONVENCIONAL DE LA UNIÓN 
En virtud del poder de concluir tratados que las normas constitutivas atribuyen a las Instituciones, el ordenamiento jurídico de la UE utiliza, de forma similar a los Estados, de manera creciente este sistema de instrumentación jurídica. En un sentido estricto es posible constatar la ausencia de una norma que establezca el sistema de recepción del derecho convencional en el orden jurídico de la Unión, pero el segundo apartado del artículo 216 TFUE, aunque de manera insuficiente, se acerca a ese tipo de normas: Los acuerdos celebrados por la Unión vincularán a las instituciones de la Unión y a los Estados miembros.

Desde luego, la fórmula no es la mejor de las posibles, ya que, aunque establece con nitidez su «obligatoriedad», no aclara su eventual integración en el orden jurídico de la Unión, ni su alcance en relación con la jerarquía ni sus posibles efectos en relación con los particulares. Una vez más, una interpretación sistemática y la jurisprudencia del TJ han permitido precisar su alcance, aunque la cuestión no está exenta de imprecisiones y ambigüedades. En relación con su integración en el Derecho de la UE, la sentencia Haegeman de 30 de abril de 1974 estableció que las disposiciones del acuerdo forman parte integrante, a partir de su entrada en vigor, del ordenamiento jurídico comunitario39. Podemos en este sentido concluir que el Derecho de la Unión se rige por un sistema de «recepción automática» de estos tratados y, en consecuencia, no sujeta a ningún acto de transformación ni de ejecución. Conviene, no obstante, tener presente que los acuerdos celebrados por la UE se introducen en el Derecho de la Unión a través de un acto jurídico interno (reglamento o decisión). Por el carácter meramente instrumental de estos actos, se entiende que no afectan a la llamada «recepción automática». El artículo 17.1.d) del Reglamento interno del Consejo obliga a la publicación de los acuerdos en el DO incluyendo su fecha de entrada en vigor.

En relación con su jerarquía, varios argumentos parecen avalar la preeminencia de los Tratados constitutivos sobre los tratados celebrados con terceros. Retengamos el que nos parece de mayor trascendencia: la previsión en el artículo 218.11 TFUE de la posible solicitud (por el PE, el Consejo, la Comisión o un Estado miembro) previa a «la entrada en vigor» de un dictamen al TJ para que éste se pronuncie sobre la compatibilidad del tratado en cuestión con los Tratados y, lo que es fundamental, que en caso de que el Tribunal determinase dicha incompatibilidad sería necesaria para su conclusión de una previa revisión de los Tratados constitutivos. Los Dictámenes 1/75, 1/91, 1/92, 2/94, 3/94, 1/00 y 2/00 precisaron el alcance de la disposición. Poco importa, a nuestro juicio, que el argumento explícito del TJ sea el de «evitar las complicaciones» que podrían derivarse de la conclusión de un acuerdo con disposiciones contrarias, si la consecuencia es su inaplicabilidad en el Derecho de la Unión salvo que haya mediado una revisión que hiciera desaparecer la contradicción. Por otra parte, en el Dictamen 3/94, aunque el Tribunal ha excluido la posibilidad de solicitud del dictamen a posteriori de la conclusión del Tratado, ha señalado que los Estados miembros o las Instituciones legitimadas disponen del recurso de anulación contra la decisión del Consejo de concluir el acuerdo e incluso de solicitar medidas cautelares, como se estudia con atención en el Capítulo 2440. Creo que es legítimo concluir sobre la supremacía de las normas constitutivas sin necesidad de asimilarlo en sentido estricto a un «control de constitucionalidad». Con relación al Derecho derivado, más que de una exigencia dogmática es susceptible de deducirse como una exigencia lógica de cumplir las obligaciones contraídas en el plano internacional y evitar las consecuencias de un incumplimiento. La misma norma que establece esta obligatoriedad y una reiterada jurisprudencia (ya aludida) sobre la obligación del respeto a las mismas es suficiente. 

Más problemático resulta la posible atribución de efecto directo a las disposiciones de un tratado de estas características. Integrada la norma convencional en el orden jurídico de la Unión es coherente que cuando presente las características de claridad, precisión e incondicionalidad pueda gozar del «efecto directo». No obstante la cuestión no es pacífica por la particular naturaleza de la norma internacional y su vinculación a reciprocidades internacionales. En una sentencia de 12 de diciembre de 1995, el TJ reiteró la jurisprudencia sentada en el asunto International Fruit Company en el sentido de que la cuestión del efecto directo de las disposiciones de un acuerdo celebrado por la Comunidad con países terceros requiere, invariablemente, un análisis del espíritu, el sistema y la letra de dicho acuerdo41. 

Los «acuerdos mixtos» presentan algunas dificultades particulares por ser instrumentos convencionales internacionales en los que participan conjuntamente la Unión y los Estados miembros por afectar simultáneamente a competencias de una y otros. Es la única solución posible cuando concurren competencias estatales y de la Unión, o cuando así lo exigen las modalidades de financiación del acuerdo en cuestión. El TJUE ha confirmado en su jurisprudencia el recurso a los acuerdos mixtos en esas circunstancias, extendiendo a ellos sus competencias de control en  lo que concierne al ámbito competencial europeo.

3.3. LAS DECISIONES DE LOS REPRESENTANTES DE LOS GOBIERNOS DE
LOS ESTADOS MIEMBROS REUNIDOS EN EL SENO DEL CONSEJO 
Las decisiones de los representantes de los Gobiernos de los Estados miembros reunidos en el seno del Consejo constituyen una categoría especial de actos que se caracterizan por ser adoptados por los miembros del Consejo colectivamente en su calidad de representantes de los Estados miembros. Desde luego, hay que aceptar la enorme dificultad de su caracterización. Existe, no obstante, un cierto grado de acuerdo para considerarlos una categoría mixta entre lo internacional y lo institucional. 

Son, por una parte, acuerdos internacionales sujetos a las normas de Derecho Internacional que determinan los aspectos relativos a la conclusión y a la eficacia de estas decisiones, aspecto que los califica mejor como elementos de instrumentación convencionales. Pero tienen, por otra parte, innegablemente, un aspecto específico de la UE que viene dado por su objeto y fin y por el marco institucional, en el que son adoptadas, no sólo porque ese hecho se produzca en el «seno del Consejo», sino también porque, a veces, se cuenta con la intervención de otras Instituciones en el proceso de su preparación. Su publicación en el DO y el que las sucesivas Actas de Adhesión establezcan que los nuevos Estados miembros se adhieren, por medio de ellas, a las decisiones y  acuerdos adoptados por los representantes de los Gobiernos de los Estados miembros en el seno del Consejo potencian su carácter de Derecho de la Unión. Pero a nuestro juicio, ello no elimina su carácter de instrumentación internacional como, por otra parte, pone de manifiesto la «adhesión» de los Estados nuevos a través del Acta de adhesión. Según la jurisprudencia del TJUE, esas decisiones no se encuentran sujetas al control de legalidad42.








\clearpage
\section{Unión Europea: Sistema Institucional}
\puntosclave{

}

Cuando se fundaron las Comunidades Europeas, los Tratados constitutivos de cada una de ellas contemplaban un sistema institucional propio, incluso con denominaciones y competencias diferentes para sus Instituciones respectivas.

Así pues, el entramado institucional estaba formado en 1957 por tres instituciones independientes de los Estados (la Alta Autoridad de la CECA y las dos Comisiones de la CEE y Euratom), tres Consejos, tres Asambleas parlamentarias y tres Tribunales de Justicia. De la función consultiva socioeconómica se encargaba un Comité consultivo (CECA) y dos Comités Económicos y Sociales (CEE y Euratom). Éstas eran las previsiones de los Tratados en 1951 y 1957.

Sin embargo, carecía de sentido esta multiplicación de instituciones y de órganos auxiliares en cada una de las tres organizaciones internacionales que, aunque persiguen objetivos económicos específicos, constituyen un proceso unitario de integración política y económica progresiva.

Por ello, la unificación de las Asambleas y de los Tribunales de Justicia se llevó a cabo el mismo día en que se firmaban los Tratados CEE y Euratom (25 de marzo de 1957), mediante la firma del «Convenio relativo a ciertas instituciones comunes». Materialmente, pues, nunca hubo tres Asambleas ni tres Tribunales. Además, ese Convenio de 1957 unificó el Comité Económico y Social (CES) de la CEE con el del Euratom, pero no con el Comité Consultivo de la CECA, que por su carácter sectorial mantuvo su singularidad orgánica hasta la extinción del Tratado CECA en julio de 2002.

Esa unificación realizada en 1957 fue operativamente fácil, ya que el Tratado CEE y el Tratado Euratom aún no habían iniciado su andadura legal y, esa unificación antes de su entrada en vigor, evitaba las dificultadas de unificar servicios, jurisprudencia, etc. La Asamblea y el Tribunal de Justicia existentes añadieron a sus competencias las previstas en los dos nuevos Tratados. Y todavía fue más fácil unificar entre sí los dos CES, pues no existía ninguno y un único CES asumió las competencias consultivas previstas en los dos Tratados.

La unificación de las Comisiones y de los Consejos se abordó varios años después en el Tratado de 8 de abril de 1965 que constituyó un Consejo único y una Comisión única de las Comunidades Europeas (también denominado Tratado de Fusión, en vigor desde 1 de julio de 1967; hoy derogado).

La unificación de la Comisión exigió notables esfuerzos: las tres instituciones ya estaban en funcionamiento, los períodos de mandato y el nombramiento de la Alta Autoridad y de los Comisarios CEE y CEEA diferían y sus competencias eran muy dispares. La solución salomónica adoptada al crear la Comisión única fue la unificación orgánica (seis únicas personalidades, nombramiento y período de mandato unificado, organización administrativa única y presupuesto único), pero mantendría las competencias bien diferentes que le otorgaban los tres tratados y las ejercerían en las condiciones y modos previstos en dichos tratados.

La unificación de los Consejos de las tres Comunidades Europeas no fue tan difícil pues debido a su composición no permanente mantuvieron siempre unos servicios administrativos comunes. El marco de unificación del Consejo se completó con la creación de un Comité de Representantes Permanentes (COREPER) como único órgano auxiliar de preparación de los trabajos del Consejo. 

Obsérvese que el sistema seguido fue el de la unificación orgánica o estructural, sin afectar a las competencias que cada Tratado atribuye a dichas instituciones y órganos, ni al modo de funcionamiento, ni a las condiciones de ejercicio de dichas competencias. No hubo, pues, unificación competencial. Las Instituciones unificadas o comunes siguieron ejerciendo los poderes y competencias que les atribuían cada uno de los Tratados y en las condiciones previstas en cada uno de ellos y en el convenio de unificación.

Con mayor razón no debe confundirse la unificación de las instituciones con la unificación de las Comunidades: esas sucesivas unificaciones institucionales no alcanzaron a las Comunidades Europeas mismas, que siguieron siendo tres organizaciones internacionales con sus tres respectivos tratados fundacionales hasta que en julio de 2002 expiró el plazo de vigencia del Tratado CECA. Ahora bien, el preámbulo del Tratado de Fusión de los Ejecutivos de 1965 consideraba la fusión de las tres Comunidades como un objetivo final al que sirve la creación de instituciones comunes. Pero el Tratado de la Unión Europea de 1992 (Maastricht), y su reforma de 1997 (Tratado de Ámsterdam), desaprovecharon la oportunidad de unificar las Comunidades en una sola y tampoco acometieron la unificación competencial. La Declaración número 23 del Tratado de Niza (2001) señalaba que la simplificación y clarificación del sistema de los Tratados era uno de los objetivos de la reforma del 2004. El fracasado Tratado Constitucional fusionaba la Comunidad Europea y la UE y este avance lo heredó el Tratado de Lisboa de 2007, logrando la unificación institucional de modo que la Unión ha sucedido a la Comunidad Europea, si bien aún se mantienen separados el Tratado Euratom y sus reglas materiales propias.

2. COMPOSICIÓN

Desde la entrada en vigor del Tratado de Reforma de Lisboa de 2007, el sistema institucional de la Unión Europea se compone de siete instituciones: Parlamento Europeo, Consejo Europeo, Consejo, Comisión Europea, Tribunal de Justicia, Banco Central Europeo y Tribunal de Cuentas (art. 13 TUE). Se ha pasado de un sistema originariamente cuatripartito a un extenso sistema institucional. Sin embargo, se observa que al margen de las presiones o recelos que han llevado a ese amplio marco institucional, el Tratado de la Unión sólo le dedica un precepto marco o de cabecera con sus roles, funciones y composición en el citado Tratado a cinco instituciones que son las fundamentales en el plano jurídico-político: Parlamento, Consejo Europeo, Consejo, Comisión y Tribunal de Justicia. En relación con las otras dos (BCE y Tribunal de Cuentas) remite plenamente al TFUE para saber qué funciones y papeles asumen. El Parlamento Europeo, el Consejo y la Comisión, en calidad de triángulo de decisión, se ven asistidos por unos órganos auxiliares, el Comité Económico y Social y el Comité de las Regiones, que tienen funciones consultivas. El Consejo dispone, además, de un destacado órgano específico que le auxilia en sus actividades, el Comité de Representantes Permanentes. A su vez, cada Institucióny algunos órganos auxiliares han creado una constelación de comités y grupos de trabajo, muchísimos permanentes, otros ad hoc.

Los Tratados han sido muy rígidos a la hora de distinguir entre unos órganos principales, para los que reservan la denominación de instituciones, y otros de carácter secundario cuya misión es asistir a las Instituciones, a los que habitualmente se les denomina órganos consultivos así como el propio TFUE (art. 300).

Entre esas dos categorías extremas hay otros organismos que ejercen con autonomía sus funciones, sin asistir propiamente a otras Instituciones. Este complicado y rígido escalonamiento de Instituciones y de órganos consultivos añade cierta complejidad al sistema de la Unión al incluir, además, un nivel ambiguo de otros organismos (el Defensor del Pueblo, la OLAF...) y varias decenas de Agencias.

Una Agencia es un organismo regulado por el Derecho de la Unión, dotado de personalidad jurídica propia; se crea mediante una disposición de Derecho derivado con el fin de que realice una labor técnica o científica muy concreta. Responden a varios objetivos: la descentralización administrativa y geográfica y la especialización técnica y científica o integrar diversos grupos de interés.

Hay agencias reguladoras con capacidad de adopción de decisiones individuales, asesoramiento técnico o científico a la Comisión y los Estados miembros y asumen la responsabilidad de actividades operativas: así la Agencia del Medicamento, FRONTEX, la Agencia Europea para la Seguridad y la Salud en el Trabajo, Agencia para el Control de la Pesca, la Agencia de los Derechos Fundamentales, etc.); otras operan en el ámbito de la Política Exterior (Agencia Europea de Defensa), otras en el ámbito de la cooperación judicial (Europol, Eurojust y la Escuela Europea de Policía).

Otras agencias son de carácter ejecutivo (se les delega la ejecución de determinados programas, así la Agencia Ejecutiva de Sanidad y Consumo o la de Investigación).

Sin embargo, el Banco Europeo de Inversiones no es ni una institución ni un órgano de la Unión, sino un establecimiento público internacional vinculado a los fines de la Unión, cuyos funcionarios lo son de la Unión, pero el BEI conserva su personalidad jurídica separada de la UE.

3. CARACTERES DEL SISTEMA INSTITUCIONAL
En los trabajos preparatorios de la reforma que condujo al Tratado de Niza (2001) hubo consenso sobre el modelo institucional. El sistema institucional no fue objeto de reflexión por parte de un grupo de trabajo específico durante el debate del fracasado Tratado constitucional en el seno de la Convención para el Futuro de Europa (2002-2004). No se han presentado alternativas al modelo institucional de la Unión y a sus principios, aceptándose el modelo establecido desde la fundación en los años cincuenta. Lo que importa en todas las reformas es hacer mejoras concretas que lo hagan más democrático, más eficaz y más transparente. En los últimos años preocupaba hacer comprensible ese sistema en su conjunto, expresar y hacer inteligibles los principios implícitos del sistema y describir las tareas y atribuciones de cada Institución: quién hace qué.

Los roles que representan y los principios que encarna cada Institución y organismos de la UE no se han visto alterados. No ha sido desmontada la complejidad del «método comunitario» en el Tratado de Lisboa (2007). Lo que consigue el TUE (Título III, «Disposiciones sobre las instituciones») es describir mejor qué hace cada Institución en el precepto de cabecera correspondiente a cada una de ellas, al tiempo que confirma el método de pesos y contrapesos o de interdependencia entre las tres principales Instituciones que encarnan la función legislativa.

3.1. LA ORIGINALIDAD
Una de las cualidades reconocidas a la Unión, desde la fundación misma de las Comunidades Europeas en los años cincuenta del pasado siglo, ha sido su original estructura de poderes que deriva de la naturaleza de las Comunidades Europeas en el pasado y de la Unión Europea en el presente: son organizaciones internacionales que, aunque participan de evidente inspiración constitucional, no admiten analogías fáciles con una estructura estatal ni tan siquiera federal. No son un Estado federal, por la evidente razón de que no son un Estado ni pretenden serlo.

La Unión Europea es una organización internacional, como lo fueron también las tres Comunidades Europeas por separado en el pasado; debe su creación a tratados internacionales y éstos regulan relaciones internacionales, aunque no de forma exclusiva. Pero no resisten tampoco fáciles comparaciones con el resto de organizaciones internacionales. Por ello, resulta difícil encajar a la Unión Europea en una organización internacional clásica —fundada en la cooperación—, por la razón evidente de que éstas no echan sus raíces en la soberanía popular y no absorben podereslegislativos, reglamentarios y judiciales de sus Estados miembros. La Unión es una organización de integración de pueblos y Estados. Asienta su fundamento en la legitimidad internacional o intergubernamental y democrática y esa doble legitimidad se traduce también en su estructura institucional.

En efecto, en las organizaciones internacionales clásicas o de cooperación los órganos están marcados por la presencia de los Estados y por la representación de intereses estatales (a excepción de los órganos jurisdiccionales internacionales y de las secretarías generales o de la peculiar composición tripartita de la OIT). La diferencia estriba en que la Unión Europea no es una organización internacional de cooperación, sino de integración basada en la atribución del ejercicio de una parte de la soberanía de cada Estado miembro.

En la Unión sólo dos de sus siete instituciones, el Consejo Europeo y el Consejo, acogen la presencia legítima y la defensa de intereses nacionales, encarnando el principio de la representación de los Estados que forman la UE.

El resto de Instituciones y otros órganos auxiliares —salvo el específico órgano auxiliar del Consejo, el Coreper— representan y defienden intereses distintos; independientemente de la forma de designación o de elección, tienen en común la obligación de ser independientes de los Gobiernos y no defender intereses nacionales.

En efecto, la Comisión asume la gestión centralizada de los asuntos comunes para lo que debe defender intereses generales con independencia respecto de los Gobiernos y de los particulares. El Parlamento Europeo encarna el principio democrático, la representación de los ciudadanos agrupados por sus concepciones políticas. El Tribunal de Justicia representa el interés social en el respeto del Derecho. El Tribunal de Cuentas representa la responsabilidad financiera de los poderes públicos en una sociedad democrática. El mismo Banco Central Europeo asume la estabilidad de los precios y la independencia del emisor de la moneda. Lo mismo sucede si analizamos la representación de intereses en los órganos auxiliares; sólo en el Comité de Representantes Permanentes (COREPER) hay presencia legítima de intereses estatales a diferencia del Comité de las Regiones o del Comité Económico y Social que representan los intereses de las regiones y de los medios socio-profesionales, respectivamente.

En definitiva, únicamente en la esfera del Consejo Europeo (Jefes de Estado o de Gobierno) y del Consejo (ministros o asimilados), incluido su órgano auxiliar, el COREPER, tiene lugar la representación y defensa de los intereses particulares de los Estados. Precisamente, los miembros del Consejo Europeo y del Consejo lo son por reunir la cualidad de miembros de los Gobiernos, a los cuales deben representar; no son independientes de los Estados, sino que representan el interés nacional, la razón de Estado.

3.2. EL «MÉTODO COMUNITARIO» DE LA UNIÓN
No se concibe un acto de la Unión sin la colaboración insustituible de la Comisión —que hace la propuesta— con el Parlamento y el Consejo —que deciden los actos legislativos ordinarios (arts. 289 y 294 TFUE)—. Sin el diálogo triangular («triálogo») entre la Comisión, el Consejo y el Parlamento Europeo no hay actos legislativos ordinarios. Lo mismo podría decirse de la colaboración en el procedimiento presupuestario y de bastantes actos de procedimiento legislativo especial y de los actos no legislativos.

Los «padres fundadores» del proceso de integración trataron de evitar cualquier connotación que hiciera aparecer a las entonces Comunidades Europeas como emulaciones estatales y querían superar cualquier reticencia que hiciera fracasar un proyecto, entonces y hoy, indispensable para la paz en Europa. Este sistema de pesos y contrapesos en el que se basa el «método comunitario» se caracteriza por poner en juego muchos intereses y factores; genera una dinámica compleja que exige negociar y dialogar entre las tres instituciones decisorias; no deja en manos de una persona ni tan siquiera de una Institución la responsabilidad de la decisión, diluyendo y repartiendo el poder entre varias Instituciones, diseminándolo en el sistema intrincado y peculiar de checks and balances que ha caracterizado al sistema europeo. Es un sistema complejo y distinto del interno en nuestras Constituciones y si cabe más próximo, en la concepción del poder y la responsabilidad, al sistema americano (contrapesos y coparticipación de varios poderes) que al tradicional sistema político interno en Europa. 

Sin embargo, siendo extraño a las tradiciones constitucionales europeas, ha evitado crear un poder definido así como competir con los modelos nacionales. Para la opinión pública es un poder ambiguo y difuso, aunque con una gran capacidad de impacto social y jurídico. También es cierto que su carácter difuso ha facilitado enormemente la integración al exigir negociar todo y tener en cuenta intereses diversos, y no imponer de forma lineal una opción; ninguna Institución partícipe del proceso de decisión legislativo es irrelevante; estas instituciones representan diversos contrapesos (de ahí el respeto por los equilibrios institucionales), se necesitan materialmente entre sí para alcanzar cualquier decisión y están obligadas a colaborar (art. 295 TFUE).

Sin mencionar al «método comunitario», la reforma de Lisboa mantiene su idea de síntesis implícita: la presencia y dialéctica de diversos intereses para identificar el interés general y la interrelación entre las instituciones. En definitiva, las competencias de la Unión se seguirán ejerciendo según un modo comunitario, lo que proclama una clara intención de continuidad histórica entre la «nueva» Unión de la reforma de Lisboa y la «antigua» fundada en las Comunidades Europeas y en los pilares intergubernamentales. 

3.3. LA DIALÉCTICA «INTERESES COMUNES-INTERESES NACIONALES» 
Aparentemente hay una mayor representación de los intereses generales de la Unión en el entramado institucional; pero no es así en la realidad. En efecto, el Consejo ha sido y sigue siendo en cierta medida el centro del poder de decisión en el sistema institucional, a pesar de los notables avances del Parlamento Europeo que ha logrado compartir con el Consejo la función legislativa y presupuestaria. Su peso político, con el valor añadido del Consejo Europeo como ejecutivo intergubernamental, y su poder normativo como legislador ha superado al resto del conjunto institucional. Especialmente, en los últimos quince años, y de forma paralela al desprestigio de la Comisión, es notorio el peso del Consejo Europeo y del Consejo no sólo en las grandes decisiones de orientación, sino en la definición de las políticas. Y, por otra parte, algunas instituciones encargadas de velar por los intereses generales, como la Comisión y, en menor medida, el Parlamento Europeo, son en la práctica diaria y, siempre, en los momentos decisivos muy permeables a los intereses gubernamentales.

Por ello, aunque la Unión se aleja del modelo formal de las organizaciones internacionales, el método de cooperación interestatal y la negociación continuada han estado y están presentes en la vida diaria de la Unión. De igual modo, las estructuras de integración siguen el rumbo marcado por los Gobiernos nacionales. No ha sido una diplomacia rígida o formal, sino una diplomacia dúctil y familiarizada con los asuntos comunes  que ha controlado el poder de decisión en los múltiples niveles en que se ha desdoblado institucionalmente (grupos de trabajo, COREPER, Consejo, Consejo Europeo).

En todo caso, es cierto que las influencias intergubernamentales
en el proceso de decisión no llegan a anular la voluntad propia de
las Instituciones. Las instituciones gozan de autonomía una vez
formada su voluntad.
El proceso de integración y su sistema institucional es original,
complejo y lleno de contradicciones y ambigüedades. Hay
inspiración y esencia federal en los moldes organizativos
internacionales; hay cooperación intergubernamental en el método
de integración. Lo «comunitario» y lo intergubernamental no son
compartimentos estancos.
3.4. DISTRIBUCIÓN DE FUNCIONES ENTRE LAS INSTITUCIONES
La analogía con los sistemas estatales de división de poderes es casi imposible. Es cierto que el viejo principio de la división de poderes tiene plasmaciones constitucionales muy diversas y que se ha exagerado una división rigurosa. Frente a la división de los tres poderes clásicos del Estado, hay siete Instituciones de las cuales al menos cuatro participan y concurren en el ejercicio de los poderes. En efecto, cualitativamente no se puede hacer una comparación entre la división de poderes, bastante precisa en las Constituciones democráticas, y la distribución y participación de varias instituciones en una misma función en la Unión Europea. Hay, como han señalado algunos autores (K. Lenaerts: 1991), una separación funcional, pero no una separación orgánica que permitiera identificar una institución con una forma determinada de poder.

De entrada, en la Unión Europea las denominaciones de las instituciones no se corresponden con las funciones que asumen a la luz de las tradiciones estatales. Desde luego, el sistema institucional de la Unión es muy complejo, como consecuencia de un proceso que integra no sólo Estados, sino también pueblos. 

Así, el Consejo de la Unión, formado por miembros de los Gobiernos, es ante todo el colegislador de la Unión de incuestionable legitimidad, no fundada en la representatividad democrática directa, sino en la representatividad internacional o territorial, al estilo de las segundas cámaras federales. En efecto, desde el Tratado de Maastricht (1992), más las posteriores reformas de los Tratados de Ámsterdam (1997), Niza (2001) y Lisboa (2007), el Consejo comparte en pie de igualdad la función legislativa con el Parlamento Europeo, así como la adopción del presupuesto. El Parlamento Europeo, además de las compartidas funciones legislativas y presupuestarias, goza de una competencia consultiva amplia, pero no siempre muy influyente. Sin embargo, Parlamento Europeo y Consejo apenas poseen iniciativa legislativa, que está depositada de forma generosa en la Comisión —como Institución independiente— pero cada vez más compartida con los Estados miembros en materias como el espacio de libertad, seguridad y justicia y la política exterior.

Pero, además de colegislador, el Consejo es el titular de amplios poderes ejecutivos, lo que significa que adopta normas generales o de base y puede adoptar en casos justificados las de ejecución o atribuir la ejecución a la Comisión bajo ciertas modalidades de seguimiento o control. Así pues, el poder ejecutivo en la Unión se reparte entre el Consejo y la Comisión, además de concretos y relevantes poderes del Consejo Europeo y del Banco Central Europeo.

Tan sólo el Tribunal de Justicia y el Tribunal General asumen las funciones propiamente jurisdiccionales equiparables a las que ejercen los altos órganos judiciales en un Estado, si bien el Tribunal de Justicia es una instancia suprema con una diversidad de competencias.

En definitiva, en el sistema de la Unión también encontramos los tres poderes clásicos; lo que sucede es que en el ejercicio de cada uno de los poderes pueden concurrir varias Instituciones. Visto desde otra perspectiva, una institución de la Unión puede ejercer poderes de muy distinta naturaleza. No se puede asociar un poder ni tan siquiera una función con una concreta Institución. Varias Instituciones pueden llegar a concurrir en una misma función. Esto es un problema de cara a la opinión pública y los medios de comunicación pues no siempre comprenden bien quién hace qué y cómo.

Este complejo tejido de interrelaciones podría subsumirse en dos poderes básicos: poderes de decisión (asumidos por la Comisión, el Parlamento, el Consejo, el Consejo Europeo y el Banco Central Europeo) y poderes de control (asumidos por la Comisión, el Parlamento, el Banco Central Europeo, el Tribunal de Justicia y el  Tribunal de Cuentas).

4. EL PRINCIPIO DEL EQUILIBRIO
INSTITUCIONAL
Una piedra clave del sistema institucional de la Unión es el principio de equilibrio institucional. El sentido de este principio fundamental reside en el ejercicio por cada institución de las competencias que le han sido atribuidas en los tratados en las modalidades y condiciones previstas por los mismos (art. 13.2 TUE). Las instituciones gozan de poderes determinados que ejercen autónomamente; poderes protegidos, a su vez, frente a invasiones de otras instituciones u órganos. Esa autonomía, sin embargo, no les permite desentenderse de las tareas de las otras y deben colaborar entre sí. La doble consecuencia que se extrae es que: — en primer lugar, los poderes atribuidos a una institución no pueden ser ejercidos por otra institución o por órganos auxiliares (con la única excepción de la delegación de poderes para adoptar actos no legislativos —art. 290 TFUE— y de ejecución, prevista en el art. 291.2 TFUE); — en segundo lugar, que los poderes atribuidos a las instituciones tienen que ser ejercidos por éstas en los límites y condiciones impuestos en los tratados.

A pesar de que los poderes de las Instituciones aparecen descompensados (piénsese en el inmenso poder del Consejo y del Consejo Europeo), con funciones entremezcladas y dispersas, los Tratados ponen un gran cuidado en la reconducción de estas situaciones hacia un diálogo y colaboración entre las Instituciones. El diálogo y la colaboración entre las tres instituciones que participan en el proceso de decisión es esencial; lo prevén y lo exigen los tratados desde su fundación (art. 295 TFUE): el Parlamento Europeo, la Comisión y el Consejo se consultarán mutuamente y organizarán de común acuerdo las modalidades de esa colaboración.

El equilibrio institucional, basado en la colaboración, no ha sido el mismo a lo largo de estos años; ha sido un equilibrio dinámico, tan evolutivo como la diversidad misma de la colaboración y diálogo emprendido entre las instituciones. También el tiempo ha demostrado que el principio del equilibrio institucional no es algo puramente formal o abstracto, no es algo sin consecuencias jurídicas prácticas; no sólo es un regulador de las relaciones entre las Instituciones, sino que afecta a los particulares.

En efecto, el equilibrio institucional, como señaló el Tribunal de Justicia en la sentencia Meroni, es una «característica de la estructura institucional» y «una garantía fundamental acordada por el tratado, especialmente a las empresas y asociaciones de empresas a las que se aplica»1. El equilibrio institucional significa tanto el respeto a las relaciones entre las instituciones como el respeto al ejercicio de sus poderes respectivos2.

Cuando una Institución debe intervenir en el proceso de decisión, el Tribunal ha estimado que la institución llamada a decidir no puede prescindir de la opinión de otras instituciones ni fundarse en sus propias prácticas para despojar a otra institución de sus competencias3. El Tribunal de Justicia afirmó refiriéndose al derecho del PE a ser consultado que su opinión consultiva (a fortiori, su participación en igualdad de condiciones con el Consejo en el procedimiento legislativo ordinario)4.

El equilibrio institucional ha sido el argumento central de importantísimas sentencias del Tribunal de Justicia con marcada incidencia constitucional. Así, el razonamiento del Tribunal, en la sentencia Chernobyl de 22 de mayo de 1990, se centró en el principio del equilibrio institucional para justificar la legitimación directa del Parlamento Europeo en el recurso de anulación (que el antiguo art. 173 TCEE —hoy, art. 263 TFUE— no contemplaba) a fin de proteger sus competencias o prerrogativas. El Tribunal declaró que El respeto del equilibrio institucional implica que cada una de las Instituciones ha de ejercer sus competencias sin invadir las de los demás. Exige también que cualquier incumplimiento de esta regla, caso de que se produzca, pueda ser sancionado. Este Tribunal, que en virtud de los Tratados ha de velar por el respeto del Derecho en su interpretación y aplicación, debe igualmente poder garantizar el mantenimiento del equilibrio institucional, y, por consiguiente, el control jurisdiccional del respeto de las prerrogativas del Parlamento […] por lo que […] una laguna procesal no puede prevalecer contra el interés fundamental en que se mantenga y respete el equilibrio institucional definido por los Tratados […]5. 

En el asunto Chernobyl, el Tribunal tuvo en cuenta el status de cada Institución y la protección efectiva de sus derechos o atribuciones, aún en el silencio del Tratado: le incumbe a este Tribunal garantizar la plena aplicación de las disposiciones de los Tratados relativas al equilibrio institucional […], en virtud del cual […] no pueden menoscabarse las prerrogativas del Parlamento Europeo […] (apdos. 24 y 25).

El Tribunal no tiene el poder de modificar el equilibrio institucional (atribuido a los Estados) pero el Tribunal acepta perfeccionar ese equilibrio mediante un recurso directo de toda Institución contra los actos que atenten contra sus prerrogativas.

En otras ocasiones el Tribunal de Justicia ha hecho respetar el reparto de competencias entre las Instituciones en relación con la base jurídica de los actos y el proceso legislativo: aunque sin mencionar este principio, el Tribunal ha estimado que una controversia sobre el fundamento jurídico no tiene un alcance meramente formal6. Cuando se utilizan dos bases jurídicas que no confieran las mismas prerrogativas a las Instituciones que participan en la adopción del acto, se debe recurrir al procedimiento que garantice la mayor participación del Parlamento (el legislativo ordinario)7.

En definitiva, el equilibrio institucional es interdependiente y multipolar; es el equilibrio establecido por la voluntad de los Estados, que evoluciona con los logros de la integración misma y es tan dinámico como el proceso mismo, y al surgir sus raíces de las reglas de los Tratados es controlado y protegido por el propio Tribunal de Justicia. En todo caso, del entendimiento y diálogo entre intereses diversos dependerá un progreso equilibrado y aceptable para los Estados y la ciudadanía de la Unión Europea. Por ello, esa dialéctica o tensión entre las Instituciones ha sido considerada como un signo de vitalidad y fuente de progreso. 

El respeto al equilibrio institucional obliga también a los Estados miembros; el Tribuna declara que «procede recordar que las normas relativas a la formación de la voluntad de las instituciones de la Unión están establecidas en los Tratados y, por tanto, no tienen carácter dispositivo ni para los Estados miembros ni para las propias»8.



El Tratado de Lisboa otorga el rango de institución europea al Consejo Europeo y al Banco Central Europeo elevando el número de instituciones a siete. Se trata de un grupo heterogéneo de instituciones dotadas de diferentes competencias para lograr los objetivos de la Unión. en el que  no es posible distinguir claramente el esquema de división tradicional de poderes.

\subsection{Consejo Europeo}
De entre las grandes modificaciones que introdujo el Tratado de Lisboa mereció especial atención el cambio de naturaleza del Consejo Europeo. Por primera vez en la larga historia de esta instancia decisoria pasó a ser considerada una Institución de la Unión. El artículo 13 TUE que enuncia las Instituciones de la Unión Europea así lo contempla. El artículo 15 TUE que contiene su regulación esencial dice que El Consejo Europeo estará compuesto por los Jefes de Estado o de Gobierno de los Estados miembros, así como por su Presidente y por el presidente de la Comisión […] Precisa asimismo que Participará en sus trabajos el Alto Representante de la Unión para Asuntos Exteriores y Política de Seguridad.

Su estatuto expresa en lo esencial una cierta continuidad con la regulación anterior y, en cierto sentido, su institucionalización formal en el Tratado de Lisboa no ha despejado las dudas que pesan sobre su confusa naturaleza, producto de una larga y no siempre pacífica evolución. El TUE describe la función de esta Institución señalando  que El Consejo Europeo dará a la Unión los impulsos necesarios para su desarrollo y definirá sus orientaciones y prioridades políticas generales. Y así, alimentando la idea de que se trata de una Institución de impulsión política que no se ha de mezclar en la actividad jurídica ordinaria del sistema de la UE, añade de manera contundente que «No ejercerá función legislativa alguna». Junto a ello, unas escuetas reglas de funcionamiento relativas a la asistencia, en su caso, a sus componentes de los Ministros de los Estados miembros y de un miembro de la Comisión, a la periodicidad de sus reuniones («dos veces por semestre» y la posibilidad de celebrar reuniones extraordinarias) y a la regla general de decisión por «consenso» perfilan los principales rasgos institucionales del Consejo Europeo enunciados en el artículo 15 TUE. Todo ello parece responder a la tradicional contención del Consejo Europeo en un mundo aparte para que no interfiera en el equilibrio esencial del sistema Parlamento, Consejo y Comisión.

Pero junto a esa continuidad esencial, la institucionalización formal del Consejo Europeo, la importancia de la figura del Presidente del Consejo Europeo y la asunción de amplias funciones de especial relevancia dentro del sistema han convertido a esta Institución en un verdadero directorio político y jurídico de la Unión Europea que, lentamente, va cambiando —si no lo ha hecho ya— la naturaleza misma del modelo europeo. Pese a ello, y a pesar incluso de la más precisa regulación del funcionamiento del Consejo Europeo que ha acompañado a esta gran reforma, resultará imposible, al menos durante un período largo de tiempo, identificar con claridad el verdadero significado político y jurídico de esta Institución y de su capacidad de condicionar el modelo institucional europeo. Bien es cierto que su desenvolvimiento en la práctica deestos primeros años nos desvela ya algunas claves de su nuevanaturaleza. La gestión de la crisis económica ha dibujado un papel del Consejo Europeo extraordinario por su importancia, inquietante desde la perspectiva del equilibrio institucional de la Unión por su efecto potenciador de los elementos más intergubernamentales del sistema y desconcertante por la ambigüedad de su actuación no siempre ortodoxa con el marco de su regulación en los Tratados y por la ausencia de una regulación clara de los límites de sus funciones o competencias. Sea como fuere sigue siendo absolutamente cierto que comprender esta compleja instancia de poder de la UE y su peculiar configuración jurídica sólo es posible si se analiza el Consejo Europeo en una doble dimensión histórica y  política. Seguiremos este camino.



[Desarrollo histórico de la formación del Consejo Europeo como Institución... (Considero que no es importante)]

el Tratado de Lisboa puso fin a la cuestión del carácter institucional del Consejo Europeo al formalizarlo como una Institución de la Unión (art. 13 TUE). Como consecuencia de ello, el 1 de diciembre de 2009 el propio Consejo Europeo aprobó la Decisión por la que adoptó su Reglamento interno en virtud de lo previsto en el artículo 235.3 TFUE2. Más allá, sin embargo, de esta formalización del carácter institucional del Consejo Europeo siguen planteándose algunas dudas sobre su significado. Aparentemente el Consejo Europeo no pareció haber sufrido en los nuevos tratados más cambios que los que se deducían de su «institucionalización» formal (la real ya venía funcionando hacía tiempo) y el de convertirse en el soporte de la nueva figura de la Presidencia del propio Consejo Europeo.


Sin embargo, en mi opinión, el nuevo encaje del Consejo Europeo en el sistema institucional de la UE ha tenido un significado y alcance más profundos y de consecuencias más amplias de lo que su mero enunciado formal traducía. Y es que lo exiguo de su regulación formal no expresa en toda su dimensión el verdadero alcance de su institucionalización ni la muy considerable ampliación de poderes que la realidad ha puesto después de manifiesto. Naturalmente esta realidad ha dibujado, guste o no, la consolidación de un espacio decisorio que profundiza los aspectos más intergubernamentales del sistema. Algunos cambios son meramente formales y no han alterado sustancialmente nada como el que puso fin a la ficción del «Consejo en su formación de Jefes de Estado y de Gobierno» asumiendo ahora el Consejo Europeo sus funciones (arts. 7.2 TUE y 142.2 TFUE). Pero otros encerraban, más allá de cierta contención reguladora e intentos de aparentar una cierta limitación del potencial invasivo del poder del Consejo Europeo, elanuncio de un verdadero cambio de naturaleza que afecta alsistema institucional completo. Así la sujeción establecida en el artículo 15.1 TUE por la que el Consejo Europeo se limita a dar los «impulsos necesarios» y a definir las «orientaciones y prioridades políticas generales», excluyendo su participación en función legislativa alguna, alimentan la idea de que el Consejo Europeo no afectará a los equilibrios tradicionales del sistema europeo, en particular los «comunitarios». Pero pese a este prudente enunciado general de su función, la regulación concreta de sus poderes y funciones previstas en el TUE y TFUE, y, eventualmente, en los Protocolos anexos, es de tal entidad que el anunciado confinamiento del Consejo Europeo en la actividad meramente impulsora no se corresponde con la realidad del papel asumido por el Consejo Europeo. Por último, tampoco parece que sea cierta la pretendida globalidad de la acción del Consejo Europeo. Es de todo punto evidente que, en la medida en que el Tratado de Lisboa se supone que habría hecho desaparecer la división en pilares, la globalidad de la acción del Consejo Europeo resultaría su consecuencia más lógica y hasta creíble. Pero claro está que esto sólo ha ocurrido en la medida en que esa desaparición de los pilares ha sido real, lo que no siempre es obvio más allá de lo meramente formal. Así que no está nada claro que el Consejo Europeo ocupe una «posición institucional» equiparable en todos los ámbitos de acción de la UE. En lo que hace a la Política Exterior y de Seguridad Común (PESC),incluida la Política Común de Seguridad y Defensa (PCSD) cuyaregulación sustancial permanece en el TUE con una cláusula de no contaminación de los procedimientos PESC (respecto de otras competencias de la UE) y viceversa (art. 40 TUE) la posición del Consejo Europeo es distinta. Este es un ámbito (arts. 23 a 46 TUE) en el que se concentran las más importantes y concretas funciones del Consejo Europeo, marcando una diferencia cierta con su función en los ámbitos que vienen del antiguo sistema comunitario (competencias de los arts. 3 a 6 TFUE). Es, a todas luces, evidente una mayor capacidad de condicionamiento del desarrollo de la Unión por el Consejo Europeo en PESC/PCSD, donde la pervivencia de procedimientos intergubernamentales y el hecho de que buena parte de la ejecución recaiga sobre los EEMM, favorece la fuerza de una Institución compuesta por la máxima representación de los Estados. Incluso en el contexto del Espacio de Libertad, Seguridad y Justicia donde la supresión del pilar es más evidente pero está muy condicionada por excepciones temporales y singularidades procesales el papel del Consejo Europeo es más incisivo que en el ámbito tradicionalmente comunitario como pueden ser las materias del mercado interior.

Sin embargo, la experiencia desde 2010 del Consejo Europeo ha puesto de manifiesto una «deriva» del mismo hacia una imprevista función de «gobierno» de la UE que algunos justifican en la crisis, que otros entienden como un verdadero atentado contra el «método comunitario» y hasta quienes lo ven como la expresión de un poder emergente fuera de todo marco de legitimidad.

Uno de los grandes conocedores de esta institución, Philippe de Schouetheete, dice que «la evolución del Consejo Europeo en el curso de estos dos últimos años —se refiere a 2010 y 2011— ha estado marcada más por la gestión de la crisis del euro que por las innovaciones del tratado» de Lisboa (2012: 18). Difícilmente rebatible esta apreciación a la luz de las Conclusiones de las diversas reuniones (de toda naturaleza) del Consejo Europeo. Cita este autor la propia aceptación de estas «circunstancias excepcionales» derivadas de la crisis para justificar la poco ortodoxa actuación del Consejo europeo en medio de la crisis por quien fuera su primer Presidente van Rompuy en la Universidad Humboldt: «en tiempos de crisis alcanzamos los objetivos dentro de los límites de las instituciones fundadas en la atribución de competencias. Cuando abordamos nuevas tierras y hay que elaborar nuevas reglas el Consejo Europeo está bien emplazado para aportar su contribución». Y recuerda que es una de las razones que avalaron su creación en los años setenta.


\subsubsection{Composición y funcionamiento}
Como hemos resaltado en la introducción, toda la estructura organizativa y funcional del Consejo Europeo es deudora de las razones que presiden su origen y la particularidad de su naturaleza jurídica. Su composición y funcionamiento son clara expresión de ello. Aparte de las normas previstas en el artículo 15 TUE y de su propio Reglamento interno, el Consejo Europeo tiene una posición muy especial dentro del sistema de la Unión que explica muchas de sus peculiares características que derivan de su misma composición. En efecto, según el artículo 15 TUE, el Consejo Europeo está compuesto por los Jefes de Estado o de Gobierno, así como por su Presidente y por el Presidente de la Comisión.

3.1.1. Composición

La composición por los Jefes de Estado o de Gobierno viene a expresar una clara intencionalidad política: que quienes lo componen sean las personas que tienen el más alto grado de representatividad y de responsabilidad política de cada Estado miembro. De ahí que la composición deba entenderse de una manera flexible que permita encajar las diversas situaciones constitucionales de los Estados miembros. En realidad, todos los Estados miembros están representados por sus Jefes de Gobierno salvo los casos francés (Presidente de la República) y chipriota (Presidente de Chipre), en que están representados por sus Jefes de Estado. El caso francés, tal vez el más complejo, se justifica por razones constitucionales habida cuenta de las altas funciones de gobierno del Jefe del Estado francés, en particular, en lo que toca a la política exterior. Incluso justifica que en los períodos de «cohabitación» en Francia la representación en el Consejo Europeo haya sido bicéfala (Presidente de la República y Jefe del Gobierno), pero no parece necesario recurrir, como a veces se ha pretendido, a la «legitimidad democrática» de la representación en el Consejo Europeo. La escueta regla sobre su composición recoge una realidad política previa que ni define ni condiciona. No obstante, el artículo 10 TUE ha querido terciar en estos problemas de la «legitimidad democrática» del sistema y, junto a la representación de los ciudadanos por el Parlamento Europeo, afirma en su párrafo 2 que Los Estados miembros estarán representados en el Consejo Europeo por su Jefe de Estado o de Gobierno y en el Consejo por sus Gobiernos, que serán democráticamente responsables, bien ante sus Parlamentos nacionales, bien ante sus ciudadanos.

El Presidente de la Comisión forma parte, asimismo, del Consejo Europeo. El artículo 15 TUE, que recoge la misma fórmula con la que el AUE formalizó una realidad aceptada desde la Conferencia de París de 1974, es contundente. Es la expresión de una seria conquista de la Comisión, que durante tiempo reivindicó no sólo su presencia, sino un papel relevante en el Consejo Europeo. La globalización de funciones del Consejo Europeo y la consiguiente imbricación de esta Institución en el sistema europeo ha consolidado su carácter de miembro del Consejo Europeo y aquilatado sus funciones en el mismo. El Tratado de Lisboa se ha limitado a formalizar esa evolución y su consecuencia. La disparidad real de su función con respecto a los altos representantes de los Estados miembros no presenta un problema grave en la medida en que el Consejo Europeo suele adoptar sus actos por consenso, como veremos, pero es evidente que cuando recurre a una votación la exclusión del Presidente de la Comisión (y del Presidente del propio Consejo Europeo) expresamente prevista en el artículo 235.1 TFUE (y 4 del Reglamento interno) acentúa el carácter interestatal de la institución y diluye, en consecuencia, el papel de la Comisión. 

Junto a esta composición «ordinaria» del Consejo Europeo, el propio Consejo Europeo, cuando «el orden del día lo exija» permite a los miembros contar con la asistencia de un ministro y al Presidente de la Comisión con un miembro de la Comisión (art. 15.3 TUE). También si el Consejo Europeo lo estima conveniente podrá invitar al Presidente del Parlamento Europeo que, salvo decisión  contraria del propio Consejo Europeo, procederá al «intercambio de pareceres» al principio de la reunión (art. 4.2 del Reglamento interno). De manera «excepcional», y previa unanimidad del Consejo Europeo, podrán celebrarse al margen de la reunión del Consejo Europeo, «encuentros» con «representantes de terceros Estados, organizaciones internacionales u otras personalidades «a iniciativa del presidente del Consejo Europeo». Da la impresión esta norma del Reglamento interno del Consejo Europeo (art. 4.2) de tratar de disciplinar y hacer más excepcional una práctica algo perturbadora de la normalidad del Consejo Europeo que había tenido cierto éxito en etapas anteriores. Sin embargo, esta práctica se ha reeditado como con la Reunión de los miembros del Consejo Europeo con el Jefe del Estado de Turquía el 18 de marzo de 2016, formando parte del Consejo Europeo de 17-18 de marzo de 2016 en el que se adoptó la Declaración UE-Turquía de 18 de marzo de 2016 sobre la crisis migratoria…

3.1.2. Reuniones del Consejo Europeo

En cuanto al funcionamiento podría decirse que, históricamente, éste se ha distinguido por una acusada flexibilidad e incluso que su actividad se ha desarrollado tradicionalmente en una atmósfera de informalidad generalizada. La falta de rigidez y de excesivas formalidades ha sido parte fundamental de la esencia del Consejo Europeo y explican en buena medida su éxito. El Tratado de Lisboa y el Reglamento interno podrían haber alterado esta flexibilidad, pero da la impresión de que las nuevas reglas del Consejo Europeo no han afectado a esa esencial flexibilidad que exige su propia naturaleza.

El artículo 15 TUE prevé que el Consejo Europeo se reunirá «dos veces por semestre» convocado por su Presidente. Son las reuniones ordinarias. Para las reuniones ordinarias, la Presidencia del Consejo Europeo un año antes del comienzo de un «semestre», «en estrecha cooperación con el Estado miembro que vaya a ejercer la Presidencia rotatoria del Consejo» establecerá las fechas, para dos sesiones lógicamente. Será después, el Consejo de Asuntos Generales el que, de conformidad con la Decisión que establece su Reglamento interno, prepare las reuniones y garantice su actuación subsiguiente «en contacto con el Presidente del Consejo Europeo y la Comisión».

El artículo 15 TUE prevé, asimismo, que el Consejo Europeo podrá reunirse de manera extraordinaria por convocatoria de su Presidente cuando «la situación lo exija».

En realidad, en relación a las reuniones del Consejo Europeo, el Tratado de Lisboa se ha limitado a trasladar una práctica que en los últimos años había elevado el número de reuniones ordinarias previsto en la regulación anterior (dos veces al año) hasta cuatro y dejado abierta la posibilidad de convocatorias extraordinarias. Una práctica que el propio Consejo Europeo de Sevilla había ya previsto en sus «Normas para la organización de los trabajos del Consejo Europeo», recogidas por la Secretaría General del Consejo junto a la modificación del Reglamento interno del Consejo de 20063 que se añadieron a la modificación del Reglamento del Consejo.

Se simplifica, pues, en principio, el modelo reducido a las ordinarias «dos por semestre» y a las extraordinarias que la situación aconseje. No hay motivo para pensar que haya ninguna diferencia jurídica entre unas y otras ya que el Tratado no las diferencia.

La formalización del sistema hacía presagiar la desaparición de las reuniones «informales» y, desde luego, las «especiales» como la de 19 de noviembre de 2009 en la que se nombró al Presidente del Consejo Europeo y al Alto Representante de la Unión para Asuntos Exteriores y Política de Seguridad ante la inmediata entrada en vigor del Tratado de Lisboa, pero la práctica posterior ha refrendado las reuniones informales. Como vengo insistiendo, una cierta «flexibilidad» parece irrenunciable para el Consejo Europeo. Así, pese a todo, se sigue la práctica de las reuniones informales del Consejo Europeo, que se suele matizar —no en todos los idiomas— con que son reuniones informales de los «Jefes de Estado y de Gobierno» (11 de febrero de 2010 sobre cambio climático) o de los «Miembros del Consejo Europeo» (30 de enero de 2012 sobre empleo y crecimiento) o, más confusamente aun, con fórmulas tan peculiares como las que siguieron a la «cena informal de los Jefes de Estado y de Gobierno de 7 de mayo de 2014 recogidas como «observaciones» del Presidente del Consejo Europeo sobre los resultados de las últimas elecciones al PE o recogidas como «Declaración de los Jefes de Estado o de Gobierno sobre Ucrania» en un documento del Consejo. Fórmulas, en fin, no cerradas que tratan de encauzar la actividad más libre y, eventualmente, urgente que reclama esta alta Institución. Caso excepcional es el del Consejo Europeo del 28 de junio de 2016 al que asistieron los 28 Jefes de Estado y de Gobierno en el que el Premier británico presentó el resultado del referendo relativo a la retirada del Reino Unido de la UE (Brexit) y que continuó al día siguiente como «reunión informal de los 27 Jefes de Estado o de Gobierno» para debatir las implicaciones políticas y prácticas del «Brexit» y el debate sobre el futuro de la Unión Europea de 27 Estados miembros. Lo dicho: la flexibilidad es la característica fundamental de esta instancia de poder y gobierno político de la UE.

Por último, reténgase que los Jefes de Estado o de Gobierno de los países que tienen como moneda el euro, se reúnen en las llamadas «Cumbres del Euro». Quede claro que, en ningún caso, es el Consejo Europeo que, como hemos visto, junto a otros factores de su composición contempla a la totalidad de los Estados Miembros de la UE. Pero la tendencia tras la crisis del euro ha sido la de mantener lógicamente separadas ambas formaciones de Jefes de Estado y de Gobierno hasta el punto de iniciar, al margen de los Tratados una cierta «institucionalización» paralela de las «Cumbres del Euro». Se inició esta «formalización» con la adopción de algunas reglas en la Declaración de la Cumbre de los Jefes de Estado o de Gobierno de la zona del euro de 26 de octubre de 20114 y se ha perfeccionado en el Tratado de Estabilidad, Coordinación y Gobernanza en la Unión Económica y Monetaria de 2 de marzo de 20125 y las Conclusiones de los Consejos Europeos de 18 y 19 de octubre de 2012 y de 13 y 14 de diciembre de 20126. Estas normas prevén la elección del Presidente de las Cumbres del Euro en el mismo momento que la del Consejo Europeo, lo que permite, como es el caso, que se elija al mismo. Esta es la situación actual en la que el Presidente de las Cumbres del euro es el mismo que el del Consejo Europeo. Las reglas están muy inspiradas en las del Consejo Europeo si bien la existencia del «Eurogrupo» (con un Presidente elegido por dos años y medio7) hace que las decisiones se canalicen básicamente a través del Eurogrupo dejando a las Cumbre del euro la adopción de declaraciones con posiciones o decisiones adoptadas por consenso.

Los lugares de reunión, tema tradicionalmente pacífico, que solían ser capitales o ciudades del Estado que ejercía la Presidencia rotatoria se vio complicado por la impresionante movilización de recursos humanos y económicos que parecían no compensar el «prestigio» de la organización de las cumbres. Aparecieron adicionales complicaciones de seguridad debidas a que las reuniones del Consejo Europeo, cuya notoriedad como poder político de la Unión y de sus Estados era cada vez más evidente, se habían convertido en el lugar y momento elegido para la expresión de la contestación política. Una Declaración aneja al Tratado de Niza estableció que: A partir de 2002, una reunión del Consejo Europeo por Presidencia se celebrará en Bruselas. Cuando la Unión cuente con dieciocho miembros, todas las reuniones del Consejo Europeo se celebrarán en Bruselas. 

Esta declaración tiene su origen en un intento de compensación a Bélgica por la pérdida de su tradicional equiparación con los Países Bajos en la ponderación del voto en el Consejo. Con la Presidencia española del primer semestre de 2002, el asunto comenzó con un incumplimiento de sus previsiones (las dos reuniones tuvieron lugar en territorio español). Sin embargo ahora el Reglamento interno del Consejo Europeo establece con contundencia en su artículo 1.2 que se reunirá en Bruselas, salvo «circunstancias excepcionales», apreciadas por unanimidad del Consejo de Asuntos Generales o el COREPER.

3.1.3. El Presidente del Consejo Europeo 
Previsto en el artículo 15 TUE, el Presidente del Consejo Europeo es una de las novedades más conocidas del Tratado de Lisboa. Responde a la vieja preocupación de otorgarle a la UE, en particular en su proyección exterior, una mayor visibilidad y permanencia que viene a sumarse a la que ya se intentó con el Alto Representante de PESC y vuelve de nuevo con la remodelada y reforzada figura del Alto Representante de la Unión para Asuntos Exteriores y Política de Seguridad y con la Presidencia del Consejo Europeo. El Presidente —dice el artículo 15— es elegido por el propio Consejo Europeo por un período de dos años y medio, con posibilidad de una sola reelección (lo que hace coincidir el máximo de su nombramiento con una legislatura europea). Según ese artículo, dicha elección se adoptará por «mayoría cualificada» del Consejo Europeo. El primer Presidente del Consejo Europeo fue designado por unanimidad en una reunión especial del Consejo Europeo, el 18 de noviembre de 2009 antes de la entrada en vigor del Tratado de Lisboa y fue reelegido el 1 de marzo de 2012, por lo que su mandato se extendió hasta el 30 de noviembre de 2014. El 1 de diciembre el nuevo Presidente inició su mandato.

El Presidente del Consejo Europeo, que «no podrá ejercer mandato nacional alguno», podrá ser removido de su cargo por el propio Consejo Europeo por mayoría cualificada sólo por «impedimento» o «falta grave» (art. 15.5 TUE). En caso de enfermedad, fallecimiento o como consecuencia de una decisión del artículo 15.5 TUE será sustituido por el Miembro del Consejo  Europeo del Estado que esté ejerciendo la Presidencia rotatoria del Consejo hasta la elección del sucesor.

Sus funciones, no demasiado bien definidas, tienen una dimensión interna y una externa. La interna está dirigida esencialmente a organizar la propia actividad del Consejo Europeo ya su relación con el Parlamento Europeo. A tenor del artículo 15.6 TUE, el Presidente del Consejo Europeo:
a) presidirá e impulsará los trabajos del Consejo Europeo;
b) velará por la preparación y continuidad de los trabajos del Consejo Europeo, en cooperación con el Presidente de la Comisión y basándose en los trabajos del Consejo de Asuntos Generales;
c) se esforzará por facilitar la cohesión y el consenso en el seno del Consejo Europeo;
d) al término de cada reunión del Consejo Europeo, presentará un informe al Parlamento Europeo.

La dimensión externa tiene que ver con las carencias de visibilidad en la representación exterior de la Unión. A tales efectos el propio artículo 15.6 establece que: El Presidente del Consejo Europeo asumirá, en su rango y condición, la representación exterior de la Unión en los asuntos de política exterior y de seguridad común, sin perjuicio de las atribuciones del Alto Representante de la Unión para Asuntos Exteriores y Política de Seguridad (cursiva añadida).

No parece que esta cuestión haya quedado muy bien resuelta. El «rango y condición» no parecen criterios de fácil objetivación y como era previsible han aparecido algunas disfunciones. Por si el embrollo no estuviera asegurado, conviene recordar que, en lo que atañe a la representación exterior de la Unión, hay que sumar, en primer lugar, la Presidencia rotatoria del Consejo que, aunque no rige para el Consejo de «Asuntos Exteriores» (que está presidido por el Alto Representante), sí para el resto de las formaciones en virtud de lo previsto en el artículo 16.9 TUE8 (distinción ésta que opera sobre la cada vez más artificiosa distinción de lo interior y lo exterior) y, en segundo lugar, que deberá convivir, a su vez, con la acción de la Presidencia de la Comisión que, según dispone en el artículo 17.1 TUE, con excepción de la política exterior y de seguridad común y de los demás casos previstos por los Tratados, asumirá la representación exterior de la Unión. 

Si la distinción entre lo interno y externo es artificiosa, la línea que separa la acción en PESC de las clásicas relaciones internacionales comunitarias es, a veces, simplemente imposible de percibir.

El Reglamento interno (art. 2.3) añade alguna función más de especial interés, en la medida en que le otorga la responsabilidad de establecer «una estrecha cooperación y coordinación» con el Presidente de turno del Consejo (se deduce que con la Alta Representante en asuntos exteriores) y con el Presidente de la Comisión, estableciendo encuentros regulares. Es evidente que esta función afecta a la dimensión interna y externa.

3.1.4. La organización de los trabajos y procedimiento de votación
del Consejo Europeo

El sistema tradicional establecía que el Consejo Europeo estaría asistido en sus reuniones por los Ministros de Asuntos Exteriores y un miembro de la Comisión, pero la profundización institucional actual parece tratar de limitar las funciones de éstos: Cuando el orden del día lo exija —dice el artículo 16.3 TUE—, los miembros del Consejo Europeo podrán decidir contar, cada uno de ellos, con la asistencia de un ministro y, en el caso del Presidente de la Comisión, con la de un miembro de la Comisión. Es evidente que parece hacer algo excepcional de las funciones que los Ministros y del miembro de la Comisión que, hasta ahora, venían siendo esenciales en la organización del trabajo del Consejo Europeo.


El artículo 235.4 TFUE prevé que «El Consejo Europeo estará asistido por la Secretaría General del Consejo». La práctica de los últimos años antes de su institucionalización dejó sentado que, si el primer borrador del «orden del día» de la reunión del Consejo Europeo es elaborado por el COREPER, las «Normas...» aludidas del Consejo Europeo de Sevilla establecían que el definitivo correspondería al «Consejo de Asuntos Generales y Relaciones Exteriores» (denominación de esta formación del Consejo decidida asimismo por el Consejo Europeo de Sevilla). La cuestión se complicó con el Tratado de Lisboa. La función es atribuida al Consejo de Asuntos Generales «en contacto con el Presidente del Consejo Europeo y la Comisión» (segundo párrafo del art. 16.6 TUE). Al Consejo de Asuntos Generales corresponde, pues, la preparación de los trabajos del Consejo Europeo (dos semanas antes de su celebración) y la elaboración (cuatro semanas antes) de un proyecto de orden del día adoptado a propuesta del Presidente del Consejo Europeo, «en estrecha colaboración» con el miembro del Consejo Europeo al que corresponda el ejercicio de la Presidencia semestral del Consejo y del Presidente de la Comisión (art. 3.1 Reglamento interno).

En consecuencia, en este marco de autoorganización del Consejo Europeo, el «orden del día», elemento de especial importancia, debe ser el Consejo de Asuntos Generales (con la colaboración de los órganos señalados) el que prepara el proyecto de orden del día «comentado» en el que, en la práctica actual, se desglosan los siguientes puntos:
— Los puntos a ser aprobados o ratificados sin debate.
— Los puntos a debatir para definir orientaciones políticas generales.
— Los puntos que requieran una decisión del Consejo Europeo.

Si hasta ahora se referían solamente al debate y derivación posterior del Consejo para que éste adoptase las decisiones oportunas según las reglas de los Tratados, tras Lisboa, incluye aquellas vigentes para muchos ámbitos en que el Consejo Europeo debe decidir por sí mismo.
— Los puntos a debatir que no han de ser objeto de
conclusiones. 
Este proyecto, que debe ser elaborado con cuatro semanas de antelación, se eleva a «orden del día» provisional por el Presidente del Consejo Europeo a la luz del debate celebrado por el propio Consejo de Asuntos Generales en los cinco días anteriores a la celebración del Consejo Europeo. El orden del día definitivo lo adopta el propio Consejo Europeo «al inicio de su reunión».

Ciertas dificultades que en la práctica se planteaban para que al Consejo Europeo llegaran las propuestas de algunas formaciones del Consejo más especializadas (UEM) empujaron a facilitar la presencia de sus componentes (a la que se añadía la del Parlamento y el Comisario correspondiente) haciendo de las reuniones del Consejo Europeo algo engorroso y difícil de hacer funcionar con agilidad, dado el escaso margen de tiempo en el que se celebran. Por ello, el cambio introducido por el Tratado de Lisboa que obligó a una revisión de estas normas y el Reglamento interno del Consejo Europeo previene que «Las contribuciones de las demás formaciones del Consejo a los trabajos del Consejo Europeo se transmitirán al Consejo de Asuntos Generales a más tardar dos semanas antes de la reunión del Consejo Europeo» (art. 3.1 Reglamento). En línea con esta tendencia simplificadora de las reuniones, y con el fin de evitar interferencias de última hora, el Reglamento interno del Consejo Europeo determina que entre esta última reunión del Consejo de Asuntos Generales que ha adoptado el orden del día provisional y el inicio de la reunión del Consejo Europeo «excepto por razones imperativas e imprevisibles» ninguna formación del Consejo u órgano preparatorio podrán debatir un asunto sometido al Consejo Europeo (art. 3.2).

El papel del Parlamento Europeo, pese a la práctica de la intervención del Presidente del Parlamento Europeo, previa al inicio formal de la reunión del Consejo Europeo (opción prevista en el art. 235.2 TFUE), que le permite expresar la opinión de la Institución sobre las cuestiones más importantes del orden del día, es poco relevante, al menos desde un punto de vista formal. El TFUE prevé que el Consejo Europeo «podrá invitar al Presidente del Parlamento Europeo a comparecer ante él» (art. 235.2 TFUE) y mantiene la obligación de presentarle un «informe» después de cada reunión (función ésta asignada al Presidente del Consejo Europeo por el art. 15.3 TUE). El Reglamento interno ha añadido una obligación para el Estado miembro que ejerza la Presidencia de turno de presentar al Parlamento sus prioridades para el semestre con anterioridad y, transcurridos los seis meses, los resultados alcanzados. Es un intento más o menos logrado de asociar más al Parlamento Europeo con una lectura tal vez poco coherente con el papel del Presidente y que, en cualquier caso, no consigue ocultar el escaso papel del PE. Por su alta función representativa su relación con el Consejo Europeo habría merecido un mejor tratamiento.

El procedimiento decisorio en el seno del Consejo Europeo, sin necesidad de disposición expresa y como consecuencia de ese  carácter «internacional» que lo imbuye, es el consenso excepto cuando los tratados dispongan otra cosa, según señala el artículo  15.4 TUE, entendiéndose por tal la unanimidad de los Jefes de Estado o de Gobierno o, al menos, la no oposición frontal de alguno de ellos. Cuando a lo largo de los Tratados se establecen los procedimientos de adopción de decisiones del Consejo los tres sistemas son: la unanimidad, la mayoría cualificada y la mayoría simple. Para proceder a una votación se requiere como quorum la presencia de dos tercios de los miembros (se excluyen Presidente del propio Consejo Europeo y Presidente de la Comisión) y se procede a votación a iniciativa del Presidente (art. 6 del Reglamento interno).

La unanimidad exige el voto favorable de la totalidad de los Estados miembros. El TUE permite que cada miembro del Consejo Europeo pueda actuar en «representación» de uno solo de los demás miembros, sin que la abstención de los miembros presentes o representados obste a la adopción de la decisión (art. 235 TFUE). Igualmente especifica que, respecto a cuestiones de procedimiento y aprobación del reglamento interno, el Consejo Europeo se pronuncia por «mayoría simple», lo que debemos entender como «la mayoría de los miembros que lo componen»9.

La mayoría cualificada para el Consejo Europeo se define como la prevista en los artículos 16.4 TUE y 238.2 TFUE, ambos relativos a la votación en el Consejo. Este sistema de mayoría cualificada previsto es aplicado a partir del 1 de noviembre de 2014. Pero, entre el 1 de noviembre de 2014 y el 31 de marzo de 2017, tampoco se aplicará si cualquier miembro del Consejo (Europeo) solicita que se siga aplicando el sistema transitorio. Y, en consecuencia, mientras tanto, se regirá por el sistema recogido en el Protocolo 36 sobre disposiciones transitorias que, sobre la ponderación de voto, prevé en su artículo 3.3, que la mayoría cualificada requerirá, habida cuenta que aquí los actos no se adoptan a propuesta de la Comisión, al menos 260 votos que representen la votación favorable de dos tercios de los miembros como mínimo. Debe añadirse que en estos casos hay un requisito adicional si cualquier miembro del Consejo Europeo solicita que se compruebe que los Estados miembros que constituyen la mayoría cualificada representan como mínimo el 62 por 100 de la población total de la Unión, lo que, de no darse, impediría la adopción del acto.

Antes de abandonar el asunto de los sistemas de votación conviene tener presente que en el Consejo Europeo juega asimismo lo previsto para el Consejo en el artículo 238.3 TFUE para cuando el Consejo Europeo adopte decisiones en las que no participan todos sus miembros. Si en unanimidad y mayoría simple no ofrece problema, ya que simplemente no cuenta el Estado o Estados no participantes, en el caso de la mayoría cualificada el tema es algo más complicado10.

\subsubsection{Poderes y funciones}
La historia misma del Consejo Europeo es una historia de crecimiento continuo de su poder. Los motivos están claros. Las circunstancias políticas y económicas han contribuido de manera decisiva a esta imparable evolución. Tampoco resulta ajena a ello la permanente revisión de las normas constitutivas del sistema (Maastricht, Ámsterdam, Niza, Lisboa) que ha durado algo más de tres lustros y que por su excepcionalidad llama a una intervención de los más altos poderes de los EEMM, es decir, los que se reúnen en el Consejo Europeo. Pero como hemos visto, desde su origen, el poder del Consejo Europeo se ha debatido siempre entre su mayor presencia en el proceso europeo, convertida ya en habitual (bastever la evolución del número de reuniones por año —11 en 2015—), y su ubicación en los ámbitos de decisión excepcionales o, si se prefiere, no ordinarios del sistema. El Tratado de Lisboa no escapa a esta tensión y el resultado es un difícil compromiso entre ambas tendencias. Por un lado, la institucionalización del Consejo Europeo expresa su incorporación inequívoca al modelo decisorio europeo, mientras que, por otro, el TUE y TFUE tratan de ubicar al Consejo Europeo en los márgenes de lo ordinario, así como en la excepcionalidad y muy cerca de una especie de instancia modificadora del Derecho originario.

Las funciones que los Tratados adjudican al Consejo Europeo son realmente amplias y de extraordinaria importancia. No es, desde luego, nuestra intención detallar aquí todas y cada una de ellas, cuyo alcance habrá sido explicado o lo será en los diferentes temas de este libro, pero sí que trataremos de sistematizarlas por su alcance e incidencia. En este sentido, en mi opinión, podemos hablar de: a) funciones constitutivas, b) funciones de orientación y dirección general, c) funciones institucionales especiales y d) funciones como instancia de solución de conflictos sistémicos. 

Las funciones constitutivas contienen el conjunto de funciones que permiten al Consejo Europeo modificar las disposiciones y procedimientos establecidos en los propios Tratados, esto es, en las normas constitutivas. Son muchas y de esencial importancia. Las especiales circunstancias en que se desarrolló la CIG 2007 que condujo al Tratado de Lisboa, apremiada por plazos muy cortos ante el temor a un fracaso más, el intento de evitar el riesgo de nuevos controles de los Parlamentos nacionales o de referendos y de flexibilizar el duro blindaje de los procedimientos de revisión de los Tratados, llevaron a conceder a esta instancia poderes verdaderamente excepcionales. Así, entre otros, deben retenerse, establecer las formaciones del Consejo y las reglas de rotación de sus Presidencias (art. 236 TFUE), la adopción del sistema de rotación de la Comisión a partir de que tenga menos componentes que Estados miembros la Unión (si esto llegare a suceder alguna vez (art. 17.5 TUE)), recomendar la adopción de una «defensa común» (art. 42.2 TUE), en el procedimiento de revisión simplificado de la parte tercera del TUE (art. 48.6 TUE), dictar las orientaciones de negociación del tratado de retirada de un Estado miembro y, en su caso, prorrogar los dos años de la retirada de un Estado desde su notificación (art. 50 TUE), decidir el cambio de adopción de decisiones del Consejo en cualquier ámbito (incluida la PESC con la excepción de las que tengan repercusiones militares o en el ámbito de la defensa) de unanimidad a mayoría cualificada, así como el paso en el TFUE del procedimiento legislativo especial al ordinario (art. 48.7 TUE), ampliaciones de la competencia de la Fiscalía Europea modificando lo dispuesto en los apartados 1 y 2 del artículo 86 TFUE (art. 86.4 TFUE) en su caso, permitir al Consejo que adopte el marco financiero plurianual por mayoría cualificada y no por unanimidad (art. 312.2 TFUE), por iniciativa del Estado de que se trate (en el ámbito de la aplicación territorial de los Tratados) puede modificar el estatuto de la Unión para algunos de los territorios daneses, franceses u holandeses de los apartados 1 y 2 (art. 355.6 TFUE). Súmense a ello las múltiples funciones previstas en los distintos Protocolos (por ejemplo, los números 4, 9, 32, etc.) y se tendrá una idea clara del papel conferido al Consejo Europeo. Adviértase que, en razón de la prohibición expresa de emitir actos legislativos que pesa sobre el Consejo Europeo, los actos resultantes de estas funciones evidencian sus rasgos excepcionales, no reconducibles a una legalidad ordinaria. Las funciones de orientación y dirección general de la Unión son, sin duda, herederas de su función esencial de impulsión y orientación. Destacan la prevista en el artículo 22.1 TUE para la totalidad de la acción exterior de la UE fijando los «intereses y objetivos estratégicos», las orientaciones generales en PESC, incluida la PCSD (art. 26 TUE) o las orientaciones estratégicas de la programación legislativa y operativa del ELSJ (art. 68 TFUE), entre otras.

Por funciones institucionales especiales entiendo aquellas que responsabilizan al Consejo Europeo de los nombramientos de las más altas instancias de la Unión (su propio Presidente, el Alto Representante, propuesta del Presidente de la Comisión (arts. 15, 17 y 18 TUE), nombramiento del Presidente, Vicepresidente y los Consejeros del BCE (art. 283 TFUE) o de las cuestiones más delicadas políticamente. En estas últimas le corresponde constatar la existencia, en su caso, de una violación grave y persistente por un Estado miembro de los valores del artículo 2 TUE (democracia y derechos humanos) que pone en marcha el mecanismo sancionador de la Unión del artículo 7 TUE, o decidir en la UEM cuándo se cumplen los criterios de convergencia y, en consecuencia, se han de suprimir las excepciones (art. 140 TFUE), entre otras.

Por último conviene recordar que el Consejo Europeo asume funciones como instancia de solución de conflictos sistémicos. No me refiero a la tradicional función de dar salida a conflictos aparentemente irresolubles que cuenta con tradición como la, cuanto menos peculiar, solución a la cuestión danesa de 1992 todavía recogida en el Protocolo 22, la planteada por Irlanda tras el resultado negativo del referéndum sobre el Tratado de Lisboa en200911, y la no menos peculiar pirueta jurídica del Protocolo sobre la República Checa aceptado para salvar in extremis el último escollo para la entrada en vigor del Tratado de Lisboa12 que no encuentran un fundamento en los Tratados, sino a las previsiones constitucionales que convierten al Consejo Europeo en una verdadera instancia política de «apelación». Entre ellas, la capacidad de parar un proyecto de acto legislativo en el campo de la seguridad social cuando un miembro declare que le perjudica en determinados aspectos y acuda al Consejo Europeo (art. 48 TFUE), en la cooperación judicial si un miembro considera que una directiva armonizadora afecta a aspectos fundamentales de su sistema penal (art. 82.3 TFUE), en cooperación policial si fracasa el Consejo en la adopción de actos legislativos especiales en cooperación operativa y un grupo de al menos nueve Estados solicitan la intervención del Consejo Europeo (art. 87.3 TFUE) o las cruciales que tiene en el Protocolo sobre el acervo Schengen integrado (arts. 5.4 y 5 del Protocolo n.º 19).

\subsubsection{LOS ACTOS}

El tipo de actos ha respondido tradicionalmente a la flexibilidad con que se concebía la propia acción del Consejo Europeo y Declaraciones, Orientaciones o Decisiones que se adoptaran en el Consejo Europeo se incluían en bajo la fórmula de «Conclusiones de la Presidencia» que recogía los resultados de estas reuniones del Consejo Europeo, eventualmente acompañadas de «Anexos» cuando el Consejo Europeo deseaba expresar con precisión su consenso.

La institucionalización del Consejo Europeo desde el Tratado de Lisboa ha permitido, sin perder demasiado la flexibilidad, precisar y adaptar el modelo de actos del Consejo a las exigencias procesales y de forma propias de esa institucionalización. Así, el Consejo Europeo, a tenor del Acta que se redacta tras su reunión adopta tres tipos de actos. Las Conclusiones que albergan distintas fórmulas para expresar el consenso de los Jefes de Estado y de Gobierno («insta», evalúa», «pide al Consejo…») propias de su actividad impulsora y de orientación general13. Las Declaraciones, en los casos de adopción posiciones importantes de principio o de especial importancia tanto en la dimensión interna como externa de la UE (o cuando así se solicita por alguno de los EEMM a efectos de expresar alguna posición particular). Y Decisiones, expresión de un poder regulador más concreto del Consejo Europeo. En este último caso, salvo que se trate de «decisiones» de procedimiento interno, van firmadas por el Presidente del Consejo Europeo y el Secretario General del Consejo. Cuando no indican un destinatario concreto se publican en el DOUE14 y cuando identifican al destinatario se «notifican» por la Secretaría General del Consejo mediante el procedimiento general. Ha de recordarse que ningún acto del Consejo Europeo tiene naturaleza legislativa.

Esta mayor precisión en la regulación de los actos ha llevado a que tanto el acceso a sus documentos como la forma de los actos queden referidos (con las salvedades derivadas de la distinta naturaleza de las instituciones) a las normas previstas en el Reglamento interno del Consejo15.







Aunque desde 1974 los Jefes de Estado y de Gobierno de los Estados miembros se han reunido
periódicamente como '·Consejo Europeo" y los Tratados posteriores lo mencionaban, solo el
Tratado de Lisboa ha convertido al Consejo Europeo en una institución de pleno derecho de la
Unión y le ha dotado de un Presidente estable. El Presidente del Consejo Europeo lo es por un
periodo de dos años y medio (renovable una vez). No puede ejercer un mandato nacional en
paralelo. pero sí podría ejercer un mandato en el seno de otra institución europea (por lo que el
Tratado no descarta que la Presidencia del Consejo Europeo y la de la Comisión se fusionaran
algún día, si los Estados miembros así lo decidieran). El Consejo Europeo es un órgano de máximo nivel político, sin competencias legislativas expresas, que desempeña funciones muy relevantes:
Define las orientaciones y prioridades políticas de la Unión. Ha ido asumiendo de forma creciente funciones de iniciativa política (rivalizando la Comisión).
Define los principios y los objetivos de la política exterior y de seguridad común. Desempeña un papel capital a la hora de acordar el Marco Financiero Plurianual ( que después adopta el Consejo con arreglo a un procedimiento legislativo especial) y el sistema de recursos propios.
En el marco del Semestre Europeo, avala las orientaciones de política económica y aprueba las recomendaciones específicas por país.
Toma decisiones institucionales: fija la composición del Parlamento (número de escaños por EM). puede modificar el número de miembros de la Comisión, decide la formación del Consejo y el calendario de las presidencias rotatorias.
En cuanto a nombramientos: elige a su propio Presidente, propone al Presidente de la Comisión, nombra al Alto Representante, y nombra al Presidente del BCE (y al Comité Ejecutivo). También nombra oficialmente al Colegio de Comisarios en su conjunto, tras el voto de aprobación del Parlamento. 
Puede detenninar que existe un caso de violación grave y persistente de los valores de la Unión por parte de un Estado miembro (con sus consiguientes sanciones), de acuerdo con el artículo 7 del TUE. Pueden modificar la reglas de decisión del Consejo ("'cláusulas pasarela'· generales). 

Para llevar a cabo estas funciones. las reuniones del Consejo Europeo tienen lugar en Bruselas dos veces por semestre (en marzo, junio, octubre y diciembre). Asimismo, el Presidente puede convocar reuniones extraordinarias si es necesario. Asisten los miembros del Consejo Europeo (es decir, los 27 Jefes de Estado y de Gobierno, el Presidente del Consejo Europeo y el Presidente de la Comisión) y, en función de los temas, se invita al Alto Representante o al Presidente del BCE.\footnote{%
    El Presidente del 13CE siempre asiste cuando los miembros del Consejo Europeo se reúnen como Cumbre del Euro. La Cumbre del Euro se instituyó oficialmente mediante el Tratado de Estabilidad. Coordinación y Gobcmanza. En principio. aprovechando la celebración de algún Consejo Europeo. la Cumbre del Euro reúne dos veces al año a los Jefes de Estado y de Gobierno de los países del euro (pudiéndose ampliar la convocatoria a todos los Estados m iembros para tratar temas de arquitectura de la UEM). El Presidente de la Cumbre del Euro es nombrado por los Jefes de Estado y de Gobierno de los Estados miembros de la zona del euro. El nombramiento se produce al mismo tiempo que el del Presidente del Consejo Europeo y tiene la misma duración. Los dos cargos pueden ser desempeñados por la misma persona. como h a ocurrido tanto en el caso de Donald Tusk como el de Charles Michel.
} En cada reunión, se aprueban unas conclusiones que reflejan los principales mensajes resultantes de los debates y hacen un balance de las decisiones adoptadas. Las reuniones se preparan en el seno del Consejo de Asuntos Generales y del COREPER 11. En general, el Consejo Europeo decide por consenso (si bien en algunos casos, como el de los nombramientos, aplica la mayoría cualificada). Cuando el Consejo Europeo decide por votación, solo pueden votar los Jefes de Estado y de Gobierno.

\subsection{Consejo de la Unión Europea}
El Consejo es una institución en la que están representados los intereses nacionales y, por ello, encarna el principio de la representación de los Estados integrados en la Unión. Hasta la entrada en vigor de la reforma de Lisboa era la única institución que encarnaba el principio de la legitimidad intergubernamental o representatividad territorial, pero desde entonces comparte esa representación de los intereses nacionales con el Consejo Europeo. Pero este último es un ejecutivo intergubernamental a diferencia del Consejo que es fundamentalmente un legislativo intergubernamental. El Consejo asume los más importantes poderes de decisión legislativa y presupuestaria, así como de definición de políticas y de coordinación en la Unión (art. 16 TUE). Asume tan importantes y diversas funciones porque la creación, primero, de las Comunidades Europeas y, más tarde, de la Unión, así como su dinámica misma, está condicionada a la existencia y permanencia de los Estados miembros. El Consejo se regula en los artículos 16 TUE y 237 a 243 TFUE.

El Consejo no es una conferencia internacional ni sus actos son acuerdos internacionales, a pesar de ser una institución compuesta por miembros de los Gobiernos de los Estados miembros, quienes, en consecuencia, representan y defienden legítimamente intereses nacionales1.

Es cierto que los miembros del Consejo hacen valer intereses nacionales, pero el resultado de sus deliberaciones es armonizar esos intereses particulares en el superior interés de la Unión. El Consejo es una institución que, después de un complejo proceso de confrontación de intereses nacionales, define el interés común.

El Consejo, aunque no es independiente de los Estados, expresa una voluntad propia, distinta de la voluntad de cada uno de los Estados y la decisión del Consejo se impone a todos los Estados miembros, hayan votado a favor o en contra o se abstuvieran. Como la voluntad del Consejo es distinta a la de sus miembros, un Estado, aun habiendo votado a favor, podría impugnar un acto del Consejo, pues el voto de un Estado no tiene relevancia autónoma; es un elemento de la voluntad colegiada. Un Estado miembro no pierde sus derechos (la legitimación activa ante el Tribunal de Justicia) por el hecho de haber contribuido a formar la voluntad de una institución de la Unión.

\subsubsection{Composición y funcionamiento}
El Consejo está compuesto por un representante de cada Estado miembro, de rango ministerial, facultado para comprometer al Gobierno del Estado miembro al que represente (art. 16.2 TUE).

Una interpretación coherente con el Derecho internacional y con el Derecho constitucional nos lleva a unas deducciones que, en lo básico, coinciden con la composición tradicional del Consejo. En efecto, del artículo 16.2 TUE se deduce, en primer lugar, que habrá que atenerse al sistema constitucional de cada Estado para saber quiénes son miembros del Gobierno de cada Estado miembro, pues sólo éstos pueden formar parte de pleno derecho del Consejo.

En segundo lugar, no debe interpretarse literalmente la alusión al «rango ministerial» en el sentido de que sólo los «ministros», como componentes del Consejo de Ministros de un Estado, podrían formar el Consejo. El rango ministerial puede serlo en el nivel regional o autonómico, siempre que el Gobierno del Estado miembro faculte a esa persona para comprometer su voluntad caso por caso. 

En tercer lugar, si no pudiera asistir un miembro del Gobierno podría hacerse representar por otra persona; ahora bien, el voto habrá de delegarse en un Estado miembro que esté representado por un miembro del Gobierno (art. 239 TFUE); esta peculiar exigencia es, indirectamente, una regla de quorum que asegura una presencia mínima de la mitad de los Estados representados por miembros de los Gobiernos para proceder a una votación. En la práctica, todas las instancias del Consejo suspenden sus trabajos en presencia de menos de la mitad de las delegaciones. En cuarto lugar, la composición personal del Consejo es variable según los temas a tratar en el orden del día de cada Consejo; incluso podrían participar varios ministros como titulares en una  misma formación del Consejo (con un solo turno de palabra y un solo voto).

2.1. LAS FORMACIONES DEL CONSEJO
Hasta la entrada en vigor de la reforma de Lisboa había un total de nueve formaciones y en la actualidad son diez2. Es una lista propia de la Unión, que no interfiere con la organización interna de cada gobierno. 

El artículo 16.6 TUE prevé que el Consejo se reúna en diferentes formaciones. Atribuye al Consejo Europeo la competencia para crear formaciones del Consejo, al tiempo que reconoce al menos dos que ya existían unificadas pero separadas de facto (órdenes del día y reuniones en fechas distintas): el Consejo de Asuntos Exteriores y el de Asuntos Generales. La formación más relevante y la que se reúne con más frecuencia es el Consejo de «Asuntos Exteriores» al que asisten los ministros de Asuntos Exteriores y es presidido por el Alto Representante de la Unión para los Asuntos Exteriores y la Política de Seguridad.

Otra formación relevante, y que se ha desgajado de la anterior a partir de la reforma de Lisboa, es el Consejo de Asuntos Generales (CAG). Se ocupa de los temas de organización y funcionamiento de la UE, cuestiones que no competan a un Consejo especializado o temas multidisciplinares o la cohesión económica y social, las relaciones con las demás Instituciones o las Conferencias Intergubernamentales o la adhesión de nuevos miembros. Dependiendo de cada Gobierno, a esta formación acuden los Ministros de Asuntos Europeos o los de Exteriores. El artículo 16.6 TUE le encarga velar por la coherencia de los trabajos de las otras formaciones y preparar las reuniones del Consejo Europeo y garantizar la continuidad de sus trabajos junto al Presidente del Consejo Europeo y de la Comisión.

Una formación muy destacada y tradicional es el Consejo «EcoFin» (asuntos económicos y financieros), encargado de seguir la política económica y presupuestaria de los Estados miembros para establecer su conformidad con las orientaciones de la Unión Económica y Monetaria (UEM): se interesa también por la política fiscal, de empleo, social, etc.; examina las previsiones financieras a medio plazo, ligadas a su vez con la PAC y la política de cohesión; hace el seguimiento de los aspectos financieros de las actuaciones políticas exteriores: ayudas macrofinancieras para terceros Estados, situaciones de graves crisis financieras, o la representación en el exterior de la moneda única.

Una consecuencia de la composición variable de las formaciones del Consejo es la falta de homogeneidad y de continuidad en sus trabajos, que ha tratado de ser paliada con un órgano auxiliar del Consejo, el Comité de Representantes Permanentes (COREPER), que prepara detalladamente los trabajos del Consejo, y mediante la infraestructura administrativa que comprarte con el Consejo Europeo y aglutina la Secretaría General del Consejo (alrededor de 3.100 funcionarios). 

Cabe señalar que en las sesiones del Consejo —y por ende de toda su infraestructura, como el COREPER y sus grupos de trabajo — está siempre presente la Comisión, a través de su Presidente o los Comisarios competentes.

3. LA PRESIDENCIA 

Hasta 1995 la rotación en la Presidencia del Consejo se seguía el orden alfabético3. Pero a partir de 1996 se acordó que no hubiera un orden rígido: se mantuvo el principio de rotación y la duración semestral. Con la entrada en vigor de la reforma de Lisboa se han introducido otros cambios: habrá un sistema dual, con presidencias rotatorias para el conjunto de formaciones y, por el contrario, con presidencia permanente en el Consejo de Asuntos Exteriores a cargo del Alto Representante de la Unión (y Vicepresidente de la Comisión) a fin de lograr una mayor visibilidad, coherencia y eficacia en materia de acción exterior.

Una segunda excepción al sistema general de rotación es el Eurogrupo (Consejo de los Estados que tienen como moneda el euro).

Salvo para las dos excepciones citadas, la presidencia rotatoria se ejerce, de forma sucesiva, por cada Estado miembro, durante un período de seis meses. Pero el orden preciso se acuerda por el Consejo por mayoría cualificada y puede ser modificada como consecuencia de circunstancias especiales que afecten a un Estado (procesos electorales, etc.)4.

El sistema rotatorio semestral en la Presidencia del Consejo se aplica también a su órgano auxiliar (COREPER) y a los dos centenares de grupos de trabajo (permanentes o temporales) dependientes de éste. De este modo, todas las presidencias en la esfera del Consejo se ejercen por personas de una misma nacionalidad que aseguran así una cierta coordinación y continuidaden las actividades del Consejo. La presidencia del Consejo es una función política —y no meramente un rango honorífico— que da al Estado que la ejerce un poder de coordinación y dirección de los trabajos, de impulso y conciliación de los intereses nacionales y de representación. En efecto, la presidencia hasta ahora ha sido «el motor político de la UE», asumiendo una función de impulso y programación de los trabajos, velando por la conclusión de expedientes en plazos razonables y esforzándose por conciliar los intereses de todos. Se estima que la rotación da lugar a una obligación de balance que constituye una presión dinámica (y positiva) para el avance de los asuntos. La rotación semestral tiene muchas ventajas: ofrece la misma oportunidad de influir a cada Estado, moviliza las energías de todos, aproxima la Unión y sus asuntos a la ciudadanía y garantiza  coherencia en los trabajos gracias a la unicidad del control político del conjunto de la maquinaria de la Unión.

Los inconvenientes de la rotación son, en especial, la falta de continuidad en los trabajos de grupos y comités, y calendarios presidenciales demasiado coincidentes con intereses nacionales. Desde 2007, anticipándose a la reforma de Lisboa, se practica una presidencia mediante grupos predeterminados de tres Estados durante dieciocho meses y cada miembro del grupo la ejercerá durante seis meses.

3.1. LA ORGANIZACIÓN, LA COORDINACIÓN Y LA DIRECCIÓN DE LOS
TRABAJOS

El calendario de sesiones se programa con una antelación de siete meses al del ejercicio mismo de la presidencia de forma concertada con los otros Gobiernos y la Comisión. Los programas no son propiamente semestrales sino programas anuales confeccionados conjuntamente por dos presidencias en función de directrices estratégicas a medio plazo que determina el Consejo Europeo; estos programas anuales, elaborados tras mantener las consultas pertinentes, tienen la doble ventaja de ser más equilibrados, al estar menos condicionados por inquietudes nacionales determinadas, y más operativos ya que la duración de un año se corresponde más con el ritmo legislativo normal de la Unión
(calendario legislativo acordado entre PE y Comisión). 

La convocatoria y el orden del día se hacen siempre por el presidente en ejercicio, ya sea a su iniciativa, ya sea a petición de algún Estado miembro o de la Comisión. Lo habitual es que sea a iniciativa del presidente en ejecución del calendario de sesiones programado, si bien puede convocar nuevas reuniones si lo estima necesario o urgente. El orden del día se elabora propiamente por el presidente del COREPER, que es el Jefe de la Representación Permanente del Estado miembro que en ese semestre ejerce la presidencia, el cual consulta a las otras Representaciones Permanentes y a la Comisión por si estimasen pertinente la inclusión de otros temas y se envía con una antelación de dos semanas (catorce días) acompañado de la documentación necesaria sobre cada punto; sólo se podrán incluir nuevos puntos por unanimidad. Es obvio que en la elección concreta de los temas a tratar hay cierta dosis de elección política.

Es importante destacar que en el orden del día provisional se indicarán los puntos respecto de los cuales la presidencia, las delegaciones de los Estados miembros o la Comisión han solicitado una votación y, por tanto, se realizará en la sesión convocada. La ordenación de los debates se dirige por la presidencia; lo usual es que siga un turno de intervenciones. La agilidad de las discusiones y la conducción del debate hacia los problemas importantes dependen del buen conocimiento que tenga la presidencia de los puntos en discusión.

El recurso a la votación es el poder concreto más importante y eficaz que se asigna a la presidencia. Frente al abuso del consenso hasta 1986, el Reglamento interno del Consejo permite a su presidente en ejercicio decidir por sí solo cuándo un punto del orden del día será objeto de votación formal. Los miembros del Consejo sabrán así que, agotada la deliberación sin acuerdo, tales divergencias no serán obstáculo para que, si se reúne la mayoría cualificada, la propuesta normativa sea adoptada. Al señalarse en el orden del día, esta indicación acelera los contactos previos entre las delegaciones y estimula una transacción. 

También pueden solicitar una votación los Estados miembros o la Comisión; esta institución puede solicitar una votación en razón de que goza de paridad frente al Consejo y, además, éste delibera sobre sus propuestas normativas y vela por el cumplimiento de los Tratados (que prevén para todas las decisiones del Consejo el recurso al voto). Pero, a diferencia del poder otorgado al Presidente del Consejo, cuando la Comisión o un Estado miembro solicitan la votación, se precisa que estén a favor de proceder a la votación la mayoría de los miembros del Consejo.

La asistencia de la Secretaría General del Consejo cumple funciones de gran importancia para el Consejo y su Presidente en ejercicio. Está formada por funcionarios formalmente independientes y estructurada de la siguiente forma: el Secretario General y el Secretario General Adjunto (nombrados, por mayoría cualificada), su Gabinete (protocolo, prensa, seguridad, etc.), el servicio jurídico (análisis de los actos, contenciosos ante el Tribunal de Justicia, etc.) y varias Direcciones Generales.

Además de asumir las funciones propias de secretaría en las reuniones del Consejo (también del Consejo Europeo, del COREPER, de los comités y grupos de trabajo) y de preparar la documentación que acompaña a cada punto del orden del día, son el apoyo imparcial de la Presidencia: un apoyo técnico propio, de carácter colectivo y común, frente al soporte nacional de la Representación Permanente y de la Administración nacional. Los funcionarios de la Secretaría General aseguran también, además de calidad técnica, una coordinación a la dispersa y heterogénea estructura del Consejo.

3.2. LA LABOR DE IMPULSO, CONCILIACIÓN E IMPARCIALIDAD DE LA PRESIDENCIA

Finalizar el mandato semestral habiendo hecho progresar en su conjunto las propuestas pendientes, culminar algunas y desbloquear otras, constituiría un gran éxito para una presidencia. Igualmente, si un asunto se presentase muy controvertido el presidente puedeoptar por evitar un debate que ahonde las diferencias; propondrá un estudio más detallado en niveles inferiores, o asumirá la responsabilidad de la promoción de acuerdos o mayorías cualificadas tan amplias como sea posible, conociendo los obstáculos que plantean los Estados y proponiendo compromisos. 

Como no basta ser neutral sino que hay que parecerlo, el Presidente en ejercicio no encabeza la delegación de su país durante el semestre (igual sucede en el COREPER y en sus grupos de trabajo) ni toma la palabra en defensa del interés nacional. Otro miembro del Gobierno de ese Estado defenderá los intereses nacionales. Esto requiere un notable esfuerzo para doblar la presencia en toda la infraestructura del Consejo.

3.3. LA REPRESENTACIÓN

El Presidente en ejercicio no preside la Unión Europea ni representa formalmente a ésta. Sólo preside y representa al Consejo mismo.

La presidencia elabora un programa de prioridades e intenciones para el semestre que se presenta al Parlamento Europeo por el jefe del gobierno o el ministro de Asuntos Exteriores; Ministros sectoriales asisten a todas las sesiones plenarias del PE; también suelen tener algún encuentro con el Comité Económico y Social y unas estrechas y continuas relaciones con el Presidente de la Comisión.

Sin embargo, al haberse creado una presidencia estable para el Consejo Europeo, la del Consejo pierde la función de representación, coordinación y portavocía en las relaciones diplomáticas y protocolarias con terceros Estados y organizaciones internacionales.

En definitiva, la presidencia da a los Estados la oportunidad de ofrecer y mejorar la imagen de su país engrandecido por su pertenencia a la Unión y, desde luego, a los Estados medios y pequeños el ejercicio de la presidencia los convierte en «grandes» Estados.


\paragraph{Votación}
Los Tratados prevén varios sistemas: mayoría simple, mayoría cualificada y unanimidad. 

6.1. MAYORÍA SIMPLE
Se prevé en pocas ocasiones; significa que cada Estado posee un voto igual y la mayoría se obtiene cuando votan a favor la mayoría de los Estados que lo componen11.

6.2. MAYORÍA CUALIFICADA 

La mayoría cualificada es el procedimiento generalizado para la formación de la voluntad de la Unión (arts. 16.1 TUE y 238.2 TFUE)12; se vota por mayoría cualificada siempre, salvo cuando el Tratado dispone de otra forma. Desde el Tratado de Maastricht y en cada reforma posterior han aumentado las materias que son objeto de votación por mayoría cualificada.

La nueva mayoría cualificada, introducida por las reformas del Tratado de Lisboa, se aplica desde el 1 de noviembre de 2014. El sistema antiguo de voto ponderado sólo se podrá utilizar excepcionalmente a petición de cualquier Estado hasta noviembre de 2017 (Protocolo n.º 36 sobre las disposiciones transitorias). Siempre que se vota por mayoría cualificada, ya sea el sistema de Lisboa, ya fuera el excepcional y temporal de Niza, el Consejo puede aprobar la propuesta de la Comisión por mayoría cualificada, pero no puede modificar el texto más que por unanimidad, ya sea para desviarse de la propuesta de la Comisión, ya para aceptar enmiendas del PE no asumidas por la Comisión (art. 293 TFUE). Es el llamado «privilegio» de la Comisión. Luego, si en el Consejo no hay unanimidad para enmendar una propuesta de la Comisión o aceptar las enmiendas del PE, el Consejo tiene dos opciones: adoptarla por mayoría cualificada, o rechazarla, es decir, no adoptar decisión alguna en esa materia (más ampliamente, Cap. 11, epígrafe 8). Únicamente la Comisión puede modificar su propuesta en todo momento —antes de la votación en segunda lectura— y facilitar — preservando en lo esencial el interés general que inspira su acción — el acuerdo en el seno del Consejo (art. 293.2 TFUE).

Desde el 1 de noviembre de 2014 para aprobar un acto por mayoría cualificada, a propuesta de la Comisión, se precisará:
— el 55 por 100 de los Estados que representen el 65 por 100 de la población.

Esto significa que el peso del voto de cada Estado miembro depende de su población13. Deberán concurrir al menos dieciséis Estados miembros (quince, en caso de retirada británica) y los Estados concurrentes en esa mayoría deben reunir al menos el 65 por 100 de la población de la Unión.

La minoría perdedora (distinta a la de bloqueo)14 serán un grupo de doce o menos Estados miembros o un grupo de Estados que reúna menos del 35 por 100 de población.

La minoría de bloqueo (impedir por poco la mayoría cualificada cuando la propuesta perjudique seriamente al grupo de bloqueo) consistirá en reunir a trece Estados en contra; o, alternativamente, un mínimo de cuatro Estados con algo más del 35 por 100 de población. Se exige que en una minoría de bloqueo haya al menos cuatro Estados, dado que el 35 por 100 de la población podría ser reunida por tres grandes Estados. Fue un contrapeso exigido por España para contrarrestar a los grandes y que aceptó la CIG de 2007.

Cuando no se requiere propuesta de la Comisión (decisiones en materia de PESC y JAI, sanciones, retirada, etc.) entonces se exige un 72 por 100 de Estados y un 65 por 100 de población (21 Estados sobre 28; caso de retirada del Reino Unido, 20 sobre 27). La reforma de Lisboa hereda los contrapesos que España exigió en la elaboración de esta parte del fracasado Tratado constitucional, si bien fue Polonia la que propuso reintroducir el «compromiso de Ioannina», consistente en que la posible minoría perdedora de una votación pida una reconsideración temporal de una propuesta legislativa antes de su adopción. Sin embargo, hasta noviembre de 2017 podría recurrirse al sistema anterior basado en la mayoría ponderada tradicional si un Estado miembro así lo solicita15.


6.3. UNANIMIDAD
La unanimidad se prevé para decisiones de especial importancia. La unanimidad responsabiliza a uno o varios Estados del «veto», pero también puede influir sobre un Estado aislado para aceptar una transacción. En todo caso, las abstenciones no impiden la adopción de decisiones cuando se requiere la unanimidad (art. 238.4 TFUE). La unanimidad se mantiene para determinadas materias como fiscalidad, recursos propios, régimen lingüístico, política exterior, defensa, seguridad social, cultura, imprevisión de competencias, número de diputados del Parlamento Europeo, revisión y simplificación de los Tratados, admisión de nuevos Estados miembros, pasaportes, documentos de identidad y permisos de residencia, derecho de familia con implicaciones transfronterizas, aspectos de cooperación en materia penal y policial, determinados aspectos en varias políticas (trasportes, propiedad intelectual, política monetaria, política comercial), etc. 

En un mismo ámbito, pero diferenciando las materias concretas, pueden preverse dos procedimientos de adopción y votación diferentes.

En relación con esa serie de ámbitos sometidos todavía a la unanimidad se debe destacar que se prevén unas pasarelas para facilitar el abandono progresivo de la unanimidad mediante un procedimiento en el que se requiere iniciativa y la unanimidad del Consejo Europeo, la aprobación del Parlamento Europeo y la no oposición de los Parlamentos nacionales (art. 48.7 TUE).


\subsubsection{ATRIBUCIONES DEL CONSEJO}
La asignación de funciones al Consejo no ha variado por la reforma del Tratado de Lisboa. Se sistematizan y se clarifican mejor en su precepto de cabecera (art. 16 TUE). En coherencia con la redacción dada a las misiones del Parlamento Europeo, se repiten de forma paralela las dos misiones que comparten: la función legislativa y presupuestaria («conjuntamente con el Parlamento Europeo»). Por el contrario, son funciones específicas del Consejola definición de políticas y la coordinación. El Consejo encarna la influencia del control permanente e institucionalizado de los Estados miembros sobre las decisiones políticas y sobre la legislación de la Unión.

5.1. LA FUNCIÓN LEGISLATIVA DEL CONSEJO

Muchos años antes de las reformas hechas por el Tratado de Lisboa, ya se calificaba al Consejo por la doctrina y la jurisprudencia como «poder legislativo»; también su reglamento interno, desde hacía años, utilizaba la noción de actos legislativos para adjetivar los actos de codecisión.

El Tribunal de Justicia, en la importantísima sentencia Simmenthal, se refirió al «poder legislativo de la Comunidad»5. El antiguo juez del Tribunal de Justicia, Pierre Pescatore, decía que el Consejo es un verdadero legislador; un legislador de carácter no parlamentario, y cuya base de legitimidad es, desde luego internacional (la representatividad) y no directamente popular. Sin embargo, ese poder legislativo del Consejo tiene un contrapeso en favor de la Comisión, al gozar ésta del cuasi-monopolio de la iniciativa legislativa, al menos en el procedimiento legislativo ordinario. El poder de decisión del Consejo en los ámbitos legislativo y ejecutivo está condicionado, no obstante, por el poder de iniciativa de la Comisión y por el acuerdo del Parlamento Europeo. Este sistema de pesos y contrapesos constituyen la cuidadosa articulación del principio del equilibrio institucional o equilibrio de poderes.
La función legislativa del Consejo, compartida con la del Parlamento Europeo, se traduce fundamentalmente:
— En la aprobación en pie de igualdad de actos legislativos por el Consejo y el Parlamento, adoptados en procedimiento ordinario. Este procedimiento es el previsto de forma generalizada para la gran mayoría de los ámbitos cubriendo el conjunto del mercado interior (dado el carácter conjunto del procedimiento se expone en Cap. 11, epígrafe 8, sobre el Parlamento Europeo).
— En la aprobación de actos legislativos por el Consejo, adoptados previa consulta, en unos casos, y con la aprobación en otros por el Parlamento Europeo, mediante un procedimiento especial.

A su vez, se podrá delegar el poder legislativo del Consejo y del Parlamento Europeo en la Comisión para la adopción de actos «delegados» (pero ya serán actos no legislativos, art. 290.1 TFUE). Estos actos añadirán ese calificativo (reglamento delegado, directiva delegada, decisión delegada; véase Cap. 10, epígrafe 7.2.1). 

5.2. FUNCIÓN LEGISLATIVA DEL CONSEJO MEDIANTE PROCEDIMIENTO ESPECIAL
El procedimiento legislativo especial se distingue por la preponderancia, según los casos, del Consejo o por la delParlamento Europeo, frente a la paridad de Parlamento y el Consejo en el ordinario. La utilización de un procedimiento u otro viene predeterminado en cada base jurídica, caso por caso, y, por tanto, las Instituciones no tienen libertad de elección. Son actos con valor legislativo aunque no sean aprobados mediante el procedimiento legislativo ordinario (art. 289.3 TFUE). El Consejo está habilitado para adoptar actos mediante un procedimiento especial en los casos específicamente previstos en los Tratados (art. 289.2 TFUE).

Hay que constatar que se prevén muchísimos más actos legislativos del Consejo que del Parlamento Europeo. Ahora bien, siempre en todo acto legislativo de procedimiento especial, ya sea del Consejo o del Parlamento, hay una participación de la otra institución. Por ejemplo, los actos legislativos del Consejo requieren en varias ocasiones de la aprobación del Parlamento6 y, por tanto, la legitimidad democrática directa se asegura en lo esencial; pero en la gran mayoría se prevé sólo la mera consulta7 al Parlamento. En tanto que actos legislativos sus caracteres sustantivos son  idénticos: tienen su base jurídica directamente en el Tratado,  enuncian el marco jurídico fundamental de un ámbito y las opciones políticas básicas.

Ahora bien, la legitimidad democrática cambia, o al menos es distinta, cuando se trata de actos legislativos del Consejo adoptados mediante procedimiento especial (sólo concurriría una legitimidad internacional o territorial y sólo indirectamente democrática); también cuando ese acto legislativo del Consejo requiera la aprobación del Parlamento, pues su aportación a la formación del acto no es constructiva y su mayor poder es el veto y, desde luego, su aportación es ínfima cuando se trata de la mera consulta. No cabe una caracterización general del procedimiento legislativo especial, a diferencia de lo que sucede con el procedimiento legislativo ordinario; es un sistema plenamente casuístico que se detalla en cada base jurídica en la que se prevé la adopción de actos por este procedimiento. Sus variantes se producen ya sea en materia de iniciativa8, ya de consultas o de votación; caso por caso, se marca el itinerario que seguirá cada acto, frente a la uniformidad del procedimiento ordinario.

La utilización del procedimiento legislativo especial tiene un carácter de excepción o temporal, pues el TUE prevé, mediante el procedimiento simplificado de reforma de los Tratados, su progresiva eliminación a favor del procedimiento legislativo ordinario (art. 48.7.2.º párrafo TUE). De todos modos, todavía no se ha utilizado esta vía o pasarela para eliminar o reducir los procedimientos legislativos especiales en favor del ordinario.


5.3. ATRIBUCIONES NORMATIVAS NO LEGISLATIVAS. FUNCIÓN DE
DEFINICIÓN Y COORDINACIÓN DE POLÍTICAS
El Consejo goza del poder de definición y coordinación de las políticas en forma muy amplia. Es frecuente que se atribuya la coordinación al Consejo, o a los Estados miembros en el marco de la Unión, por lo tanto en el Consejo, o junto a otras instituciones. Así, tiene atribuciones de coordinación en diversos casos tales como los siguientes:
— supervisar la evolución económica de cada uno de los Estados miembros y de la Unión, así como la coherencia de las políticas económicas con las orientaciones generales (arts. 120 y 121 TFUE);
— reforzar la coordinación y supervisión de su disciplina presupuestaria (art. 136 TFUE);
— coordinar la estrategia para el empleo (arts. 145 y 146 TFUE);
— coordinar en materias diversas como salud pública (art. 168 TFUE), redes transeuropeas (art. 171.2), industria (art. 173), cohesión económica, social y territorial (art. 175), investigación y política espacial (arts. 181 y 189), políticas nacionales de cooperación al desarrollo (art. 210), en caso de recurrir a la cláusula de solidaridad (art. 222), lucha contra el fraude (art. 325), etc. La política económica, a diferencia de la monetaria, no es política exclusiva. Los artículos 119 a 126 TFUE y varios protocolos sientan las bases de esa política de coordinación, incluidas las normas de  disciplina financiera y presupuestaria que deben respetar los Estados miembros para mantener la unión monetaria.

En el marco de la UEM, el Consejo tiene la máxima importancia, pues la Comisión ejerce en ocasiones su iniciativa mediante una recomendación en lugar de una propuesta y el Consejo puede modificarla sin requerir la unanimidad (arts. 121.2 y 4, 126.7 TFUE), aunque también se prevé el procedimiento legislativo ordinario (art. 121.6 TFUE). El Consejo se encarga de la vigilancia de las políticas económicas y presupuestarias de los Estados miembros; significa que puede advertir y sancionar a un Estado miembro en caso de déficit presupuestario (arts. 126.8 a 11 TFUE). Igualmente, no es de descartar la posibilidad de que adopte recomendaciones y así armonizar los diversos casos en que el Tratado permite adoptar «medidas» sin referirse a concretas normas vinculantes.

El Consejo es una de las instituciones habilitadas para adoptar actos no legislativos siempre que caso por caso se prevea su adopción. La noción de actos no legislativos es residual: si el precepto no menciona un procedimiento legislativo, las demás bases jurídicas que prevean la adopción de actos dan lugar a actos no legislativos.

Resulta obvio que, desde el punto de vista sustantivo, un acto no legislativo no debe contener opciones políticas fundamentales, ni elementos esenciales del régimen jurídico en cuestión, ni definir objetivos generales ni principios. Nunca debe regular o alterar el disfrute de derechos humanos. Sin embargo, tiene que tener necesariamente su base jurídica directamente en el Tratado (principio de atribución de competencias)9.

Los actos no legislativos atribuidos al Consejo se adoptan con base directamente en los Tratados o para desarrollo de actos legislativos10. Se puede observar que el Consejo y la Comisión son las dos instituciones que por excelencia están llamadas a producir actos no legislativos con fundamento directo en el Tratado, cuando así se prevea expresamente.

Los actos no legislativos se rigen para su adopción caso por caso por el precepto que establece su base jurídica (lo que también sucedía con los actos legislativos «especiales»). Por tanto, la casuística sobre los derroteros de cada procedimiento no legislativo es innumerable y con toda clase de variantes, ya sea en cuanto a la iniciativa (en general, la Comisión, pero a veces el BCE, los Estados miembros, etc.), ya sea a las consultas institucionales (al Comité Económico y Social, al Comité de las Regiones, al BCE, etc.), ya sea al sistema de votación.

5.4. ATRIBUCIONES EN MATERIA DE RELACIONES EXTERIORES
(REMISIÓN)

Debe velar por la unidad, la coherencia y la eficacia de la acción de la Unión. Por ello le corresponde al Consejo ejercer en lo esencial la responsabilidad de la acción exterior, tanto en materia de relaciones económicas externas como en materia de Política exterior y de seguridad común (PESC) (véanse Caps. 23 a 25).










El Consejo es el órgano formado por los representantes con rango ministerial de los gobiernos de cada uno de los 27 Estados m iembros. Dichos representantes han de estar facultados para comprometer al gobierno del Estado al que representen y ejercer el derecho a voto. Los representantes del Consejo varían en función del tema a debatir, dando lugar a distintas formaciones (diez en la actualidad).\footnote{
Cuatro recaen en el ámbito del COREPER 11 (Embajadores Representantes Permanentes): Asuntos Económicos y F 0 inancieros (ECO FIN): Asuntos Exteriores: Asuntos Generales: Justicia y Asuntos de Interior. Seis recaen en el ámbito del COREPER I (Embajadores Representantes Permanentes Adjuntos): Agricultura y Pesca: Competitividad; Medio Ambiente: Empleo. Política Social, Sanidad y Consumidores: Educación. Juventud. Cultura y Deporte: Transporte. Telecomunicaciones y Energía.
} Aunque el Consejo no es un órgano fijo, sí es un órgano único. Cualquiera de las formaciones del Consejo puede adoptar actos del ámbito de otra, por lo que en ningún acto legislativo adoptado por el Consejo se menciona la formación. La Presidencia del Consejo cambia cada seis meses en un sistema de rotación entre los Estados miembros, a excepción del Consejo de Asuntos Exteriores, que preside de forma permanente el Alto Representante.

Cabe señalar que el Eurogrupo, que reúne a los m inistros de Economía y Hacienda de los Estados miembros de la zona del euro, tiene carácter informal y no puede tomar decisiones oficiales (como sí hace el ECOFIN), lo que no le resta relevancia política. El Eurogrupo está establecido en virtud del Protocolo 14 del TFUE para promover el crecimiento económico y la estabilidad financiera de la zona del euro mediante la coordinación de las políticas económicas. Cuenta con un Presidente elegido por sus miembros por un mandato de dos años y medio. Las principales funciones del Consejo son las siguientes:
Adopta la legislación europea sobre la base de las propuestas presentadas por la Comisión, en la mayor parte de los casos junto con el Parlamento Europeo con arreglo al procedimiento legislativo ordinario. En los casos en los que aplica un procedimiento legislativo especial (de aprobación o de consulta), el Consejo es el único legislador.
Adopta conclusiones, resoluciones y declaraciones, sin efectos jurídicos, que establecen compromisos o posiciones políticas.
Firma oficialmente los acuerdos entre la Unión y otros países. o con organizaciones internacionales. Estos acuerdos también están sujetos a la aprobación del Parlamento en ámbitos en los que posee poderes de codecisión. 
Coordina las políticas de los Estados miembros. En materia económica, se realiza tanto en el seno del ECOFIN como del Eurogrupo.

Para el buen cumplimiento de sus fu nciones. las reuniones del Consejo son preparadas por el Comité de Representantes Permanentes (COREPER) formado por los Embajadores que cada uno de los Estados miembros tiene al frente de su Representación Permanente ante la UE. El COREPER se reúne semanalmente y cuenta con dos formaciones (1 y II según los temas). Entre otras cosas, negocia las cuestiones que no han podido ser acordadas previamente en los Comités y Grupos de Trabajo.\footnote{%
    Los lista de órganos preparatorios de los trabajos del Consejo se acerca a los trescientos. Algunos de los Comités han sido creados por los propios Tratados (caso, por ej emplo. del Comité Económico y Financiero y el Comité de Política Comercial), otros por decisión intcrgubemamental. por un acto del Consejo. o incluso por el COREPER para tratar asuntos muy concretos.
} Además, existe una Secretaría General del Consejo que apoya tanto al Consejo y a su Presidencia rotatoria (facilitando la continuidad de los trabajos). como al Consej o Europeo y su Presidente.


\subsection{Comisión Europea}
La Comisión es una institución creada desde el origen mismo de las Comunidades Europeas en los años cincuenta, si bien se denominaba Alta Autoridad en la extinta Comunidad Europea del Carbón y del Acero.

La Comisión es la institución que se encarga de la gestión centralizada de los asuntos comunes a la Unión. Frente a la falsa creencia de medios de comunicación y políticos, la Comisión  Europea es una institución dotada de plena legitimidad democrática pues se forma al obtener la confianza mayoritaria del Parlamento Europeo tras las elecciones europeas, de forma similar a los gobiernos nacionales tras las elecciones generales. La Comisión es independiente de los Estados y encarna el interés general (art. 17.1 TUE).

Es una institución independiente de los Estados miembros, de las otras Instituciones de la Unión y de los intereses privados. La independencia es una cualidad exigida con carácter previo a su nombramiento: deben ofrecer «plenas garantías de independencia», pues los Tratados exigen que los miembros de la Comisión ejerzan sus responsabilidades con absoluta imparcialidad y en «interés general de la Unión» (art. 17.1 y 3 TUE). Aunque los miembros de la Comisión, obviamente, tienen que ser nacionales de los Estados miembros y son propuestos por los Gobiernos, ni individual ni colegiadamente los representan y tampoco pueden defender intereses de los Estados ni de los particulares (art. 245 TFUE). Esa independencia exige de los miembros de la Comisión que no soliciten ni acepten instrucciones de ningún Gobierno ni de ningún organismo ni realicen actos incompatibles con el carácter de sus funciones, lo que significa que no pueden ejercer ninguna actividad profesional, remunerada o no, durante su mandato. No dependen del mantenimiento de la confianza de sus Gobiernos ni del Consejo, que en ningún caso podrán cesarles. Únicamente podrán ser cesados de sus cargos por el Tribunal de Justicia, a instancia del Consejo o de la propia Comisión, si se constatase que un miembro de la Comisión no reúne ya las condiciones para el ejercicio de su cargo o si ha cometido una falta grave (arts. 245 y 247 TFUE). Aunque deberán conservar la confianza del Presidente de la Comisión (véase epígrafe 3).

Gozan de una amplia libertad de expresión para enjuiciar problemas europeos (aunque no deben hacer uso de informaciones que estén bajo secreto, art. 339 TFUE) o asuntos nacionales, expresándose bajo su responsabilidad personal, pues sólo las tomas de postura de la Comisión la comprometen en tanto que Institución.

Por otra parte, para reforzar esa independencia se les exige el compromiso solemne de respetar, en tanto permanezcan en sus funciones y aun después del cese de éstas, las obligaciones derivadas de su cargo, especialmente los deberes de honestidad y discreción en cuanto a la aceptación, después de este cese, de ciertas funciones o ciertas ventajas (art. 245 TFUE). Si alguno violase este compromiso estando en funciones, el Tribunal de Justicia, a instancia del Consejo o de la Comisión, podría obligarle a dimitir o cesarle y si ya hubiese cesado, retirarle la pensión. Los deberes que se imponen a los Comisarios son, a su vez, una obligación para los Estados miembros que se comprometen a respetar esa independencia y a no tratar de influir sobre los miembros de la Comisión en la realización de sus tareas (art. 245 TFUE). 

Otra de las características definitorias es la colegialidad. En palabras del Tribunal, el principio de colegialidad se basa en la igualdad de los miembros de la Comisión, en cuanto a la participación en la adopción de decisiones, e implica en particular que se delibere colectivamente sobre las decisiones y que todos sus miembros sean responsables en forma colectiva, en el plano político, de todas las decisiones adoptadas1. 

Las decisiones o acuerdos se adoptan por mayoría absoluta y la Institución asume la responsabilidad colectivamente: son acuerdos de la Institución en su conjunto y no individualmente de sus miembros.

Como los Tratados confían a la Comisión la iniciativa de toda  acción normativa, la gestión y el control del cumplimiento del Derecho de la Unión, parece coherente que la Comisión dé cuenta de tales responsabilidades ante el Parlamento Europeo. Por ello, la Comisión tiene que presentar ante el Parlamento un Informe General anual sobre las actividades de la Unión (art. 233 TFUE), dando lugar a un amplio debate y a una publicación continuada desde 1958. También presenta informes monográficos anuales o periódicos, entre otras materias, sobre la situación social, la política agrícola, la competencia leal entre las empresas y el control del cumplimiento del Derecho de la Unión (véase el capítulo sobre el Parlamento Europeo, epígrafe 4.2). Además, los parlamentarios pueden dirigir interpelaciones verbales y escritas a la Comisión en cualquier momento (art. 239 TFUE).

La Comisión asume responsabilidad política colegiada; en efecto, el Parlamento puede aprobar una moción de censura y hacer dimitir colectivamente a la Comisión (arts. 17 TUE y 234 TFUE), si bien, hasta su reemplazo, la Comisión dimisionaria continuará gestionando los asuntos corrientes (sobre la moción de censura, ver más ampliamente en el capítulo sobre Parlamento Europeo). No cabe responsabilidad política individual de un Comisario ante el PE. Ya se han señalado las condiciones (art. 16.6 TUE) para que, en caso necesario, el Presidente, previa autorización por mayoría del Colegio de comisarios, puede solicitar la dimisión de un Comisario y éste deberá dimitir («procedimiento Prodi»).

\subsubsection{Composición y funcionamiento}
La Comisión —desde 2004— está compuesta por tantos Comisarios como Estados miembros (28; en caso de consumarse el proceso de retirada del Reino Unido, 27). En ese número se incluye  a la Presidencia y al Alto Representante de la Unión para Asuntos Exteriores y Política de Seguridad, que es uno de sus vicepresidentes (art. 17.4 TUE).

Si bien sus miembros no representan a los Estados ni los intereses nacionales, comprende, pues, un nacional de cada Estado miembro. Los miembros de la Comisión Europea son propuestos y nombrados en razón de su competencia general, su compromiso europeo y han de ofrecer garantías plenas de independencia (art. 17.3 TUE).

El peligro de haber reducido su composición a un Comisario por país lleva inevitablemente a una relativa intergubernamentalización, poniendo en duda su tradición de supranacionalidad y la representatividad del interés en general en sus decisiones. Más europeísta e integrador hubiera sido que el número de Comisarios fuera siempre inferior al de Estados miembros mediante un número fijo y predeterminado, con una fórmula como la que sugiere para el futuro el artículo 17.5 TUE: un número de miembros igual a dos tercios del número de Estados miembros. Dicho de otro modo: un tercio de los Estados miembros no tendrían un nacional en la Comisión y habría rotaciones en estricto pie de igualdad para todos los Estados y la composición global deberá satisfacer la diversidad demográfica y geográfica del conjunto de los Estados miembros (art. 244 TFUE). Pero esta previsión ad futurum no tiene claros apoyos actuales. Tampoco a los Estados grandes les convence la rotación que les puede dejar sin su nacional durante una legislatura. Por ello, la rotación tiene limitadas posibilidades de llegar a producirse.

1.1. EL NOMBRAMIENTO DE LA COMISIÓN
La Comisión goza de un mandato de cinco años (art. 17.3 TUE), al igual que el de los miembros del Parlamento Europeo, de modo que su proceso de nombramiento se inicia tras cada renovación parlamentaria. Por otra parte, los Tratados toman algunas precauciones para que siempre se mantenga ese paralelismo en los mandatos del Parlamento y de la Comisión: se prevé que, en caso de moción de censura aprobada por el Parlamento Europeo contra la Comisión, los nuevos miembros de la Comisión solo desempeñarían su cargo durante el período que restase a la Comisión censurada y obligada a dimitir. De forma similar se resuelve en caso de fallecimiento, cese o dimisión de un Comisario (art. 246 TFUE).

Su nombramiento ha experimentado una evolución cada vez más democratizadora. En la práctica, desde la Declaración Solemne sobre la Unión Europea, adoptada en el Consejo Europeo de Stuttgart de 19 de junio de 1983, el nombramiento de la Comisión sigue un sistema bastante similar a la investidura de un Gobierno por un Parlamento. Esa práctica, basada en la necesidad de una doble y separada votación de confianza en el Parlamento europeo, fue incorporada a los Tratados de Maastricht y de Ámsterdam dotando al nombramiento del Presidente de una superior base jurídica, mayor transparencia y legitimidad democrática.

A tenor de las modificaciones introducidas por los Tratados citados y el Tratado de Lisboa, el procedimiento es el siguiente (art. 17.7 TUE):
— Tras cada elección al Parlamento Europeo y a la vista de los «resultados electorales» (y no de la lista más votada en la UE), el Consejo Europeo, por mayoría cualificada (frente al común acuerdo, anterior a Niza) propondrá la personalidad candidata a nombrar como Presidente o Presidenta de la Comisión. Claro que esta politización o parlamentarización de la Comisión, al hacer depender el nombramiento de la presidencia de la mayoría parlamentaria o de los partidos mayoritarios, puede ocasionar alteraciones en el futuro sobre la independencia de la Comisión. La Comisión es un ejecutivo en el sentido de una administración federal, representa la gestión centralizada de los asuntos comunes, pero no es un Gobierno, no es un poder político, sino un poder moderador e independiente. Su excesiva dependencia de los dictados de la mayoría parlamentaria puede quebrar su independencia, la defensa de los intereses generales y afectar a su papel motor.
— El pleno del PE decidirá por mayoría de los miembros que lo componen si aprueba o no, con carácter vinculante, la propuesta del Consejo en torno al candidato a la Presidencia. En caso de no aceptar, el Consejo Europeo tendrá que proponer otra candidatura. 
— El nuevo Presidente inicia consultas con los Gobiernos a fin de proponer a las demás personalidades que formarán la Comisión. Los miembros de la Comisión serán propuestos por mayoría cualificada por el Consejo, de común acuerdo con el Presidente designado. La persona nombrada por el Consejo Europeo —con la  aprobación del Presidente de la Comisión— como Alto Representante de la Unión para Asuntos Exteriores y Política de Seguridad asumirá la vicepresidencia de la Comisión. Cuando los Gobiernos han alcanzado un acuerdo sobre el resto de miembros de la Comisión, el Parlamento invita a cada una de las personalidades propuestas, en función de sus competencias previsibles, a comparecer ante las comisiones parlamentarias pertinentes y éstas informan al Presidente de sus conclusiones sobre la capacidad de los propuestos para los cargos. Este examen del Parlamento Europeo se inspira en el sistema de audiencias (hearings) del Senado norteamericano y no está previsto directamente en los Tratados pero sí en su Reglamento interno.
— Todos los designados, incluido la Alta Representante (hasta ahora siempre han sido mujeres) y el Presidente, se someterán colegiadamente a un voto de aprobación del Parlamento Europeo (por mayoría de los votos emitidos). Entonces, ya podrán ser formalmente nombrados por el Consejo Europeo por mayoría cualificada.


\paragraph{Presidencia}
El Presidente de la Comisión marca la orientación política y decide la organización interna a fin de garantizar la coherencia, la eficacia y la colegialidad de su acción (arts. 17.6 TUE y 248 TFUE). Debe gozar, pues, de amplia libertad para conferir funciones dentro  del Colegio de Comisarios pero el respeto a estas previsiones  introducidas por los Tratados de Ámsterdam y de Niza depende siempre de la personalidad del Presidente.

Se reconocen poderes al Presidente para estructurar y repartir las responsabilidades tanto al inicio del mandato como su reordenación a lo largo del mismo, recordando a los Comisarios —a veces aquejados de un fuerte afán de protagonismo— que ejercerán sus funciones bajo la autoridad del Presidente (art. 248 TFUE). También puede nombrar Vicepresidentes distintos del Alto Representante (art. 17.6 TUE).

Además, el Tratado, desde la reforma de Niza, recogió el denominado «procedimiento Prodi», o compromiso de dimisión del comisario al que se lo solicite el Presidente (art. 17.6 TUE), si bien en el caso del Alto Representante y Vicepresidente de la Comisión se requiere además el acuerdo del Consejo Europeo (art. 18.1 TUE).

La figura del Alto Representante, que es al tiempo vicepresidente de la Comisión y Presidente del Consejo (de ministros) de Asuntos Exteriores, introduce ciertas disfunciones en la colegialidad e independencia de la Comisión. Como Alto Representante obedece a las instrucciones del Consejo Europeo. El Presidente de la Comisión puede pedir unilateralmente la dimisión de cualquier Comisario pero no la del Alto Representante para la que precisará el acuerdo del Consejo Europeo. Si hubiera una moción de censura le afectará ineluctablemente, aunque en caso de censura triunfante el Alto Representante sólo cesará en sus funciones de Comisario, rompiendo de hecho la colegialidad de la Comisión. Se incrusta, por primera vez, un elemento claramente intergubernamental en el seno del colegio mismo de la Comisión.

La Comisión aprueba su propio reglamento interno2. Sin perjuicio de su carácter colegial, la Comisión distribuye los sectores de actividad —salvo las relaciones exteriores que se asumen por el Alto Representante— entre los distintos Comisarios: esta técnica de departamentalización significa que cada Comisario se encarga de varios ámbitos, no siempre homogéneos ni comparables con las tradiciones gubernamentales de las «carteras ministeriales». Ahora bien, su carácter colegial se ve atenuado al prever en su reglamento interno la posibilidad de adoptar decisiones, al margen del debate y aprobación en sesión formal, mediante procedimiento escrito o por habilitación o por delegación.

Se sigue el procedimiento escrito para aligerar el orden del día de las reuniones de asuntos poco controvertidos (como medidas de aplicación o de ejecución de actos ya adoptados en sesiones anteriores). El acto se prepara por el Comisario que tiene la competencia en ese ámbito, aunque con conocimiento del resto de los gabinetes de los Comisarios y de los distintos servicios de la Comisión. De todos modos, cualquier miembro de la Comisión puede pedir que el texto sometido al procedimiento escrito sea  objeto de debate en una sesión de la Comisión.

Además, la Comisión puede habilitar a un Comisario para que adopte medidas de gestión o de administración claramente delimitadas. El Tribunal de Justicia ha reconocido la legalidad del sistema de habilitaciones en favor de un determinado Comisario; «las Decisiones adoptadas por habilitación deben ser consideradas como Decisiones de la Comisión» de forma que «no vulneran el principio de colegialidad»3.

Y cabe que la Comisión delegue, asumiendo siempre su responsabilidad colegiada, la adopción de medidas de gestión o de administración en los Directores Generales y Jefes de Departamento dentro de los límites y condiciones que establezca.

De la Comisión depende una vasta infraestructura administrativa (aproximadamente unos 24.000 empleados entre funcionarios y empleos temporales, excluidos los del Servicio Europeo de Acción Exterior) compuesta por unos Servicios generales e internos (Secretaría General, los Gabinetes de los Comisarios, servicio jurídico, oficina de estadística, oficina europea de lucha antifraude, oficina de publicaciones, comunicación, servicios de traducción e interpretación, etc.) y una veintena de direcciones generales dedicadas a las Políticas y a las Relaciones Exteriores.

\subsubsection{Atribuciones}
La Comisión tiene asignada una misión esencial: promover el interés general de la Unión (art. 17 TUE). Para ello asume importantes atribuciones.

El papel de la Comisión Europea no ha sufrido variaciones sensibles tras las reformas introducidas en el Tratado de Lisboa.

Mantiene su cuasi-monopolio de iniciativa normativa, si bien comparte más casos con otros actores como los Estados miembros u otras instituciones.

Conserva sus poderes de guardiana de los Tratados y los ve ligeramente reforzados por el artículo 260.3 TFUE que le faculta para proponer al Tribunal de Justicia una sanción al Estado infractor en el propio recurso por incumplimiento en caso de no informar sobre las medidas de transposición de las directivas. Mantiene los poderes de gestión y decisión, incluida la ejecución del presupuesto. Además, esos poderes se ven facilitados al prever la delegación de capacidad legislativa del Parlamento y del Consejo en favor de la Comisión —actos delegados, art. 290 TFUE—. Y puede seguir recibiendo poderes de ejecución de los Estados— actos de ejecución— (art. 291 TFUE).

Sin embargo, la Comisión ha perdido peso político en el sistema institucional en las recientes reformas y se ha debilitado tras las ampliaciones del siglo XXI, aunque también por la desigual aptitud de sus miembros y la falta de autoridad de sus últimos Presidentes. El papel de la Comisión debe seguir siendo importante y decisivo para evitar la parálisis o los enfrentamientos o su dependencia de facto del Consejo Europeo y del Consejo, pero posiblemente su forma de ejercer su papel de motor no deba ser la misma que en el pasado. La Comisión tiene que encontrar su papel en el nuevo  contexto de la Europa ampliada, del ascenso del Consejo Europeo y del propio Parlamento Europeo y marcar un liderazgo sobre la base de grandes objetivos comunes (cambio climático, aprovisionamientos energéticos y de materias primas, innovación tecnológica...).

5. DERECHO DE INICIATIVA NORMATIVA
5.1. SIGNIFICADO Y ALCANCE DE LA PROPUESTA
Uno de los poderes que singularizan a la Comisión es, sin duda, su poder de iniciativa normativa. Mediante el derecho de propuesta, casi general, la Comisión participa en la formación de los actos legislativos y no legislativos de la Unión (arts. 17.2 TUE y 289 TFUE) e impulsa la acción normativa. 

Hay excepciones a este derecho de propuesta, pero son casos  específicos y detallados en los Tratados en los que se confía la iniciativa a un grupo de Estados miembros, al Parlamento Europeo,al Banco Central Europeo o al Tribunal de Justicia o al Banco europeo de inversiones (art. 289.4 TFUE).

El derecho de iniciativa revela que se ha confiado en la Comisión la concepción y la preparación de la legislación de la Unión y que le corresponde, en la inmensa mayoría de los casos, aportar en la concepción de las propuestas una visión de conjunto inspirada en el interés de la Unión; en consecuencia, asume el papel de órgano «motor» de la Unión.

El procedimiento legislativo ordinario, consistente en la adopción conjunta por el Parlamento Europeo y el Consejo, exige con carácter general la propuesta previa de la Comisión (art. 289 TFUE). Igualmente, una amplia mayoría los actos no legislativos y, obviamente, la Comisión elabora la propuesta de los de carácter delegado y ejecutivo que ella misma decide.

Este derecho significa:

1) que por voluntad expresa de los Tratados (en definitiva, de los Estados mismos), la gran mayoría de los actos legislativos del Parlamento Europeo y del Consejo, así como los actos no legislativos del Consejo sólo pueden ser adoptados si ha habido la correspondiente propuesta normativa de la Comisión;

2) que en la gran mayoría de los casos, el Consejo puede aprobar o rechazar el acto propuesto por mayoría cualificada, pero si quiere enmendarlo o modificarlo, incluso a instancia del Parlamento Europeo, necesita la unanimidad (art. 293 TFUE). Es el privilegio de la Comisión.

La razón de ser de este privilegio reside en que la Comisión encarna el interés general del conjunto y se presume que obra con independencia de todo interés nacional. De ahí que, como ha señalado J. V. Louis, el derecho de propuesta de la Comisión sea un contrapeso al poder legislativo y normativo del Consejo y del Parlamento Europeo.

Ese privilegio es una garantía de protección de los intereses generales, especialmente de los intereses de los Estados medios y pequeños frente a las grandes potencias de la UE. Al no poder modificar la propuesta de la Comisión más que por unanimidad (salvo en los casos previstos en el art. 293.1 TFUE), se parte de la presunción jurídica de que tiene en cuenta el interés general, y que éste sólo puede ser desplazado cuando todos los Estados demuestran con la unanimidad la existencia de una alternativa alinterés general definido por la Comisión. El interés general es una idea-fuerza que entronca y se origina en la legalidad de la Unión: es la actuación de un sistema dinámico que deriva de las finalidades, principios y reglas establecidos directa e indirectamente en los Tratados.

Es cierto que el privilegio de la Comisión se justificaba plenamente en la etapa anterior a la reforma del Tratado de Maastricht al disponer el Consejo hasta entonces de todo el poder de decisión. Pero al haberse reformado para que Parlamento y Consejo compartan el poder legislativo, ese privilegio ha ido sufriendo ciertas limitaciones lógicas y derivadas del mismo procedimiento legislativo ordinario. Los supuestos en los que no sería preciso que el Consejo se pronunciara por unanimidad cuando hay una modificación a la propuesta de la Comisión son los siguientes:

— cuando se recurre al Comité de Conciliación en la segunda y tercera lectura en el procedimiento legislativo ordinario (art. 294.10 y 13 TFUE);
— en la adopción del presupuesto anual, del marco financiero plurianual y prórroga del presupuesto (arts. 310, 312, 314 y 315, párrafo segundo, TFUE). 

Excepcionalmente, siempre ha habido situaciones en las que la iniciativa era exclusiva del PE, por ejemplo, la propuesta sobre la composición del PE, el procedimiento electoral, el estatuto de los diputados, comisiones de investigación o el estatuto del defensor del pueblo europeo. O se compartía con los Estados miembros y las otras instituciones afectadas en materia de revisión de los tratados o de sus estatutos respectivos. Finalmente, se observa que, tras la reforma de Lisboa, la Comisión comparte la iniciativa en algunos ámbitos, por ejemplo, con los Estados miembros (en materia de cooperación judicial penal y policial y en materia de ayudas públicas), o con la Alta Representante en materia de PESC-PESD. La Comisión ejerce este derecho de propuesta de forma discrecional con carácter general4.


La Comisión sólo puede presentar una propuesta normativa si hay una base jurídica que reconozca la competencia de la UE y su derecho de propuesta o de iniciativa (principio de atribución de la competencia). Cuando la Comisión estima que debe actuar en un determinado ámbito, intervienen en esa decisión criterios de oportunidad, de necesidad, o de urgencia y, como es lógico, el mandato que, en cada caso, se prevea en los Tratados. En esta fase, en la que se forma la «idea» de responder a la realización de objetivos de los Tratados, puede decirse que la Comisión actúa en solitario, con independencia, y se presume que actúa en interés general de la Unión.

La Comisión y sus servicios técnicos tienen canales de información y muchos contactos con medios nacionales y europeosdesde donde puede haber surgido alguna sugerencia, y si ésta está fundada en razones objetivas, la Comisión no es insensible a tales sugerencias. Pero aun en estos casos, cuando desde fuera de la Comisión surge el estímulo para que actúe, la Comisión tamiza esas invitaciones o consejos y acuerda ponerse en acción desde la reflexión interna y con independencia.

Cuando considera que debe elaborar una propuesta normativa, las Direcciones Generales afectadas preparan informes o estudios de base, solicitan datos y estadísticas, hacen estudios de legislación comparada, etc.; si se trata de propuestas de modificación de normas de la UE, se tienen en cuenta las consecuencias de la legislación objeto de cambio, o se apuntan las eventuales repercusiones, etc.

En esta fase puramente interna en la que se idea la propuesta, el método de trabajo en la Comisión es fundamentalmente «funcionalista». Son los servicios técnicos de la Comisión, los que inician las discusiones y los que identifican los problemas. Este método funcionalista se caracteriza por un profundo racionalismo y utilitarismo: se identifican unos intereses comunes, se fija un objetivo común y se comparte la necesidad de una acción común. El estudio del «problema-solución» es progresivo y los acuerdos se van intensificando conforme el «problema-solución» va ascendiendo desde los grupos de base de los servicios técnicos a los escalones superiores de la Comisión.

\begin{specialblock}{\textbf{Influencias en el proceso legislativo de la Comisión}}
    6.1. LA INFLUENCIA DE LOS EXPERTOS NACIONALES
Con frecuencia, a la par que se elaboran estudios por los servicios de la Comisión, o bien ya al término de una reflexión de conjunto sobre la futura propuesta, la Comisión, a través de sus servicios, entra en contacto con las administraciones nacionales y convoca a los expertos propuestos por éstas. Los funcionarios expertos acuden a estas reuniones, no en su calidad de funcionarios nacionales unidos por un vínculo de fidelidad a su administración, sino en calidad de «expertos» sin representación nacional y sin instrucción o mandato.

Estos encuentros con los expertos nacionales son muy útiles para la Comisión. Obtiene datos más precisos, conoce con más detalle la normativa y la práctica administrativa de cada Estado en la materia en cuestión. La Comisión, por sí sola, no estaría ensituación de abarcar pormenorizadamente las necesidades y los problemas de cada Estado miembro o de una región determinada, así como las diferencias de tradiciones o de concepciones. La Comisión reconoce que, gracias a esas consultas, se asegura una  completa información sin estar necesariamente vinculada por las conclusiones obtenidas.

Además de esas ventajas inmediatas en la elaboración de la propuesta, la Unión se beneficia porque, gracias a esa coordinación de las actividades de la Comisión con los servicios de las Administraciones nacionales, ve indirecta y eficientemente aumentada su capacidad administrativa. La Comisión tiene la oportunidad de observar las eventuales dificultades que inicialmente puede encontrar en algún Estado, el grado de aceptación o las reticencias que pueden surgir.

Para las Administraciones nacionales estos encuentros son una suerte de preaviso del contenido de la futura propuesta y ya les alerta para ir preparando la posición nacional en el seno del Consejo y tratar de influir discretamente en este estadio inicial. En cualquier caso, esta asociación a los trabajos preparatorios permite mantener unas estrechas relaciones entre las Administraciones nacionales y la de la Unión.
6.2. LA INFLUENCIA DE LOS GRUPOS SOCIO-ECONÓMICOS
Con vistas a la elaboración de una propuesta o a perfilar la orientación de un sector, la Comisión también entra en contacto con los medios sociales, económicos y profesionales afectados. A veces, la Comisión, apenas ha pergeñado un proyecto, lo lanza al debate público, o simplemente lo suscita mediante Libros Verdes repletos de preguntas y alternativas en los que se estimula la reflexión de la sociedad haciendo diversos planteamientos con alternativas, cuestionarios y direcciones para remitir la opinión ciudadana. Las informaciones y las inquietudes que transmiten los grupos económicos y sociales revisten gran importancia para la Comisión pues desea conectar su acción con tales fuerzas. Por otra parte, la Comisión sabe que si su futura propuesta goza de un nivel de aceptación en esos medios tendrá un mayor peso en el momento de la deliberación en el Consejo. 

La facilidad del acceso a los servicios de la Comisión es un aliciente para las asociaciones socio-económicas nacionales y transnacionales y para los «lobbistas» en general: el momento más adecuado y eficaz de influir en una determinada orientación normativa (o en la adopción de medidas de defensa comercial o ser incorporado a líneas de financiación con fondos de la Unión) es el estadio inicial de elaboración. Cuanto más esperen a influir en un proyecto, más elaborado e impermeable será a sus intereses. Estos grupos de interés económico o socio-profesional se han organizado a escala de la Unión y gozan de una gran fuerza que despliegan como grupos de presión, especialmente, ante la Comisión y también ante el Parlamento y el Comité Económico y Social. La Comisión, junto con el Parlamento Europeo, ha establecido un registro público de estos grupos o lobbies con todos sus datos6.

Estos encuentros son vistos, en ocasiones, con algún recelo por parte del Consejo. Pero cuando la Comisión entra en contacto con los medios socio-económicos o con los expertos, afirma que evita formalmente entrar en una negociación de sus iniciativas, no utiliza «borradores» de propuestas, sino que intercambian puntos de vista sobre un problema o sector en general, tratando de obtener las mejores informaciones. La Comisión niega que solicite de los interlocutores socio-económicos una negociación o un apoyo. Este método de apertura y permeabilización de la Comisión ha sido aceptado desde hace tiempo y confirmado por el Acuerdo interinstitucional entre el PE, el Consejo y la Comisión (2016) sobre la mejora de la legislación. En efecto, se acepta que la Comisión, antes de la adopción de su propuesta, efectuará consultas públicas de forma abierta y transparente, garantizando que las modalidades y plazos de dichas consultas públicas, permitan una participación lo más amplia posible7.

En síntesis, en esta fase de elaboración de la propuesta, la Comisión concilia un método funcionalista de análisis racional y técnico en la aproximación a los objetivos comunes, una influyente presión transnacional y una inquietante permeabilidad a los elementos intergubernamentales.
\end{specialblock}


5.2. MODIFICACIONES DE LA PROPUESTA POR LA COMISIÓN 
La Comisión es «dueña» de su propuesta, pues puede modificarla en todo momento (con algunas salvedades referidas al procedimiento legislativo ordinario, art. 294.10 y 13 TFUE) y, llegado el caso, retirarla con adecuada motivación5. El hecho de que la Comisión pueda modificar su propuesta, en tanto no haya decisión del Consejo, significa que la Comisión dispone de un verdadero poder de negociación con el Consejo y con el Parlamento Europeo, articulándose en torno a este derecho de iniciativa el llamado diálogo triangular Comisión-Consejo-Parlamento Europeo (art. 295 TFUE). Con esa relación estrecha entre propuesta-decisión puede comprenderse fácilmente que la colaboración triangular se materialice en un diálogo permanente mediante «consultas recíprocas» entre tales Instituciones (art. 295 TFUE).

Los Tratados dan una gran fuerza a la Comisión en los procedimientos decisorios que tienen como base su iniciativa legislativa. En efecto, la propuesta está privilegiada en los mecanismos de votación en el Consejo ya que se puede adoptar la propuesta por mayoría cualificada, pero sólo se puede enmendarpor unanimidad en el Consejo, a lo que se añade la facultad de la Comisión de modificarla en todo momento (salvo en la fase del comité de conciliación) y de retirarla. Igualmente, las enmiendas del Parlamento serán más fácilmente aceptadas por el Consejo si la Comisión las hace suyas; por ello, se requiere que haga un informe tras cada participación del Parlamento y del Consejo.

Este poder de negociación y de conciliación en el seno del Consejo y del Parlamento europeo puede favorecer la formación demayorías y desbloquear el Consejo cuando éste no reúne la mayoría necesaria para aprobar ni la unanimidad para modificar. En esos casos, sólo la Comisión, velando por el interés general, puede hacer los cambios que estime aceptables para lograr esa mayoría  suficiente, si estima que una concesión a tiempo —sin afectar a los principios y al interés general— permiten formar la mayoría cualificada o, a veces, lograr una mayoría más amplia o el consenso de todos.

5.3. LOS MOTORES SUBSIDIARIOS

El Parlamento Europeo y el Consejo no pueden sustituir a la Comisión en su derecho de propuesta, aunque sí pueden pedirle que le someta todas las propuestas adecuadas (arts. 225 y 241 TFUE). Ahora bien, si la Comisión no ejerce su derecho de propuesta, además de explicar las razones de su negativa, el Consejo o el Parlamento Europeo pueden interponer ante el Tribunal de Justicia un procedimiento por omisión (art. 265 TFUE).

Por otra parte, el Parlamento puede utilizar, proporcionadamente, los instrumentos parlamentarios de control, tales como las interpelaciones (art. 230 TFUE) o la moción de censura, si estima que la falta de diligencia de la Comisión hace merecer su dimisión (art. 234 TFUE).

Hay que tener en cuenta que el derecho de iniciativa (para actos legislativos y no legislativos) exige que la Comisión no se desentienda de la misma, que defienda su propuesta en todas las instancias donde va a ser debatida (Parlamento, Comité de las Regiones, Comité Económico y Social). Por ello, tiene derecho a estar presente en todas las sesiones de las Instituciones, órganos auxiliares y comités o grupos de trabajo en los que se delibere sobre su propuesta.


7. ATRIBUCIONES DE DECISIÓN NORMATIVA NO
LEGISLATIVA
La Comisión Europea no dispone de ningún poder de carácter legislativo ni tampoco del poder general de ejecución de los actos vinculantes de la Unión que corresponde a los Estados miembros (art. 291 TFUE).

La Comisión tiene asignadas funciones de coordinación, ejecución y gestión estrictamente fijadas por los Tratados para las que se le atribuye poder de decisión normativo no legislativo. El desarrollo de los Tratados por la Comisión puede ser una atribución propia, es decir, reconocida directamente por los Tratados y da lugar a actos no legislativos en la forma de reglamentos, directivas o decisiones de la Comisión. Además también puede adoptar actos no legislativos: — bien por atribución del Parlamento Europeo y del Consejo cuando le delegan aspectos parciales y limitados de poder legislativo, dando lugar a reglamentos delegados, directivas delegadas y decisiones delegadas (art. 290 TFUE); — bien por atribución mediante actos de la Unión cuando se requieran condiciones uniformes de ejecución. Adoptará, entonces, reglamentos de ejecución, directivas de ejecución y decisiones de ejecución (art. 291 TFUE).

7.1. LOS PODERES DE EJECUCIÓN Y GESTIÓN ATRIBUIDOS POR LOS
TRATADOS
El poder de decisión propio de la Comisión, es decir, por atribución directa del Tratado, está muy limitado, caso por caso, y es de naturaleza ejecutiva o reglamentaria. A tal fin, puede adoptar reglamentos, directivas y decisiones. Hay que tener en cuenta que si el poder de decisión autónomo no es intenso o esencial, sin embargo, debido a los ámbitos económicos y sociales tan extensos e importantes a los que afecta, sus atribuciones son importantes, aún siendo siempre de naturaleza ejecutiva.

La Comisión se presenta con una clara vocación de «administración» que asume la gestión de la UE. La Comisión, como gestora de la Unión, adopta actos de naturaleza reglamentaria y administrativa, tanto de alcance general como individual.\footnote{
Entre otros:
-La Comisión también está encargada de gestionar importantes sectores de la actuación de la Unión a través de la administración de los fondos estructurales, como es el caso del Fondo Social Europeo (FSE), el Fondo Europeo de Desarrollo Regional (FEDER), el Fondo Europeo de Garantía Agraria (FEGA), el Fondo Europeo de Desarrollo (FED), el Fondo de Cohesión, etc., si bien está asistida por diversos comités intergubernamentales. — Se le confía el poder de ejecución del presupuesto de la Unión (art. 317 TFUE).
}


7.2. PODERES DE DECISIÓN POR DELEGACIÓN O POR ATRIBUCIÓN
El Tratado prevé la delegación en la Comisión de poderes para adoptar actos no legislativos de alcance general que completen o modifiquen determinados elementos no esenciales del acto legislativo (art. 290 TFUE). También prevé la posibilidad de conferirle competencias de ejecución de actos jurídicamente vinculantes (legislativos y no legislativos) (art. 291 TFUE).

7.2.1. Por delegación del Parlamento y del Consejo (art. 290 TFUE) 
Se trata de materias que competen al poder legislativo y que podrían ser reguladas exhaustivamente por el Parlamento y el Consejo, pero que renuncian a hacerlo y delegan su poder en la Comisión (art. 290 TFUE)8. La finalidad es agilizar el proceso normativo, de modo que el Consejo y Parlamento Europeo se concentren en los aspectos esenciales de la legislación (cuestiones con vocación de permanencia o largo plazo), mientras que los aspectos más técnicos, o temporales, se desarrollarán por el poder ejecutivo. También se busca facilitar en determinados aspectos la adaptación técnica y circunstancial del acto legislativo. Se parece pero no es el sistema tradicional de atribución de competencias de ejecución a favor de la Comisión.

El legislador (Parlamento y Consejo) fijará en el acto legislativo de base (un reglamento o una directiva o una decisión) los objetivos, contenido y alcance de la delegación. No tienen la obligación jurídica de atribuir sus poderes en favor de la Comisión; es una opción si bien sólo podrán delegar ese poder en la Comisión. Y se podrán establecer mecanismos de control de la delegación: a) pueden revocarlo, b) vetar el acto concreto, y c) fijarle una duración.
a) La revocación: Sin duda el más importante y definitorio es que el PE o el Consejo pueden revocar la delegación separadamente, recuperando la competencia legislativa. La revocación se adoptará por mayoría de los miembros en el Parlamento Europeo y mayoría cualificada en el Consejo.
b) Veto: Parlamento y Consejo pueden impedir, por decisión autónoma y separada, la entrada en vigor del acto delegado. Expresan su confirmación mediante el silencio dejando pasar el plazo o bien pueden formular objeciones en dicho plazo.
c) Duración: Siempre cabe su prórroga. La fijación de su duración
(sunset clause) significará que las disposiciones del acto delegado
tendrían una duración limitada; una vez pasado este plazo, el
legislador podría renovar la habilitación.
Los actos así adoptados se denominan actos «delegados»
(reglamento delegado, directiva delegada, decisión delegada).
7.2.2. Poderes conferidos por los Estados miembros (art. 291.2
TFUE)
Este poder ha experimentado una importante evolución jurídica. Hasta el Acta Única Europea en 1987, el Consejo podía atribuir, caso por caso, poderes de decisión de naturaleza ejecutiva y sólo en favor de la Comisión (antiguo art. 155.4 TCE). El Acta Única Europea estableció la obligación general de atribuir a la Comisión las competencias de ejecución (antiguo art. 202 TCE).

Por el contrario, la nueva redacción del Tratado de Funcionamiento parte de que la ejecución de los actos legislativos y no legislativos corresponde a los Estados miembros (art. 291.1 TFUE). Ahora bien, si se necesitaran condiciones uniformes de ejecución, los propios actos de la Unión podrán conferir competencias de ejecución en la Comisión y, excepcionalmente, en el Consejo cuando se trate de materia de Política Exterior y de Seguridad Común (PESC, arts. 24 y 26 TUE). Estos actos se denominarán actos «de ejecución» (reglamento de ejecución, directiva de ejecución…).

La naturaleza de los poderes que pueden ser atribuidos a la Comisión es únicamente de ejecución y no de naturaleza legislativa.

El Tribunal de Justicia ha estimado, en relación con el antiguosistema de atribución de competencias, que la noción de ejecución comprende tanto la elaboración de normas de aplicación como la aplicación de las normas a casos particulares mediante actos de alcance individual9.

Ni el Parlamento Europeo ni el Consejo pueden revocar las competencias de ejecución conferidas (sólo podría revocar un acto concreto cuando concurran determinadas circunstancias). Las modalidades de control se regulan en un reglamento adoptado en procedimiento legislativo ordinario (Reglamento «Comitología»)10.

El control consiste en la consulta a comités que seguirán en unos casos un procedimiento consultivo y en otros un procedimiento de examen.

El procedimiento de examen se aplica para actos de alcance general destinados a la ejecución de actos de base y de actos de ejecución que puedan tener una repercusión importante, además de los relativos a la política agrícola, pesquera, comercial y fiscalidad así como programas con implicaciones presupuestarias importantes o dirigidos a terceros países; su informe favorable se adoptará por la mayoría cualificada prevista en el Tratado para el Consejo (art. 16.4 y 5 TUE) y si fuera desfavorable requiere la modificación del acto por la Comisión o acudir a un comité de apelación. De manera excepcional, la Comisión puede adoptar un acto de ejecución, incluso en el caso de dictamen desfavorable de un comité, cuando tal acto sea necesario para evitar perturbaciones significativas en los mercados, en el sector de la agricultura o un riesgo para los intereses financieros de la Unión.

En el resto de los casos debe recurrirse al procedimiento consultivo (informe por mayoría simple y no vinculante para la Comisión).

La composición de estos comités es siempre la misma:representantes de los Estados miembros (funcionarios), ponderación del voto conforme a las reglas del Consejo (art. 205 CE) y presidencia por un funcionario de la Comisión11.



8. ATRIBUCIONES DE CONTROL DEL CUMPLIMIENTO DEL DERECHO
El artículo 17.1 TUE confía a la Comisión la vigilancia del cumplimiento de las disposiciones de los Tratados y del Derecho derivado. Ciertamente, la Comisión no tiene un poder exclusivo de control del cumplimiento del Derecho de la Unión, pero su misión de control o «guardiana» del Derecho de la Unión tiene unos caracteres propios que le diferencian del control judicial que, por ejemplo, lleva a cabo el Tribunal de Justicia, o del control político del Parlamento Europeo, o del control financiero ejercido por el Tribunal de Cuentas. En el campo de la Unión Económica y Monetaria los poderes de supervisión y control se confieren al Consejo y al Banco Central Europeo (arts. 127 y 132 TFUE)12. También los Estados miembros deben controlar, a través de sus servicios administrativos y de inspección, la correcta aplicación del Derecho de la Unión.

Para desempeñar esta atribución, la Comisión está dotada de medios adecuados tanto desde el punto de vista de infraestructura administrativa y canales oficiales u oficiosos de información como de instrumentos jurídicos para conocer y exigir la aplicación de las normas de la Unión. Esta misión exige una gran dedicación, una minuciosa labor de vigilancia e información, una sabia combinación del uso de su responsabilidad como guardiana de los Tratados y de sus dotes de flexibilidad y prudencia a la hora de exigir el respeto del Derecho al eventual infractor, especialmente si es un Estado miembro.

Debido a la multiplicidad de los ámbitos de competencia de la Unión, a la variedad y la sutilidad de las formas que puedan revestir las violaciones del Derecho de la Unión, esta tarea es uno de los aspectos más importantes dentro del conjunto de las actividades de la Comisión. Y ello es así porque esa misión de vigilancia se ejerce tanto sobre los comportamientos de los Estados como, en algunos casos, sobre los particulares

8.1. EL CONTROL SOBRE LOS ESTADOS MIEMBROS
El fundamento jurídico general de este poder de control se prevé en el artículo 17.1 TUE. De modo específico, en virtud del principio de cooperación leal, los Estados miembros se comprometen a facilitar a la Unión el cumplimiento de sus misiones (art. 4.3 TUE). Además de estos dos preceptos, reconocidos por el Tribunal de Justicia como marco jurídico del deber de informar, otras disposiciones de los Tratados exigen para casos concretos que los Estados miembros tengan informada a la Comisión.

El procedimiento por incumplimiento se desenvuelve en dos fases consecutivas, primeramente, en una fase administrativa (o precontenciosa) y, después, la fase contenciosa ante el Tribunal de Justicia. Nos interesa ahora la fase administrativa.

La Comisión tiene el deber de dirigirse al Estado miembro si tiene conocimiento de que un Estado miembro ha incumplido susobligaciones de la Unión: no ha transpuesto una directiva, no respeta o no cumple obligaciones derivadas del Tratado o del Derecho derivado, mantiene en vigor normas nacionales contrarias que impiden o pueden impedir el cumplimiento del Derecho de la Unión, etc. A tal fin le envía una carta de emplazamiento (denominada también requerimiento) en la que denuncia el retraso en el cumplimiento, haciéndole las observaciones pertinentes y dando al Estado miembro un plazo para que conteste ofreciendo explicaciones de su conducta o rectificando la misma. Cuando el Estado miembro no contesta a la carta de emplazamiento o, en el caso de que responda, sus explicaciones no den satisfacción a la Comisión y persista la infracción, entonces, esa fase preliminar se completa con la adopción por la Comisión de un dictamen motivado en el que conmina al Estado infractor a poner término al origen de la infracción en un plazo dado. En el dictamen motivado se tiene que fijar un plazo para el cumplimiento de la obligación denunciada. 

[...]

Mediante estas dos iniciativas de la Comisión, la carta de emplazamiento y el dictamen motivado, se trata de dar ocasión al Estado miembro de justificar su posición formulando adecuadamente las alegaciones que, en su defensa, estime pertinentes frente a las imputaciones de la Comisión —como así lo ha señalado el Tribunal de Justicia— y, llegado el caso, permitirle conformarse voluntariamente a las exigencias del tratado. En el caso de que este esfuerzo no se corone con el éxito, el dictamen motivado tiene por función definir el objeto del litigio21. Al ser un  procedimiento previo, no produce ningún efecto jurídico vinculante frente al destinatario del dictamen motivado.

Una vez transcurrido el plazo marcado por la Comisión para la eliminación voluntaria del hecho ilícito por el Estado mismo, la Comisión puede demandar al Estado miembro ante el Tribunal de Justicia, fundada en el artículo 258 TFUE, iniciándose así el recurso por incumplimiento. Las decisiones de emitir un dictamen motivado y de interponer un recurso no son medidas de administración o gestión y no pueden ser objeto de delegación en ningún Comisario, aplicándoseles estrictamente el principio de colegialidad23. Por otra parte, no hay un derecho de acceso del público a los requerimientos y dictámenes motivados enviados por la Comisión a los Estados miembros en el marco de los procedimientos de infracción24.

Si la Comisión solicita, además, al Tribunal la imposición de una multa coercitiva en el marco del procedimiento por incumplimiento del artículo 260, ya sea por inejecución de sentencias o porincumplimiento de la obligación de informar sobre las medidas de transposición, deberá indicar el importe que considere adecuado a las circunstancias

8.2. EL CONTROL SOBRE LOS PARTICULARES
El poder de investigación y control puede dirigirse sobre los particulares, principalmente empresas asentadas en la Unión. Aunque el fundamento jurídico general radica en el artículo 337 TFUE, sin embargo las actuaciones de la Comisión se han basado siempre en disposiciones particulares relacionadas con los sectores energético, transportes, competencia y aquellos sectores que tienen contingentada la producción (en los que la Comisión tiene poderes de intervención directa).

En el marco del derecho de la competencia leal entre empresas, la Comisión hace sus investigaciones y recibe denuncias de otros particulares. En una primera fase, tras la presentación, la Comisión recaba los datos que le permitirán apreciar el curso que dará a la denuncia. En la segunda fase, la Comisión indica a la parte  denunciante las razones por las que no le parece justificado dar curso favorable y le da un plazo para que presente, llegado el caso, sus observaciones el denunciante. En una tercera fase, o bien inicia el procedimiento con la parte denunciada, o bien adopta una decisión definitiva de archivo de la denuncia, que ha de comunicar al denunciante, y que puede ser objeto de un recurso de anulación. Si se abstiene de iniciar el procedimiento y de adoptar la decisión de archivo, entonces cabe el recurso por omisión
































Los Tratados asignan unas características y competencias a la Comisión tratando que, en su actuación como órgano colegiado, prevalezca la supranacionalidad y la independencia frente a los Estados miembros. Es por tanto, en gran parte, el motor de la integración. La Comisión desempeña cuatro funciones principales: (i) iniciativa legislativa; (ii) guardiana de los Tratados; (iii) órgano ejecutivo; y (iv) poder de negociación internacional.

El monopolio de la iniciativa legislativa otorga a la Comisión un poder fundamental. que ha logrado mantener en el Tratado de Lisboa.\footnote{
En algunos casos detcnninados . otros actores pueden pedir a la Comisión que presente una propuesta: El Parlamento, por mayoría de los miembros que lo componen El Consejo. actuando por mayoría simple Por iniciativa de una cuarta parte de los EEMM (cooperación judicial en materia penal o cooperación policial) Por recomendación del 13anco C<!ntral Europeo A petición del Tribunal de Justicia de la Unión Europea (para la creación de tribunales especializados) A petición del Banco Europeo de Inversiones Por iniciativa ciudadana europea (reuniendo al menos un millón de finnas)
} Antes de adoptar una determinada propuesta legislativa, la Comisión ha podido anticiparla desarrollando sus ideas y fomentando el debate mediante comunicaciones, hojas de ruta, libros blancos, libros verdes (menos específicos) y procesos de consulta pública. 

En lo que concierne al control y vigilancia del cumplimiento de los Tratados, la Comisión puede llevar a cabo distintas acciones, como exigir información a empresas y Estados miembros, iniciar un procedimiento de infracción, e imponer sanciones. Algunos ámbitos en los que la Comisión se mantiene especialmente vigilante es el de trasposición de directivas, la política de competencia, el Pacto de Estabilidad y Crecimiento, o el respeto al Estado de Derecho. 

En el ámbito ejecutivo, el Tratado de Lisboa ha reservado a la Comisión por primera vez la potestad de adoptar dos tipos de actos encaminados a la aplicación de la legislación europea. La Comisión adopta estos actos asistida por los Estados miembros a través de un comité (sistema de "comitología"). 
• Los actos delegados son actos de alcance general que completan o modifican determinados elementos no esenciales de un acto legislativo para adaptarlos a los avances científicos o a la evolución del mercado. 
• Los actos de ejecución son medidas muy detalladas que garantizan que los actos legislativos se aplican de manera uniforme en toda la Unión.

En los ámbitos de acción exterior distintos a la política exterior y de seguridad, la Comisión desempeña un papel principal, en particular en relación a la política comercial y la ayuda humanitaria. Por mandato del Consejo, la Comisión negocia los acuerdos y los tratados internacionales sobre materias que sean competencia de la Unión. En cuanto a la política exterior y de seguridad, el Tratado de Lisboa coloca al  Alto Representante de la Unión para Asuntos Exteriores y Política de Seguridad al frente del nuevo Servicio Europeo de Acción Exterior.

En lo que concierne a la composición de la Comisión, el Tratado de Lisboa preveía que, a partir del I de noviembre de 20 14, el número de Comisarios fuera el correspondiente a los dos tercios del número de Estados miembros.\footnote{%
    El término '"Comisión" no solo hace referencia al conjunto de miembros del ·'Colegio de Comisarios". que dirige la institución y toma sus decisiones, y en el que se centra el texto, sino que también se refiere a la propia institución y todo su personal. El personal de la Comisión (unos 32.000 trabajadores) se organiza en "direcciones generales" (DG de Economía y Finanzas. DG de Comercio, DG de Competencia, etc.) y ··servicios·• (Servicio Jurídico. Servicio de Traducción, Eurostat. etc.).
} Sin embargo, el Consejo Europeo adoptó en 2013 la decisión (ya prevista en 2009 para vencer reticencias de Irlanda) de mantener el principio de un Comisario por país, incluyendo al Presidente y al Alto Representante. Esta composición se ha mantenido para la Comisión Juncker 20 14-20 19 y para la Comisión von der Leyen 2019-2024. No es previsible que cambie en el futuro.

El mandato de cada Comisión es de cinco años. Su designación se produce en un plazo de seis meses tras las elecciones al Parlamento Europeo. Es el Parlamento quien elige al Presidente de la Comisión por mayoría, a propuesta del Consejo Europeo. Posteriormente, de acuerdo con el Presidente electo, el Consejo aprueba por mayoría cualificada la lista del resto de candidatos a miembros de la Comisión, quienes deberán someterse al examen individual y al voto en bloque del Parlamento.\footnote{%
    Para vincular más a la Comisión con los resultados de las elecciones al Parlamento Europeo. se ha defendido en las últimas ocasiones un proceso de ··SpitzenkandidaC (compromiso do.: que el candidato que presentó el principal grupo político a la vista de las elecciones sea el nominado a Presidente de la Comisión). Este método se empleó en 2014. En 2019, se tem1inó no siguiendo. lográndose el equilibrio político en la reunión del Consejo Europeo dedicada a discutir sobre nombramientos tras las elecciones.
} El Presidente electo de la Comisión define unas prioridades políticas con arreglo a las cuales la Comisión planificará su trabajo durante su mandato. Durante éste, ha ido ganado  relevancia la comparecencia anual del Presidente de la Comisión ante el Pleno del Parlamento de septiembre (el llamado debate sobre el estado de la UE, o SOTEU), cuando anuncia con marcado tono político las principales iniciativas para el año siguiente.

También se ha reforzado notablemente la posición del Presidente dentro del Colegio de Comisarios, determinando éste la distribución de carteras y la organización del Colegio (por ejemplo, la Comisión von der Leyen cuenta con un Vicepresidente primero, dos Vicepresidentes ejecutivos y cinco Vicepresidentes, uno de los cuales ha de ser el Alto Representante) que puede cambiar a lo largo de su mandato. También exigir la dimisión a un miembro de la Comisión.

\subsection{Parlamento Europeo}
El Parlamento Europeo es la asamblea de los representantes de los ciudadanos de los Estados miembros de la Unión Europea, elegida mediante sufragio universal directo. El Parlamento Europeo encarna el principio democrático en la estructura institucional de la Unión Europea. La ciudadanía está directamente representada en la Unión a través del Parlamento Europeo (art. 10.1 TUE).

El Parlamento Europeo (PE) no es una simple «asamblea» o «conferencia general» de una organización internacional; por el contrario, ya desde su fundación fue una novedad en las relaciones entre Estados por sus atribuciones, en especial, por su poder colegislador y sus atribuciones de control político, extraños y difícilmente ubicables en los modelos organizativos internacionales. 

La decisión de dotar de una institución parlamentaria corrobora la inspiración política del proceso iniciado en los años cincuenta del pasado siglo. La opción por el principio democrático se debió a la creencia, firmemente compartida y mantenida por los Estados miembros, de que la integración europea tiene que encontrar sus cimientos políticos en una «unión cada vez más estrecha entre los pueblos europeos» (Preámbulo del antiguo TCE y del vigente TFUE) que hiciese realidad, a través de la participación de los ciudadanos en la vida política, el establecimiento de «una Comunidad más amplia y más profunda entre pueblos durante largo tiempo opuestos por divisiones sangrientas» (Declaración Schuman).

La participación de la ciudadanía en el Parlamento Europeo hace de la Unión Europea una organización fundada en la democracia representativa y en el pluralismo político. La vinculación de la Unión Europea al pluralismo político se consagró por el Tratado de Maastricht de 1992, al reconocer que «los partidos políticos a escala europea constituyen un importante factor para la integración» al contribuir «a la formación de la conciencia europea y a expresar la voluntad política de los ciudadanos de la Unión» (hoy, art. 10.4 TUE) y se proclama que el funcionamiento de la Unión se basa en la democracia representativa, principio que se encarna de forma central en el Parlamento Europeo.

El Parlamento Europeo es la institución que más cambios ha experimentado desde el inicio del proceso de integración. Su denominación en los Tratados fundacionales (CEE y CEEA) era la de «Asamblea» o «Asamblea común» (CECA), pero decidió en 1962 autodenominarse «Parlamento europeo», término que aceptó el Acta Única Europea de 1987. De ser inicialmente una institución consultiva pasó a ejercer conjuntamente el poder presupuestario (desde 1977) y legislativo (desde 1993) con el Consejo (art. 14.1 TUE), además de mantener funciones de control político y consultivas. La redacción dada al precepto de cabecera sobre el Parlamento Europeo en el Tratado de Lisboa (art. 14 TUE) ayuda a la visibilidad del Parlamento ante la opinión pública y los medios de comunicación al enunciar qué papel desempeña.

\subsubsection{Composición y funcionamiento} 
Desde 1952 a 1979 los diputados al PE eran designados en el seno de los Parlamentos nacionales, de forma que su representación era de segundo grado, y obligaba, por tanto, a la acumulación del doble mandato: la duración del mandato europeo se vinculaba a su permanencia como parlamentario nacional. Pero los Tratados fundacionales preveían su propia evolución hacia un sistema de elección directa por sufragio universal, mediante la elaboración por el propio Parlamento del proyecto que posibilitase su elección.

En efecto, el Consejo adoptó en 1976 una Decisión y un Acta aneja relativa a la elección de los representantes de la Asamblea mediante sufragio universal directo (en adelante, AEE) por la que se han regido, desde 1979, todas las elecciones al Parlamento. Ha sido objeto de importantes modificaciones mediante la Decisión de 23 de septiembre de 20021. El Acta, que contiene las reglas de fondo aplicables en las elecciones europeas, es un acuerdo internacional y constituye Derecho originario.

Con fundamento en dicha Decisión y AEE se han ido celebrando las sucesivas elecciones mediante sufragio universal desde junio de 1979 y, con posterioridad, cada cada cinco años en torno a mayojunio. La media de participación popular en las elecciones bajó ostensiblemente al aproximarse las grandes ampliaciones del siglo XXI y, después de producidas, por la alarmante abstención de la ciudadanía de los países del Este.

El artículo 223 TFUE establece que la elección será de acuerdo con un procedimiento uniforme en todos los Estados miembros o de acuerdo con principios comunes a todos los Estados miembros.

La explicación de esa segunda opción o disyuntiva, añadida en la reforma operada por el Tratado de Ámsterdam, se debe a la falta de acuerdo sobre una ley electoral única y uniforme.

Estos cambios han facilitado la reforma del Acta Electoral Europea en el 2002, renunciando a una «ley» uniforme, pero estableciendo principios comunes y aproximando la regulación de diversos aspectos electorales. Esta nueva Decisión rige los comicios europeos desde 2004; cualquier cambio futuro se hará mediante un procedimiento legislativo especial del Consejo, cuya iniciativa y previa aprobación corresponderá al Parlamento Europeo y la ratificación a los Parlamentos nacionales (art. 223 TFUE). Los principios comunes predeterminados son: — la elección es por sufragio universal, directo y secreto; — el modo de escrutinio proporcional; — las circunscripciones electorales pueden ser según las características nacionales: únicas (España, Dinamarca, Luxemburgo, Países Bajos, Grecia y Portugal), o mediante subdivisiones, por ejemplo, regionales o de otro tipo (Alemania, Bélgica, Francia, Italia, Irlanda, Polonia, Reino Unido), sin que una u otra opción desvirtúen el carácter proporcional del sistema electoraleuropeo; — se permite la votación de listas con voto de preferencia: el elector puede expresar sus preferencias (Dinamarca, Italia, Bélgica,Países Bajos) o mezclar de diversas listas (Luxemburgo —panachage—); por el contrario, en algunos Estados la lista está bloqueada por los partidos políticos (España, Francia, Grecia y Portugal); — podrán establecer un umbral mínimo para la atribución de escaños; a escala nacional, ese umbral no podrá ser superior al 5 por 100 de los votos emitidos; — se pueden fijar límites a los gastos de la campaña electoral de los candidatos; — nadie puede votar más de una vez; — las fechas de las elecciones se predeterminan en el Acta.

Todos los ciudadanos de la Unión tienen el derecho a ser electores y elegibles en las elecciones al Parlamento Europeo en su lugar de residencia (arts. 20 y 22.2 TFUE y 39 de la Carta de Niza). Los nacionales de la Unión tienen garantizado este derecho de sufragio, cualquiera que sea el Estado miembro en el que residan. Pero ¿el derecho de sufragio activo y pasivo es un derecho exclusivo de los nacionales de los Estados miembros? ¿Un Estado miembro puede reconocer ese derecho de sufragio a nacionales de terceros países vinculados por residencia al Estado miembro? En la controversia sobre estas cuestiones entre España y el Reino Unido, a propósito del derecho de voto de los ciudadanos de la colonia británica en Gibraltar, el Tribunal estimó que no se deduce «de forma expresa y precisa quiénes son titulares del derecho de sufragio activo y pasivo en las elecciones al Parlamento Europeo». [Esas] disposiciones no excluyen, en sí mismas, que a una persona que no tenga la condición de ciudadano de la Unión, […] se le reconozca el derecho de sufragio activo y pasivo (fund. 70 ss.)2.

Por tanto, en lo que no esté regulado por los Tratados y el Acto electoral (edad, residencia, incompatibilidades, etc.), corresponde a cada Estado decidir quién puede ser elector y elegible y, por tanto, puede ampliar esa participación a nacionales de terceros países. Ninguna disposición del derecho vigente ni del Tratado de Lisboa excluye la posibilidad de que los extranjeros residentes puedan participar en las elecciones ni se opone, según el Tribunal, a que los Estados miembros reconozcan ese derecho de sufragio activo y pasivo a determinadas personas que tengan un estrecho vínculo con ellos y que no sean sus propios nacionales o los ciudadanos de la Unión residentes en su territorio» (funds. 77 y 78). Como dijera el Abogado General A. Tizzano, los nacionales de los Estadosmiembros son al menos los titulares necesarios de este derecho.

Este derecho se desarrolló mediante una Directiva3 que prevé que todos los ciudadanos de la Unión, tengan o no la nacionalidad de su Estado miembro de residencia, puedan ejercer en él, en igualdad de condiciones con los nacionales, su derecho a ser electores y elegibles. Esto significa que se puede optar por votar o presentarse como candidato en el Estado del que se es nacional o en el de residencia. Si se opta por el Estado de residencia, deberá inscribirse en el censo electoral y disfrutará del derecho en igualdad de condiciones con los nacionales, en particular, en lo que afecta a la duración y la prueba del período de residencia aplicables a los nacionales, o en caso de exclusión del derecho de sufragio activo o pasivo debido a resolución civil o penal.

Diversas cuestiones electorales se remiten al derecho nacional de cada Estado miembro por el AEE o por la Directiva, con el límite de que las particularidades propias no deberán desvirtuar el principio básico del régimen electoral europeo —el carácter proporcional del modo de elección—, por lo que se rigen en todo lo restante por el derecho nacional; así,
— la edad mínima para ejercer el derecho de sufragio activo y pasivo (que varía entre los dieciocho y los veinticinco años), que en España es a los dieciocho años;
— el voto por correo o en los consulados (permitido en España) o por poder; o su prohibición absoluta como optan algunos Estados miembros;
— la renovación de los puestos en caso de dimisión o fallecimiento, etc.

En España se regula por la Ley Orgánica del Régimen Electoral General (LOREG), cuyo Título VI «Disposiciones especiales para las Elecciones al Parlamento Europeo» fue introducido en 1987 para regular las elecciones europeas en España. Esta Ley ha sido objeto de varias modificaciones: así, en 1994 para adaptar nuestro Derecho electoral al artículo 22.2 TFUE (antiguo art. 19.2 TCE) y a la Directiva 93/109/CE y permitir, así, el derecho de sufragio activo y pasivo de los ciudadanos de la Unión residentes en nuestro país que han optado previamente por ejercer su derecho de voto en España4. Diversas disposiciones administrativas facilitan su inscripción en la Oficina del Censo Electoral mediante el cumplimiento de ciertos requisitos.

España ha optado en la LOREG por listas bloqueadas y circunscripción electoral única para todo el territorio nacional (si bien los partidos nacionalistas desearían circunscripciones autonómicas); prohíbe la acumulación de mandatos parlamentarios y establece causas de inelegibilidad aplicando las incompatibilidades que se establecen para los parlamentarios nacionales, en abierta contradicción con el AEE —que sólo admite incompatibilidades—.

Los diputados al Parlamento Europeo son elegidos por un período de cinco años.

El mandato europeo tiene carácter representativo: formalmente significa que los diputados votan individual y personalmente y no pueden estar vinculados por instrucción ni recibir mandato imperativo alguno.

Gozan de los privilegios e inmunidades siguientes: libertad de desplazamiento sin restricción alguna, facilidades aduaneras y de control de cambios, no pueden ser objeto de persecución, arresto o detención en razón del ejercicio de su función y gozan de inmunidad de jurisdicción penal que sólo podrá ser levantada por el propio Parlamento.

Es importante señalar y comprender que los Tratados reiteran que el Parlamento está compuesto por representantes de los ciudadanos de los Estados miembros (art. 14.2 TUE); no es una representación de los pueblos europeos o de las regiones, sino de los ciudadanos que forman los Estados miembros; incluso en caso de dividir un país en varias circunscripciones, los diputados al PE «seguirán siendo... los representantes de los pueblos de los Estados miembros reunidos en la Comunidad y no los representantes de sus circunscripciones» (como recordaba un antiguo proyecto de acta electoral aprobado en 1998 por el PE). Las regiones y sus pueblos tienen otro foro de representación (el Comité de las Regiones); en el Parlamento Europeo, por el contrario, debe estar representada la ciudadanía articulada en los Estados miembros que han formado la Unión.

El reparto de escaños es decrecientemente proporcional. El Tratado establece que su número no excederá de 751 diputados, con un tope máximo de 96 diputados y mínimo de seis por Estado (art. 14.1 TUE).
— Alemania, 96; Francia, 74; Reino Unido, 735; Italia, 73;
— España, 54; Polonia, 51; Rumanía: 32;
— Países Bajos, 26; Bélgica, 21; Grecia, 21; Portugal, 21; República Checa, 21; Hungría, 21;
— Suecia, 20; Austria, 18; Bulgaria, 17;
— Eslovaquia, 13; Dinamarca, 13; Finlandia, 13;
— Irlanda, 11; Lituania, 11; Croacia, 11;
— Letonia, 8; Eslovenia, 8; Estonia, 6; Luxemburgo, 6; Chipre, 6; Malta, 6.

4. ORGANIZACIÓN INTERNA

El Parlamento Europeo aprueba su propio Reglamento interno en el que se desarrolla el modo de ejercicio de sus competencias6. El PE cuenta con algo más de 5.700 empleos permanentes y un centenar de agentes temporales. El Parlamento Europeo elige a su propio presidente por mayoría absoluta (a partir de la cuarta votación, por mayoría simple), a los vicepresidentes y a los cuestores (quienes se ocupan de las cuestiones administrativas y económicas de los diputados), todos ellos por un período de dos años y medio renovables. El Presidente del Parlamento Europeo dirige el conjunto deactividades parlamentarias. Dispone de todos los poderes cuando preside las deliberaciones (da la palabra, somete a votación los asuntos, proclama los resultados, hace observar el Reglamento, etc.) y representa al Parlamento en el interior de la Unión Europea y en las relaciones internacionales. Si participase en algún debate, debe dejar la presidencia hasta que termine la discusión de la cuestión.

Cada período anual comprende doce sesiones; la sesión anual se inicia el segundo martes de marzo de cada año y se denomina período parcial de sesión a la reunión que tiene lugar cada mes (de tres a cinco días). También puede reunirse en sesión extraordinaria a petición de la mayoría de sus miembros, del Consejo o de la Comisión.

4.1. LOS GRUPOS POLÍTICOS

Los diputados se agrupan por afinidades políticas y no por su nacionalidad. Todo grupo político estará integrado por diputados elegidos en al menos una cuarta parte de los Estados miembros. El número mínimo de diputados necesario para constituir un grupo  político es de veinticinco. No está permitido constituir un grupo político con diputados de una misma nacionalidad. Los diputados que no desean adherirse a un grupo no precisan, como es obvio, de un número mínimo para estar en el «grupo» de los «no inscritos». Sólo se puede pertenecer a un grupo político. 

La existencia de los grupos tiene su importancia a efectos de la formación de las Comisiones parlamentarias y de la Mesa (y de la Conferencia de Presidentes) del Parlamento. Son, a su vez, el resultado de las coaliciones electorales y de la formación reciente de partidos políticos a nivel europeo, si bien hay que reconocer que, pese al desideratum expresado por el artículo 10.4 TUE en torno al fortalecimiento de los partidos políticos europeos, éstos son aún construcciones artificiales.

4.2. LA ORGANIZACIÓN DE LOS TRABAJOS PARLAMENTARIOS

La Mesa del PE está formada por el Presidente, los vicepresidentes y los cuestores (sin derecho de voto). Le corresponde a la Mesa regular las cuestiones económicas, de organización y administrativas que afecten a los diputados, así como a la secretaría y a los funcionarios y otros agentes del PE. La Conferencia de Presidentes está compuesta por el Presidente del PE y por los presidentes de los grupos políticos. Es el órgano de acción del PE y le corresponde a la Conferencia tratar las cuestiones referidas a la organización de los trabajos parlamentarios, el orden del día, comparecencias, la programación legislativa, etc. y las relaciones con otras instituciones de la Unión y organismos externos.

La Conferencia de Presidentes de Comisión reúne a todos los presidentes de las comisiones parlamentarias temporales o permanentes; la Conferencia de Presidentes de Delegación está integrada por los presidentes de todas las delegaciones interparlamentarias permanentes.

El PE, además de desarrollar sus trabajos en sesiones plenarias, también trabaja en comisiones parlamentarias.

El PE puede crear comisiones parlamentarias permanentes o temporales, generales o especiales, y comisiones de investigación (art. 226 TFUE). Las comisiones se forman teniendo en cuenta una representación equitativa por Estados y de las ideologías políticas (lo que significa que la adscripción no compete siempre al diputado de forma personal); se puede pertenecer a más de una comisión parlamentaria. 

En el seno de las comisiones parlamentarias se examinan las cuestiones que le solicita el Pleno y se confrontan las posiciones de los grupos. Se encarga a un ponente la elaboración de un Informe y un proyecto de resolución o de Dictamen, según la naturaleza deliberante o legislativa del acto informado; si se aprueban en la comisión parlamentaria se eleva al Pleno para su discusión y aprobación. Además, el Pleno puede delegar en la comisióncompetente la elaboración del dictamen, cuando son asuntos técnicos sin alcance general.

La Comisión Europea tiene derecho a asistir y participar en las sesiones plenarias y es costumbre arraigada que asistan uno ovarios Comisarios o el Presidente mismo de la Comisión (art. 230 TFUE). También tiene derecho a asistir a las comisiones parlamentarias, estando representada por altos funcionarios (a veces, también por algún Comisario al presentar una propuesta). La razón de la presencia de la Comisión en los trabajos parlamentarios estriba en su derecho de iniciativa normativa: de un lado, porque el Parlamento Europeo debate sus proyectos normativos y puede recoger sus sugerencias y, de otro, porque, en tanto el Consejo no haya decidido, sólo la Comisión puede modificar su propuesta y tener en cuenta el dictamen parlamentario. Por otra parte, es obvio que la Comisión, como el Consejo, está obligada a asistir al Pleno cuando deben responder a las interpelaciones de los diputados.


4.3. EL SISTEMA DE VOTACIONES Y QUÓRUM
El sistema de votaciones previsto con carácter general es el de mayoría absoluta de los sufragios expresados. Sin embargo, para cuestiones específicas los tratados o el reglamento interno pueden exigir otras mayorías. Por ejemplo, el rechazo en el procedimiento legislativo ordinario, la aprobación de enmiendas y el rechazo del presupuesto, exigen una mayoría absoluta de los diputados que componen el PE; la moción de censura a la Comisión prevé una doble mayoría de dos tercios de los votos expresados y la mayoría absoluta de los miembros que componen el PE; la consolidación de las enmiendas al presupuesto exigen una mayoría de tres quintos en la tercera lectura. El quórum fijado por el Reglamento interno es de un tercio de los miembros efectivos —con mandato en curso—; sólo es exigible su constatación si cuarenta diputados lo piden previamente a una votación. Pero si en la Sala hay menos de cuarenta diputados el propio Presidente puede constatar la falta de quórum.


\subsubsection{Atribuciones}
5. ATRIBUCIONES DE CONTROL POLÍTICO
El Parlamento Europeo es informado asiduamente por el Presidente de la Comisión y por los Comisarios; se pronuncia sobre el programa de trabajo de la Comisión y pacta con ésta el calendario legislativo. También es informado por la Presidencia del Consejo. Controla con carácter general la conducción de las distintas políticas de la Unión.

El Parlamento puede ejercer el control político mediante las interpelaciones o preguntas, la discusión del informe general anual de la Comisión y la moción de censura a ésta; también mediante el descargo o aprobación de la ejecución del presupuesto (objeto del Cap. 13, epígrafe 6). Con algunos matices, estos instrumentos de control tienen un gran parecido a los utilizados por los Parlamentos nacionales. Igualmente, el control democrático se articula mediante las comisiones de investigación y su participación en la designación de las personalidades que componen las Instituciones.

Además, mediante acuerdos específicos, el Parlamento ha ido extendiendo el control democrático a ámbitos diversos; así, el control sobre el Mecanismo Único de Resolución (MUR), facultado para liquidar Bancos con dificultades7.

5.1. LAS INTERPELACIONES Y EL TURNO DE PREGUNTAS 
Las preguntas son interpelaciones en las que los parlamentarios recaban de la Comisión Europea o del Consejo o del Banco Central Europeo una opinión, una información o una explicación sobre un asunto determinado que afecta a la competencia de esa Institución o en relación con las actividades de la Unión Europea en su conjunto.

El Tratado confirma que la Comisión Europea comparecerá ante el Parlamento Europeo si así lo solicita (art. 230 TFUE). Las preguntas pueden ser de varias clases: orales —con o sin debate—, escritas y el turno de preguntas. La Conferencia de Presidentes decidirá sobre su inclusión en el orden del día y sobre el orden de las mismas.

Para las preguntas orales la iniciativa corresponde a una comisión parlamentaria, un grupo político o cuarenta diputados. Las preguntas escritas se pueden plantear por cualquier diputado y no precisan el acuerdo de la Conferencia de Presidentes. Las preguntas, con la respuesta de la Comisión o del Consejo, se publican en el Diario Oficial (serie C). El turno de preguntas tiene lugar en todos los períodos parciales de sesiones y se presentarán una semana antes. Se trata de diálogos breves entre los parlamentarios y la Presidencia de la Comisión o del Consejo o los Comisarios sobre aspectos específicos de actualidad y de interés general.

5.2. EL INFORME GENERAL ANUAL Y LOS INFORMES SECTORIALES

Además, los tratados han previsto que la Comisión presente anualmente (cada mes de febrero) al PE un Informe General sobre todas las actividades de la Unión y sobre el desarrollo de las diferentes políticas (art. 233 TFUE).

El Informe es estudiado en sus diferentes partes por las comisiones parlamentarias permanentes y da lugar en sesión plenaria (en marzo, al abrirse el período anual de sesiones) a un debate sobre el mismo y a la aprobación de una resolución. Este Informe General relata tanto las actividades pretéritas como proyectos o programas de futuro.

Un informe de especial transcendencia es el Informe anual sobre Control de la aplicación del Derecho de la Unión que realiza la Comisión desde 1988; se publica en el Diario Oficial de la UE. Este estudio meticuloso es analizado por diversas comisiones parlamentarias, encargándose a la de asuntos jurídicos la elaboración de un informe al Pleno; la resolución del PE y el informe parlamentario se transmiten al Consejo, a la Comisión, a los Gobiernos y a los Parlamentos nacionales.


5.3. LAS COMISIONES TEMPORALES DE INVESTIGACIÓN
La creación en el seno del PE de comisiones temporales de investigación tiene como objetivo examinar las alegaciones de infracción o de mala administración en la aplicación del Derecho de la Unión.

Estas comisiones podían formarse con anterioridad al TUE, por lo que tampoco fue una novedad su regulación en el Tratado de Maastricht. El Reglamento interno del PE regulaba su constitución y sus funciones. Ahora, su fundamento no es exclusivamente de Derecho derivado como hasta 1993, sino de Derecho originario al estar previstas en el artículo 226 TFUE.

La finalidad de la comisión temporal es elaborar un informe que se presenta al Pleno; éste decidirá si se hace público. Las modalidades de ejercicio del derecho de investigación se regulan en un Reglamento del PE, de su propia iniciativa, mediante un procedimiento legislativo especial, previa aprobación del Consejo y de la Comisión9.

5.4. LA MOCIÓN DE CENSURA Y LA VOTACIÓN DE CONFIANZA A LA
COMISIÓN
Otro gran instrumento de control político es la moción de censura, la cual se dirige únicamente contra la Comisión a fin de exigir responsabilidades políticas por su gestión (art. 234 TFUE).

La moción de censura se puede plantear por la décima parte de los diputados. Deberá ser planteada por escrito y estar motivada. El Presidente del PE deberá anunciarla inmediatamente a los diputados y notificarla a la Comisión Europea. Además de algunos obstáculos de «enfriamiento» en cuanto al momento del debate y de la votación (en la práctica suponen de una a cuatro semanas), se exige un procedimiento de votación muy riguroso que hace de la moción de censura un instrumento de difícil éxito: la votación es nominal y pública y precisa una doble mayoría de dos tercios de los votos expresados y la mayoría absoluta de los miembros que componen el PE.

Si la moción de censura lograse su aprobación, los miembros de la Comisión y su Presidente tendrían que dimitir colectivamente si bien proseguirían despachando los asuntos corrientes hasta su reemplazo por una nueva Comisión; el Alto Representante de la Unión (AR) dimitirá del cargo de Vicepresidente. La moción de censura sólo puede ser planteada en relación con la Comisión en su conjunto; no se admiten reprobaciones de algún Comisario en particular. Hasta ahora se han presentado en varias ocasiones mociones de censura contra la Comisión. Otras mociones han sido presentadas pero fueron retiradas antes de su votación.

5.6. EL DERECHO DE PETICIÓN ANTE EL PARLAMENTO EUROPEO

5.7. CONTROL DEMOCRÁTICO DE LA LEGALIDAD: EL IUS STANDI ANTE EL
TRIBUNAL DEL JUSTICIA DE LA UNIÓN
El Parlamento tiene legitimación activa y pasiva ante el Tribunal de Justicia, por tanto puede ser demandante o demandado por sus actos y omisiones ilegales. La reiterada petición del Parlamento Europeo de serle reconocida la legitimación activa para recurrir actos de otras instituciones en el recurso de anulación (art. 263 TFUE) fue satisfecha por el Tratado de Niza (2001). Se le reconoce la misma capacidad que al Consejo y a la Comisión de modo que, como demandante privilegiado, ya no tiene que justificar el interés para recurrir en la violación de sus prerrogativas.

Igualmente, le ha sido reconocida capacidad para solicitar al Tribunal de Justicia dictámenes sobre la compatibilidad con los Tratados de los acuerdos internacionales con terceros que proyecte suscribir la Unión. Se le reconoce con plenitud que es un defensor objetivo de la legalidad de la Unión en pie de igualdad con las Instituciones que concurren en el proceso legislativo y tanto en la acción legislativa interna (art. 263 TFUE, recurso en anulación) como en la externa (art. 218.11 TFUE, competencia consultiva).


6. LA INICIATIVA LEGISLATIVA Y EL PARLAMENTO EUROPEO

En el sistema europeo de integración, la iniciativa normativa es una atribución tradicional de la Comisión. Los Tratados siempre han reconocido con un carácter muy amplio el derecho de iniciativa a la Comisión, aunque haya disminuido o se vea más compartida desd el Tratado de Lisboa. Se ha considerado siempre ese cuasimonopolio del derecho de iniciativa como un elemento fundamental del equilibrio institucional. Por ello, las reiteradas peticiones del Parlamento Europeo de compartir el derecho de propuesta han chocado con la oposición de los Estados miembros a modificar los Tratados en ese punto esencial de la dinámica de la Unión. 

La solución acordada en Maastricht en torno al actual artículo 225 TFUE sigue el modelo existente para el Consejo (art. 241 TFUE)10. Este precepto, en su segundo párrafo, prevé que el PE pueda solicitar a la Comisión que [le] presente las propuestas oportunas sobre cualquier asunto que a juicio de aquél requiera la elaboración de un acto de la Unión para la aplicación de los Tratados.

En esta cuestión, los Estados han igualado plenamente al Parlamento con el Consejo. Pero no se le reconoce derecho deiniciativa al PE como tampoco al Consejo. No podrá subrogarse a la Comisión en caso de inacción de ésta. Tampoco el Parlamento podrá prejuzgar la libertad de apreciación de la Comisión sobre las prioridades normativas y el momento de su utilización. Ahora bien, la obligación de colaboración entre las instituciones afecta a la iniciativa legislativa y al conjunto del proceso legislativo (art. 295 TFUE)11.

En todo caso el Parlamento tiene reconocida alguna iniciativa normativa restringida a supuestos muy concretos que no siempre se corresponden con la clásica iniciativa legislativa, así: — puede tomar la iniciativa para iniciar una revisión ordinaria de los Tratados (art. 48.2 TUE) o la revisión simplificada de la parte III del TFUE (art. 48.6 TUE); — o la iniciativa de presentar al Consejo Europeo una propuesta de decisión por la que se fije la composición del Parlamento Europeo, que tras su paso por el Consejo Europeo requiere la aprobación del Parlamento (art. 15.2 TUE). Más parecido con la iniciativa legislativa tienen las previsiones yatradicionales de reconocerle el derecho a presentar los proyectos relativos:
— al procedimiento para la elección de los diputados al PE,
— al estatuto y condiciones generales de ejercicio de las funciones de los diputados. El Estatuto se aprueba por procedimiento legislativo especial del Consejo, con posterior aprobación del PE. Las condiciones generales es un procedimiento legislativo especial del propio Parlamento con posterior aprobación del Consejo (art. 223.1 y 2 TFUE);
— a las modalidades del ejercicio del derecho de investigación (se aprueba mediante procedimiento legislativo especial del PE, previa aprobación del Consejo por mayoría cualificada y de la Comisión, art. 226 TFUE);
— al Estatuto y las condiciones generales de ejercicio de las funciones del Defensor del Pueblo (se aprueba mediante procedimiento legislativo especial del PE y previa aprobación del Consejo, art. 228.4 TFUE).

Lo esencial del derecho de iniciativa sigue estando en poder de la Comisión. Claro que si la Comisión se mostrara poco receptiva a las peticiones del Parlamento de presentar propuestas legislativas en las materias que le sugiera, aquélla tendrá que asumir el riesgo de una moción de censura de los desairados parlamentarios. Otra opción menos drástica, si el Parlamento Europeo estima que hay una carencia de la Comisión en materia de iniciativa legislativa, sería demandar a la Comisión conforme al procedimiento por omisión previsto en el artículo 265 TFUE.


7. LA FUNCIÓN LEGISLATIVA CONJUNTA DEL PARLAMENTO Y DEL CONSEJO

Desde las reformas introducidas por el Tratado de Maastricht de 1992, el Parlamento logró ser un legislador, junto al Consejo, participando en el proceso de decisión de normas básicas. Una puntualización importante es que la codecisión sólo se preveía para determinadas materias del antiguo Tratado de la Comunidad Europea. En la posterior reforma introducida por el Tratado de Ámsterdam se incrementaron en una quincena los ámbitos que pasaron a ser decididos conjuntamente por el Consejo y el Parlamento Europeo. Tras las reformas introducidas por el Tratado de Lisboa se ampliaron de nuevo los ámbitos materiales y el número de normas que se aprobarán por el procedimiento legislativo ordinario.

7.2. CARACTERES BÁSICOS DE LA FUNCIÓN LEGISLATIVA 

El Tratado de la Unión Europea, tal como ha sido modificado por el Tratado de Lisboa, ha abordado la distinción entre la función legislativa, que atribuye de forma conjunta al Parlamento Europeo y al Consejo (arts. 14 y 16 TUE), y la ejecutiva o no legislativa, que atribuye, por separado, al Consejo, a la Comisión y a otras instituciones, entendiendo que en lo no establecido corresponde a los Estados miembros (art. 291.1 TFUE). Esa distinción explícita, entre función legislativa y ejecutiva (aunque sin definirlas), aporta transparencia a las funciones de las instituciones y al sistema de actos. 

La función legislativa no se describe en el Tratado; sin embargo, se deduce del contexto general de los Tratados de la Unión que se refiere a la capacidad de decidir los elementos esenciales de un régimen jurídico en un ámbito dado, estableciendo las opciones políticas, los principios que lo rigen, los objetivos generales y los contenidos esenciales que conllevan las opciones políticas fundamentales.

Desde el punto de vista sustantivo, un acto legislativo debe expresar las opciones políticas fundamentales, los elementos esenciales del régimen jurídico en cuestión, los objetivos generales y los principios que lo rigen. Necesariamente su base jurídica tiene que estar prevista directamente en el Tratado (principio de atribución de competencias)12.
El Tratado proclama que son actos legislativos los adoptados mediante un procedimiento legislativo, ya sea ordinario, ya sea especial (art. 289.3 TFUE). La definición es formal y esa categorización viene condicionada por el procedimiento de adopción. Este procedimiento rige para la adopción de actos en las materias que el TUE y el TFUE reservan expresamente a este procedimiento. Emulando la práctica interna hay, pues, materias de «reserva de acto legislativo».

La distinción entre funciones conlleva racionalizar los  procedimientos de la Unión a este respecto; permite una mayor lógica de la toma de decisiones y, por tanto, una mayor transparencia de dicho proceso para los ciudadanos responsabilizando al Parlamento y al Consejo de las decisiones más importantes. La tarea de los colegisladores debería concentrarse en la definición de grandes objetivos y principios de la legislación de la Unión más que en disposiciones detalladas y de carácter altamente técnico, que se corresponderían con la función ejecutiva.

Se denomina procedimiento legislativo ordinario al proceso que
se inicia con la propuesta de la Comisión y se concluye con la
adopción conjunta por el Parlamento y el Consejo (arts. 289.1 y 294
TFUE); los actos así aprobados tendrán la categoría genérica de
actos legislativos.

Excepcionalmente, y en los casos detallados por los Tratados  expresamente, los actos legislativos podrán ser a propuesta de los Estados miembros, del Parlamento Europeo, del Banco Central Europeo, del Banco Europeo de Inversiones y del Tribunal de Justicia.

La utilización del calificativo «legislativo» no debe extrañar, aunque se utilice por primera vez de forma clara en el articulado del Tratado13. Se aprueban, pues, conjuntamente por el Parlamento Europeo y por el Consejo con participación paritaria mediante un sistema de doble lectura o debate. Este procedimiento se caracteriza por su transparencia y carácter público tanto en el Parlamento como en el Consejo (art. 294 TFUE).

Desde el Tratado de Maastricht es una característica determinante de los actos legislativos de la Unión el hecho de que no se pueda aprobar una norma legislativa sin el acuerdo del Parlamento Europeo. La garantía de la legitimidad democrática de las normas de codecisión o del procedimiento legislativo ordinario significa que la ciudadanía de la Unión tiene asegurado que no se puede adoptar ninguna norma legislativa que, además de la mayoría cualificada del Consejo —en determinados casos, la unanimidad—, no tenga la aprobación mayoritaria del Parlamento. Los actos legislativos ordinarios se aprueban, con carácter general, por mayoría del Parlamento y por la mayoría cualificada del Consejo, salvo que se establezca expresamente la unanimidad (art. 16.3 TUE).

Los ámbitos sujetos al procedimiento legislativo ordinario son muy numerosos; más de noventa preceptos contienen las bases jurídicas para adoptar actos legislativos ordinarios en materias muy diversas y fundamentales como no discriminación, ciudadanía, mercado interior, espacio de libertad, seguridad y justicia, aproximación de legislaciones, transportes, empleo, cooperación aduanera, política social, igualdad de oportunidades y de trato, protección de consumidores, redes transeuropeas, cohesión económica y social, investigación y desarrollo, medio ambiente, cooperación al desarrollo, transparencia, lucha contra el fraude, etc. 

Conviene recordar que, a veces, en un mismo ámbito material pueden coexistir varios procedimientos de decisión, de modo que la adopción de los principios o líneas básicas de un sector se aprueben en procedimiento legislativo ordinario, mientras que acciones concretas se adoptan mediante procedimientos legislativos especiales o mediante actos no legislativos con consulta al Parlamento Europeo.

8.1. LA PRIMERA LECTURA
[...]


8.2. LA SEGUNDA LECTURA
[...]

En esta segunda fase, tanto el Parlamento Europeo como el Consejo tienen un plazo de tres meses ampliable a un mes a petición de cualquiera de dichas instituciones.

8.3. TERCERA LECTURA EN CASO DE ACUERDO EN EL COMITÉ DE CONCILIACIÓN
[...]


9. LA FUNCIÓN LEGISLATIVA MEDIANTE EL PROCEDIMIENTO LEGISLATIVO ESPECIAL

El procedimiento legislativo especial se diferencia del ordinario, basado en la adopción conjunta del Parlamento Europeo y del Consejo, porque se adoptan por una de estas dos instituciones con aprobación o consulta, según los casos, de la otra (art. 289.2 TFUE). El procedimiento legislativo especial, con excepción del procedimiento presupuestario, se distingue por la preponderancia del Consejo o por la del Parlamento Europeo frente a la paridad de Parlamento y Consejo en el ordinario. Son actos con valor legislativo aunque no sean aprobados mediante el procedimiento legislativo ordinario (art. 289.3 TFUE).

En tanto que actos legislativos sus caracteres sustantivos son idénticos: tienen su base jurídica directamente en el Tratado, enuncian el marco jurídico fundamental de un ámbito y las opciones políticas básicas. Ahora bien, la legitimidad democrática es un distintivo evidente tanto en los actos que se aprueban en el procedimiento ordinario como en los aprobados por procedimiento especial por el Parlamento Europeo.

Hay que constatar que se prevén muchísimos más actos legislativos del Consejo que del Parlamento Europeo. Ahora bien, siempre en un acto legislativo de procedimiento especial, ya sea del Consejo o del Parlamento, hay una participación de la otra institución. 

Cabe resaltar que los procedimientos legislativos especiales son muy casuísticos, ya sea en materia de iniciativa, ya de consultas o de votación; caso por caso, se marca el itinerario que seguirá cada acto, frente a la uniformidad del procedimiento ordinario. 

Se prevén tres actos legislativos del Parlamento Europeo en los siguientes preceptos y con estas características procedimentales:
— Estatuto y condiciones generales de ejercicio de las funciones de los diputados. El Parlamento Europeo se pronunciará por propia iniciativa, previo dictamen de la Comisión y previa aprobación del Consejo por mayoría cualificada. El Consejo se pronunciará por unanimidad sobre toda norma o condición relativa al régimen fiscal de los diputados o de los antiguos diputados (art. 223.2 TFUE).
— Modalidades del ejercicio del derecho de investigación. El Parlamento Europeo se pronunciará por propia iniciativa, previa aprobación del Consejo por mayoría cualificada y de la Comisión (art. 226 TFUE).
— Estatuto y las condiciones generales de ejercicio de las funciones del Defensor del Pueblo. El Parlamento Europeo se pronunciará por propia iniciativa, previo dictamen de la Comisión y previa aprobación del Consejo por mayoría cualificada (art. 228 TFUE).

10. LA PARTICIPACIÓN EN LA CONCLUSIÓN DE ACUERDOS INTERNACIONALES Y EN LA ADHESIÓN DE NUEVOS ESTADOS

El Acta Única Europea en 1986 introdujo un nuevo procedimiento  de participación del Parlamento Europeo en el poder de decisión del Consejo mediante la emisión de dictámenes conformes. El Tratado de Lisboa ha cambiado esta denominación, algo opaca, por la de aprobación. Las diversas reformas han ampliado los supuestos de participación del Parlamento en el proceso de prestación del consentimiento de la Unión en materia de tratados internacionales.

Así, de conformidad con el artículo 218.6.a) TFUE necesitarán la previa aprobación del Parlamento Europeo, por mayoría de los votos emitidos19, para los siguientes acuerdos o convenios internacionales:
— los acuerdos internacionales de asociación; 
— el acuerdo de adhesión de la Unión al Convenio Europeo para la Protección de los Derechos Humanos y de las Libertades Fundamentales;
— los acuerdos que crean un marco institucional específico organizando procedimientos de cooperación;
— los acuerdos con implicaciones presupuestarias importantes para la Unión;
— los acuerdos que implican una modificación de un acto adoptado según el procedimiento legislativo ordinario o el especial cuando precisa de la aprobación del Parlamento;
— los tratados de adhesión de nuevos Estados miembros (por mayoría de los miembros que componen el Parlamento Europeo, es decir, mayoría reforzada, art. 49 TUE).

La naturaleza jurídica de la aprobación es la de un acto de autorización para celebrar un acuerdo internacional, similar al que se precisa en el Derecho interno por parte de los Parlamentos nacionales para la prestación del consentimiento del Estado. El Parlamento Europeo, al igual que los Parlamentos nacionales, tiene que aceptar o rechazar el tratado como un todo.

Mediante la aprobación participa del poder de decisión en materia de acción exterior en la medida en que no podrá haber acuerdo sin su autorización y que su negativa supondría impedir la adopción del acuerdo en cuestión.

El artículo 218.6.a) exige una mayoría muy rebajada (mayoría de los votos emitidos en materia de acción exterior) y, por tanto, fácil de obtener. Pero también es un arma de doble filo, ya que esa misma mayoría fácil se requiere para rechazarlo, es decir para el dictamen no conforme o veto al acuerdo en cuestión. En materia de aprobación de acuerdos internacionales no se ha seguido el método asimétrico del procedimiento legislativo ordinario en el que se requieren mayorías fáciles para aprobar y mayorías reforzadas para rechazar.

Para el resto de acuerdos, el Parlamento Europeo deberá ser consultado previamente (consulta ordinaria o dictamen consultivo). En el mismo sentido, es reconfortante la extensión de la simple consulta (no vinculante) al resto de acuerdos (nada habitual en el interior de los Estados con ese carácter previo)20.


11. LA PREVIA APROBACIÓN DEL PARLAMENTO
EUROPEO EN OTRAS MATERIAS Y
PARTICIPACIÓN EN LA FUNCIÓN LEGISLATIVA
DEL CONSEJO
Los Tratados han ido extendiendo la necesidad de obtener la previa aprobación del Parlamento Europeo (lo que antes del Tratado de Lisboa se denominaba el dictamen conforme) en materias distintas a los acuerdos internacionales. Sin duda, es un avancepositivo en la medida en que acrece los poderes de influencia ydecisión del Parlamento mediante este derecho de veto: en los ámbitos de la previa aprobación sólo se podrán aprobar esas normas si tienen el apoyo mayoritario del Parlamento Europeo.

Tiene trascendencia el reconocimiento del derecho a participar en  la revisión de algunos aspectos de los Tratados, como es la necesidad de obtener su previa aprobación para el paso del procedimiento especial al procedimiento legislativo ordinario y de la unanimidad a la mayoría cualificada (art. 48.7 TFUE).

Sin embargo, la previa aprobación del PE es un sistema un tanto anómalo en materias legislativas pues carece de la opción de presentar enmiendas. Es un progreso democrático, pero no en la dirección deseada de la plena coparticipación normativa. Sólo permite al Parlamento Europeo aprobar el texto normativo que le presenta el Consejo o rechazarlo, lo que no impide que el Consejo y el Parlamento puedan dialogar y negociar directamente.


12. LAS FUNCIONES CONSULTIVAS Y DE DELIBERACIÓN

12.1. LA CONSULTA ORDINARIA: EL DICTAMEN CONSULTIVO O DE LECTURA ÚNICA
Hasta 1986 la competencia consultiva del PE era el núcleo de su participación en el proceso de decisión. Como declaró el Tribunal de Justicia en las sentencias sobre los asuntos Maizena y Roquette (isoglucosas), la consulta es el medio que permite al Parlamento participar de forma efectiva en el proceso legislativo de la Comunidad. Esta competencia representa un elemento esencial del equilibrio institucional querido por el Tratado. Es el reflejo a nivel de la Comunidad, aunque limitado, de un principio fundamental según el cual los pueblos participan en el ejercicio del poder a través de una asamblea representativa21. Los Tratados establecen que el Consejo y la Comisión antes de adoptar determinadas decisiones deben solicitar y conocer el dictamen del PE (dictamen consultivo preceptivo). El dictamen consultivo no es vinculante para la institución que lo solicita.

Sin embargo, en el caso de un dictamen consultivo preceptivamente exigido por los Tratados, si el Consejo o la Comisión no solicitasen o no esperasen a conocer su dictamen, la decisión así adoptada puede ser objeto de recurso de anulación ante el Tribunal de Justicia por violación de formas sustanciales (art. 263 TFUE). Así, el Tribunal de Justicia anuló un Reglamento del Consejo en los citados asuntos Roquette y Maizena (en los que el Parlamento hizo uso de su derecho de intervención) debido, entre otras alegaciones, a que el Consejo no esperó su dictamen y era preceptivo conocerlo para adoptar el reglamento en cuestión. Para el Tribunal, el derecho de consulta del Parlamento Europeo no queda satisfecho por la simple petición de dictamen sino que su respeto implica que el Parlamento Europeo exprese su opinión22.

Además, es práctica del Parlamento pedir formalmente a la Comisión que modifique sus propuestas acogiendo las enmiendas o sugerencias del dictamen del PE (incluso en la fase de comisión parlamentaria), aplazando la emisión del dictamen y recabando, en su caso, información de las razones por las que no se han tenido en cuenta.

Los ámbitos materiales de la consulta ordinaria se han ido reduciendo en las sucesivas reformas en favor del procedimiento legislativo ordinario en los que tiene plena capacidad de decisión. Pero aún así se prevé la consulta al Parlamento en unos cinco casos en el TUE y en unos cincuenta preceptos del TFUE; muchos son casos de procedimiento legislativo especial del Consejo con la previa consulta al Parlamento Europeo.


12.2. LA DELIBERACIÓN POR PROPIA INICIATIVA
Además de los dictámenes consultivos en los que se pronuncia a solicitud del Consejo o de la Comisión para los actos adoptados en el marco de los Tratados, el Parlamento es competente para deliberar y pronunciarse por propia iniciativa sobre cuestiones que entren en el marco de las actividades internas o externas, en sentido lato, de la Unión Europea o, incluso, sobre cuestiones al margen de tal actividad pero en las que el PE desea expresar su opinión (violaciones de los derechos humanos en cualquier parte del mundo, catástrofes naturales, etc.). Estos poderes de deliberación significan que el PE no se limita a expresar su opinión sobre la materia y el momento elegido por el Consejo o la Comisión, sino que a iniciativa propia puede discutir y expresar su opinión.


13. LA FUNCION PRESUPUESTARIA

Véase Capítulo 13 sobre financiación de la UE.

El Parlamento fonna junto al Consejo la autoridad presupuestaria de la Unión: Participa en el proceso presupuestario anual desde la fase de preparación; es decir, el establecimiento de las orientaciones generales y la naturaleza de los gastos. Posteriormente, aprueba el presupuesto, controla su ejecución y aprueba la gestión de tal ejecución.







El Parlamento Europeo es la institución europea que en mayor medida ha ido reforzando y aumentando sus competencias en el curso de las sucesivas refom1as institucionales, de manera que ha pasado de ser una asamblea consultiva a ser pleno "colegislador" en la mayor parte de los casos. Las principales funciones del Parlamento son las siguientes:

El Parlamento participa en la adopción de los actos legislativos de la Unión en diversa medida, según la base jurídica de cada acto.
• Adopta legislación europea con arreglo al procedimiento legislativo ordinario en pie de igualdad con el Consejo. En caso de acuerdo entre ambas instituciones, el acto se adopta en primera o segunda lectura; en caso de desacuerdo, el acto solo puede adoptarse tras un proceso de conciliación satisfactorio. El Tratado de Lisboa no ha modificado el desarrollo del procedimiento (antes conocido como ,'codecisión''), pero lo ha conver tido en el "caso nor mal" incrementándose el número de áreas en las que aplica (libertad, seguridad y justicia, comercio exterior, medio ambiente, PAC, etc.). 
• Emite dictámenes en el marco del procedimiento de consulta. En estos casos, el Consejo no queda vinculado por la posición del Parlamento, sino solo por la obligación de consultarlo. La ausencia de tal consulta convie11e el acto en ilegal. Se consulta al PE en un número limitado de ámbitos como política de competencia, armonización de los impuestos indirectos, etc. 
• Tiene derecho de veto o de aprobación en los casos en los que aplica el procedimiento de aprobación. Es decir, El PE puede aceptar o rechazar un actopropuesto, pero no modificarlo. En caso de que el PE no dé su aprobación, el actono puede adoptarse. Este procedimiento aplica, por ejemplo, a la adhesión de nuevos Estados miembros a la UE, a su retirada, a la ratificación de detenninados acuerdos internacionales, a los casos de violaciones graves de los derechos fundamentales en virtud del artículo 7 del TUE, a la aprobación del Marco Financiero Plurianual, etc.



El Parlamento dispone de competencias de control del ejecutivo. Por ejemplo, elige al Presidente de la Comisión a propuesta del Consejo Europeo, el Colegio de Comisarios se somete al voto de investidura del PE, puede presentar una moción de censura pidiendo la dimisión de toda la Comisión, formula preguntas parlamentarias a la Comisión y al Consejo, tiene la facultad de establecer comisiones de investigación, etc.

En lo que concierne a la composición del Parlamento, el Tratado de Lisboa establece que el número de miembros del Parlamento Europeo no puede exceder los 750, más el Presidente. Tras la retirada de Reino Unido, el número total de escaños se ha reducido de 751 a 705. El Consejo Europeo decide la composición exacta del Parlamento Europeo. El mínimo de escaños por Estado miembro es de seis, y el máximo de 96, asignándose el número concreto con proporcionalidad decreciente (sobre la base del tamaño de la población del Estado miembro, pero evitando la infrarrepresentación de los pequeños corrigiendo la proporción; cuanto mayor sea el país, menor será el número de escaí'los con respecto a su población).

Los eurodiputados son elegidos por sufragio universal directo cada cinco años. La mayor parte  son miembros de un partido político nacional en su país de origen. Algunos partidos políticos están presentes a escala de la UE; los cuales sirven para que, en el Parlamento Europeo, la mayor parte de los diputados se organicen en grupos políticos supranacionales en función de sus afinidades políticas y no en delegaciones nacionales. No obstante, sigue pendiente· el establecimiento de un sistema electoral unifonne y la creación de partidos políticos europeos, paso previo a que el Parlamento pueda equipararse a una asamblea propia de un Estado federal.

Los miembros del Parlamento Europeo eligen a su Presidente para un mandato de dos años y medio. Las sesiones plenarias, a las que asisten todos los diputados, se celebran en Estrasburgo (sede oficial del PE) una vez al mes, con una duración de una semana, y en Bruselas seis veces al año, de duración más corta. También pueden convocarse plenos extraordinarios. Existe una Secretaría General del Parlamento Europeo que coordina el trabajo Legislativo, y presta asistencia técnica, jurídica y especial izada, y servicios de interpretación y traducción.

\subsection{Tribunal de Justicia de la Unión Europea} 
El Tribunal de Justicia de la Unión Europea (TJUE) es la
Institución que encarna el poder judicial en la Unión Europea, ya que
el artículo 19 TUE afirma que
garantizará el respeto del Derecho en la interpretación y aplicación de los
Tratados.
Ahora bien, el TJUE comparte el ejercicio de la función
jurisdiccional con los órganos jurisdiccionales de los Estados
miembros, que son los llamados a aplicar, en primer lugar, las
normas europeas y que son los jueces ordinarios del Derecho de la
UE. La aplicación administrativa descentralizada de las normas
europeas justifica este control jurisdiccional descentralizado. El
nuevo artículo 19.1 TUE introduce un inciso que alude a esta
función esencial de los jueces nacionales, señalando que
los Estados miembros establecerán las vías de recurso necesarias para
garantizar la tutela judicial efectiva en los ámbitos cubiertos por el Derecho
de la Unión.
En el presente capítulo nos vamos a concentrar sólo, como es
lógico, en el estudio del TJUE.

La articulación del Tribunal de Justicia y de sus competencias de control jurisdiccional no ha experimentado cambios sustanciales desde su creación en 1957 (régimen competencial de atribución, carácter general de la jurisdicción del TJCE, mantenimiento del sistema de cooperación entre las instancias judiciales de la Unión y las instancias judiciales de los Estados miembros como eje del modelo, limitaciones del acceso de los particulares, no especificidad de las instancias judiciales comunitarias...). Los únicos cambios relevantes se han producido como consecuencia de la ampliación del número de países miembros y del aumento de los asuntos que ha tenido que resolver el Tribunal. Así, el Tribunal contó hasta 1988 con una única instancia, denominada Tribunal de Justicia, y todos los asuntos jurídicamente relevantes eran resueltos por la formación plenaria. Pero el aumento del contencioso hizo necesaria la creación de otra instancia, subordinada a la anterior y denominada Tribunal de Primera Instancia, y con posterioridad, el Tratado de Niza agregó otra instancia adicional, las Salas jurisdiccionales para contenciosos específicos, de las que sólo se creó en 2004 el Tribunal de la Función Pública de la Unión Europea, que ha desaparecido en septiembre de 2016, siendo sus jueces y sus competencias absorbidos por el Tribunal General.

El Tratado de Niza introdujo las modificaciones indispensables en el poder judicial de la UE para asimilar la incorporación de los doce nuevos Estados miembros y para hacer frente al aumento del volumen del contencioso debido al aumento de los países miembros y al incremento de las competencias atribuidas a la UE. Como consecuencia de ello, el TJUE ha sido la Institución en la que elTratado de Lisboa ha introducido menos cambios. Esto se debe aque este modelo de protección jurisdiccional ha constituido, sin duda, un pilar y un motor básico del proceso de integración y a que ha cumplido con éxito su función de garantizar el respeto del derecho en la interpretación y aplicación de las normas de la Unión.Ahora bien, en los últimos años este modelo ha comenzado a mostrar signos de agotamiento y han surgido propuestas doctrinales sobre modificaciones sustanciales de la denominada «arquitectura jurisdiccional comunitaria», completadas por el propio Tribunal mediante la publicación de un documento de reflexión en junio d 1999, titulado El futuro del sistema jurisdiccional de la Unión Europea1. Sin duda, el Tratado de Lisboa ha sido muy conservador en relación con el TJUE y no ha incorporado ninguna modificación sustancial que adecue el modelo de protección jurisdiccional a las nuevas exigencias de una UE ampliada y con mayores y más heterogéneas competencias. Se ha aplicado un planteamiento exclusivamente referido a la funcionalidad o efectividad del sistema judicial de la UE, partiendo, en consecuencia, de la base sobreentendida de que el esquema de racionalidad jurídica del modelo es apropiada. Un enfoque cuantitativo, el aumento de asuntos y del tiempo para su resolución, ha provocado undiagnóstico institucional cargado de realismo y de tratamientos prudentes. Tal vez demasiado realistas, tal vez demasiado prudentes.

El Tratado de Lisboa mantiene, por tanto, la misma articulación del poder judicial en la UE, introduciendo sólo cambios en la denominación de las instancias jurisdiccionales y algunas modificaciones en el procedimiento de designación de los jueces y abogados generales, así como en la articulación de algunos recursos jurisdiccionales. La desaparición de la estructura de pilares conlleva, también, una extensión de la jurisdicción del TJUE debido a la desaparición del sistema de control jurisdiccional limitado, existente en el tercer pilar de la antigua UE, relativo a la cooperaciónpolicial y judicial penal, aunque se mantiene la PESC al margen del control del TJUE, con la salvedad de lo previsto en el artículo 275 TFUE.

Continuando con la simplificación de las normas de Derecho originario aplicables al Tribunal, el Tratado de Lisboa incluye al TJUE en el artículo 13.1 TUE como una de las Instituciones de la UE y su regulación básica se contiene en el artículo 19 TUE, que es desarrollado ampliamente por los artículos 251 a 281 TFUE. Las disposiciones formalmente constitucionales, pero no materialmente constitucionales, sobre el TJUE se recogen en el Estatuto del Tribunal de Justicia de la Unión Europea, que constituye el Protocolo número 3, anejo al TUE, al TFUE y al TCEEA. Según el artículo 281.2 TFUE, el Parlamento Europeo y el Consejo, con arreglo al procedimiento legislativo ordinario, podrán modificar las disposiciones del Estatuto, a excepción de su Título I (estatuto de jueces y abogados generales) y su artículo 64 (régimen lingüístico). El Parlamento Europeo y el Consejo se pronunciarán bien a peticióndel Tribunal de Justicia y previa consulta a la Comisión, bien a propuesta de la Comisión y previa consulta al Tribunal de Justicia. Además, la regulación procesal complementaria la encontramos en los reglamentos de procedimiento de las tres instancias que componen en la actualidad el TJUE, cuya modificación ha propuesto el TJUE a raíz de la entrada en vigor del Tratado de Lisboa, juntocon algún cambio del Estatuto2.

De conformidad, con el nuevo artículo 19 TUE, el Tribunal de Justicia de la Unión Europea es el nuevo nombre de la Instituciónque asume la función jurisdiccional en la arquitectura institucional de la UE. Pero el TJUE, que es una Institución única, está integrado por dos instancias diferentes, que con el Tratado de Lisboa cambiaron también de denominación, a saber:
— La instancia superior seguirá denominándose «Tribunal de Justicia».
— La instancia intermedia se llama «Tribunal General»,desapareciendo por inadecuada la denominación anterior de Tribunal de Primera Instancia. 

El TFUE permite la creación de instancias inferiores, denominadas «tribunales especializados». En 2004 se creó elTribunal de la Función Pública de la Unión Europea, pero lamodificación operada por el Reglamento 2015/24223 del Estatuto del Tribunal de Justicia de la Unión Europea ha supuesto la desaparición de este tribunal especializado y la absorción de sus jueces y de su contencioso por el Tribunal General desde el 1 de septiembre de 2016.

Esta controvertida reforma de 2015 se está desarrollando en tres etapas que conducirán gradualmente hasta 2019 a la duplicación del número total de miembros del Tribunal General, que pasa de 28 a 56 jueces:
— 12 nuevos Jueces entran en funciones en 2016;
— Coincidiendo con la renovación parcial del Tribunal General en septiembre de 2016, el número de Jueces aumentará en 7 unidades, como consecuencia de la integración del Tribunal de la Función Pública en el Tribunal General;
— Con la siguiente renovación parcial del Tribunal General en septiembre de 2019, el número de Jueces aumentará en 9 miembros, alcanzándose el número total de 56 jueces.

Esta reforma pretende que el Tribunal General reduzca de forma duradera el número de asuntos pendientes, lo que tendrá como consecuencia una reducción de la duración media de los procedimientos ante él. Se elimina la triple instancia con la desaparición del único tribunal especializado (el TFPUE) y, de esta manera, se pretende simplificar la arquitectura jurisdiccional de la Unión. Esta reforma ha sido impulsada por el Tribunal de Justicia, pero con la oposición del Tribunal General, que era partidario de la creación de más tribunales especializados. Sin duda, la incapacidad de los Estados para ponerse de acuerdo sobre el nombramiento de los jueces del TFPUE y algunas dificultades de funcionamiento deeste tribunal especializado de composición reducida han sido losmotivos que han originado la reforma de 2015, que hadesembocado en una duplicación de jueces del Tribunal General, no solicitada por el Tribunal de Justicia, pero acordada por los Estados para evitar problemas con los turnos de rotación en los nombramientos. La duplicación de jueces del Tribunal General parece excesiva y se han puesto algunos límites al incremento del gasto con la reforma.

Las dos instancias del TJUE cuentan con una estructura jurisdiccional y administrativa especialmente concebida para que puedan cumplir la función que tienen encomendada en el seno de la estructura institucional europea. Esta misma idea inspira las normas procedimentales que se aplican a los diferentes tipos de recursos sustanciados ante ellas.

\subsubsection{Composición y funciones}

1. EL TRIBUNAL DE JUSTICIA
La instancia superior del TJUE continúa denominándose, como hasta ahora, Tribunal de Justicia (TJ, en adelante). En su análisis, hay que diferenciar la composición y su organización y funcionamiento.
1.1. COMPOSICIÓN

Los miembros del TJ son los jueces y los abogados generales, que designan, asimismo, un secretario para que les asista en el ejercicio de sus funciones.

El artículo 19.2 TUE acoge expresamente el principio de un juez por Estado miembro. Como consecuencia de ello, el TJ ha pasado a  tener veintiocho jueces desde la incorporación de Croacia el 1 de julio de 2013. Entre los Estados ha existido un amplio acuerdo sobre la conveniencia de mantener el principio de un juez por Estado miembro por la necesidad de que todos los sistemas jurídicos nacionales estén representados en el Tribunal y puedan, así, contribuir a la conformación de su jurisprudencia, y porque ello favorece la aceptación de dicha jurisprudencia en los derechos internos. Lógicamente, el inconveniente de este principio es que impide al TJ disponer de una formación plenaria con capacidad de trabajo jurisdiccional, dado el notable incremento del número de jueces.


2. EL TRIBUNAL GENERAL 

El Tribunal General (TG) es la instancia intermedia del TJUE, cuya denominación estableció ex novo el Tratado de Lisboa, ya que esta instancia se llamaba anteriormente Tribunal de Primera Instancia (TPI). Esta denominación ya no era adecuada, porque el TG conoce muchos asuntos en primera instancia, pero es tribunal de apelación respecto a las sentencias de los tribunales especializados.

La creación de una instancia inferior al Tribunal de Justicia se llevó a cabo con la adopción por el Consejo de la Decisión de 24 de octubre de 19885, mediante la cual se establecía el TPI, que inició su funcionamiento en noviembre de 1989. La creación del TPI perseguía un doble objetivo:
— la mejora de la protección judicial de los justiciables y, en especial, de sus derechos de defensa, sobre todo en recursos que requieren un estudio detenido de hechos complejos, como es el caso de los asuntos de competencia, dumping, ayudas de Estado y funcionarios, y
— la reducción del volumen de trabajo del TJ, para que pudiera concentrar su actividad en su labor esencial, que es velar por una interpretación uniforme del Derecho de la Unión.

El TPI se concibió, por tanto, como una jurisdicción de primera instancia, que debía ocuparse de todos aquellos asuntos de gran complejidad fáctica, para el examen de los cuales el TJ no estaba bien equipado técnicamente. El Tratado de Niza reforzó la posición del TPI, al proceder a su consolidación «constitucional» como instancia del TJUE con competencia general para el conocimiento de todos los recursos, salvo los atribuidos expresamente al Tribunal de Justicia y los pertenecientes a contenciosos específicos, que se encomendaron a las futuras salas jurisdiccionales.

El Tratado de Lisboa cambió la denominación de esta instancia para hacerla más acorde con sus funciones, pero mantiene básicamente la regulación tal como salió del Tratado de Niza.























































La misión fundamental del TJUE es hacer garantizar el respeto del Derecho europeo mediante la interpretación y aplicación de las normas europeas según criterios uniformes en todo el territorio de la Unión. En concreto, desempeña tres funciones principales: Controla de la legalidad de todos los actos, nonnativos y ejecutivos. que emanan de las instituciones de la Unión. El TJUE cotej a si el acto en cuestión regula una competencia interna de los Estados miembros. posee un defecto de forma sustancial o contradice una norm a de rango superior. El Tribunal puede ej ercer este controldirectamente, o a partir de recursos (recursos de anulación) interpuestos por un Estado miembro u otra institución europea. Un particular o empresa también puede presentar un recurso de anulación siempre que tenga interés directo (es decir, que se vea afectado por la aplicación de la norma comunitaria que cree defectuosa). Las mismas instancias pueden también interponer un recurso por omisión contra una institución europea si ésta. contradiciendo los Tratados, se ha abstenido de actuar (tras haber sido requerida previamente para que actúe). Controla a los Estados miembros para asegurar que cumplan las obligaciones que los Tratados y el Derecho derivado les imponen. La Comisión o, en casos excepcionales, un Estado miembro pueden iniciar un procedimiento de infracción si tienen razones para creer que un Estado miembro no cumple con sus obligaciones conforme al Derecho de la Unión. Interpreta el Derecho de la Unión para garantizar su correcta aplicación por los jueces nacionales. Si, al tratar de juzgar un litigio, un juez nacional tiene dudas sobre la conveniencia o no de aplicar una norma europea frente a una nacional o sobre la interpretación correcta de una norma comunitaria, éste puede plantear sus dudas al Tribunal de Justicia (cuestión prejudicial). La contestación que el Tribunal emita (sentencia prejudicial) se ceñirá estrictamente a la pregunta planteada, sin entrar en el fondo del asunto, pero tendrá carácter vinculante para el juez nacional.

Las sentencias del Tribunal (que se deciden por mayoría) son de obligado cumplimiento para los sujetos afectados, pudiendo el Tribunal imponer multas coercitivas a aquellos que no cumplan el contenido de las sentencias por las que hubieren sido condenados.

El TJUE, con sede en Luxemburgo, está integrado por dos órganos jurisdiccionales 19 :

El Tribunal de Justicia ( creado por el Tratado de París):
• Está compuesto por un Juez por cada Estado miembro y 11 Abogados Generales. Los Jueces y los Abogados Generales son designados de común acuerdo por los Gobiernos de los Estados miembros, previa consulta a un comité encargado de emitir un dictamen sobre la idoneidad de los candidatos. Su mandato es de seis años renovable. Los Jueces del TJ eligen de entre ellos al Presidente y al Vicepresidente por un período de tres años renovable. Los Abogados Generales están encargados de presentar, con toda imparcialidad e independencia, un dictamen jurídico en los asuntos que se les asignen.
• Como órgano jurisdiccional supremo de la Unión, resuelve las cuestiones esenciales para el ordenamiento comunitario, tales como los recursos interpuestos por las instituciones europeas y por los Estados miembros, así como las cuestiones prejudiciales. Además, actúa como Tribunal de Casación de las sentencias dictadas por el Tribunal General.
El Tribunal General (creado por el Acta Única Europea):
• Está compuesto por dos Jueces por cada Estado miembro. El nombramiento de los Jueces sigue las mismas normas que las del TJ.
• Atiende recursos de anulación y por omisión interpuestos por personas f isicas o jurídicas en casos de ilegalidad u omisión de actos jurídicos de la Unión, recursos de Estados contra la Comisión o el Consejo en determinados ámbitos, y recursos de indemnización por daños.

\subsection{Tribunal de Cuentas}
Además del control externo del PE, los Tratados han establecido otro sistema más de control externo no político a través del Tribunal de Cuentas. Es una de las siete instituciones de la Unión.

El Tribunal de Cuentas, calificado como la «conciencia financiera» de la Unión, fue creado mediante el Tratado de 22 de julio de 1975 al establecer un régimen de financiación íntegra mediante recursos propios. Por tanto, la autonomía financiera llevó a la intensificación del control financiero en la ejecución del presupuesto. Anteriormente fue, desde su creación en 1975 hasta el Tratado de Maastricht de 1992, un órgano auxiliar de la co-autoridad presupuestaria de la Unión (Consejo y PE).

Es hoy, pues, una institución, de naturaleza administrativa y no judicial, que ejerce funciones de control externo de las cuentas y de consulta con plena independencia en el interés general de la Unión. Ahora bien, aunque es una Institución en el sentido formal definido por el artículo 13 TUE, sin embargo ese Tratado no define sus funciones y se remite al Tratado de Funcionamiento, debido a que materialmente ejerce sus tareas de control asistiendo al Parlamento y al Consejo en tanto que diarquía presupuestaria.


\subsubsection{Composición y funcionamiento}
Se compone de un nacional por cada Estado miembro. Son propuestos por los Estados y nombrados para un período de seis años por el Consejo mediante decisión por mayoría cualificada y previa consulta (no vinculante) al Parlamento (art. 286 TFUE).

Las cualidades de carácter personal que deben reunir los miembros del Tribunal de Cuentas son las de profesionalidad e independencia, comprometiéndose a ejercer sus funciones con dedicación exclusiva y permanente. Los miembros del Tribunal de Cuentas son elegidos entre personalidades que han pertenecido en sus países respectivos a instituciones de control externo o que poseen una cualificación particular para esta función.

El Tribunal de Cuentas distribuye las atribuciones de control entre sus miembros. Cada uno tiene responsabilidad directa de control de gastos e ingresos en ciertos sectores de la Unión y sobre una parte de las actividades del Tribunal mismo. Cada miembro es responsable ante el Tribunal de Cuentas de realizar las verificaciones necesarias en los sectores que le han sido atribuidos y de estar informado de los restantes sectores de control. Los informes sectoriales se examinan colegiadamente y se decide el texto definitivo de su dictamen o informe.


\subsubsection{Atribuciones}
Su misión es fiscalizar la utilización de la masa presupuestaria por la Unión y el conjunto de sus organismos. Su función primordial es ejercer el control sobre «las cuentas de la totalidad de los ingresos y gastos» (art. 287.1 TFUE).

Se trata pues de un auditor externo de la Unión que debe contribuir a mejorar la gestión financiera; evalúa la obtención y lautilización de los fondos de la Unión, comprueba que lasoperaciones financieras se han contabilizado y presentado con corrección, se han ejecutado legal y regularmente y se han gestionado de manera que se asegura la economía, la eficiencia y la eficacia.

El control siempre es externo. El objeto de verificación son las cuentas. Pero las instituciones y los diferentes organismos siguen siendo responsables del control interno (vigilancia interna de la ejecución del presupuesto, del control de los compromisos y ordenación de los gastos y control de los ingresos). Los criterios en torno a los que debe desenvolverse el examen externo de los gastos e ingresos por el Tribunal de Cuentas deben ser la legalidad, la regularidad y la buena gestión financiera.

En relación con el criterio de la legalidad, el Tribunal examina si las operaciones respetan las normas presupuestarias de los tratados, el reglamento financiero y los actos jurídicos que les afecten. Pero debe tener en cuenta también la regularidad de la operación, es decir, si ha quedado reflejada correctamente desde un punto de vista contable. Por el contrario, el criterio de la buena gestión financiera no obedece a criterios jurídicos, sino que se trata de apreciaciones sobre el buen juicio de la operación tanto desde la perspectiva de la oportunidad política o económica, criterios éticos, la adecuada organización interna y el funcionamiento de los controles internos, la adecuación a los objetivos perseguidos, su coherencia con otras acciones, etc. El control es a posteriori, es decir, en el estadio final de la operación cuando ya hay un pago o transferencia en un ejercicio cerrado; pero no excluye que los controles puedan efectuarse antesdel cierre de las cuentas del ejercicio presupuestario, lo que facilita la rectificación de las irregularidades.

Además, el control tiene lugar sobre la documentación contable, pero, si fuera necesario también se pueden desplazar al lugar mismo de ejecución del gasto, ya sea en los servicios administrativos de la Unión, ya sea en las Administraciones públicas de los Estados miembros o en las dependencias de cualquier persona física o jurídica que perciba fondos del presupuesto. Este control en el territorio de los Estados miembros se efectuará en coordinación con las instituciones de control nacional u organismos competentes, sobre todo para no gravar las estructuras administrativas de la Unión de control y para hacerlas más eficaces. Las Instituciones y los Estados miembros se comprometen a suministrar todo documento y toda información necesaria al cumplimiento de su misión.

El resultado global de sus actividades de control se plasma en el Informe Anual que debe realizar después del cierre de cada ejercicio. La realización del Informe exige un comportamiento coordinado de las instituciones y órganos conforme a un calendario fijado por el reglamento financiero. El Tribunal de Cuentas envía a las Instituciones y órganos auxiliares las observaciones que figurarán en el Informe Anual6. Aquéllas tienen un plazo para enviar sus respuestas. El Informe Anual, junto con las respuestas de las Instituciones se publica en el Diario Oficial de la UE aproximadamente 10 u 11 meses tras el cierre del ejercicio en tiempo útil para que se exijan responsabilidades políticas o penales, a diferencia de lo que sucede en instituciones similares en algún Estado miembro, como España, cuyos informes se suelen presentar ya prescrita toda responsabilidad.

También puede emitir dictámenes. Debe ser consultado preceptivamente por el Consejo en varias ocasiones previstas por los Tratados e igualmente puede ser consultado con carácterfacultativo cuando así lo solicite alguna Institución; pero sucompetencia consultiva no depende únicamente de la petición preceptiva o facultativa por las Instituciones, pues el Tribunal de Cuentas podría, por propia iniciativa y en todo momento, expresar sus observaciones sobre cuestiones particulares.

Junto a estas competencias clásicas, los Tratados de Maastricht y de Ámsterdam acentuaron la finalidad de su función de control: el Tribunal de Cuentas deberá presentar al Parlamento y al Consejo una declaración sobre la fiabilidad de las cuentas, la regularidad y legalidad de las operaciones correspondiente, y podrá presentar informes especiales.












El Tribunal de Cuentas es un órgano independiente encargado del control de las finanzas públicas de la Unión. Se creó en 1975, comepzó a funcionar en Luxemburgo en 1977 y fue elevado al rango de institución por el Tratado de Maastricht. A pesar de su denominación, el Tribuna del Cuentas no ejerce funciones jurisdiccionales (sus miembros no son jueces) y carece de poder sancionador. Las funciones que realiza son las siguientes:

Controla la ejecución del presupuesto anual de la Unión, emitiendo una declaración anual (DAS, del francés déc/aration d ·assurance) sobre la fiabilidad de las cuentas, y la legalidad y regularidad de las operaciones subyacentes. Esto es fundamental para que la autoridad presupuestaria (Consejo y Parlamento) pueda aprobar ex-post la gestión del presupuesto realizado por la Comisión.

Puede auditar cualquier cuenta de ingresos o gastos de la Unión, por iniciativa propia o a instancia de otra institución. Tiene libertad para seleccionar sus temas de análisis, decidir el enfoque y elegir cómo presenta sus informes y dictámenes.

También examina las normas con impacto financiero. En algunos casos, esto es obligatorio (modalidades de recursos propios, medidas contra el fraude, etc.). En los que no, las demás instituciones pueden solicitarle que lo haga.

El Consejo nombra a los miembros del Tribunal de Cuentas (cada Estado miembro tiene un representante) para un mandato de seis años, a propuesta del Estado miembro en cuestión y previa consulta al Parlamento Europeo. El Tribunal de Cuentas elige a un Presidente de entre sus miembros, por un período de tres años renovable.


\subsection{Banco Central Europeo} 
Desde la entrada en vigor del TUE en 1993, el establecimiento de una Unión Económica y Monetaria (UEM) constituye uno de los objetivos principales de la Unión Europea (art. 3.4 TUE). En el presente Capítulo se pretende dar cuenta de la estructura institucional creada para gestionar el proceso de unión monetaria, que esencialmente se identifica con el Banco Central Europeo (BCE) y con el Sistema Europeo de Bancos Centrales (SEBC) (art. 282 TFUE).

Las disposiciones dedicadas a la cooperación en las políticas económica y monetaria entre los Estados miembros en la redacción inicial del Tratado de Roma, incluso tras su modificación por el AUE, resultaban muy breves y dejaban un amplio margen de discrecionalidad a los Estados miembros en la conducción de dichas políticas1. Esta regulación, no obstante, sirvió de base para l  articulación de un mecanismo denominado «serpiente monetaria», que fue el antecedente del Sistema Monetario Europeo. Se trataba en ambos casos de mecanismos de cooperación monetaria que permitieron limitar el alcance y la intensidad de las variaciones de los tipos de cambio entre las distintas monedas comunitarias desde los años setenta. Estos mecanismos de cooperación monetaria tuvieron un efecto muy beneficioso sobre el funcionamiento del mercado interior, ya que la estabilidad de los tipos de cambio favorece el comercio internacional.

Tras el fracaso de un primer proyecto de unión monetaria en la década de los setenta, sobre la base del Informe Werner (de 8 de octubre de 1970), el Consejo Europeo encargó a un Comité de Expertos presidido por Jacques Delors la preparación de un informe (presentado el 17 de abril de 1989) para la realización de una UEM por etapas. El calendario para la puesta en práctica de las mismas fue aprobado por el Consejo Europeo de Madrid de junio de 1989. La primera fase de la UEM comenzó el 1 de julio de 1990 y tenía como objetivo la completa liberalización de los movimientos de capital entre los Estados miembros. La segunda y la tercera fase de la UEM sólo podían llevarse a cabo mediante una reforma de los Tratados constitutivos, que tuvo lugar con el Tratado de Maastricht.

La segunda fase, iniciada el 1 de enero de 1994, implicaba la asunción de unos planes de convergencia que permitieran garantizar la estabilidad económica de los países que iban a compartir la futura moneda común. Desde el punto de vista institucional, para regular esta fase se creó el Instituto Monetario Europeo (IME), que se configuró como un órgano provisional, cuya función primordial era poner las bases para la creación del Banco Central Europeo (BCE) que habría de regir la tercera fase de la UEM.

La tercera, y última, fase comenzó el 1 de enero de 1999 y trajo consigo la fijación definitiva de los tipos de cambio de las monedas de los Estados que participaron desde su inicio en la fase definitivade la unión monetaria. Esta fase se identifica por la creación de una moneda común para todos estos países: el euro. A partir del 1 de enero de 2002, el euro sustituyó a las monedas nacionales de los primeros doce países que participaron en este proceso. Las antiguas monedas nacionales carecen hoy de curso legal.

Esta tercera fase supuso, desde la perspectiva institucional, la creación, como hemos visto, del BCE, que tiene como misión principal la elaboración de una política monetaria única para los Estados que comparten esa moneda común. Desde noviembre de 2014, el BCE asume también la supervisión prudencial de la mayoría de las entidades de crédito de los Estados miembros participantes en el Mecanismo Único de Supervisión (MUS).

El BCE fue creado el 1 de junio de 1998. Su regulación jurídica se encuentra en los artículos 127 a 133 y 282 a 284 TFUE, y en un Protocolo adicional anexo al TFUE «sobre los Estatutos del SEBC y del BCE», que como tal, forma parte de dicho Tratado. De esta manera, el estatuto jurídico del BCE forma parte del Derecho originario, lo que le proporciona una gran estabilidad normativa. Por otro lado, esta «constitucionalización» de los Estatutos puede generar dificultades si surge la necesidad de modificarlos. Precisamente por ello, se ha previsto un mecanismo de revisión de las disposiciones más técnicas de los Estatutos mediante el procedimiento legislativo ordinario (arts. 129.3 TFUE y 41 de los Estatutos), que evite la necesidad de recurrir a los mecanismos de revisión de los Tratados previstos en el artículo 48 TUE2.

El BCE tiene su sede en Francfort (Alemania), en virtud de la Decisión tomada en tal sentido por acuerdo entre los Jefes de Estado y de Gobierno de los Estados miembros el 29 de octubre de 19933 y goza de personalidad jurídica propia (arts. 282.3 TFUE y 9.1 de los Estatutos). Desde la reforma introducida por el Tratado de Lisboa, el BCE ha sido incluido formalmente entre las «instituciones» de la UE (art. 13 TUE). No obstante, su autonomía financiera (su presupuesto y recursos son independientes de los de la UE) y su personalidad jurídica propia explican que la Unión no sea responsable de los daños causados por el BCE en el curso de sus actuaciones, de las que exclusivamente responde este último (art. 340 TFUE). 

Desde un punto de vista monetario, el BCE aparece inmerso en un conjunto institucional más amplio, el SEBC, del que forman parte los bancos centrales de todos los Estados miembros. Ello responde a la necesidad de aprovechar la experiencia y el entramado orgánico de los bancos centrales nacionales en el sistema monetario de la Unión. Paralelamente, el BCE y los bancos centrales de los Estados miembros cuya moneda es el euro constituyen el «Eurosistema». Por tanto, los bancos centrales nacionales continúan existiendo y mantienen una personalidad jurídica propia, distinta de la del BCE. Por su parte, el SEBC como  conjunto carece de dicha personalidad. Ahora bien, en todo sistema en el que existen diversos órganos deben establecerse unos principios de relación entre ellos. En el SEBC estos principios son los de jerarquía y descentralización.

En virtud del principio jerárquico, el BCE ocupa la cúspide de esta estructura institucional, ya que la actuación de los bancos centrales nacionales en materia monetaria debe ajustarse a las orientaciones e instrucciones del BCE (art. 14.3 de los Estatutos). Asimismo, el BCE puede llevar ante el TJUE a un banco central nacional cuando considere que éste «ha incumplido algunas de las obligaciones que establecen los presentes Estatutos» (art. 35.6 de los Estatutos). El hecho de que el procedimiento precontencioso y la demanda ante el Tribunal de Justicia se dirijan directamente contra el banco central nacional, y no contra el Estado miembro, reafirman el carácter jerárquico de la relación entre el BCE y los bancos centrales nacionales, y se explica por la independencia de los bancos centrales respecto de sus gobiernos. Por otro lado, el carácter descentralizado del sistema se pone de manifiesto cuando el artículo 12.1, in fine, de los Estatutos establece que «el BCErecurrirá a los bancos centrales nacionales para ejecutar las operaciones que correspondan a las funciones del SEBC». En cualquier caso, esta descentralización no puede poner en peligro el carácter unitario de la política monetaria. De ahí que la ejecución por parte de los bancos centrales nacionales de las medidas de política monetaria decididas por el BCE tenga un límite previsto en el mismo artículo 12.1: la salvaguarda de las competencias de los órganos del propio BCE, que garantizan la uniformidad de la política monetaria de la Unión.

Esta estructura jerárquica y descentralizada se reproduce, aunque con significativas variaciones, en el diseño institucional del MUS. Desde noviembre de 2014 el BCE ha asumido directamente la  supervisión prudencial de las entidades de crédito de dimensiones sistémicas («significativas», en la terminología comunitaria) en los Estados miembros participantes en el Mecanismo. Ahora bien, las autoridades nacionales competentes en materia de supervisión  prudencial deben desarrollar sus competencias de control prudencial sobre el resto de las entidades de crédito de acuerdo con los reglamentos, directrices o instrucciones generales que les sean dirigidas por el BCE4. Si éste considera que las autoridades nacionales están incumpliendo sus obligaciones, puede en cualquier momento absorber sus funciones y ejercer por sí mismo las competencias de supervisión en relación con las entidades de crédito afectadas [art. 6.5.b) del Reglamento 1024/2013].


\subsubsection{Composición y funcionamiento}
El SEBC está compuesto por el BCE y por los bancos centrales nacionales de todos los Estados miembros (incluidos, por tanto, los bancos centrales de los Estados miembros «sometidos a unaexcepción», que son los que no participan en la zona euro), es decir, por veintinueve entidades con personalidad jurídica propia. En esta situación, los bancos centrales nacionales desarrollan dos tipos de funciones: las que les son asignadas por el derecho europeo (particularmente, en el ámbito monetario, en el seno del SEBC), y las que adicionalmente le corresponden a cada banco central nacional en función del Derecho interno. Como ya se ha dicho, el SEBC carece de personalidad jurídica propia, aunque actúa siguiendo el principio de unidad de acción, gracias a su estructura jerarquizada. Los órganos del BCE son al mismo tiempo los órganos del SEBC.

a) El Consejo de Gobierno está compuesto por los miembros del Comité Ejecutivo y por los Gobernadores de los bancos centrales nacionales de los Estados miembros cuya moneda es el euro (art. 283.1 TFUE). Este órgano debe reunirse al menos diez veces al año (aunque suele hacerlo cada dos semanas) y dichas reuniones tienen carácter confidencial, aunque el propio Consejo de Gobierno puede decidir hacer públicos los resultados de las mismas (art. 10.4 y 5 de los Estatutos), sobre los que habitualmente se proporcionan explicaciones en conferencias de prensa periódicas. Se trata del órgano más importante del SEBC puesto que entre sus misiones se encuentra la adopción de las «orientaciones y decisiones necesarias para garantizar el cumplimiento de las funciones» del SEBC, además de la «formulación de la política monetaria de la Unión» (art. 12.1 de los Estatutos). Adicionalmente, es el Consejo de Gobierno el órgano competente para adoptar «el Reglamento interno que determinará la organización interna del BCE y de sus órganos rectores» (art. 12.3 de los Estatutos)5, emprender acciones en nombre del BCE ante el TJUE (art. 35.5 de los Estatutos), tomar las principales decisiones de naturaleza financiera previstas en el Capítulo VI de los Estatutos o realizar las recomendaciones del BCE previstas en el procedimiento de revisión simplificado de esos mismos Estatutos (art. 40 de los Estatutos). También es el órgano que adopta formalmente las decisiones en el ámbito de la supervisión prudencial (art. 26.8 del Reglamento 1024/2013). De todas formas, si lo estima conveniente, el Consejo de Gobierno  puede delegar algunos de sus poderes en el Comité Ejecutivo (art.12.1 de los Estatutos). 

El procedimiento de votación en el seno del Consejo de Gobierno es complejo y varía en función de la cuestión a tratar. La regla general es la mayoría simple, lo que acentúa la independencia de la Institución y hace que prime su naturaleza de órgano de integración. Cada miembro del Consejo de Gobierno dispone en principio de un voto (aunque el Presidente goza de un voto de calidad en caso de empate —art. 10.2 de los Estatutos—). No obstante, desde enero de 2015 los gobernadores se dividen en grupos (dos grupos en la actualidad, mientras el número de gobernadores se sitúe entre quince y veintiuno6, y tres grupos a partir del momento en que haya veintidós gobernadores), en función del PIB y del balance agregado total de las instituciones financieras monetarias del Estado miembro que representen. En tales circunstancias, el número máximo de votos se ha fijado en veintiuno. Los miembros del Comité Ejecutivo conservan un voto cada uno (seis votos) y, por tanto, el número de gobernadores con derecho a voto es de quince. Ese derecho de voto va rotando mensualmente entre todos los gobernadores representados en el Consejo de Gobierno, en cada uno de los grupos creados, en función de los criterios previstos en el artículo 10.2 de los Estatutos y en el Reglamento interno del BCE. En cualquier caso, todos los miembros del Consejo de Gobierno podrán seguir asistiendo a sus reuniones, en las fases en las que carezcan temporalmente de derecho de voto.

Adicionalmente, para determinadas decisiones financieras se establece un sistema de voto ponderado, en el que el valor de los votos de los distintos miembros del Consejo de Gobierno depende de la participación de cada banco central nacional en el capital suscrito del BCE. Los miembros del Comité Ejecutivo no tienen derecho de voto en este sistema, en el que se establecen dos variantes: la mayoría simple (con el voto ponderado) y la mayoría cualificada (para la que se requieren un número de votos favorables que representen al menos dos tercios del capital suscrito y representen al menos la mitad de los accionistas) (art. 10.3 de los Estatutos). Por último, el Estatuto prevé algunos supuestos en los que la decisión debe adoptarse por mayoría de dos tercios del Consejo de Gobierno (p. ej., art. 20 de los Estatutos), y un caso en el que se precisa la unanimidad de los mismos (art. 40.3 del Estatuto).

b) El Comité Ejecutivo está compuesto por un Presidente, un Vicepresidente y otros cuatro miembros, todos con dedicación exclusiva, que son nombrados para un período de ocho años por el Consejo Europeo por mayoría cualificada (sobre la base de una recomendación del Consejo y previa consulta al Parlamento Europeo y al Consejo de Gobierno) de entre personalidades de experiencia y reconocido prestigio profesional en asuntos monetarios o bancarios (art. 11.2 de los Estatutos). El largo período de tiempo que dura el mandato (ocho años) pretende garantizar la independencia de los miembros de este órgano; sin embargo, en sentido contrario, el hecho de que este mandato no sea renovable muestra una preferencia de los Estatutos por la rotación de los miembros del Comité Ejecutivo entre los Estados miembros antes que aprovechar la experiencia de los mismos en un segundo mandato. La separación de un miembro del Comité Ejecutivo, por dejar de reunir los requisitos exigidos para desempeñar sus funciones o si en su conducta se observara una falta grave, ha de ser decidida por el TJ a petición del Consejo de Gobierno o del Comité Ejecutivo (art. 11.4 de los Estatutos).

Los poderes del Comité Ejecutivo comprenden la puesta en práctica de la política monetaria de acuerdo con las orientaciones y decisiones adoptadas por el Consejo de Gobierno, impartiendo en ese contexto las instrucciones necesarias a los bancos centrales nacionales (art. 12.1 de los Estatutos), así como la preparación de las reuniones del Consejo de Gobierno (art. 12.2 de los Estatutos) y la gestión ordinaria del BCE (art. 11.6 de los Estatutos). Además, como ya se ha indicado, el Comité Ejecutivo puede desempeñar las funciones que le sean delegadas por el Consejo de Gobierno (art. 12.1 de los Estatutos). La adopción de decisiones en el Comité Ejecutivo se rige por el principio de mayoría simple, disfrutando cada miembro del órgano de un voto. En caso de empate, corresponde al Presidente el voto decisivo (art. 11.5 de los Estatutos).

c) El tercer órgano rector del BCE, en tanto existan Estados miembros acogidos a una excepción, es el Consejo General. Este órgano está compuesto por el Presidente y el Vicepresidente del BCE, así como por los Gobernadores de todos los bancos centrales nacionales de los Estados miembros. Los demás miembros del Comité Ejecutivo pueden participar en las reuniones del Consejo General, aunque sin derecho a voto (art. 44 de los Estatutos). El Presidente del Consejo y un miembro de la Comisión también pueden participar en las reuniones de este órgano sin derecho de voto (art. 45.2 de los Estatutos). La misión del Consejo General es establecer un foro de cooperación que permita mantener vinculados a la política monetaria de la Unión a los Estados miembros acogidos a una excepción, y desarrollar las funciones que antes realizaba el IME en relación con los mismos. La existencia por un período de tiempo indefinido de un órgano de carácter transitorio, en el que los Estados miembros acogidos a una excepción discutan cuestiones relativas a la política monetaria de la Unión, no estuvo exenta de polémica. Sin embargo, teniendo en cuenta el carácter limitado de las competencias del Consejo General, este órgano parece más bien estar diseñado para inducir a esos Estados miembros acogidos a una excepción a participar plenamente en la tercera fase de la UEM.

De hecho, las funciones decisorias del Consejo General selimitan en lo esencial a las tareas del antiguo IME que deban seguir desempeñándose durante la tercera fase en relación con los Estados acogidos a una excepción (art. 43 de los Estatutos), a la decisión sobre el desembolso que estos últimos deberán hacer sobre su porcentaje en el capital del BCE como contribución a sus costes operativos (art. 47 de los Estatutos) y a la adopción del Reglamento interno del propio Consejo General (art. 45.4 de losEstatutos). El resto de las funciones atribuidas a este órgano tienenun carácter consultivo, preparatorio o técnico y están taxativamente previstas en el artículo 46 de los Estatutos. En cualquier caso, esas funciones no pueden menoscabar las competencias atribuidas al Consejo de Gobierno ni al Comité Ejecutivo. Los Estatutos no establecen ningún sistema de votación para el Consejo General7 por lo que es el Reglamento interno del órgano el que regula esta cuestión8.

d) El Consejo de Supervisión. Este órgano está compuesto por un Presidente, un Vicepresidente (elegido de entre los miembros del Comité Ejecutivo), cuatro representantes del BCE (que no pueden ejercer funciones directamente relacionadas con la política monetaria del BCE) y un representante de la autoridad nacionalcompetente en materia de supervisión de cada Estado miembro participante en el MUS. El Presidente y el Vicepresidente son nombrados por el Consejo por mayoría cualificada (sólo votan los Estados miembros participantes en el MUS), a propuesta del BCE y previa aprobación del Parlamento Europeo (art. 26 del Reglamento 1024/2013). Si el Presidente deja de reunir las condiciones necesarias para el ejercicio de su cargo o comete una falta grave, puede ser destituido mediante una decisión ejecutiva del Consejo, a propuesta del BCE y previa aprobación del Parlamento Europeo. El Vicepresidente, en tanto que miembro del Comité Ejecutivo, sólo puede ser cesado mediante el procedimiento previsto a tal efecto para los miembros de ese órgano, aunque el Parlamento Europeo oel Consejo pueden informar al BCE de que a su juicio se cumplen las condiciones para la destitución, quedando entonces el BCE obligado a pronunciarse sobre el asunto.

El Consejo de Supervisión es el órgano encargado de elaborar y ejecutar las funciones de supervisión prudencial atribuidas al BCE. No obstante, el órgano formalmente decisorio en materia de supervisión es el Consejo de Gobierno, ya que el Consejo deSupervisión se limita a presentarle «proyectos completos dedecisiones para su adopción por éste» (art. 26.8 del Reglamento 1024/2013). En principio, lo habitual es que las propuestas del Consejo de Supervisión se aprueben mediante el procedimiento de no objeción (que el Consejo de Gobierno no se oponga a la decisión en un plazo que varía dependiendo del tipo de decisión, pero que no será superior a diez días hábiles). Para los casos en que el Consejo de Gobierno objete una decisión del Consejo de Supervisión, se ha creado una Comisión de mediación, compuesta por un representante por cada Estado miembro participante, elegido por cada Estado miembro entre los miembros del Consejo de Gobierno y del Consejo de Supervisión. Esta Comisión tomará sus decisiones por mayoría simple y en ella cada miembro tendrá un voto. Con la intención de favorecer su operatividad, se ha previsto que el Consejo de Supervisión adopte sus decisiones por mayoría simple, disponiendo cada miembro de un voto (aunque el Presidente goza de un voto de calidad en caso de empate).
 
A pesar del funcionamiento colegiado del Consejo de Gobierno, del Comité Ejecutivo y del Consejo General, no puede dejar de resaltarse la figura del Presidente del BCE. Además de presidir las reuniones de estos tres órganos, goza de un voto de calidad para resolver los empates, como ya se ha señalado con anterioridad. Es también la persona que puede comprometer legalmente al BCE frente a terceros (art. 38 de los Estatutos) y quien representa a esta Institución en el exterior (art. 13.2 de los Estatutos).

La cuestión de la representación del SEBC en las relaciones monetarias y financieras internacionales no quedó definitivamente resuelta en la reforma de Maastricht. El artículo 138 TFUE establece que es el Consejo (sólo votan los Estados miembros cuya moneda es el euro), a propuesta de la Comisión y previa consulta al BCE, quien decide sobre las posiciones comunes y sobre la representación (que podrá ser única) de la Unión en las Instituciones financieras internacionales. Por su parte, el artículo 6 de los Estatutos indica que en el ámbito de la cooperación monetaria internacional, «el BCE decidirá cómo estará representado el SEBC», añadiendo a continuación que «el BCE y, siempre que éste lo apruebe, los bancos centrales nacionales podrán participar en instituciones monetarias internacionales». Corresponde por tanto al Consejo de Gobierno decidir sobre la forma de representación del SEBC en el exterior (art. 12.5 de los Estatutos), aunque dentro del marco establecido por el Consejo sobre la base del artículo 138 TFUE (art. 6.3 de los Estatutos). Los bancos centrales nacionales de los Estados miembros acogidos a una excepción conservan sus competencias externas «con arreglo a la legislación nacional» (art. 42.1 y 2 de los Estatutos).

El Consejo Europeo de Viena de diciembre de 1998 aprobó un informe presentado por el Consejo EcoFin en el que se diseñaba con carácter general la representación exterior de la zona euro, sobre la base de las disposiciones del Tratado antes comentadas. Sin embargo, la cuestión sigue abierta, en la medida en que se ha adoptado una actitud flexible y pragmática en las distintas Instituciones financieras internacionales, acomodando la representación comunitaria en cada momento a las competencias de cada uno de dichos organismos, a su contexto jurídico, y a la posición de los terceros Estados en el seno de las mencionadas Instituciones9. Se ha articulado así una complejísima variedad de fórmulas de representación, todavía en evolución, en la que la resistencia de los Estados miembros a perder su estatus en las Instituciones financieras internacionales redunda, a veces, en perjuicio de la eficacia de esta acción exterior y de la proyección internacional del euro.



3. EL ESTATUTO DE LOS ESTADOS MIEMBROS NO PARTICIPANTES EN LA EUROZONA
Para poder adoptar el euro como moneda común, los Estados miembros debieron cumplir unos criterios jurídicos y económicos de convergencia, que aparecen reflejados en el artículo 140 TFUE y son detallados en un Protocolo anexo a los Tratados. En la actualidad diecinueve Estados comparten el euro como moneda común y otros nueve mantienen su moneda nacional. El artículo 140 TFUE establece el procedimiento que debe seguirse para que nuevos Estados miembros puedan incorporarse al Eurosistema, adoptando también el euro como moneda Los Tratados prevén un régimen específico (arts. 139 a 144 TFUE) para los Estados miembros que conservan su moneda nacional, denominándolos «Estados miembros acogidos a una  excepción». Este régimen se aplica a todos los Estados miembros que no comparten el euro como moneda, aunque el Reino Unido y Dinamarca disfrutan de regímenes jurídicos particulares establecidos en sendos Protocolos anexos a los Tratados.

Examinaremos, en primer lugar, el régimen general para los Estados miembros acogidos a una excepción, en segundo lugar, estudiaremos las particularidades de las disposiciones aplicables a Dinamarca y al Reino Unido y, posteriormente, haremos mención a la situación de los Estados miembros acogidos a una excepción en el MUS.

El Consejo Europeo de Ámsterdam adoptó una Resolución, de 16 de junio de 199731, mediante la que se decidió crear un nuevo mecanismo de tipos de cambio para la tercera fase de la UEM, para limitar las fluctuaciones entre el euro y las monedas de los Estados miembros acogidos a una excepción32. Este instrumento, conocido como MTC II, sustituyó al SME desde el 1 de enero de 199933. El órgano encargado de gestionar y administrar el sistema de cambios y el mecanismo de intervención y financiación del MTC II es el Consejo General del BCE.

3.1. «ESTADOS MIEMBROS ACOGIDOS A UNA EXCEPCIÓN»

Según lo dispuesto en el artículo 139.2 TFUE, no son aplicables a los Estados miembros acogidos a una excepción las orientaciones generales de política económica que afecten a la zona euro (art. 121.2 TFUE), o las disposiciones relativas al procedimiento de sanción en los casos de déficit público excesivo (art. 126.9 y 11 TFUE), a los objetivos y funciones del SEBC (art. 127.1, 2, 3 y 5 TFUE), a la emisión de billetes de banco y moneda metálica (art. 128 TFUE), al poder normativo del BCE (art. 132 TFUE), al uso del euro (art. 133 TFUE), a la política cambiaria con las monedas de terceros países (arts. 219 y 138 TFUE) y a la elección de los miembros del Comité Ejecutivo (art. 283.2 TFUE). Estos Estados también carecen de derecho de voto en el Consejo cuando se adoptan medidas sobre la base de disposiciones que no les son deaplicación, adaptándose para tales supuestos la regla de la mayoría cualificada [art. 238.3.a) TFUE]. A cambio, los bancos centrales de los Estados miembros acogidos a una excepción conservan sus competencias en materia monetaria con arreglo a la legislación nacional (art. 42.2 de los Estatutos) y la autoridad nacional competente también mantiene su autonomía en cuanto a la políticacambiaria, con el límite establecido en al artículo 142 TFUE (la política cambiaria nacional se considera como una cuestión de interés común).

Desde un punto de vista institucional, los bancos centrales de los Estados miembros acogidos a una excepción forman parte del SEBC (art. 282 TFUE). Su participación se articula a través del Consejo General, cuya composición y funciones ya han sido mencionadas. El artículo 140 TFUE establece el procedimiento mediante el cual los Estados miembros acogidos a una excepción pueden acceder a la moneda única, una vez alcanzada la convergencia jurídica y económica necesaria con arreglo a los criterios establecidos en ese mismo artículo. Para adoptar ese paso es necesaria una decisión positiva del Consejo (en su composición íntegra), tras recibir una recomendación por mayoría cualificada de los Estados miembros cuya moneda es el euro [de acuerdo con el sistema de votación previsto en el art. 238.3.a) TFUE].

3.2. LOS REGÍMENES ESPECIALES DE DINAMARCA Y EL REINO UNIDO

En cuanto a los regímenes jurídicos específicos, el Protocolo 16 aplicable a Dinamarca es relativamente simple ya que indica que en el caso de que este país decida no participar en la tercera fase de la UEM se le considerará como un «Estado miembro acogido a una excepción». Por tanto, a este país le es de aplicación el régimen general que se acaba de describir, con la particularidad de que el procedimiento de ingreso en la tercera fase previsto en el artículo 140 sólo puede ponerse en marcha a petición de Dinamarca.

Por lo que se refiere al Reino Unido, la situación es diferente, ya que el Protocolo 15 anexo a los Tratados que regula la situación de este país establece un régimen jurídico especial. En principio, elReino Unido no tiene ninguna obligación de participar en la fase definitiva de la unión monetaria (art. 1 del Protocolo), aunque si este país desea participar en la moneda única y cumple los requisitos de convergencia, no tiene más que solicitarlo, siguiendo elprocedimiento previsto en el artículo 140 TFUE (art. 9 del Protocolo). En realidad, la mayor parte de las disposiciones del Protocolo sobre el Reino Unido reproducen las condiciones establecidas para los Estados miembros acogidos a una excepción. No obstante, existen algunas excepciones significativas. Así, por ejemplo, el Reino Unido no asume una obligación jurídica de evitar los déficit públicos excesivos o de consultar al BCE sobre los proyectos legislativos nacionales en materia monetaria (art. 4 del Protocolo).

3.3. EL MUS Y LOS ESTADOS MIEMBROS ACOGIDOS A UNA EXCEPCIÓN

Todos los Estados miembros de la Eurozona participan  automáticamente en el MUS (son «Estados miembros participantes» en la terminología del Reglamento 1024/2013). Sin embargo, los Estados miembros que conservan su moneda nacional pueden optar entre participar en el MUS (con lo que se convertirían en «Estados miembros participantes») o mantener su sistema independiente de supervisión prudencial (con el régimen previsto para los «Estados miembros no participantes»).

Para integrarse en el MUS, los Estados miembros de fuera de la Eurozona deben establecer una «cooperación estrecha»34 entre el BCE y la autoridad nacional competente en materia de supervisión. La cooperación estrecha se establece mediante una Decisión del BCE, previa petición del Estado miembro interesado, en la que éste se comprometa a que su autoridad nacional competente cumpla las orientaciones o solicitudes formuladas por el BCE y proporcione toda la información necesaria sobre las entidades de crédito establecidas en su territorio, una vez adoptada la legislación nacional pertinente (art. 7.2 del Reglamento 1024/2013). A partir del momento en que un Estado miembro establece una cooperación estrecha con el BCE, los poderes de este último en relación con las entidades de crédito establecidas en su territorio son los mismos que si se tratase de un Estado miembro de la Eurozona. No obstante, en la medida en que los Estados miembros de fuera de la Eurozona no están representados en el Consejo de Gobierno del BCE, que es el órgano que tiene la última palabra sobre las decisiones supervisoras de esta institución, se ha articulado un mecanismo de salida que adopta varias modalidades y que permitiría a un Estado miembro cuya moneda no sea el euro retirarse (o ser expulsado) del MUS si se produce un desencuentro fundamental entre sus autoridades supervisoras y el BCE. Así, los Estados miembros participantes cuya moneda no es el euro pueden dar por terminada la cooperación estrecha una vez transcurridos tres años desde su inicio, también tienen la posibilidad de notificar que no están de acuerdo con una decisión del Consejo de Supervisión confirmada por el Consejo de Gobierno y dar por terminada la cooperación estrecha con efecto inmediato o, por último, pueden notificar que no se considerarán vinculados por una decisión que deba modificarse tras el rechazo del Consejo de Gobierno a una propuesta del Consejo de Supervisión, en cuyo caso el BCE podría suspender o hacer cesar la cooperación estrecha (párrafos 6, 7 y 8 del art. 7 del Reglamento 1024/2013). 

Por su parte, los Estados miembros de fuera de la Eurozona no participantes en el MUS conservan plenamente sus competencias supervisoras sobre las entidades de crédito establecidas en su territorio, aunque deben coordinar dicha actividad con el BCE y con las autoridades supervisoras de los Estados miembros participantes en el seno de la Autoridad Bancaria Europea en las condiciones establecidas por el Reglamento (UE) n.º 1093/201035.


\subsubsection{Atribuciones}
a) Tanto los artículos 127.1 y 282.2 TFUE como el artículo 2 de los Estatutos destacan que el objetivo primordial del SEBC es «mantener la estabilidad de precios». Éste es el signo de identidad que preside las decisiones del BCE. Esta opción jurídica por un determinado modelo de política monetaria pretende configurar el euro como una moneda fuerte que tenga la misma estabilidad que las antiguas monedas nacionales de los Estados miembros menos inflacionistas. En esta línea, la independencia del BCE pretende facilitar la labor del SEBC a la hora de adoptar las medidas que permitan mantener esa estabilidad de precios.En consecuencia, corresponde al BCE determinar lo que puede considerarse un nivel aceptable de inflación y adoptar las medidas oportunas para corregirlo en el caso de que se estime excesivo. Aunque nos hallamos sin duda ante una obligación legal establecida en el Derecho originario, parece difícil articular un control jurisdiccional de esa obligación.

b) El segundo objetivo del SEBC, que debe desarrollarse sin perjuicio del primero (la estabilidad de precios), es el apoyo a «las políticas económicas generales en la Unión» (arts. 127.1 TFUE y 2  de los Estatutos). No se trata únicamente de contribuir a la política económica de la Unión, ya que la «política económica», en sentido estricto, permanece esencialmente en manos de los Estados miembros, aunque es objeto de coordinación, más estrecha entrelos miembros de la zona euro10. En consecuencia, el SEBC no sólo debe tener en cuenta las «orientaciones generales para las políticaseconómicas de los Estados miembros» que aprueba el Consejo, según el procedimiento previsto en el artículo 121 TFUE, y los objetivos que se establezcan en el marco de la coordinación económica articulada en la Eurozona sobre la base del artículo 136 TFUE, sino que debe intentar apoyar las grandes líneas de la política económica de los Estados miembros, siempre y cuando ello no comprometa su objetivo de garantizar la estabilidad de precios. En cualquier caso, el BCE goza de total libertad de opinión cuando emite sus dictámenes sobre la base del artículo 127.4 TFUE.

En este contexto, se han planteado debates sobre los límites jurídicos de una política monetaria expansiva del BCE que pretenda contribuir al éxito de las políticas contra el desempleo en la Unión, facilite la financiación de la deuda pública de los Estados miembros o estimule el crecimiento11. El TJ ha señalado que las repercusiones económicas indirectas de las medidas monetarias del BCE no alteran su naturaleza monetaria y están amparadas por esta obligación de prestar apoyo a las políticas económicas generales de la Unión12. 

c) Por último, los artículos 127.1 TFUE y 2 de los Estatutos señalan que el SEBC actuará según el principio de una economía de mercado abierta y de libre competencia, favoreciendo una eficiente asignación de recursos y conforme a los principios establecidos en el artículo 119 TFUE. Estos principios son los siguientes: precios estables, finanzas públicas y condiciones monetarias sólidas y balanza de pagos estable.


2.2.2. Funciones

Las funciones del SEBC aparecen recogidas en los artículos 3 y siguientes de los Estatutos y en el artículo 127 TFUE. Estas funciones corresponden tanto al BCE como a los bancos centrales nacionales, a través de los cuales actúa el SEBC. Cuando las funciones son diferentes, el TFUE y los Estatutos se refieren expresamente al BCE o a los bancos centrales nacionales (p. ej., en el art. 128.1 TFUE en relación con la emisión de billetes de euro).


a) Sin duda, la primera de esas funciones es también la principal. El SEBC tiene por misión «definir y ejecutar la política monetaria de la Unión». La redacción del artículo 3.1, primer guión, de los Estatutos es incondicionada y proporciona una enorme libertad al Eurosistema a la hora de elegir los objetivos monetarios y adoptar las medidas que considere más oportunas para conseguirlos13, configurándose como una competencia exclusiva de la Unión [art. 3.1.c) TFUE]. Aunque no existe una definición en el Tratado de lo que debe entenderse por política monetaria, un repaso a los poderes del SEBC permite comprobar que se encuentran incluidos  dentro de esa expresión, entre otros, la apertura de cuentas en el BCE y en los bancos centrales nacionales a entidades de crédito o a entidades públicas (art. 17 de los Estatutos), las operaciones en mercado abierto tanto de cambio como crediticias (art. 18 de los Estatutos), el establecimiento de reservas mínimas para las entidades de crédito establecidas en los Estados miembros (art. 19 de los Estatutos) o la realización de determinadas operaciones  financieras exteriores (art. 23 de los Estatutos).

Estas funciones aparecen incluidas dentro de los preceptos de los Estatutos que son susceptibles de modificación a través del procedimiento de revisión simplificada (art. 40 de los Estatutos), por lo que estas competencias pueden verse aumentadas si la evolución de los mercados monetarios internacionales aconseja la atribución de nuevos poderes para que el SEBC siga cumpliendo eficazmente sus objetivos.

b) El artículo 119 TFUE coloca al mismo nivel la política monetaria y la política de cambios cuando indica que esta última será «única» para la Unión. En ese contexto, la segunda función del SEBC consiste en realizar operaciones de cambio de divisas, aunque con un estricto respeto al reparto de competencias realizado sobre la base del artículo 219 TFUE (arts. 127.2 TFUE y 3 de los Estatutos). La competencia para realizar operaciones de cambio de divisas incluye también las operaciones con monedas de aquellos Estados miembros que no participan en la moneda única.

Durante la negociación del Tratado de Maastricht se desató una fuerte polémica entre quienes deseaban que la política de cambio fuese decidida por las autoridades políticas (el Consejo) y quienes preferían que dicha política fuese competencia del BCE, para reforzar así su papel de guardián de la estabilidad de precios. El compromiso que refleja la redacción final del artículo 219 TFUE es complejo, puesto que su párrafo 1.º otorga al Consejo competencia para celebrar (por unanimidad) acuerdos internacionales14 en los que se cree un sistema de tipos de cambio en el que participe el euro junto con monedas de terceros Estados, pudiendo esta misma Institución decidir (por mayoría cualificada) sobre la adopción, el ajuste o el abandono del tipo central del euro en ese sistema.

Adicionalmente, el párrafo 2.º del artículo 219 establece que el Consejo (por mayoría cualificada) podrá formular «orientaciones generales» para la política de cambios comunitaria, que en ningún caso pueden atentar contra el objetivo preferente de la estabilidad de precios. Ciertamente, estas orientaciones generales no tienen un carácter vinculante para el BCE (art. 130 TFUE); sin embargo, tampoco parece razonable que el BCE pueda ignorar dichas orientaciones generales. En todas estas cuestiones (art. 219.1 y 2 TFUE), el Consejo decide sobre la base de una recomendación del BCE, o de la Comisión y previa consulta del BCE.

c) La tercera de las competencias del SEBC se refiere a su capacidad para «poseer y gestionar las reservas oficiales de divisas de los Estados miembros»15 (arts. 127.2 TFUE y 30.1 de los Estatutos). La utilización de las expresiones «poseer y gestionar», unidas a la referencia a los Estados miembros, parecen sugerir que la propiedad de las reservas oficiales permanece en manos de las autoridades nacionales. En cualquier caso, esa propiedad teórica de las reservas no afecta a la libertad de acción del BCE en su manejo, ya que éste tiene «pleno derecho a poseer y gestionar las reservas exteriores que le sean transferidas» (art. 30.1 de los Estatutos).

En realidad, las reservas exteriores de divisas de los Estados miembros se dividen en tres grupos: a) una parte se transfiere al BCE, determinándose la aportación de manera proporcional a la participación de cada banco central nacional en el capital del BCE  (art. 30.2 de los Estatutos); b) el resto de los activos exteriores de reserva que queden en manos de los bancos centrales nacionales sólo pueden ser utilizados previa aprobación del BCE, para garantizar así la coherencia de la política monetaria y de cambio de la Unión (art. 31.2 de los Estatutos); c) los Gobiernos de los Estados miembros se reservan el derecho a la tenencia y gestión de «fondos de maniobra oficiales» en divisas, necesarios para el desarrollo cotidiano de su actividad, que podrán utilizarse libremente en operaciones por debajo de un determinado límite establecido por el Consejo de Gobierno, y con la autorización del BCE en las operaciones de un montante por encima de dicho límite (art. 31.2 de los Estatutos).

d) Los artículos 127.2 TFUE y 3.1 de los Estatutos identifican como cuarta función del SEBC la promoción de un «buen funcionamiento del sistema de pagos» en la Unión. El artículo 22 de los Estatutos concreta el significado de esta misión señalando que «el BCE y los bancos centrales nacionales podrán proporcionar medios y el BCE dictar reglamentos, destinados a garantizar unos sistemas de compensación y liquidación eficientes y solventes dentro de la Unión, así como con otros países». La articulación de sistemas de compensación es una función tradicional de los bancos centrales que también es asumida por el SEBC. El buen funcionamiento del sistema de pagos es una condición indispensable para el éxito de la moneda común. En este contexto, el BCE ha creado el sistema TARGET216, que conecta los sistemas de compensación y liquidación en euros de los bancos centrales nacionales17, conformándose así un sistema de pagos europeo que permite realizar transacciones en tiempo real entre los sistemas de pago nacionales y sirve de vehículo a las decisiones de política monetaria adoptadas por el propio BCE.

e) Los artículos 127.5 TFUE y 3.3 de los Estatutos confieren al SEBC la responsabilidad de «contribuir» a la «buena gestión de las políticas que lleven a cabo las autoridades competentes con respecto a la supervisión prudencial de las entidades de crédito y la estabilidad del sistema financiero». Más específicamente, los artículos 127.6 TFUE y 25.2 de los Estatutos exigen la unanimidad en el Consejo para encomendar al BCE «tareas específicas» sobre supervisión prudencial de las entidades financieras con excepción de las compañías de seguros. La redacción de estos preceptos es poco clara, ya que refleja un difícil compromiso entre quienes deseaban que el BCE asumiera un papel central en la definición de la supervisión prudencial en la Unión, y quienes estimaban que esa misión podía perjudicar su independencia y entrar en colisión con su objetivo de mantener la estabilidad de precios sin conflictos de intereses. No obstante, la necesidad de reforzar la credibilidad y la independencia de los mecanismos de supervisión en la UE ha llevado al Consejo a atribuir al BCE competencias de supervisión de las entidades de crédito de carácter sistémico en la Eurozona y en los demás Estados miembros que establezcan una «cooperación estrecha» con el BCE a tal efecto, en virtud del Reglamento (UE) n.º 1024/2013, todo ello en el contexto más amplio de una unión bancaria18. En virtud de este Mecanismo Único de Supervisión (MUS), el BCE ejerce la supervisión directa sobre esas entidades sistémicas19 (significativas, en el lenguaje del mencionado Reglamento) desde noviembre de 2014 en los Estados participantes, quedando las demás entidades de crédito bajo la supervisión de las autoridades nacionales. En cualquier caso, en los Estados participantes, las autoridades nacionales deberán desarrollar sus funciones respetando los reglamentos, directrices o instrucciones generales emanados del BCE (art. 6.5 del Reglamento n.º 1024/2013).

Desde una perspectiva más general, la crisis financiera ha propiciado la creación en 2010 de un conjunto orgánico denominado «Sistema Europeo de Supervisión Financiera», con el propósito de mejorar la coordinación de la supervisión prudencial en toda la UE, y en el que el BCE tiene una participación variada, dependiendo del órgano20.

f ) Además de las funciones principales asignadas al SEBC por el artículo 127 TFUE que acaban de mencionarse, es posibleidentificar algunas otras funciones del sistema que aparecen recogidas en diversas disposiciones del Tratado y de los Estatutos. Entre ellas cabe destacar el control sobre la emisión de billetes de banco y de moneda metálica en la Unión, lo que se identifica como uno de los principales instrumentos de la política monetaria y como una de las competencias tradicionales de un banco central (arts. 128 TFUE y 16 de los Estatutos), así como determinadas funciones de recopilación de información estadística (art. 5 de los Estatutos), que es beneficiosa tanto para los Gobiernos como para los operadores de los mercados financieros y el propio BCE. Igualmente, el BCE asume una función consultiva, de acuerdo con lo dispuesto en los artículos 127.4 TFUE y 4 de los Estatutos.

g) Los problemas experimentados por algunos Estados miembros para financiarse en los mercados y la consiguiente articulación de instrumentos de asistencia financiera (entre otros, el Mecanismo Europeo de Estabilidad, que es formalmente una organización internacional independiente), han llevado a asignar al BCE funciones que no estaban contempladas en los Tratados constitutivos. Esta atribución de funciones en instrumentos internacionales ha sido autorizada mediante Decisiones de los representantes de los Gobiernos de los Estados miembros de la UE21. En el asunto Pringle, el TJ ha señalado que la asunción de esas funciones por el BCE era compatible con el derecho originario mientras éstas no desvirtúen las competencias que los Tratados atribuyen a esta institución y concuerden con ellas22. Como ha señalado Hinojosa (2014, 225), de esta manera se ha flexibilizado significativamente el alcance del principio de atribución institucional.

En el desarrollo de todas estas funciones, las personas que integran los órganos rectores y el personal del SEBC tienen acceso a numerosa información confidencial procedente de las distintas entidades financieras. Esto explica la exigencia de confidencialidad que el artículo 37.1 de los Estatutos exige a los miembros del BCE y de los bancos centrales nacionales en relación con la información sometida a secreto profesional, «incluso después de cesar en sus  funciones». La obligación de confidencialidad se extiende a todas las personas que «tengan acceso a datos amparados por la legislación de la Unión» (art. 37.2 de los Estatutos).

El Capítulo VI de los Estatutos contiene las disposiciones que regulan el régimen financiero del SEBC. Este régimen específico implica la exclusión de las cuentas financieras del SEBC del presupuesto de la Unión (lo que acentúa, una vez más, la independencia del BCE). 

\paragraph{EL PODER NORMATIVO Y SANCIONADOR DEL BCE }

En una primera aproximación, puede sorprender el importante poder normativo que los Tratados han otorgado al BCE, sobre todo si se tienen en cuenta la independencia con la que actúa, y la  circunstancia de que sus actos no necesitan por lo general de la intervención de ninguna otra Institución para producir efectos obligatorios.

En efecto, los artículos 132 TFUE y 34 de los Estatutos conceden al BCE la capacidad para dictar reglamentos, decisiones, recomendaciones y dictámenes, tal y como estos actos son definidos en el artículo 288 TFUE. En consecuencia, las reglas establecidas en los artículos 296 a 299 TFUE son aplicables a los reglamentos y las decisiones emanados del BCE (normas sobre motivación, publicación, transparencia y ejecutoriedad de los actos típicos obligatorios), aunque éstos deban calificarse como «actos no legislativos» por exclusión de la definición que realiza el artículo 289 TFUE. En definitiva, el BCE puede adoptar toda la tipología de actos recogidos en el artículo 288 TFUE con excepción de las directivas.

El alcance del poder normativo del BCE, sin embargo, se encuentra con dos límites materiales importantes. En primer lugar, con el principio de atribución de competencias, reiterado en el ámbito institucional por el artículo 13.2 TUE. En segundo lugar, los artículos 132.1 TFUE y 34.1 de los Estatutos limitan taxativamente los ámbitos materiales en los que el BCE podrá elaborar reglamentos23, que es el instrumento normativo más poderoso. Junto a los actos típicos mencionados en los artículos 132 TFUE y 34 de los Estatutos, el BCE también puede adoptar las orientaciones e instrucciones mencionadas en los artículos 12.1 y 14.3 de los Estatutos. El artículo 14.3 indica que el BCE formulará orientaciones e instrucciones para los bancos centrales nacionales que son de obligado cumplimiento para los mismos. Paralelamente, el artículo 12.1 establece que el Consejo de Gobierno podrá adoptar orientaciones que vincularán al Comité Ejecutivo y que éste, a su vez, impartirá instrucciones a los bancos centrales nacionales para la puesta en práctica de la política monetaria de la Unión24. Mientras las orientaciones se configuran como una norma más general, que puede dejar un cierto margen de discrecionalidad a sus destinatarios, las instrucciones tienen un contenido más concreto, o incluso directamente ejecutivo. Nos hallamos, por tanto, ante actos que persiguen configurar el SEBC como un sistema coherente, que debe actuar con arreglo a criterios unitarios. Ello no tiene por qué impedir que de dichos actos puedan derivarse derechos para terceras personas (fuera del Sistema) en determinadas circunstancias, como ocurre por ejemplo en el caso de las directivas (Smits, 1997: 105; Zilioli/Selmayr, 2001: 108).

En cualquier caso, los actos del BCE están sometidos al control jurisdiccional del TJUE, en la medida en que producen efectos jurídicos obligatorios (arts. 263 TFUE y 35 de los Estatutos).

Por lo que se refiere a la potestad para adoptar los actos jurídicos que menciona el artículo 34 de los Estatutos, esta disposición habla del BCE, en general, de lo que se puede deducir que tanto el Consejo de Gobierno como el Comité Ejecutivo podrían tener capacidad para elaborar reglamentos y decisiones en el desarrollo de sus competencias. Aunque el artículo 17.1 del Reglamento interno del BCE atribuye en exclusiva al Consejo de Gobierno la potestad de adoptar reglamentos, el artículo 12.1 de los Estatutos  permite a este órgano delegar sus competencias en el Comité Ejecutivo. Más específicamente, el artículo 17.3 del Reglamento interno del BCE permite al Consejo de Gobierno delegar sus competencias normativas en el Comité Ejecutivo, aunque sólo para la ejecución de sus reglamentos y orientaciones.

El BCE tiene también un importante poder sancionador,
reconocido por el artículo 34.3 de los Estatutos. Esta disposición
establece que \textit{el BCE estará autorizado a imponer multas y pagos periódicos coercitivos a las empresas que no cumplan con sus obligaciones respecto de los reglamentos y decisiones del mismo}.

Este poder sancionador debe ejercerse «dentro de los límites y en las condiciones establecidas por el Consejo». El Reglamento 2532/98 del Consejo establece el necesario marco normativo para el desarrollo de esta competencia sancionadora25, que resulta también de aplicación en el ámbito de las competencias del BCE sobre supervisión prudencial26. En este Reglamento se articula un mecanismo de investigación en el que participan los bancos centrales nacionales, que tiene un carácter contradictorio, y que es susceptible de revisión judicial. Además del régimen general previsto en el Reglamento 2532/98, existen algunas especificidades para la aplicación de sanciones contenidas en algunas otras disposiciones adoptadas sobre la base del artículo 41 de los Estatutos (relativo a la legislación complementaria que debe ser adoptada por el Consejo)27.


\paragraph{LA INDEPENDENCIA DEL BCE Y LA LEGITIMIDAD DEMOCRÁTICA DE SUS DECISIONES}

El TFUE apuesta por la independencia del SEBC como un eje fundamental de su funcionamiento. La opción por un banco centralindependiente como mejor fórmula para garantizar la estabilidad deprecios no es unánime entre los economistas y, aunque ese debatecontinúe en el plano teórico, el propio Derecho originario realiza una clara elección por un alto nivel de independencia. Los mercados financieros reciben así el mensaje de que la política antiinflacionista no estará condicionada por la coyuntura política.

El principio de independencia es formulado en el artículo 130 TFUE (art. 7 de los Estatutos) como una prohibición de «solicitar oaceptar instrucciones de las Instituciones, órganos u organismos de la Unión, ni de los Gobiernos de los Estados miembros, ni de ningún otro órgano», aplicable tanto al BCE como a los bancos centrales nacionales y a «los miembros de sus órganos rectores», a título personal. Este mandato es complementado por la obligación asumida por las Instituciones y órganos de la Unión y los Gobiernos de los Estados miembros de «no tratar de influir» en los miembros de los órganos rectores del BCE y de los bancos centrales nacionales en el ejercicio de sus funciones.

Esta norma central se ve completada por otras disposiciones particulares que buscan reforzar aspectos concretos de la independencia del SEBC. Siguiendo la clasificación de Louis (1995: 64), pueden distinguirse:

a) Independencia institucional, que se plasma en la ausencia de instrucciones externas, la personalidad jurídica propia del BCE y en una importante capacidad de decisión autónoma. A pesar de todo, esta independencia no es absoluta, ya que el Presidente del Consejo y un miembro de la Comisión pueden participar sin derecho de voto en las reuniones del Consejo de Gobierno del BCE, y el Presidente del Consejo puede llegar a «someter una moción a la deliberación del Consejo de Gobierno» (art. 284.1 TFUE). Igualmente, un representante de la Comisión puede participar como observador en las reuniones del Consejo de Supervisión (art. 26.11 del Reglamento 1024/2013).

b) Independencia funcional, que se traduce en la capacidad de decisión autónoma del BCE en el marco de sus competencias, sin perjuicio de los límites que ciertos actos del Consejo pueden establecer en determinados supuestos (p. ej., la legislación complementaria adoptada por el Consejo sobre la base del art. 129.4 TFUE).

c) Independencia personal, que se explica por la duración del mandato de los componentes del Comité Ejecutivo —ocho años norenovables (art. 283.2 TFUE)— y los Gobernadores de los bancoscentrales nacionales —mandato no inferior a cinco años (art. 14.2 de los Estatutos)—, así como el carácter excepcional y grave de las causas que deben concurrir para que estas personas sean separadas de sus cargos (arts. 11.4 y 14.2 de los Estatutos). El presidente del Consejo de Supervisión también goza de un mandato de cinco años no renovable, aunque puede ser destituido mediante una decisión de ejecución del Consejo, adoptada por mayoría cualificada, previa propuesta del BCE aprobada por el Parlamento Europeo (art. 26.4 del Reglamento 1024/2013).

d) Independencia financiera, en la obtención y gestión de susrecursos36. El artículo 123 TFUE refuerza esta autonomía prohibiendo el acceso privilegiado al crédito (del BCE o de los bancos centrales nacionales) de las entidades públicas europeas o nacionales.

e) Independencia de gestión interna, en especial en lo que se refiere a la gestión de recursos humanos. Como puede apreciarse, estas garantías de independencia afectan tanto al BCE como a los bancos centrales nacionales. Esto resulta imprescindible para la independencia del SEBC, ya que de otra manera los Gobiernos conservarían una posibilidad de influencia indirecta en el BCE a través de su capacidad de interferencia en el banco central nacional. Esto explica que el artículo 131 TFUE obligue a los Estados miembros a adecuar los Estatutos de los respectivos bancos centrales nacionales a estas exigencias.

La necesidad de salvaguardar la independencia de la política monetaria ha llevado a establecer una división interna dentro del BCE, de manera que el procedimiento de adopción de decisiones en el ámbito de la supervisión prudencial se prepare en un órganoajeno a la política monetaria (el Consejo de Supervisión), y que cuando el Consejo de Gobierno tenga que adoptar decisiones en el ámbito prudencial, lo haga en reuniones diferenciadas (art. 25.4 del Reglamento 1024/2013).

La gran cantidad de garantías de independencia de las que goza el SEBC ha sido criticada por un sector de la doctrina, que consideraque se incurre en un déficit democrático a la hora de adoptar las decisiones de política monetaria en la Unión, y que el Tratado debería reformarse para permitir que las autoridades políticas  democráticamente elegidas pudieran revocar las decisiones del BCE (Gormley/Haan, 1996: 109). En cualquier caso, lo que esta polémica pone de manifiesto es que los redactores de los Tratados prefirieron apostar por la solidez de la nueva moneda común (a través de la independencia del SEBC), antes que poner el énfasis en el control democrático de la toma de decisiones monetarias. 

A pesar de todo, se crean tres mecanismos que permiten a las autoridades políticas ejercer un cierto control sobre la actividad delBCE en el plano monetario:

a) El Presidente del Consejo y un miembro de la Comisión pueden participar, sin derecho de voto, en las reuniones del Consejo de Gobierno del BCE. El Presidente del Consejo puede incluso someter una moción para su deliberación en el Consejo de Gobierno. Paralelamente, el Presidente del BCE será invitado a las reuniones del Consejo en las que se delibere sobre cuestiones relativas a los objetivos y funciones del SEBC (art. 284.1 y 2 TFUE). Se propicia así una colaboración institucional y una cierta influencia entre el BCE y el Consejo.

b) El BCE debe presentar una serie de informes periódicos sobre sus actividades. En concreto, estas obligaciones incluyen la elaboración de un informe anual sobre las actividades del SEBC y sobre la política monetaria del año precedente y del año en curso (art. 284.3 TFUE, art. 15 de los Estatutos), la publicación de un informe anual sobre las cuentas del BCE (art. 26.2 de los Estatutos), la publicación de informes sobre las actividades del SEBC con una periodicidad al menos trimestral (art. 15.1 de los Estatutos), y la publicación semanal de un estado financiero consolidado del SEBC (art. 15.2 de los Estatutos). Los informes y estados mencionados en el artículo 15 se pondrán gratuitamente a disposición de los interesados (art. 15.4 de los Estatutos). Se potencia así la transparencia del Sistema.

c) El Presidente del BCE presenta su informe anual ante el Parlamento Europeo (y el Consejo), que puede proceder a un debate general sobre el mismo. Igualmente, el Presidente del BCE y los restantes miembros del Comité Ejecutivo, a petición del Parlamento Europeo o por iniciativa propia, pueden comparecer ante las comisiones competentes del Parlamento Europeo (art. 284.3 TFUE). En estos debates los miembros del Comité Ejecutivo deberán explicar los motivos de sus decisiones y podrán hacerse una idea clara de lo que desean los representantes de los ciudadanos.

En el contexto del control de la actividad de supervisión prudencial del BCE también se han establecido diversos mecanismos de rendición de cuentas: 

a) El BCE elaborará un informe anual sobre la ejecución de sus funciones supervisoras que será presentado públicamente por el Presidente del Consejo de Supervisión ante el Eurogrupo y el Parlamento Europeo (art. 20.3 del Reglamento 1024/2013).

b) Tanto el Eurogrupo como el Parlamento Europeo pueden solicitar al Presidente del Consejo de Supervisión que participe en una audiencia sobre el desempeño de sus funciones37. El BCE responderá oralmente o por escrito a las preguntas que le sean formuladas por el Parlamento Europeo o por el Eurogrupo, con arreglo a sus propios procedimientos (párrafos 5.º y 6.º del art. 20.6 del Reglamento 1024/2013). El BCE y el Parlamento Europeo han celebrado un Acuerdo interinstitucional que regula todos los procedimientos de rendición de cuentas democrática en relación con las funciones de supervisión prudencial38.

c) El informe anual mencionado en el apartado a) también se enviará a los parlamentos nacionales. Éstos podrán dirigir al BCE sus observaciones motivadas, solicitarle que responda por escrito a cualquier observación o pregunta e, incluso, invitar a un representante del Consejo de Supervisión (junto con un representante de la autoridad nacional competente) a participar en un cambio de impresiones sobre la supervisión prudencial de las entidades de crédito de ese país (art. 21 del Reglamento 1024/2013).

Al margen de la virtud que puedan demostrar estos mecanismos de control, se echa de menos una participación más decisiva del Parlamento Europeo en la adopción de la legislación complementaria (el art. 129.4 TFUE prevé la mera consulta), o en el nombramiento de los miembros del Comité Ejecutivo (art. 283.2  TFUE), en los que desempeña un rol meramente consultivo. Estas reformas podrían redundar en un incremento de la legitimidad democrática del entramado orgánico creado en el proceso de integración monetaria, sin cuestionarse por ello la independencia del BCE en el desarrollo de sus funciones.



5. EL CONTROL JURISDICCIONAL

El artículo 35 de los Estatutos regula el sistema de control jurisdiccional de los actos adoptados en el marco de la actividad del SEBC. Esta norma se complementa con las disposiciones del TFUE relativas a la competencia general del TJUE, configurándose así un sistema exigente de revisión judicial, equiparable al que se aplica a los actos de las demás Instituciones europeas. Adicionalmente, en el ámbito de la supervisión prudencial, se ha creado un Comité Administrativo de Revisión y el procedimiento correspondiente para que cualquier persona física o jurídica a la que concierna o afecte directa e individualmente una decisión del BCE en ese terreno pueda solicitar que ésta se reexamine, sin perjuicio de su derecho a interponer un recurso ante el TJUE.

El párrafo 1.º del artículo 35 indica que tanto los actos como las omisiones del BCE están sujetos a la revisión (control de validez) y a la interpretación del TJUE. Esta legitimación pasiva del BCE en el recurso de anulación se ve acompañada por la legitimación activa que en el mismo tipo de recurso le reconoce el artículo 263 TFUE con el fin de salvaguardar sus prerrogativas. Si un acto del BCE es anulado por el TJUE, el primero tiene la obligación de «adoptar las medidas necesarias para la ejecución de la sentencia» (art. 266 TFUE), como cualquier otra Institución. En el recurso por omisión, el BCE goza igualmente de legitimación activa y pasiva (art. 265 TFUE). Además, el TJUE puede pronunciarse sobre las cuestionesprejudiciales que se planteen sobre la validez o la interpretación de los actos del BCE [art. 267.b) TFUE], y cabe igualmente alegar la excepción de ilegalidad contra un acto de alcance general del BCE (art. 277 TFUE).

La responsabilidad extracontractual del BCE se rige, en tanto que Institución europea, por el artículo 340 TFUE, es decir, que sigue los mismos principios ya enunciados por el TJUE en relación con las demás Instituciones39. No obstante, la UE no responde de los actos del BCE, cuyas consecuencias asume él en solitario. En cambio, la responsabilidad extracontractual de los bancos centrales nacionales debe regirse por el Derecho interno, ya que éstos son órganos de los Estados miembros. Esta cuestión puede plantear problemas complejos en los casos en los que sea preciso delimitar la responsabilidad de los daños causados por un banco central nacional cuando actúa siguiendo las órdenes del BCE. En la sentencia Krohn, el TJ entendió que, cuando la administración nacional realiza los actos como mero agente de la Comisión, siguiendo sus órdenes, la responsabilidad extracontractual puede corresponder a la Comisión, aunque en la misma sentencia se indica que la admisibilidad de la acción de indemnización «puede estar subordinada, en determinados casos, al agotamiento de las vías de recurso internas» si éstas «aseguran de un modo eficaz la protección de los particulares interesados»40. Sin embargo, el TJUEno tiene por qué aplicar el mismo criterio de responsabilidad extracontractual en el contexto del SEBC. El carácter cohesionado y unitario de la acción de los bancos centrales nacionales a las órdenes del BCE podría fundamentar sólidamente una acción indemnizatoria directa contra el BCE, sin necesidad de recurrir a las jurisdicciones nacionales. En cualquier caso, la naturaleza  económica de las decisiones del BCE parece asegurar un amplio margen de apreciación a esta Institución, lo que dificulta en gran medida el éxito de esas acciones indemnizatorias41 (Smits, 1997: 110). Por su parte, el BCE asume una función similar a la de la Comisión en el recurso por incumplimiento en relación con los bancos centrales nacionales. En efecto, el artículo 35.6 de los Estatutos establece que cuando el BCE considere que algún banco central nacional ha incumplido alguna de las obligaciones establecidas en los Estatutos, emitirá un dictamen motivado al respecto, después de haber dado a dicho banco central nacional la posibilidad de presentar sus alegaciones. Si el banco central nacional no se atuviere a este dictamen en el plazo establecido por el BCE, éste podrá recurrir al TJUE. Se pone así de manifiesto la independencia del SEBC, ya que la acción se dirige directamente contra el banco central nacional y no contra el Estado miembro, aunque se trata de un órgano nacional. Por tanto, no es posible utilizar en este contexto el mecanismo sancionador previsto en el artículo 260 TFUE.

El Tribunal ha señalado en una jurisprudencia ya consolidada que la forma que adopta un acto es indiferente a la hora de dilucidar si puede ser atacado mediante un recurso de anulación. Enconsecuencia, los actos del BCE sujetos a la revisión o lainterpretación del TJUE son todos los actos que produzcan «efectos jurídicos»42, y no sólo los reglamentos y las decisiones. En tales circunstancias, el BCE debe respetar escrupulosamente las obligaciones que establecen los Tratados en materia de motivacióno de transparencia, por ejemplo, para evitar que alguno de sus actos sea anulado por el TJUE.








El Banco Central Europeo (BCE) tiene el rango de institución oficial de la Unión desde la entrada en vigor del Tratado de Lisboa. Siguiendo lo dispuesto en el Tratado de Maastricht, el BCE se creó el I de junio de 1998 en Fráncfort y, desde el I de enero de 1999, se hace cargo de la instrumentación de la política monetaria de la Unión Económica y Monetaria. El BCE se sitúa en el núcleo del Sistema Europeo de Bancos Centrales (SEBC), que agrupa al BCE y a los Bancos Centrales nacionales de todos los Estados miembros de la UE.

El Consejo General del SEBC: Está fonnado por los Gobernadores de los Bancos Centrales de los 27 Estados miembros, junto con el Presidente y el Vicepresidente del BCE. El Comité Ejecutivo del BCE: Está compuesto por el Presidente, el Vicepresidente y otros cuatro miembros, todos nombrados por el Consejo Europeo por mayoría cualificada para un mandato de ocho años. El Consejo de Gobierno del BCE: Está fonnado por los seis miembros del Comité Ejecutivo del BCE y los Gobernadores de los Bancos Centrales nacionales de los Estados miembros de la zona del euro: todos ellos forman el .. Eurosistema". El Consejo de Gobierno es el principal órgano de toma de decisiones del BCE.

































\clearpage
\section{Unión Europea: Gestión de Crisis}
\puntosclave{

}


Puede decirse que, desde la entrada en vigor del Tratado de Lisboa, la Unión Europea se ha encontrado inmersa en un terreno pantanoso cercado por el encadenamiento de crisis (económica y financiera, migratoria, el shock de Brexit, la pandemia de COVID- 19, la guerra en Ucrania, etc.). Frente a estos retos, la Unión E uropea ha tratado de ser parte de la solución utilizando todos los mecanismos a su disposición. A continuación, repasamos brevemente algunos de los desarrollos institucionales de esta fase, que sigue abierta.

\subsubsection{Crisis Políticas} 
Por primera vez, el Tratado de Lisboa reconoció expresamente el derecho de sus Estados miembros a tomar la decisión unilateral de retirarse de la l)nión, evitando así cualquier asimilación superficial de la UE a una estatalidad. Al poco tiempo de que se abriera este cauce. la situación con Reino Unido puso a prueba el artículo 50 TUE en la práctica. Las tensiones entre el Reino Unido y sus socios europeos en demanda de un encaje especial habían sido una constante a lo largo de sus años como Estado miembro de la UE2.\footnote{%
    Por ejemplo: reducción contribución al presupuesto mediante el "cheque b1itánico". autoexclusión del espacio Schengen. cláusula de exención de la Unión Económica y Monetaria. recelo respecto al proyecto de Constitución, oposición al Tratado de Estabilidad. Cooperación y Gobemanza y su Fiscal Compact. etc.
} En un último pulso, el Primer Ministro David Cameron negoció en febrero de 20 16 unas nuevas condiciones especiales. Sobre la base de este paquete. convocó un referéndum sobre la permanencia de Reino Unido en la Unión, celebrado el 23 de junio de 20 16. El resultado de la consulta, favorable a la salida de Reino Unido, abrió un periodo de incertidumbres económicas e institucionales.

Reino Unido notificó oficialmente su deseo de abandonar la Unión el 29 de marzo de 20 17. La fecha de notificación es relevante ya que, a partir de ella, se abre el periodo de dos años que establece el artículo 50 .3 para la negociación de la retirada. En el caso de cumplirse el plazo sin alcanzar un acl!erdo, puede haber graves efectos y disrupciones, si bien no existe un vacío total ya que las relaciones comerciales pasan a estar regidas por las reglas de la OMC. También existe la posibilidad de prorrogar el plazo por unanimidad del Consejo Europeo. Esto es lo que sucedió tres veces debido a los problemas parlamentarios internos que tuvo Reino Unido para aceptar el acuerdo negociado.

Finalmente, el Acuerdo de Retirada fue suscrito el 2 4 de enero de 2020 por la Presidenta de la Comisión, el Presidente del Consejo Europeo y el Primer Ministro Boris Johnson, tan solo una semana antes de agotarse la última prórroga.\footnote{%
    El Acuerdo de Retirada garantiza la retirada ordenada del Reino Unido de la Unión. Abarca los derechos de los ciudadanos. la liquidación financiera. un período transitorio. sendos protocolos sobre Irlanda e Irlanda del No11e. Chipre y Gibraltar. la gobemanza y otras cuestiones relativas a la separación.
} Previamente, había sido aprobado por el Parlamento británico el 23 de enero. El Parlamento Europeo lo hizo el 29 de enero (como acuerdointernacional entre la UE y el Estado que se retira, el Acuerdo no es mixto y no requirió la autorización de los Parlamentos Nacionales). Por su parte, el Consejo adoptó la decisión relativa a la celebración del Acuerdo de Retirada en nombre de la UE el 30 de enero. Reino Unido dejó de ser Estado miembro a las 00 :00 del 1 de febrero de 20 20 .

Hasta esa fecha. Reino Unido había seguido estando vinculado por todas las obligaciones de cumplimiento del Derecho de la UE, así como conservaba sus derechos y su presencia en las instituciones. Hasta el 31 de diciembre de 20 20, hubo un periodo transitorio en el que Reino Unido siguió vinculado por determinadas obligaciones jurídicas y financieras. pero ya sin participar en las instituciones. El Acuerdo de Retirada incluía una declaración política relativa al marco de relaciones futuras, las cuales se empezaron a negociar tras la retirada. El Acuerdo de Comercio y Cooperación entre la UE y el Reino Unido, que establece un acuerdo de libre comercio y una cooperación estrecha en una variedad de áreas, empezó a aplicarse provisionalmente el I de enero de 20 2 1 .

Desde el punto de vista jurídico e institucional de la UE, la retirada de Reino Unido solo ha supuesto algunos ajustes en la composición de las instituciones europeas. En la decisión del Consejo Europeo de j unio de 20 18 que establecía la composición del Parlamento Europeo tras las elecciones de 20 19, se previeron 100 alteraciones que se se producirían cuando la retirada de Reino Unido (que contaba con 73 escaños) fuera jurídicamente efectiva. Por un lado, algunos Estados miembros, incluida España, han visto aumentada su representación.\footnote{%
    Incremento en el número de MEPs de algunos EEMM: +5. ES, fR: +3, IT. NL: +2. IE: + l . PL. RO. SE, AT. DK. SK. FL HR, EE.
} Por otro, el número total de escaños se ha reducido de 751 a 705. El número de escaños hasta alcanzar el máximo previsto por el TUE queda en reserva ante futuras ampliaciones.

En el Consejo y en el Consejo Europeo, Reino Unido perdió su representación cuando dejó de ser miembro de la UE. No obstante, desde el momento de la notificación, ya no participaba en las deliberaciones que afectaban a las negociaciones de retirada, tal y como establece el artículo 50 (por lo que el Consejo y el Consejo Europeo se reunían a 27 en formato del artículo 50). Reino Unido renunció a su presidencia rotatoria del Consejo del segundo semestre de 20 17, provocando la ampliación del turno correspondiente al gobierno estonio.

El Consejo Europeo mantuvo la composición de la Comisión con un miembro por cada Estado. Tras el referéndum, el que era Comisario europeo de Estabilidad Financiera, Jonathan Hill, dimitió, siendo sustituido por Julian King como Comisario de seguridad entre 20 16 y 20 19. La Comisión von der Leyen ya comenzó sin miembro de nacionalidad británica.





\subsubsection{Crisis Económicas y Financieras}
Poco después de la entrada en vigor del Tratado de Lisboa, la larga y grave crisis económica y financiera iniciada en 2008 exigió utilizar varios de los mecanismos de flexibilidad previstos en el Tratado para reforzar la gobernanza económica de la UE. En concreto:

El procedimiento simplificado de. revisión de los Tratados se utilizó en 201 1 para pem1itir la creación el Mecanismo de Estabilidad. Así, por Decisión del Consejo Europeo de 25 de marzo de 20 1 1, se añadió un tercer párrafo al artículo 136 del TFUE, que dice lo siguiente: "los Estados miembros cuya moneda es el euro podrán establecer un mecanismo de estabilidad .. . ", pero " ...l a concesión de toda ayuda. financiera necesaria con arreglo al mecanismo se supeditará a condiciones estrictas".  Gracias a esta disposición, el Tratado Constitutivo del Mecanismo Europeo de Estabilidad (MEDE), que constituye un acuerdo intergubernamental de los países de  la zona del euro, fue suscrito en febrero de 20 12. De forma complementaria (como requisito para recibir asistencia del MEDE, pero abierto a todos tos Estados miembros), se firmó el Tratado de Estabilidad, Coordinación y Gobernanza (TECG) en la Unión Económica y Monetaria· en marzo de 20 12, del que se autoexcluyeron Reino Unido y República Checa. Este instrumento intergubernamental, que incluye el llamado Pacto Fiscal, entró en vigor en 20 13 haciendo uso de la flexibilidad que permite la "cooperación reforzada". Al mismo tiempo, durante la crisis, se utilizó en toda su extensión el artículo 136 del TFUE, que permite adoptar medidas específicas en la zona del euro que vayan más allá de las disposiciones aplicables a todos los Estados miembros para reforzar la coordinación y la supervisión de su disciplina presupuestaria. Esto ha permitido, por ejemplo, establecer la posibilidad de imponer sanciones a los países de la zona euro en el marco del Pacto de Estabilidad y Crecimiento (en virtud del paquete legislativo "Six-Pack", que fue aprobado en 20 1 1) , o exigirles que presenten su Borrador de Presupuesto cada año para que la Comisión lo ·evalúe antes de aprobarlo (en virtud del paquete legislativo "Two-Pack", que fue aprobado en 20 13).


\subsubsection{Reflexiones sobre el futuro de la Unión Europea}
El resultado del referéndum británico en 20 16 se unió a una serie de circunstancias (consecuencias sociales de la crisis económica y financiera, aumento del nacionalismo y la xenofobia, nuevo papel geopolítico de China y Rusia, crisis climática y medioambiental, etc.) que generaron cierto sentimiento de crisis existencial de la UE.

Con el fin de responder a estas cuestiones y con motivo del 60 aniversario de los Tratados de Roma, la Comisión Juncker presentó un Libro Blanco en marzo de 20 l ? ' estableciendo cinco escenarios posibles para el futuro de la Unión, así como cinco documentos de reflexión para las áreas social, de defensa, Unión. Económica y Monetaria, globalización, y Marco Financiero Plurianual 20 2 1 -2027 (antes de que la Comisión presentara su propuesta en mayo de 20 18). A partir del Libro Blanco se celebraron una serie de diálogos con los ciudadanos entre 20 17 y 20 19, cuyo resultado también sirvió para informar a los Jefes de Estado y de Gobierno (reunión informal en Sibiu el 9 de mayo de 20 19; Agenda Estratégica adoptada en las conclusiones del Consejo Europeo de junio de 20 19, a la vista del nuevo ciclo institucional).

El énfasis en la participación de los ciudadanos en el proyecto europeo viene a responder a las críticas de déficit democrático que ha recibido tradicionalmente la UE, cuyo funcionamiento, a pesar de la creciente relevancia del Parlamento, se ha seguido percibiendo corno opaco. Siguiendo esta estela, la Presidenta de la Comisión von der Leyen hizo suya, desde su discurso de investidura. la propuesta del Presidente Macron de poner en marcha una Conferencia sobre el Futuro de Europa con el objetivo de involucrar más a los ciudadanos en el proceso de torna de decisiones y en el establecimiento de las prioridades de la Unión.

La Conferencia sobre el Futuro de Europa tuvo lugar entre marzo de 2021 y el 9 de mayo de 2022, presidida por la Presidenta de la Comisión. el Presidente del Parlamento y el Presidente del Consejo. La Conferencia se articuló a través de una plataforma digital rnultilingüe, paneles europeos de ciudadanos, paneles nacionales y plenos de la Conferencia, dando lugar a un infonne final con propuestas. En junio de 2022, la Comisión presentó una Comunicación con una primera valoración de las propuestas, el Consejo Europeo tomó nota de ellas y declaró que las instituciones de la UE deben garantizar una actuación consecutiva eficaz al informe.

Por parte del Parlamento, existen demandas de refonna de los Tratados (para, por ejemplo. convertirse en colegislador del Marco Financiero Plurianual, reforzar el procedimiento del artículo 7. o reforzar el Pilar Europeo de Derechos Sociales). No obstante, no es probable ninguna modificación de los Tratados en el futuro próximo. Respecto a la posibilidad de utilizar las cláusulas pasarela previstas en el Tratado de Lisboa para pasar de la votación por unanimidad a la votación por mayoría cualificada, la Comisión es partidaria de hacerlo en algunos ámbitos (energía, imposición, sanciones en política exterior, etc.), pero tampoco es probable que pueda alcanzarse la unanimidad de los Estados miembros que ello requiere.































\clearpage
\section{Conclusión} 


\subsection{Recapitulación}


\subsection{Extensiones y relación con otras partes del Temario}



\subsection{Opinión}

\subsubsection{Idea final (Salida o cierra)}

\clearpage
\pagenumbering{gobble}
\MyBibliography



\MyBibItem{DAWID y GATTI}{Chapter 2 - Agent-Based Macroeconomics}{2018}{https://doi.org/10.1016/bs.hescom.2018.02.006}


% ANEXOS
\clearpage
\anexo{\textbf{Preguntas Test}}
%\input{\TemaCode/questions.tex} 



\clearpage
\anexo{Bloque de la Unión Europea: Cuadro de asociación}\label{anx:bloque_ue}
Este es el primer tema del bloque corresponiente a la Unión Europea, que se expande en diez temas individuales. El temario de la oposición a TCEE tiene un carácter eminentemente económico, sin embargo, el Bloque de temas de la Unión Europea exige, en aras de la contextualización y la \textit{raison d'être} de éstos, un ejercicio de vinculación de éstos con los objetivos y obligaciones que establecen tanto las normas originarias del Derecho de la Unión como su normativa derivada: no se puede comprender y exponer un tema sin poder relacionarlo y contextualizarlo dentro de los objetivos de la Unión.































\end{document}